\documentclass{MastersThesis}

% ============================================================================
% ADDITIONAL PACKAGES
% ============================================================================

% Math (\text, \overline, aligned environments used in IR and rendering sections)
\usepackage{amsmath}
\usepackage{amssymb}

% Map literal Unicode characters used in the source to LaTeX equivalents
% (the T1/lmodern fonts otherwise silently drop them).
\usepackage{textcomp}
\usepackage{newunicodechar}
\newunicodechar{—}{\textemdash}
\newunicodechar{–}{\textendash}
\newunicodechar{§}{\S}
\newunicodechar{°}{\textdegree}
\newunicodechar{→}{$\rightarrow$}
\newunicodechar{✓}{\checkmark}
\newunicodechar{✗}{$\times$}

% Bibliography with natbib/bibtex (tectonic runs bibtex internally; no external biber needed)
\usepackage[numbers,sort&compress]{natbib}

% Code listings for IR/YAML examples
\usepackage{listings}
\usepackage{inconsolata}

% Graphics
\usepackage{graphicx}

% Tables
\usepackage{booktabs}
\usepackage{tabularx}
\usepackage{longtable}

% Captions
\usepackage{caption}

% For inline code and verbatim
\usepackage{fancyvrb}

% ============================================================================
% LISTINGS CONFIGURATION (for YAML/IR examples)
% ============================================================================

\lstset{
  basicstyle=\ttfamily\small,
  breaklines=true,
  breakatwhitespace=true,
  frame=single,
  framerule=0.5pt,
  rulecolor=\color{gray!50},
  backgroundcolor=\color{gray!5},
  numbers=left,
  numberstyle=\tiny\color{gray},
  numbersep=8pt,
  tabsize=2,
  showstringspaces=false,
  captionpos=b
}

\lstdefinelanguage{yaml}{
  keywords={true, false, null},
  sensitive=false,
  comment=[l]{\#},
  morestring=[b]',
  morestring=[b]"
}

\lstdefinelanguage{javascript}{
  keywords={const, let, var, function, return, import, from, export, async, await, if, else, for, while, new, class, extends, this, typeof, null, undefined, true, false},
  sensitive=true,
  comment=[l]{//},
  morecomment=[s]{/*}{*/},
  morestring=[b]',
  morestring=[b]",
  morestring=[b]`
}

% ============================================================================
% DOCUMENT METADATA
% ============================================================================

% ============================================================================
% OURDIAGRAM — macro for dogfooded figures rendered by our own compiler
% ============================================================================
% Usage: \ourdiagram[<width>]{<basename>}{<caption text>}
%   <width>    optional, defaults to \linewidth
%   <basename> filename stem in design/figures/ (no extension, no path)
%   <caption>  figure caption
\newcommand{\ourdiagram}[3][\linewidth]{%
  \begin{figure}[htbp]
  \centering
  \includegraphics[width=#1]{figures/#2.png}
  \caption{#3}
  \par\vspace{0.4em}
  \textit{\small Rendered by the diagram compiler this document describes
    (source: \texttt{design/figures/src/#2.mmd}, built with \texttt{make figures}).}
  \end{figure}%
}

\title{A Mermaid-Superset Diagram Compiler\\[0.5em]\Large Beautiful, Deterministic, Agent-Authorable\\[0.3em]\normalsize Architecture \& Specification}
\author{Timeline Compiler Project}
\date{June 2026}

% ============================================================================
% DOCUMENT
% ============================================================================

\begin{document}

\maketitle

\begin{quote}
\textbf{This document is dogfooded.}
Every diagram in this specification is rendered by the compiler the document describes.
Diagram sources live in \texttt{design/figures/src/} as \texttt{.mmd} files authored in our Mermaid-superset DSL;
running \texttt{make figures} from the \texttt{design/} directory invokes the CLI
(\texttt{timeline render <name>.mmd --format png --scale 3}) to produce high-DPI PNGs that are
then included via \texttt{\textbackslash{}ourdiagram}.
This reusable recipe works for any \LaTeX{} document: author diagrams in our DSL,
render to PNG, include with \texttt{\textbackslash{}includegraphics}.
Vector PDF output is a future enhancement (see \S\ref{sec:backends}).
\end{quote}

\tableofcontents
\newpage

% ============================================================================
% PART I — THESIS & POSITIONING
% ============================================================================

\part{Thesis \& Positioning}
\label{part:thesis}

\section{Problem Definition}
\label{sec:problem}

\subsection{The Communication Gap}

Executives, product managers, consultants, and program leads communicate plans through visual timelines. These artifacts appear in board presentations, investor decks, quarterly reviews, transformation roadmaps, and architecture evolution documents. The quality of these visuals directly affects stakeholder comprehension and decision-making.

Yet producing presentation-quality timelines remains surprisingly difficult. The gap exists because available tools fall into two categories, each failing half the requirement:

\begin{description}
  \item[Documentation Tools] (Mermaid, PlantUML, text-based DSLs) — machine-readable, version-controllable, agent-friendly, but produce diagrams suited for technical documentation rather than executive communication.
  
  \item[Design Tools] (PowerPoint, Visio, Lucidchart, think-cell, Figma) — produce polished visuals, but require significant manual effort, resist automation, and cannot be reliably generated by AI agents.
\end{description}

Timeline Compiler bridges this gap: a compiler target for structured plan data that produces presentation-quality visual output.

\subsection{Why Existing Tools Are Insufficient}

\subsubsection{Mermaid and Documentation-Grade Diagramming}

Mermaid~\cite{mermaid2023} has become the de facto standard for diagrams-as-code in technical documentation. Its Gantt syntax illustrates the limitation:

\begin{lstlisting}[language=yaml,caption={Mermaid Gantt — functional but not presentation-grade}]
gantt
    title Product Roadmap Q3
    dateFormat YYYY-MM-DD
    section Platform
    Core API :a1, 2026-07-01, 30d
    Mobile SDK :a2, after a1, 45d
    section Integrations  
    Partner A :b1, 2026-08-01, 60d
\end{lstlisting}

Mermaid excels at documentation — syntax is simple, output is deterministic, rendering integrates with GitHub, GitLab, and Notion. But the output is not presentation-grade:

\begin{itemize}
  \item Fixed visual style optimized for legibility in documentation contexts
  \item No theming beyond basic color overrides
  \item Limited typographic control
  \item Dependency-centric (task scheduling) rather than communication-centric (timeline narrative)
  \item No built-in support for swimlanes as first-class visual groupings
  \item No annotations, callouts, or emphasis mechanisms for executive communication
\end{itemize}

The result: Mermaid diagrams are useful in READMEs and wikis but inappropriate for board decks and investor presentations.

\subsubsection{Project Management Tools}

Tools like Microsoft Project, Jira, Asana, Monday.com, and Linear are optimized for \emph{managing work}, not \emph{communicating plans}. Their timeline views:

\begin{itemize}
  \item Export poorly or not at all to static visual formats
  \item Embed project-management semantics (resources, dependencies, critical paths) that clutter executive communication
  \item Assume ongoing interaction rather than snapshot communication
  \item Produce visuals tied to the tool's aesthetic, not the organization's brand
\end{itemize}

A quarterly roadmap exported from Jira looks like a Jira export, not like a McKinsey slide.

\subsubsection{Design and Presentation Tools}

PowerPoint, Keynote, Visio, Figma, and Lucidchart produce beautiful results — but at the cost of manual effort. Each timeline is a bespoke artifact:

\begin{itemize}
  \item No structured data input; changes require manual redrawing
  \item Not version-controllable in meaningful ways (binary or JSON blobs)
  \item Not agent-accessible; an AI cannot reliably produce a Visio file
  \item Theming requires manual application
  \item Determinism is impossible; two designers will produce different artifacts from the same brief
\end{itemize}

This creates a maintenance burden. When scope shifts, the timeline must be manually updated. When quarterly plans roll forward, someone rebuilds the slide.

\subsubsection{AI Agent Blindspot}

Modern AI agents (LLMs with tool use, planning agents, copilots) can generate structured plans — sprints, phases, milestones, dependencies. But they have no widely-adopted open-source target for rendering those plans into visual artifacts.

An agent can output JSON or YAML describing a roadmap. But then what? Without a compiler that accepts structured input and produces presentation-grade output, the agent's plan stops at structured text.

\subsection{Who This Is For}

Timeline Compiler serves two complementary audiences:

\subsubsection{Human Authors}

\begin{description}
  \item[Executives and Chiefs of Staff] — Need polished quarterly plans, transformation roadmaps, and board-ready timelines without designer dependency.
  
  \item[Product Managers] — Need roadmaps that communicate priority and sequencing to stakeholders, not dependency graphs for engineers.
  
  \item[Consultants] — Need branded deliverables (BCG-style roadmaps, McKinsey program timelines) that are data-driven yet visually refined.
  
  \item[DevRel and Technical Writers] — Need version-controlled timelines that embed in documentation but also export to slides.
  
  \item[Program Managers] — Need milestone views and phase summaries that roll up operational detail into executive narrative.
\end{description}

\subsubsection{AI Agents}

AI agents require:
\begin{itemize}
  \item A well-defined intermediate representation (IR) they can generate
  \item Deterministic rendering — the same IR always produces the same visual
  \item Theming separation — agents produce structure, not style
  \item Validation — agents can verify their output conforms to the schema
  \item Embeddability — agents can invoke rendering as a tool
\end{itemize}

Timeline Compiler's IR is designed as a \emph{compiler target}: a stable, documented, validatable representation that agents can emit and humans can inspect, version, and override.

\subsection{What Success Looks Like}

\subsubsection{Scenario: Product Roadmap}

A product manager writes (or an agent generates):

\begin{lstlisting}[language=yaml,caption={Product Roadmap IR Example}]
timeline:
  title: "Platform Roadmap 2026"
  range: { start: 2026-Q1, end: 2026-Q4 }
  tracks:
    - id: platform
      label: "Core Platform"
      activities:
        - label: "API v2"
          range: { start: 2026-01, end: 2026-03 }
          status: complete
        - label: "SDK Release"
          range: { start: 2026-04, end: 2026-06 }
          status: in-progress
          progress: 0.6
    - id: mobile
      label: "Mobile"
      activities:
        - label: "iOS App"
          range: { start: 2026-05, end: 2026-08 }
          status: planned
  milestones:
    - label: "Public Beta"
      date: 2026-06-15
      tracks: [platform, mobile]
\end{lstlisting}

The compiler, given a theme (e.g., \texttt{corporate-blue}), produces a PDF suitable for a board presentation: clean typography, proper spacing, consistent color palette, no visual noise.

\subsubsection{Scenario: Transformation Plan}

A management consultant defines a multi-year digital transformation:

\begin{lstlisting}[language=yaml,caption={Transformation Timeline IR Example}]
timeline:
  title: "Digital Transformation — 3 Year Plan"
  range: { start: 2026-01, end: 2028-12 }
  sections:
    - id: foundation
      label: "Foundation"
      range: { start: 2026-01, end: 2026-12 }
    - id: scaling
      label: "Scaling"  
      range: { start: 2027-01, end: 2027-12 }
    - id: optimization
      label: "Optimization"
      range: { start: 2028-01, end: 2028-12 }
  tracks:
    - id: technology
      label: "Technology"
      activities:
        - label: "Cloud Migration"
          range: { start: 2026-03, end: 2026-09 }
          section: foundation
        - label: "Platform Modernization"
          range: { start: 2027-01, end: 2027-09 }
          section: scaling
    - id: organization
      label: "Organization"
      activities:
        - label: "Team Restructuring"
          range: { start: 2026-06, end: 2026-12 }
          section: foundation
  annotations:
    - label: "Board Checkpoint"
      date: 2026-12-15
      style: callout
\end{lstlisting}

The output: a three-year view with clear phase demarcation, swimlanes for workstreams, and a visual hierarchy that communicates strategic sequencing rather than tactical dependencies.

\subsubsection{Success Criteria}

Timeline Compiler succeeds when:

\begin{enumerate}
  \item A human can write the IR by hand in under 10 minutes for a typical roadmap
  \item An AI agent can generate valid IR without special prompting beyond the schema
  \item The same IR, rendered with the same theme, produces byte-identical output across runs (determinism)
  \item A non-designer can produce board-ready visuals by selecting a theme
  \item The IR is diff-friendly — changes to the plan produce readable diffs
  \item The output is accepted in contexts where PowerPoint slides are currently required
\end{enumerate}

\subsection{What This Is Not}

To prevent scope creep and misunderstanding:

\begin{itemize}
  \item \textbf{Not a project management tool} — does not track work, assign resources, or compute schedules
  \item \textbf{Not a Gantt chart replacement} — though it can render timeline views, it is not optimized for dependency visualization
  \item \textbf{Not a Jira/Asana competitor} — does not replace work-tracking systems; may consume data from them
  \item \textbf{Not a design tool} — does not provide a canvas, drag-and-drop, or WYSIWYG editing
  \item \textbf{Not interactive} — output is static (SVG, PNG, PDF, PPTX); interactivity is out of scope
\end{itemize}

Timeline Compiler is a \emph{compiler}: structured input in, presentation-quality visual output out. The value is in the deterministic, themeable, agent-accessible transformation — not in the tooling around authoring or managing plans.

\section{Central Thesis: A Deterministic Diagram Compiler}
\label{sec:central-thesis}

This document specifies a \emph{deterministic, themeable, agent-authorable diagram compiler}: a system that transforms small declarative intermediate representations into presentation-quality visual output through a layered architecture of grammars, a shared kernel, and pluggable rendering backends.

\subsection{The Engine-as-Asset Insight}

The strategic insight underlying this specification is a reframe: \textbf{the real asset is not ``a timeline tool.''} It is a \emph{compiler architecture} capable of hosting a family of visual grammars, each small enough to be reliably emitted by an LLM, each compiling down to a shared Scene IR that guarantees determinism and enables theming across all diagram types.

Timeline is \emph{grammar \#1}---the first proof of the engine, not the whole product. The architecture generalizes to:

\begin{itemize}
  \item \textbf{Flow/pipeline diagrams}: ordered icon+label nodes with arrows, loops, and legends
  \item \textbf{Node-link graphs}: trees, networks, hub-and-spoke, architecture diagrams
  \item \textbf{Comparison matrices}: N columns with checkmarks, stats, or feature rows
  \item \textbf{Stat callouts}: big number + caption compositions
  \item \textbf{Step-card sequences}: numbered ordinal + icon + bullet lists
  \item \textbf{Multi-panel posters}: compositions that arrange multiple grammars into a single visual
\end{itemize}

The value proposition is the \emph{compounding payoff}: animation, themes, new backends, and quality-gates added once in the kernel are inherited by every grammar. A new grammar (e.g., flow diagrams) gains determinism, theming, SVG/PNG/PDF output, and validate-before-render for free.

\subsection{The Pipeline Model}

Every diagram in the compiler follows the same pipeline:

\begin{figure}[h]
\centering
\begin{Verbatim}[frame=single,fontsize=\small]
  Domain IR (grammar-specific)
       |
       v  [Grammar-specific Layout Engine]
       |
  Scene IR (universal primitives)
       |
       v  [Rendering Backend]
       |
  Output (SVG / PNG / PDF / PPTX)
\end{Verbatim}
\caption{The universal pipeline: Domain IR → Layout Engine → Scene IR → Backend → Output}
\label{fig:central-pipeline}
\end{figure}

\paragraph{Validate-before-render.}
The pipeline enforces a strict contract: no diagram reaches the renderer until it passes schema validation and semantic checks. This is not defensive coding; it is a \emph{specification-level guarantee}. An invalid IR produces a diagnostic, never a corrupt visual.

\paragraph{Byte-identical golden determinism.}
The same IR plus the same theme produces byte-identical output across runs, platforms, and time. This property enables CI/CD integration, golden-image testing, and reliable agent round-trips.

\subsection{Why ``Diagram Compiler,'' Not ``Diagram Tool''}

The term ``compiler'' is precise and intentional:

\begin{description}
  \item[Structured input, structured output] The system does not accept freeform drawings or WYSIWYG manipulation. It accepts a declarative specification (the IR) and produces a deterministic artifact.
  
  \item[Separation of concerns] Content (IR), styling (theme), and rendering (backend) are cleanly separated. The IR describes \emph{what} to communicate; the theme describes \emph{how} it looks; the backend describes \emph{where} it goes.
  
  \item[Typed intermediate representations] Like a programming-language compiler, this system uses typed IRs at multiple levels (Domain IR, Scene IR) with well-defined transformations between them.
  
  \item[Determinism by construction] The compiler is a pure function from (IR, theme, backend) to output. No random state, no user interaction, no system clock.
\end{description}

This is not a design tool (Figma, Canva), not a charting library (D3, Plotly), and not a project-management tool (Jira, MS Project). It is a \emph{rendering compiler} for technical explainer diagrams.

\subsection{The Target Niche: Technical Explainer Content}

The inspiration corpus (Section~\ref{sec:corpus-taxonomy}) reveals a coherent niche: \emph{AI-engineering / developer technical-explainer infographics}. These artifacts share visual vocabulary (flows, trees, matrices, step-cards, stat callouts) and share an audience (technical practitioners, executives receiving technical briefings).

This niche is defensible:

\begin{itemize}
  \item \textbf{Small grammars are the feature.} A narrow grammar is both shippable and reliably LLM-emittable. A god-schema that tries to express everything is unmaintainable and hostile to agent generation.
  
  \item \textbf{Automatic layout is required.} The author (human or agent) should not hand-place nodes. The grammar must be constrained enough that a deterministic layout algorithm can position elements without ambiguity.
  
  \item \textbf{The author is often an LLM.} Every design decision must ask: can an agent emit this reliably? If the answer is ``only with special prompting,'' the design is wrong.
\end{itemize}

The explicit \emph{non-goal} is competing with general-purpose freeform canvas tools (Canva, Figma, Excalidraw). Those tools optimize for human creativity and manual placement. This compiler optimizes for deterministic, automated, agent-driven generation of \emph{constrained} visual vocabularies.

\subsection{Elevating the Timeline Grammar Thesis}

Section~\ref{sec:grammar-abstraction-old} (now Part~III) established the ``Timeline Grammar vs.\ Task Scheduling Grammar'' distinction. That thesis remains valid but is now a \emph{special case} of a broader principle:

\begin{quote}
\emph{Every grammar in the compiler adopts a visual-communication abstraction, not an operational/computational abstraction. The IR describes what to show, not what to compute.}
\end{quote}

For timelines, this means: no dependency scheduling, no resource leveling, no critical-path analysis. For flow diagrams, this means: no execution semantics, no data-flow analysis, no simulation. For graphs, this means: no graph algorithms beyond layout---the IR specifies nodes and edges, not traversal results.

The compiler renders \emph{narratives} about technical concepts. It does not execute those concepts.

\subsection{The Kernel/Grammar/Composition Architecture}

The architecture has three layers (detailed in Part~\ref{part:architecture}):

\begin{description}
  \item[Kernel (Part~\ref{part:kernel})] The shared infrastructure: Scene IR, rendering backends, theme system, determinism contract, icon registry, asset loader, lint/quality-gate framework. The kernel knows nothing about timelines, flows, or graphs---it renders primitives.
  
  \item[Grammars (Part~\ref{part:grammars})] Peer modules atop the kernel. Each grammar defines a Domain IR and one or more layout engines that compile Domain IR → Scene IR. Timeline is grammar \#1; future grammars include flow, graph, comparison, stat, step-cards.
  
  \item[Composition (Part~\ref{part:composition})] A thin layer that arranges multiple panels into a poster. Each panel references a grammar's Domain IR. This is how the multi-section infographics are assembled.
\end{description}

\subsection{Summary of the Strategic Evolution}

This specification captures a major strategic evolution:

\begin{enumerate}
  \item \textbf{From ``timeline tool'' to ``diagram compiler.''} The product is the engine, not the first grammar.
  
  \item \textbf{From one IR to two IR layers.} Domain IRs are grammar-specific and small; all compile to the shared Scene IR.
  
  \item \textbf{From SVG-first to Scene-first.} SVG is one backend, not the universal root. The Scene IR is the render contract.
  
  \item \textbf{From timeline-only to grammar family.} Flows, graphs, comparisons, stats, and compositions are in scope.
  
  \item \textbf{From static-only to animation-aware.} Animation is additive, declarative, and backend-conditional.
  
  \item \textbf{From monolith to layered kernel.} The kernel is extractable; grammars are thin consumers.
\end{enumerate}

The following parts specify each layer in detail.

\input{sections/03-principles}
\section{Scope Boundaries}
\label{sec:scope}

This section defines explicit boundaries for the Mermaid-superset diagram compiler.

\subsection{Governing Principle}

The scope boundary is defined by the \emph{Mermaid-superset contract}: this compiler accepts
all valid Mermaid diagrams and our extensions, and produces deterministic visual output.
It does not execute, simulate, schedule, or interactively manipulate diagrams.

\begin{quote}
\emph{If Mermaid renders it, we render it (better). Beyond Mermaid, we add composition,
rich theming, animation, and agent IR. We do not add interactivity, data management,
or operational semantics.}
\end{quote}

\subsection{In Scope}

\subsubsection{Diagram Coverage}

All 22 Mermaid diagram types (Section~\ref{sec:mermaid-compat}), organized into five families,
plus our superset extensions (composition, animation, cross-references).

\subsubsection{Input Paths}

\begin{description}
  \item[Mermaid-superset DSL] Text syntax compatible with Mermaid plus our extensions.
  \item[Structured IR] JSON/YAML conforming to per-grammar Domain IR schemas.
\end{description}

\subsubsection{Output Formats}

\begin{description}
  \item[SVG] Primary vector output (source-of-truth backend).
  \item[PNG] Rasterized output via resvg/Skia.
  \item[PDF] Print-quality vector output for documents and decks.
\end{description}

\subsubsection{Key Capabilities}

\begin{description}
  \item[Theming] Complete visual customization via theme files.
  \item[Determinism] Byte-identical output guarantee.
  \item[Composition] Multi-diagram posters and grid layouts.
  \item[Animation] Declarative SMIL animation in SVG output.
  \item[Agent integration] MCP server, CLI, library API.
  \item[Schema validation] Per-grammar JSON Schema validation with clear diagnostics.
\end{description}

\subsection{Out of Scope}

\begin{description}
  \item[Interactive output] No clickable regions, tooltips, or drill-down. Output is static.
  \item[WYSIWYG editing] No drag-and-drop authoring. This is a compiler, not an editor.
  \item[Data management] No persistent storage, collaboration, or version history.
  \item[Operational semantics] No dependency scheduling, resource leveling, critical-path
        analysis, or simulation.
  \item[General data visualization] Not competing with D3/Plotly for arbitrary charts.
        The chart family covers Mermaid's four chart types only.
  \item[Native integrations] Does not natively read from Jira, Azure DevOps, or other
        systems. Consumes only its IR format or Mermaid text.
  \item[Live collaboration] No real-time multi-user editing.
\end{description}

\subsection{Scope Summary}

\begin{table}[h]
\centering
\caption{Scope Boundaries Summary}
\begin{tabularx}{\textwidth}{lX}
\toprule
\textbf{Category} & \textbf{Status} \\
\midrule
All 22 Mermaid diagram types & In Scope \\
Mermaid-superset DSL parsing & In Scope \\
Structured IR (agent path) & In Scope \\
Beautiful themed output (SVG, PNG, PDF) & In Scope \\
Deterministic byte-identical rendering & In Scope \\
Multi-diagram composition/posters & In Scope \\
Declarative animation & In Scope \\
CLI + library + MCP server & In Scope \\
\midrule
Interactive/clickable output & Out of Scope \\
WYSIWYG/canvas editing & Out of Scope \\
Operational/scheduling semantics & Out of Scope \\
Data storage and collaboration & Out of Scope \\
Native data source integrations & Out of Scope \\
General-purpose data visualization & Out of Scope \\
\bottomrule
\end{tabularx}
\end{table}

\section{Positioning: Why a Mermaid Superset}
\label{sec:comparison}

This compiler is positioned as a \textbf{full Mermaid superset}: every valid Mermaid diagram
parses and renders correctly, but the output is dramatically better-looking, the architecture
is deterministic and composable, and agents get a first-class structured IR path alongside the
text DSL.

\subsection{The Competitive Landscape}

\begin{table}[h]
\centering\small
\caption{Positioning matrix: diagram-as-code tools}
\begin{tabularx}{\textwidth}{lXXXXX}
\toprule
\textbf{Axis} & \textbf{Mermaid} & \textbf{PlantUML} & \textbf{D2} & \textbf{Vega-Lite} & \textbf{Ours} \\
\midrule
Diagram types      & 22 types & UML-centric & General graph & Charts only & 22+ (superset) \\
Visual quality     & Documentation-grade & Dated & Modern, limited & Good charts & \textbf{Presentation-grade} \\
Theming            & Basic color overrides & Skins (limited) & Themes & Spec-level & \textbf{Full theme system} \\
Determinism        & Mostly & No (Graphviz) & No & Yes & \textbf{Byte-identical} \\
Agent IR           & Text only & Text only & Text only & JSON spec & \textbf{Structured IR + DSL} \\
Composition        & None & None & None & Concat/layer & \textbf{Grid posters} \\
UML/software focus & Partial & Primary & No & No & \textbf{Tier-1 priority} \\
\bottomrule
\end{tabularx}
\label{tab:positioning-matrix}
\end{table}

\subsection{Why Full Mermaid Compatibility}

Mermaid~\cite{mermaid2023} is the de facto standard for diagrams-as-code.
It renders natively on GitHub, GitLab, Notion, Obsidian, and hundreds of platforms.
Over 88{,}000 GitHub stars and a \$7.5M raise confirm its ubiquity.
But Mermaid has clear limitations:

\begin{itemize}
  \item \textbf{Ugly output.} Mermaid diagrams look like documentation artifacts, not
        presentation-grade visuals. Fixed styles, poor typography, no coherent design system.
  \item \textbf{No structured IR.} Agents must generate raw text DSL; no JSON/YAML path
        for reliable constrained generation.
  \item \textbf{No composition.} Cannot combine multiple diagrams into posters or multi-panel layouts.
  \item \textbf{Non-deterministic edge cases.} Layout can vary across versions and renderers.
  \item \textbf{Limited theming.} Color overrides only; no typographic or geometric control.
\end{itemize}

Our strategy: \textbf{parse every Mermaid grammar faithfully}, then differentiate on
aesthetics, determinism, composition, and the agent IR. Users can drop in existing Mermaid
files and immediately get better output with zero migration cost.

\subsection{Differentiation Pillars}

\begin{description}
  \item[Beautiful, coherent output] Explicitly beating Mermaid's visual quality.
        Professional themes, proper typography, consistent spacing, dark/light modes.
        Aesthetics is a first-class architectural pillar (Part~\ref{part:aesthetics}).

  \item[UML/software diagram line] Tier-1 investment in class, state, ER, and C4 diagrams---
        the diagrams software teams use daily. PlantUML owns this niche with dated visuals;
        we deliver modern, beautiful UML.

  \item[Agent-authorable structured IR] Dual front-end: humans write
        Mermaid-superset text; agents emit/consume a structured IR (JSON/YAML).
        Both compile to the same domain IR, then to Scene IR, then to backends
        (Part~\ref{part:frontend}).

  \item[Deterministic composition] Byte-identical output; multi-diagram posters;
        cross-diagram references. No other diagram-as-code tool offers this.
\end{description}

\subsection{vs.\ PlantUML}

PlantUML~\cite{plantuml} is Java-based, GPL-licensed, Graphviz-dependent (non-deterministic),
and produces visually dated output. Its UML coverage is comprehensive but its rendering is a
liability in 2026. We capture PlantUML's UML breadth with modern visuals and deterministic output.

\subsection{vs.\ D2}

D2~\cite{d2lang} is a modern diagram language with good aesthetics but no Mermaid compatibility,
no structured IR for agents, limited diagram types, and non-deterministic layout (ELK/dagre).
We offer broader coverage, Mermaid migration path, and byte-identical determinism.

\subsection{vs.\ Vega-Lite}

Vega-Lite~\cite{vegalite2017} is the gold standard for declarative statistical charts---but
it has no concept of node-link graphs, sequence diagrams, or flowcharts. Our chart family
(Section~\ref{sec:chart-family}) draws on Vega-Lite's grammar-of-graphics principles for the
quantitative subset, while covering the full diagram-type space Vega-Lite cannot reach.
            % Sharpened: Mermaid/PlantUML/D2/Vega positioning

% ============================================================================
% PART II — THE FRONT-END (Input)
% ============================================================================

\part{The Front-End}
\label{part:frontend}

\section{Dual Front-End Architecture}
\label{sec:frontend}

The compiler exposes two input paths that converge on the same internal pipeline.
This dual front-end is the architectural key to serving both human authors
(who prefer text DSLs) and AI agents (who prefer structured data).

\subsection{The Two Input Paths}

\ourdiagram[0.95\linewidth]{dual-frontend}{%
  Dual front-end: Mermaid DSL text (Path~A) passes through a PEG/Nearley parser;
  structured IR JSON/YAML (Path~B) passes through schema validation.
  Both paths converge at the same Domain IR, then proceed through the shared
  layout and rendering pipeline to SVG, PNG, and PDF outputs.}

\paragraph{Path A: Mermaid-Superset DSL.}
Human authors write diagrams in Mermaid syntax (or our extensions).
A parser converts the text to the grammar's Domain IR.
This path preserves the ``diagrams-as-code'' ergonomics that made Mermaid popular.

\paragraph{Path B: Structured IR.}
Agents emit JSON or YAML conforming to the Domain IR's JSON Schema.
No parsing step is needed---only schema validation.
This path enables grammar-constrained generation, round-trip editing, and tool integration.

\subsection{Convergence at Domain IR}

Both paths produce the same Domain IR. From that point forward, the pipeline is identical:
Domain IR $\to$ Layout Engine $\to$ Scene IR $\to$ Backend $\to$ Output.

This convergence guarantees:
\begin{itemize}
  \item A Mermaid file and its equivalent YAML produce \textbf{byte-identical output}.
  \item Agents can consume human-authored diagrams (parse DSL $\to$ IR $\to$ emit structured).
  \item Humans can inspect agent-generated diagrams (structured IR $\to$ pretty-print DSL).
  \item Tooling (validation, linting, diffing) works on the common Domain IR regardless of input path.
\end{itemize}

\subsection{The Agent IR as First-Class API}

The structured IR is not a secondary or internal representation---it is a \textbf{first-class
public API}. Design commitments:

\begin{description}
  \item[Published JSON Schema] Every grammar's Domain IR has a versioned JSON Schema,
        suitable for OpenAI Structured Outputs, XGrammar~\cite{dong2024xgrammar},
        or llama.cpp GBNF constrained generation.
  \item[Stable contract] The IR schema follows semver. Breaking changes require a major version bump.
  \item[Round-trip fidelity] DSL $\to$ IR $\to$ DSL preserves semantics (formatting may differ).
  \item[Diff-friendly] Changes to one entity do not cascade to unrelated IR fields.
  \item[Small schemas] Each grammar's schema fits comfortably in an LLM context window
        (typically under 200 lines of JSON Schema).
\end{description}

\subsection{Parser Architecture}

The DSL parser is grammar-specific. Each grammar defines:
\begin{itemize}
  \item A PEG or Nearley grammar specification for its Mermaid-compatible syntax
  \item A CST-to-Domain-IR lowering pass
  \item Extension hooks for superset syntax (new keywords, extended directives)
\end{itemize}

The parser layer is the \textbf{new architectural component} introduced by the Mermaid-superset
positioning. It sits above the existing Domain IR / Layout Engine / Scene IR pipeline and does
not affect the proven kernel or grammar internals.

\subsection{Implementation Status}

\textbf{Built:} The structured IR path (Path B) is fully implemented for all four shipped
grammars (timeline, flow, sequence, tree). Schema validation, layout, and rendering work
end-to-end.

\textbf{Planned:} The Mermaid-superset DSL parsers (Path A) are the primary Tier-0 deliverable.
Parsers for flowchart, sequence, gantt/timeline, and mindmap syntax will wire the existing
Domain IRs to Mermaid text input.
              % Dual front-end: Mermaid-superset DSL + structured IR
\section{Full Mermaid Compatibility}
\label{sec:mermaid-compat}

The compiler targets \textbf{complete compatibility with all 22 Mermaid diagram types}.
Every valid Mermaid file must parse, validate, and render---producing output that is
semantically equivalent to Mermaid's but visually superior.

\subsection{The 22 Diagram Types}

Mermaid supports 22 diagram types as of 2026. We organize them into five families
(Section~\ref{sec:family-taxonomy}) and assign coverage tiers:

\begin{table}[h]
\centering\small
\caption{Mermaid diagram types: families and coverage tiers}
\begin{tabular}{llll}
\toprule
\textbf{Mermaid Type} & \textbf{Family} & \textbf{Tier} & \textbf{Status} \\
\midrule
flowchart       & Node-link / graph     & 0 & Built (Flow grammar) \\
sequence        & UML / software        & 0 & Built (Sequence grammar) \\
gantt           & Timeline / project    & 0 & Built (Timeline grammar) \\
timeline        & Timeline / project    & 0 & Built (Timeline grammar) \\
mindmap         & Tree / hierarchy      & 0 & Built (Tree grammar) \\
\midrule
class           & UML / software        & 1 & Planned \\
state           & UML / software        & 1 & Planned \\
erDiagram       & UML / software        & 1 & Planned \\
C4Context       & Node-link / graph     & 1 & Planned \\
C4Container     & Node-link / graph     & 1 & Planned \\
C4Component     & Node-link / graph     & 1 & Planned \\
C4Dynamic       & Node-link / graph     & 1 & Planned \\
\midrule
pie             & Charts                & 2 & Planned \\
xychart-beta    & Charts                & 2 & Planned \\
quadrantChart   & Charts                & 2 & Planned \\
radar           & Charts                & 2 & Planned \\
\midrule
sankey-beta     & Node-link / graph     & 3 & Planned \\
requirement     & Node-link / graph     & 3 & Planned \\
gitGraph        & Node-link / graph     & 3 & Planned \\
block-beta      & Node-link / graph     & 3 & Planned \\
architecture    & Node-link / graph     & 3 & Planned \\
treemap         & Tree / hierarchy      & 3 & Planned \\
kanban          & Timeline / project    & 3 & Planned \\
journey         & Timeline / project    & 3 & Planned \\
packet          & Node-link / graph     & 3 & Planned \\
\bottomrule
\end{tabular}
\label{tab:mermaid-coverage}
\end{table}

\subsection{Compatibility Strategy}

\subsubsection{Syntax Fidelity}

Each Mermaid grammar is parsed faithfully. Our parsers accept the exact syntax documented at
\texttt{mermaid.js.org}~\cite{mermaiddocs2024}, including:

\begin{itemize}
  \item Indentation-sensitive formats (timeline, mindmap, gantt)
  \item Arrow syntax variants (\texttt{-->}, \texttt{-.->}, \texttt{==>}, etc.\ in flowchart)
  \item Section/participant declarations
  \item Frontmatter (\texttt{---}$\ldots$\texttt{---}) and init directives (\texttt{\%\%\{init\}\%\%})
  \item All node shape syntaxes (round, diamond, hexagon, etc.)
\end{itemize}

\subsubsection{Semantic Mapping}

Mermaid's syntax maps to our Domain IRs via grammar-specific lowering rules:

\begin{description}
  \item[flowchart $\to$ Flow IR] Nodes become Flow nodes; edges become Flow edges;
        subgraphs become lanes/clusters.
  \item[sequence $\to$ Sequence IR] Participants, messages, activations, and fragments
        map directly.
  \item[gantt/timeline $\to$ Timeline IR] Sections become tracks; tasks become activities;
        milestones map directly.
  \item[mindmap $\to$ Tree IR] Indentation hierarchy becomes parent-child tree structure.
\end{description}

\subsubsection{Directive and Frontmatter Hooks}

Mermaid's configuration mechanisms become our extension points:

\begin{lstlisting}[language=yaml,caption={Mermaid frontmatter as theme/extension hook}]
---
title: Architecture Overview
theme: professional-dark
config:
  flowchart:
    nodeSpacing: 60
    rankSpacing: 80
---
flowchart LR
    A[Client] --> B[API Gateway]
    B --> C[Service Mesh]
\end{lstlisting}

The \texttt{theme} field selects our theme system. The \texttt{config} block maps to
grammar-specific theme extensions. Init directives (\texttt{\%\%\{init: \{...\}\}\%\%})
provide inline configuration.

\subsection{Reuse of Existing Kernels}

Most new diagram types reuse existing layout kernels:

\begin{description}
  \item[Node-link kernel] (Sugiyama/layered layout) serves: flowchart, C4 variants,
        architecture, block, requirement, gitGraph, sankey.
  \item[Tree kernel] (Buchheim--J\"unger--Leipert) serves: mindmap, treemap.
  \item[Timeline kernel] (track-based layout) serves: gantt, timeline, journey, kanban.
  \item[Sequence kernel] (order-deterministic columns) serves: sequence diagrams.
  \item[Chart kernel] (NEW grammar-of-graphics layer) serves: pie, xychart, quadrant, radar.
\end{description}

The chart family is the only genuinely new kernel required (Section~\ref{sec:chart-family}).
All other new types are primarily new \emph{parsers} wiring Mermaid syntax to existing
Domain IR shapes and layout engines.
        % Full Mermaid compatibility strategy (22 types)
% =============================================================================
% §16b — Extended Timeline Syntax
% Authored: Leslie (Spec Architect), 2026-06-15
% =============================================================================

\section{Extended Timeline Syntax}
\label{sec:extended-timeline}

Mermaid's \texttt{timeline} keyword is intentionally minimal: a sequence of
\texttt{section} groupings each containing \texttt{period : event} lines.
That simplicity makes it portable but leaves the full power of the Timeline
Domain IR~(Section~\ref{sec:timeline-grammar}) unreachable from the text
surface.

The \emph{Extended Timeline} is a strict superset of that syntax.  It uses
the same \texttt{timeline} header keyword and accepts every legal Mermaid
timeline file unchanged; files that use any extended construct are
superset-only and carry a compiler warning (see
Section~\ref{sec:et-degradation}).  The design principle is \emph{one IR,
many diagrams}: a single extended-timeline source compiles to the same
\texttt{IRDocument} regardless of which of the six layout families the author
selects (Section~\ref{sec:et-layouts}), and it inherits all seven contract
themes (Section~\ref{sec:themes}) without modification.

% ---------------------------------------------------------------------------
\subsection{Two-Tier Portability Model}
\label{sec:et-tiers}
% ---------------------------------------------------------------------------

\begin{description}
  \item[Tier 1 — Mermaid-faithful \texttt{timeline}] (unchanged, portable).
        Uses only \texttt{section} and \texttt{period : event} syntax.
        Any valid Mermaid timeline is valid here and renders identically in
        vanilla Mermaid.  The compiler maps these files to the Timeline IR
        via the existing \texttt{parseTimeline} path and defaults the layout
        to \texttt{timeline-columns} to match Mermaid's visual presentation.

  \item[Tier 2 — Extended Timeline] (superset-only, not renderable in
        vanilla Mermaid).  Uses any construct defined in this section: named
        tracks, duration activities, \texttt{@}-directives, sections with
        explicit time ranges, annotations, legend, or axis breaks.  The
        compiler emits a diagnostic warning and proceeds normally; the same
        \texttt{IRDocument} is the output regardless of which layout is
        chosen.
\end{description}

The two tiers share \emph{exactly} the same IR target.  There is no new IR
introduced by the Extended Timeline --- every construct maps to an existing
\texttt{IRDocument} field defined in \texttt{packages/core/src/types.ts} and
validated by \texttt{packages/core/src/schema.ts}.

% ---------------------------------------------------------------------------
\subsection{Frontmatter Configuration}
\label{sec:et-frontmatter}
% ---------------------------------------------------------------------------

The YAML frontmatter block (\texttt{---}\ldots\texttt{---}) configures the
document at parse time.  All keys are optional; unknown keys are warned and
ignored.  The full config surface (Section~\ref{sec:superset-extensions},
\S17.1) is available; the timeline-relevant subset is:

\begin{lstlisting}[language=yaml,caption={Extended timeline frontmatter}]
---
title: "Platform Roadmap 2026"
subtitle: "Engineering org — Q1..Q4"
theme: executive          # any of 7 contract themes
layout: roadmap           # horizontal | vertical-spine | serpentine |
                          # roadmap | gantt | timeline-columns
density: compact          # compact | normal | comfortable
spineSpacing: even        # time | even  (vertical-spine only)
---
\end{lstlisting}

\begin{table}[h]
\centering\small
\caption{Frontmatter keys and their IR targets}
\label{tab:et-frontmatter}
\begin{tabular}{lll}
\toprule
\textbf{Key} & \textbf{IR field} & \textbf{Notes} \\
\midrule
\texttt{title}       & \texttt{metadata.title}        & Already supported in Tier 1 \\
\texttt{subtitle}    & \texttt{metadata.subtitle}     & Already supported in Tier 1 \\
\texttt{theme}       & \texttt{metadata.theme}        & 7 contract theme names valid \\
\texttt{layout}      & \texttt{metadata.layout}       & 6 values; see Table~\ref{tab:et-layouts} \\
\texttt{density}     & Theme-resolution config        & Not an IR \texttt{Metadata} field; \\
                     &                                & passed to \texttt{resolveContractTheme} \\
\texttt{spineSpacing}& \texttt{RenderOptions.spineSpacing} & \texttt{time} (proportional) or \\
                     &                                & \texttt{even} (infographic); \\
                     &                                & only affects \texttt{vertical-spine} \\
\bottomrule
\end{tabular}
\end{table}

\noindent
\textbf{Note on \texttt{layout} schema gap.}  As of 2026-06-15,
\texttt{metadata.layout} in \texttt{schema.ts} accepts only four values
(\texttt{horizontal}, \texttt{vertical-spine}, \texttt{serpentine},
\texttt{roadmap}); the TypeScript type in \texttt{types.ts} correctly
includes all six.  The Zod schema must be extended to add \texttt{gantt} and
\texttt{timeline-columns} before those values can be round-tripped through
JSON Schema validation.  This is a known schema gap, not a design gap.

% ---------------------------------------------------------------------------
\subsection{Grammar Sketch}
\label{sec:et-grammar}
% ---------------------------------------------------------------------------

The following EBNF-style grammar covers the Extended Timeline body (after
frontmatter extraction).  Whitespace between tokens is insignificant except
that indented lines nest under the nearest enclosing construct.  Two-space
indentation is conventional; the parser accepts any consistent depth $\ge
1$.

\begin{lstlisting}[caption={Extended Timeline EBNF sketch}]
extended-timeline ::=
    'timeline' [ 'extended' ] EOL
    stmt*

stmt ::=
      title-stmt
    | track-stmt
    | section-stmt
    | annotation-stmt
    | legend-stmt
    | break-stmt

(* ---- top-level declarations ---- *)
title-stmt    ::= 'title' text EOL

track-stmt    ::= 'track' track-name EOL
                  INDENT track-item*

track-item    ::= activity-stmt | milestone-stmt

(* ---- activities ---- *)
activity-stmt ::= label ':' date-expr EOL
                  ( INDENT directive )*

date-expr     ::= irdate '..' irdate   (* duration form *)
               |  irdate               (* single-date (span) form *)

(* ---- milestones ---- *)
milestone-stmt ::= '@milestone' label ':' irdate EOL
                   ( INDENT directive )*

(* ---- @-directives (sub-lines under activity or milestone) ---- *)
directive     ::= '@status'      status-value EOL
               |  '@progress'    integer      EOL   (* 0-100 *)
               |  '@shape'       shape-token  EOL
               |  '@icon'        icon-name    EOL
               |  '@color'       css-color    EOL
               |  '@description' text         EOL

(* ---- sections ---- *)
section-stmt  ::= 'section' label
                  ( '[' irdate '..' irdate ']' )? EOL

(* ---- annotations ---- *)
annotation-stmt ::= 'annotation' DQUOTED-TEXT
                    '@' annotation-target EOL

annotation-target ::= irdate '..' irdate    (* period annotation *)
                    | irdate                (* point annotation  *)

(* ---- legend ---- *)
legend-stmt   ::= 'legend' ( 'show' | 'hide' )? EOL

(* ---- axis breaks ---- *)
break-stmt    ::= 'break' irdate '..' irdate EOL

(* ---- primitives ---- *)
irdate   ::= (* IR date string per §21 §4 date model;
               examples: 2026-Q2  2026-06-01  FY26-Q3  now  tbd *)

status-value ::=
    'planned' | 'in-progress' | 'done' |
    'at-risk'  | 'blocked'    | 'cancelled' | 'tentative'

shape-token  ::= (* theme-resolved icon name, e.g. 'diamond' 'star' 'flag' *)

label        ::= (* non-empty text, no leading colon *)
track-name   ::= (* non-empty text *)
\end{lstlisting}

\paragraph{Keyword \texttt{extended}.}  The optional \texttt{extended}
modifier on the opening line is a forward hint.  When present, the compiler
suppresses the ``extended features detected'' warning (the author has
explicitly opted in).  When absent, the warning is emitted on the first use
of any Tier-2 construct.

\paragraph{Indentation and nesting.}
A \texttt{track} header owns all activity and milestone lines at the next
indentation level until a dedent or another \texttt{track}/\texttt{section}/
\texttt{annotation} at the same level appears.  \texttt{@}-directives attach
to the \emph{immediately preceding} activity or milestone and must be
indented one level deeper than that activity.

% ---------------------------------------------------------------------------
\subsection{Construct-to-IR Mapping}
\label{sec:et-ir-mapping}
% ---------------------------------------------------------------------------

Every Extended Timeline construct maps to an existing \texttt{IRDocument}
field.  Table~\ref{tab:et-ir-mapping} gives the authoritative mapping using
the canonical TypeScript field names from \texttt{types.ts} in
\texttt{packages/core/src/}.

\begin{table}[h]
\centering\small
\caption{Extended Timeline syntax $\rightarrow$ IRDocument field mapping}
\label{tab:et-ir-mapping}
\begin{tabular}{p{3.8cm}p{4.2cm}p{5.0cm}}
\toprule
\textbf{DSL construct} & \textbf{IR field} & \textbf{Notes} \\
\midrule
\texttt{track Name}
  & \texttt{Track \{ id, label \}}
  & \texttt{id} = sanitized kebab-case slug of Name \\
\addlinespace
\texttt{Label: start..end}
  & \texttt{Activity \{ track, label,} \texttt{start, end \}}
  & Both \texttt{start} and \texttt{end} are \texttt{IRDate} strings \\
\addlinespace
\texttt{Label: date}
  & \texttt{Activity \{ track, label, span \}}
  & Single-date form; \texttt{Activity.span} is the period identifier \\
\addlinespace
\texttt{@status value}
  & \texttt{Activity.status} or \texttt{Milestone.status}
  & Must be one of the 7 \texttt{Status} enum values \\
\addlinespace
\texttt{@progress N}
  & \texttt{Activity.progress}
  & Author writes integer 0--100; parser stores $N \div 100 \in [0,1]$ \\
\addlinespace
\texttt{@milestone Label: date}
  & \texttt{Milestone \{ id, label, date, track \}}
  & \texttt{track} is the enclosing \texttt{track} block's ID \\
\addlinespace
\texttt{@shape token}
  & \texttt{Milestone.icon}
  & No \texttt{shape} field in IR; \texttt{icon} is the closest proxy.
    Exact marker geometry is theme-resolved.
    \emph{(Gap: see \S\ref{sec:et-gaps})} \\
\addlinespace
\texttt{section Name [s..e]}
  & \texttt{Section \{ id, label, time\_range \}}
  & \texttt{time\_range.start} and \texttt{time\_range.end} from bracket \\
\addlinespace
\texttt{annotation "text" @ date}
  & \texttt{Annotation \{ type:'callout', text, date \}}
  & Point annotation; \texttt{type} = \texttt{'callout'} \\
\addlinespace
\texttt{annotation "text" @ s..e}
  & \texttt{Annotation \{ type:'period', text, start, end \}}
  & Range annotation; \texttt{type} = \texttt{'period'} \\
\addlinespace
\texttt{legend show}
  & \texttt{Legend \{ show: true \}}
  & Renderer auto-builds entries from \texttt{Status}/\texttt{category} data \\
\addlinespace
\texttt{break s..e}
  & \texttt{Metadata.axis\_breaks: [\{ from, to \}]}
  & Collapses dead calendar gaps to a \texttt{//} notch \\
\bottomrule
\end{tabular}
\end{table}

\paragraph{The \texttt{Status} enum.}
The seven legal \texttt{@status} values (from \texttt{types.ts}) are:
\texttt{planned}, \texttt{in-progress}, \texttt{done}, \texttt{at-risk},
\texttt{blocked}, \texttt{cancelled}, \texttt{tentative}.
The theme contract maps each to a role-palette color
(Section~\ref{sec:contract-palette}).

\paragraph{The \texttt{IRDate} format.}
Both \texttt{start..end} operands and standalone \texttt{date} values must
conform to the date model defined in Section~\ref{sec:date-model} and the
regex in \texttt{schema.ts}.  Legal examples: \texttt{2026-06-01},
\texttt{2026-Q2}, \texttt{2026-H1}, \texttt{FY26-Q3}, \texttt{now},
\texttt{tbd}, \texttt{\~{}2027-Q1}.

\paragraph{\texttt{@progress} integer-to-float mapping.}
The IR stores \texttt{Activity.progress} as a decimal in $[0, 1]$ (type
\texttt{number[0..1]}).  The DSL accepts an author-friendly integer percent:
$\texttt{@progress}\ N \mapsto \texttt{progress} = N / 100$.  The parser
clamps the value to $[0, 1]$; values outside $[0, 100]$ emit a warning.

% ---------------------------------------------------------------------------
\subsection{One IR, Many Layouts}
\label{sec:et-layouts}
% ---------------------------------------------------------------------------

The canonical claim of the Extended Timeline is the \emph{reach principle}
(Section~\ref{sec:themes-reach}): one \texttt{IRDocument} can be rendered
under any of the six layout families without changing the source.  The
author sets \texttt{layout:} in the frontmatter; switching layouts requires
editing \emph{only that line}.

\begin{table}[h]
\centering\small
\caption{The six layout families and their extended-timeline affordances}
\label{tab:et-layouts}
\begin{tabular}{lp{4.2cm}p{6.5cm}}
\toprule
\textbf{\texttt{layout} value} & \textbf{Character} & \textbf{Best-fit use case} \\
\midrule
\texttt{horizontal}       & Time axis left$\rightarrow$right, track rows
  & Dense multi-track engineering schedules; high information density \\
\texttt{vertical-spine}   & Central spine, L/R alternating entries
  & Narrative timelines; investor decks; few milestones \\
\texttt{serpentine}       & Boustrophedon winding path
  & Journey maps; annual reviews with a clear chronological story \\
\texttt{roadmap}          & Band-per-track, month/quarter columns
  & Product roadmaps; portfolio views; the \emph{executive} sweet spot \\
\texttt{gantt}            & Mermaid-faithful Gantt: task bars, date axis
  & Project management; dependency-heavy schedules \\
\texttt{timeline-columns} & Section-colored header bands, evenly-spaced columns
  & Mermaid-compatible presentation; events stacked under periods \\
\bottomrule
\end{tabular}
\end{table}

\paragraph{The $7 \times 6 = 42$ look matrix.}
Any of the seven contract themes (\texttt{executive}, \texttt{midnight},
\texttt{blueprint}, \texttt{editorial}, \texttt{terminal}, \texttt{pastel},
\texttt{mono}) combined with any of the six layouts produces a distinct
visual presentation from a single source file.  Themes drive typography,
palette, and connector style (Section~\ref{sec:themes-drives-layout});
layouts drive geometry and axis orientation
(Section~\ref{sec:timeline-rendering}).  This \emph{combination matrix} is
the primary differentiator of the Extended Timeline over Mermaid's fixed
rendering.

\paragraph{The \texttt{poster} bridge.}
A \texttt{poster} (Section~\ref{sec:superset-extensions}) can reference the
same extended-timeline source under multiple layouts in adjacent cells,
producing a side-by-side showcase of the $42$-look matrix in a single
document.

\paragraph{Dimension guard.}
The layout engine enforces a dimension guard (height $\le 5000$\,px; aspect
ratio $\le 4:1$; see \texttt{gallery-dimensions.test.ts}).  Pathological
combinations --- for example, \texttt{vertical-spine} over a multi-decade
span with default \texttt{spineSpacing:time} --- trigger a render-time
warning and a concrete suggestion:

\begin{quote}
\textit{``vertical-spine with a time-proportional spine exceeds safe render
dimensions.  Consider \texttt{spineSpacing: even} or switching to
\texttt{layout: roadmap} / \texttt{timeline-columns} for multi-decade
spans.''}
\end{quote}

This is implemented in \texttt{packages/core/src/grammars/timeline/\allowbreak
vertical-spine.ts} and was validated against the full gallery in the
dimension-guard decision (2026-06-15).

% ---------------------------------------------------------------------------
\subsection{Degradation and Portability Contract}
\label{sec:et-degradation}
% ---------------------------------------------------------------------------

\begin{quote}
\textbf{Extended Timeline Portability Contract (2026-06-15).}

\begin{enumerate}
  \item A file using only Tier-1 constructs (\texttt{section} +
        \texttt{period : event}) is byte-compatible with vanilla Mermaid and
        produces no warnings from this compiler.
  \item A file using any Tier-2 construct is superset-only.  The compiler
        emits a \texttt{WARNING} at parse time:
        \textit{``Extended timeline features detected --- output is not
        renderable by vanilla Mermaid.  Use Tier-1 syntax for portable
        files.''} and continues compilation normally.
  \item The warning is \emph{suppressed} when the opening line reads
        \texttt{timeline extended}, signalling explicit opt-in.
  \item All Tier-2 constructs map to existing \texttt{IRDocument} fields.
        No new IR is introduced.  Switching from a Tier-2 file back to
        Tier-1 syntax produces a valid (reduced) document with no schema
        changes.
  \item Vanilla Mermaid renderers encountering a Tier-2 file will fail to
        parse it predictably (they do not silently corrupt output); the
        failure is a clear signal to switch to this compiler.
\end{enumerate}
\end{quote}

% ---------------------------------------------------------------------------
\subsection{Worked Examples}
\label{sec:et-examples}
% ---------------------------------------------------------------------------

\subsubsection{Full Extended Timeline}

The following example exercises every Tier-2 construct and shows the
frontmatter configuration.

\begin{lstlisting}[caption={Full extended timeline source}]
---
title: "Platform Engineering Roadmap"
subtitle: "2026 — all tracks"
theme: executive
layout: roadmap
density: normal
---
timeline extended

section Infrastructure [2026-Q1 .. 2026-Q2]
section Product [2026-Q2 .. 2026-Q4]

track Platform
  Kubernetes Upgrade: 2026-Q1 .. 2026-Q2
    @status done
    @progress 100
    @description "Migrate all clusters to k8s 1.32"
  Service Mesh Rollout: 2026-Q2 .. 2026-Q3
    @status in-progress
    @progress 40
  @milestone Hard Freeze: 2026-06-30
    @status planned
    @shape diamond

track API
  GraphQL Gateway: 2026-Q1 .. 2026-Q3
    @status in-progress
    @progress 60
  @milestone v3 Launch: 2026-09-01
    @status planned
    @shape star

track Mobile
  iOS Redesign: 2026-Q2
    @status planned
  Android Parity: 2026-Q3 .. 2026-Q4
    @status tentative

annotation "PCI-DSS audit window" @ 2026-06-01 .. 2026-06-30
annotation "Board review" @ 2026-Q3

break 2026-12-20 .. 2027-01-05

legend show
\end{lstlisting}

\subsubsection{Same Data, Multiple Layouts}

Switching \texttt{layout:} in the frontmatter (one line) re-renders the
identical \texttt{IRDocument} under a different presentation family.

\begin{lstlisting}[caption={Same source, four layout variants}]
# Layout 1: Gantt — dependency-aware task bars
layout: gantt

# Layout 2: Roadmap — band-per-track, quarter columns
layout: roadmap

# Layout 3: Serpentine — boustrophedon journey narrative
layout: serpentine

# Layout 4: Timeline-columns — Mermaid-compatible column bands
layout: timeline-columns
\end{lstlisting}

\noindent
Combined with seven contract themes, this yields 42 distinct visual
presentations from one source.  A \texttt{poster} cell grid can render
selected layout-theme combinations side by side:

\begin{lstlisting}[language=yaml,caption={Poster showcasing two layout variants of the same source}]
---
theme: executive
layout: grid 1x2
---
poster "Roadmap two ways"
  cell [0,0]: timeline extended
    # ... (roadmap layout via frontmatter override)
  cell [0,1]: timeline extended
    # ... (gantt layout via frontmatter override)
\end{lstlisting}

% ---------------------------------------------------------------------------
\subsection{Known IR Gaps and Deferred Questions}
\label{sec:et-gaps}
% ---------------------------------------------------------------------------

The following items are gaps between what the Extended Timeline syntax
expresses and what the current IR (2026-06-15) can cleanly carry.  They are
flagged here for future IR revision, not as blockers for the syntax
definition.

\begin{description}
  \item[\texttt{@shape} has no IR field.]
        \texttt{Milestone} in \texttt{types.ts} carries \texttt{icon?:
        string} and \texttt{category?: string} but no \texttt{shape} field.
        The parser maps \texttt{@shape diamond} to \texttt{Milestone.icon =
        'diamond'}, conflating pictogram identity with marker geometry.
        Future IR revision should add \texttt{Milestone.shape?: enum\{diamond,
        circle, square, star, flag\}} and separate it from \texttt{icon}.

  \item[\texttt{gantt} and \texttt{timeline-columns} missing from Zod schema.]
        \texttt{metadata.layout} in \texttt{schema.ts} lists only four values
        (\texttt{horizontal}, \texttt{vertical-spine}, \texttt{serpentine},
        \texttt{roadmap}).  The TypeScript type in \texttt{types.ts} correctly
        includes all six.  The schema must be extended to add \texttt{gantt}
        and \texttt{timeline-columns} for these layouts to survive
        JSON Schema round-trip validation.

  \item[\texttt{density} is not an \texttt{IRDocument} field.]
        Density is resolved at theme time via \texttt{resolveContractTheme}
        and never written into the IR.  A document authored at
        \texttt{density: compact} is indistinguishable from one at
        \texttt{density: normal} after IR round-trip.  Whether density should
        be promoted to \texttt{Metadata} is an open question for the Tier-2
        contract evolution (Section~\ref{sec:themes-contract}).

  \item[\texttt{legend} auto-entry generation is unspecified.]
        The IR supports \texttt{Legend.entries} (explicit \texttt{LegendEntry}
        objects), but the parser-level behaviour of \texttt{legend show}
        without explicit entries (auto-build from status/category data) is a
        rendering convention, not an IR invariant.  The spec defers the
        auto-entry algorithm to the layout engine documentation
        (Section~\ref{sec:timeline-rendering}).

  \item[\texttt{@milestone} without an enclosing \texttt{track}.]
        If \texttt{@milestone} appears outside any \texttt{track} block, the
        parser should assign \texttt{Milestone.track = undefined} (the IR
        permits this).  The layout engine then renders the milestone on the
        global axis rather than a swimlane.  This behaviour should be
        explicitly validated.
\end{description}
    % Extended Timeline Syntax: superset of Mermaid timeline
\section{Superset Extensions Beyond Mermaid}
\label{sec:superset-extensions}

The compiler is not merely Mermaid-compatible---it is a \textbf{superset} that extends
Mermaid's capabilities through well-defined extension mechanisms. Extensions use Mermaid's
own frontmatter/directive hooks plus new top-level keywords, ensuring that the base Mermaid
grammar remains parseable by standard Mermaid renderers (with graceful degradation of
extended features).

\subsection{Extension Mechanisms}

\begin{description}
  \item[Frontmatter extensions] The YAML frontmatter block (\texttt{---}$\ldots$\texttt{---})
        carries our theme selection, animation directives, and composition metadata.
        Standard Mermaid renderers ignore unknown frontmatter fields.

  \item[Init directive extensions] The \texttt{\%\%\{init: \{...\}\}\%\%} block accepts
        grammar-specific configuration that maps to our theme extensions.

  \item[New top-level keywords] For capabilities with no Mermaid equivalent
        (e.g., \texttt{poster}, \texttt{compose}), we introduce new diagram-type keywords.
        These are clearly outside Mermaid's grammar and signal superset-only features.

  \item[Inline annotations] Extended comment syntax (\texttt{\%\% @directive}) for
        per-element hints (animation targets, cross-diagram references).
\end{description}

\subsection{Composition / Posters}

A new \texttt{poster} keyword enables multi-diagram compositions directly in the DSL:

\begin{lstlisting}[language=yaml,caption={Poster composition in superset DSL}]
---
theme: professional-dark
layout: grid 2x2
---
poster "RAG Architecture"
  cell [0,0]: flowchart LR
    A[Query] --> B[Retriever] --> C[Generator]
  cell [0,1]: sequence
    User->>API: request
    API->>DB: lookup
  cell [1,0]: mindmap
    root(Components)
      Retriever
      Generator
      Router
  cell [1,1]: stat
    value: 99.2%
    label: Accuracy
\end{lstlisting}

This maps to the Composition IR (Section~\ref{sec:composition}) and produces grid-based
multi-panel posters.

\subsection{Rich Theming}

Superset theming goes beyond Mermaid's color overrides:

\begin{lstlisting}[language=yaml,caption={Rich theme configuration via frontmatter}]
---
theme: corporate-blue
themeVariables:
  fontFamily: "Inter, system-ui"
  fontSize: 14
  borderRadius: 8
  nodePadding: 16
  lineWidth: 2
darkMode: true
---
\end{lstlisting}

Our theme system (Section~\ref{sec:themes}) provides full typographic, geometric, and
chromatic control---not just color token overrides.

\subsection{Animation}

Declarative animation hints attached to elements:

\begin{lstlisting}[language=yaml,caption={Animation extension}]
---
animate:
  edges: dashflow
  duration: 2s
---
flowchart LR
    A --> B --> C
\end{lstlisting}

Animation compiles to SVG SMIL (Section~\ref{sec:animation}) with raster backends
rendering static resting frames.

\subsection{Cross-Diagram References}

Elements in one diagram can reference elements in another within the same poster:

\begin{lstlisting}[language=yaml,caption={Cross-diagram reference annotation}]
---
refs:
  - from: flow.nodeC
    to: sequence.participant.API
    style: dashed-arrow
---
\end{lstlisting}

References resolve at the composition layer and render as overlay connectors.

\subsection{Icons and Images in Nodes}

Nodes can embed icons from a registry or external images:

\begin{lstlisting}[language=yaml,caption={Icon extension in flowchart}]
flowchart LR
    A["{icon:aws-lambda}" Lambda] --> B["{icon:database}" Store]
\end{lstlisting}

The icon registry is part of the kernel; themes control icon styling (color, size).

\subsection{The Structured IR as Agent API}

The superset's most significant extension is invisible in the DSL: the structured IR path.
Agents bypass the text parser entirely and emit Domain IR directly as JSON/YAML.
This is not an ``export format''---it is a first-class input API
(Section~\ref{sec:frontend}).

\subsection{Graceful Degradation}

When a superset file is rendered by a standard Mermaid renderer:
\begin{itemize}
  \item Frontmatter extensions are silently ignored (Mermaid skips unknown fields).
  \item New keywords (\texttt{poster}, etc.) produce a parse error---clearly signaling
        that the file requires our renderer.
  \item Inline annotations in comments (\texttt{\%\% @...}) are ignored as comments.
\end{itemize}

This ensures that the base diagram content remains portable; only superset features
require our compiler.
   % Superset+ extensions beyond Mermaid

% ============================================================================
% PART III — THE KERNEL
% ============================================================================

\part{The Kernel}
\label{part:kernel}

\section{Scene IR: The Universal Render Contract}
\label{sec:scene-ir}

The Scene IR is the central abstraction of the kernel---the single handoff point between any grammar's layout engine and every rendering backend. It is a backend-agnostic, byte-deterministic description of every drawing primitive, visual treatment, and effect request.

\subsection{Design Rationale}

\paragraph{Why a Scene IR?}
A Scene IR (sometimes called a ``Render IR'' or ``Display List'') decouples the concerns of layout from the concerns of output format. This separation enables:

\begin{itemize}
  \item \textbf{Grammar-agnosticism}: Timeline, flow, and graph grammars all compile to the same Scene IR. The backends need not know which grammar produced the scene.
  
  \item \textbf{Backend-agnosticism}: SVG, PNG, PDF, and future backends consume the same input. Adding a new backend requires no changes to any grammar.
  
  \item \textbf{Determinism verification}: The Scene IR is hashable; golden-testing can verify that the same IR+theme produces byte-identical scenes.
  
  \item \textbf{Effect negotiation}: Effects (glow, drop-shadow, animation) are \emph{requested} in the Scene IR with fallback policies. Backends fulfill or degrade gracefully.
\end{itemize}

\subsection{Scene Model}
\label{sec:scene-model}

The Scene consists of four top-level structures:

\begin{Verbatim}[frame=single,fontsize=\small]
Scene {
  canvas:    { width, height, background }
  elements:  [ DrawingPrimitive ]   -- ordered; painting order preserved
  effects:   { effect_id -> EffectDefinition }
  meta:      { theme_id, fidelity_tier, scene_hash(SHA-256), version }
}
\end{Verbatim}

\subsubsection{Canvas Descriptor}

\begin{Verbatim}[frame=single,fontsize=\small]
Canvas {
  width:      number       -- logical pixels
  height:     number       -- computed from content
  background: Background   -- solid color, gradient, or pattern
}

Background := one of:
  SolidColor     { color: "#RRGGBB" | "#RRGGBBAA" }
  LinearGradient { angle_deg, stops: [(offset, color)...] }
  RadialGradient { cx, cy, radius, stops: [(offset, color)...] }
\end{Verbatim}

\subsubsection{Drawing Primitives}

The Scene IR defines a minimal set of drawing primitives---the ``assembly language'' of visual output:

\begin{Verbatim}[frame=single,fontsize=\small]
DrawingPrimitive := one of:

  Rect {
    id, x, y, width, height,
    fill, stroke, stroke_width, corner_radius,
    opacity, clip_id?, effect_refs[]
  }

  Circle {
    id, cx, cy, radius,
    fill, stroke, stroke_width, opacity, effect_refs[]
  }

  Line {
    id, x1, y1, x2, y2,
    stroke, stroke_width, stroke_style,   -- solid | dashed | dotted
    opacity, effect_refs[], animation_ref?
  }

  Path {
    id, d_commands: [(op, args)...],      -- SVG path mini-language
    fill, stroke, stroke_width, opacity, effect_refs[]
  }

  Text {
    id, x, y, content,
    font_family, font_size, font_weight, color,
    anchor: 'start' | 'middle' | 'end',
    clip_id?, effect_refs[]
  }

  MultiText {
    id, x, y, lines: [TextSpan...],
    line_height, clip_id?
  }

  Image {
    id, x, y, width, height,
    data_uri: string                      -- embedded raster or SVG
  }

  Group {
    id, children: [DrawingPrimitive...],
    clip_rect?, transform?, opacity?
  }
\end{Verbatim}

\paragraph{No semantic primitives.}
The Scene IR contains \emph{no} semantic types (Activity, Milestone, Track, Node, Edge). By the time content reaches the Scene, it has been reduced to pure geometry. This is the key to grammar-agnosticism.

\subsubsection{Effect Definitions}

Effects are visual treatments that may not be uniformly supported across backends. The Scene IR \emph{requests} effects and specifies fallback policies:

\begin{Verbatim}[frame=single,fontsize=\small]
EffectDefinition := one of:

  Glow {
    color, radius_px, intensity,
    fallback_policy: 'approximate' | 'omit' | 'embed-raster' | 'error'
  }

  DropShadow {
    offset_x, offset_y, blur_radius, color, opacity,
    fallback_policy
  }

  GradientFade {
    direction: 'left' | 'right' | 'top' | 'bottom',
    color_stops: [(offset, color)...],
    fade_px,
    fallback_policy
  }

  NoiseTexture {
    seed, scale, turbulence_type, opacity, color,
    fallback_policy
  }

  Bloom {
    radius_px, threshold, intensity,
    fallback_policy
  }
\end{Verbatim}

\paragraph{Fallback policies.}
Each effect carries a \texttt{fallback\_policy} instructing backends what to do when they cannot render the effect natively:

\begin{center}
\small
\begin{tabular}{ll}
\toprule
\textbf{Policy} & \textbf{Meaning} \\
\midrule
\texttt{approximate} & Use the closest available native mechanism \\
\texttt{omit}        & Silently omit the effect; geometry unchanged \\
\texttt{embed-raster}& Rasterize the affected layer; embed as Image \\
\texttt{error}       & Abort with a descriptive diagnostic \\
\bottomrule
\end{tabular}
\end{center}

\subsection{Animation Hints}
\label{sec:scene-animation}

Animation is modeled as \emph{optional, declarative hints} attached to primitives. Backends that support animation (SVG, HTML) honor the hints; backends that produce static output (PNG, PDF) render the resting frame.

\begin{Verbatim}[frame=single,fontsize=\small]
AnimationHint := one of:

  FlowingDashes {
    target_id,                            -- Line or Path id
    dash_length, gap_length,
    speed_px_per_sec,
    direction: 'forward' | 'reverse'
  }

  DrawOn {
    target_id,                            -- Path id
    duration_ms,
    easing: 'linear' | 'ease-in-out' | 'ease-out'
  }

  Pulse {
    target_id,
    property: 'opacity' | 'scale' | 'fill',
    from, to, duration_ms, repeat: number | 'infinite'
  }

  DotAlongPath {
    target_id,                            -- Path id
    dot_radius, dot_fill,
    duration_ms, repeat
  }
\end{Verbatim}

\paragraph{Determinism is preserved.}
An animated SVG is still byte-identical markup---the animation is declarative SMIL or CSS, not frame-by-frame. Golden tests can compare the SVG text. A rasterized frame is a static snapshot of the resting state.

\subsection{Scene Determinism Contract}
\label{sec:scene-determinism}

The Scene is byte-deterministic: two executions of any layout engine on identical IR and theme produce bitwise-identical Scene records.

\paragraph{The \texttt{scene\_hash}.}
The Scene's metadata includes a SHA-256 hash of the serialized element list and effect registry. This hash is reproducible across runs and is used for:

\begin{itemize}
  \item Golden-image testing (compare hashes before comparing pixels)
  \item Cache invalidation (skip re-render if hash matches)
  \item Audit trails (record which scene produced which artifact)
\end{itemize}

\paragraph{Layered determinism.}
Determinism applies at three levels:

\begin{enumerate}
  \item \textbf{Scene geometry---always byte-deterministic.} All coordinates, dimensions, and element ordering are bitwise-reproducible.
  
  \item \textbf{Per-backend output---deterministic given pinned backend version.} The SVG backend with version X produces identical output from identical scenes. Effect seeds for procedural effects (noise, cloud) derive from \texttt{scene\_hash + effect\_id}.
  
  \item \textbf{Cross-backend pixel identity---explicitly not promised.} The SVG and Skia backends produce visually different outputs (Skia renders art effects; SVG approximates or omits). This is correct behavior, not a defect.
\end{enumerate}

\subsection{Scene Serialization}

The Scene IR has a canonical JSON serialization for debugging, testing, and tooling interoperability. The serialization is:

\begin{itemize}
  \item \textbf{Sorted keys} at every level (for diff-friendliness)
  \item \textbf{No whitespace-only differences} (normalized indentation)
  \item \textbf{Floating-point precision} capped at two decimal places
\end{itemize}

This serialization is the basis for the \texttt{scene\_hash} computation.

\subsection{Scene IR as the Integration Point}

The Scene IR is the \emph{only} integration point between the upper layers (grammars, composition) and the lower layers (backends, output formats). This architectural choice has several consequences:

\begin{description}
  \item[Adding a grammar] requires only implementing a layout engine that emits Scene IR. No backend changes are needed.
  
  \item[Adding a backend] requires only consuming Scene IR. No grammar changes are needed.
  
  \item[Testing] can happen at the Scene level: grammar tests verify Scene output; backend tests verify rendering from canonical scenes.
  
  \item[Themes] operate entirely within the Scene---they determine fill colors, stroke widths, and effect parameters, but they do not affect layout (which is already baked into coordinates).
\end{description}

\section{Rendering Backends: SVG as Source of Truth}
\label{sec:backends}

This section defines the rendering backends that consume the Scene IR and produce final output artifacts. The key architectural decision is that \textbf{SVG is the canonical source of truth}---not one backend among equals, but the primary representation from which other formats are derived or specialized.

\subsection{The Backend Architecture}

\begin{figure}[h]
\centering
\begin{Verbatim}[frame=single,fontsize=\small]
                      Scene IR
                         |
         +---------------+---------------+
         |               |               |
         v               v               v
    SVG Backend     Skia Backend    PPTX Backend
         |               |               |
         v               v               v
       SVG file     PNG / PDF       PPTX (native shapes)
         |          (art effects)
         |
    +----+----+
    |         |
    v         v
  resvg      PDF
  (PNG)    (vector)
\end{Verbatim}
\caption{Backend architecture: SVG as the canonical path; Skia and PPTX as specialized backends}
\label{fig:backend-arch}
\end{figure}

\subsection{SVG as Canonical Source of Truth}

\paragraph{Rationale.}
SVG is the canonical representation because:

\begin{enumerate}
  \item \textbf{Text-based and diff-friendly.} SVG is human-readable XML. Golden tests can compare SVG text, not just pixel hashes.
  
  \item \textbf{Resolution-independent.} SVG scales to any resolution without quality loss.
  
  \item \textbf{Deterministic.} The same Scene produces byte-identical SVG. No font rendering variance, no anti-aliasing differences.
  
  \item \textbf{Animation-capable.} SVG supports SMIL animation and CSS animations natively. An animated SVG is still a static text file---the animation is declarative.
  
  \item \textbf{Universal embedding.} SVG embeds in web pages, Markdown, documentation, and most design tools.
  
  \item \textbf{Superset relationship.} A static SVG is an animated SVG with no animation. An animated SVG degrades gracefully to still in print contexts.
\end{enumerate}

\paragraph{PNG and PDF as exports from SVG.}
PNG output is produced by rasterizing SVG through resvg (a Rust-based SVG renderer with deterministic output). PDF output is produced by converting SVG to PDF vectors. Both are \emph{exports} of the SVG truth, not parallel primary targets.

\subsection{The Skia Backend: Art Effects}

For diagrams requiring art effects beyond SVG's native capabilities (high-quality glow, noise textures, bloom), the Skia backend~\cite{skia} renders directly from Scene IR to raster (PNG) or PDF.

\paragraph{When to use Skia.}
\begin{itemize}
  \item Effects with \texttt{fallback\_policy: embed-raster} that SVG cannot approximate
  \item Themes designated as ``Tier 2'' or ``Tier 3'' fidelity (see Section~\ref{sec:fidelity-tiers})
  \item Explicit user request for Skia rendering
\end{itemize}

\paragraph{Skia determinism.}
Skia output is deterministic given:
\begin{itemize}
  \item Pinned Skia version
  \item Fixed effect seeds (derived from \texttt{scene\_hash + effect\_id})
  \item No system font fallback (fonts embedded or specified)
\end{itemize}

\subsection{The PPTX Backend: Native Shapes}

PowerPoint output is a specialized backend that maps Scene IR primitives to PPTX's native DrawingML shapes~\cite{ooxml}. This produces editable slides, not embedded images.

\paragraph{Supported mappings.}
\begin{itemize}
  \item \texttt{Rect} → \texttt{<p:sp>} with \texttt{<a:prstGeom prst="rect">}
  \item \texttt{Circle} → \texttt{<a:prstGeom prst="ellipse">}
  \item \texttt{Line} → \texttt{<p:cxnSp>} connection shape
  \item \texttt{Text} → \texttt{<p:txBody>} with \texttt{<a:p>/<a:r>}
  \item \texttt{Group} → \texttt{<p:grpSp>}
  \item Effects → DrawingML effects (\texttt{<a:effectLst>}, \texttt{<a:glow>}, \texttt{<a:outerShdw>})
\end{itemize}

\paragraph{Limitations.}
PPTX cannot represent arbitrary SVG paths. Complex paths are rasterized and embedded as images. Animation hints are ignored (PPTX animations are a separate, complex system not in scope).

\subsection{Rejecting HTML/CSS-First Architecture}

An alternative architecture would use HTML/CSS as the primary representation, leveraging browser layout (flexbox, grid) and CSS animations. This approach was explicitly considered and rejected:

\begin{description}
  \item[Browser variance breaks determinism.] Different browsers render the same HTML/CSS differently. Font metrics, sub-pixel anti-aliasing, and layout rounding vary. This fundamentally violates the determinism contract.
  
  \item[No single-file portability.] HTML requires external CSS, fonts, and assets. SVG is self-contained.
  
  \item[Headless browser dependency.] Rasterizing HTML requires Puppeteer or Playwright---heavy, slow, and non-deterministic.
  
  \item[Layout leaks into rendering.] CSS layout (flexbox, grid) would compute positions. But determinism requires that layout happen in the Scene IR, not in the renderer. Mixing layout engines creates irreproducibility.
\end{description}

\paragraph{The honest caveat.}
The real robustness risk is the \emph{layout engine}, not the backend. Computing node positions, label placements, and connector routing is hard. CSS auto-layout is tempting precisely because it solves layout. The architectural discipline is: keep layout computed deterministically in the grammar's layout engine; the Scene IR records the result; the backend only draws.

If CSS auto-layout is ever desired (e.g., for a future HTML-interactive backend), it must be quarantined \emph{after} the Scene IR seam, with acceptance that HTML output is not golden-testable in the same way as SVG.

\subsection{Backend Selection Logic}

The compiler selects a backend based on:

\begin{enumerate}
  \item \textbf{Output format requested.} \texttt{.svg} → SVG backend; \texttt{.png} → SVG+resvg (default) or Skia (if effects require); \texttt{.pdf} → SVG+PDF conversion or Skia; \texttt{.pptx} → PPTX backend.
  
  \item \textbf{Theme fidelity tier.} Tier 1 (clean) → SVG; Tier 2/3 (effects) → Skia fallback for PNG/PDF.
  
  \item \textbf{Explicit override.} \texttt{--backend=skia} forces Skia for any format.
\end{enumerate}

\subsection{Fidelity Tiers}
\label{sec:fidelity-tiers}

Themes declare a fidelity tier that guides backend selection:

\begin{center}
\small
\begin{tabular}{llll}
\toprule
\textbf{Tier} & \textbf{Effects} & \textbf{SVG} & \textbf{Raster/PDF} \\
\midrule
Tier 1 & None (clean vectors) & Full fidelity & resvg \\
Tier 2 & Drop shadows, gradients & Approximate & Skia (recommended) \\
Tier 3 & Glow, noise, bloom & Omit/rasterize & Skia (required) \\
\bottomrule
\end{tabular}
\end{center}

The default theme (``consulting'') is Tier 1. Art-effect themes (``dark-executive'', ``neon'') are Tier 2 or 3.

\subsection{Determinism Verification}

Each backend provides determinism verification:

\begin{description}
  \item[SVG] Byte-identical file comparison. Golden SVG files are committed and compared character-by-character.
  
  \item[PNG (resvg)] Pixel-identical comparison. resvg is deterministic given the same SVG input and version.
  
  \item[PNG (Skia)] Pixel-identical comparison with tolerance. Skia's anti-aliasing may vary slightly across platforms; a small per-pixel tolerance (e.g., ±1 in each channel) is permitted.
  
  \item[PDF] Text-layer comparison (ignore binary streams). PDF metadata (timestamps) is stripped before comparison.
  
  \item[PPTX] XML comparison of shape definitions. Binary image blobs compared by hash.
\end{description}

\section{Determinism Contract}
\label{sec:determinism}

Determinism is not a feature of the diagram compiler---it is the \emph{central architectural invariant}. This section consolidates the determinism contract across all layers of the system.

\subsection{The Fundamental Guarantee}

\begin{quote}
\textbf{Determinism Contract:} The same IR document plus the same theme plus the same backend version produces byte-identical output across all runs, platforms, and time.
\end{quote}

This is a \emph{specification-level} requirement, not an implementation detail. If two conforming implementations produce different outputs from identical inputs, one of them has a defect.

\subsection{Why Determinism Matters}

\begin{description}
  \item[Version control.] Changes to the IR produce predictable visual diffs. Teams can review diagram changes in pull requests.
  
  \item[CI/CD integration.] Automated pipelines can render diagrams without human inspection. Build artifacts are reproducible.
  
  \item[Agent reliability.] AI agents can trust that their generated IR produces the expected visual. Round-trip editing (generate → render → inspect → regenerate) works. This property interacts with LLM-side grammar constraints~\cite{willard2023}: the compiler guarantees that a valid IR produces a specific visual, while grammar-constrained decoding guarantees that the LLM emits a valid IR. Together, they close the agent-authoring loop end-to-end.
  
  \item[Golden testing.] Reference images committed to the repository can be compared byte-for-byte against new renders. Regressions are detectable.
  
  \item[Audit trails.] Given an IR and theme, the exact output can be reconstructed at any future date.
\end{description}

\subsection{Sources of Non-Determinism (and How They're Eliminated)}

\subsubsection{System Clock}

\textbf{Problem:} Reading \texttt{Date.now()} or the system clock introduces time-dependent variability.

\textbf{Solution:} The compiler never consults the system clock. The ``today'' marker uses \texttt{metadata.today} from the IR. Timestamps in output metadata are either omitted or taken from IR fields.

\subsubsection{Random Number Generators}

\textbf{Problem:} Random values for layout jitter, effect seeds, or ID generation break determinism.

\textbf{Solution:} No random state is consumed. Effect seeds for procedural effects (noise textures, cloud layers) are derived deterministically from \texttt{scene\_hash + effect\_id}.

For layout engines: force-directed algorithms (Fruchterman-Reingold~\cite{fruchterman1991}, Eades~\cite{eades1984}) require a fixed initial placement. The recommended mitigation is to use stress majorization~\cite{gansnerStressMaj2004} with a deterministic initial layout (nodes sorted by ID on a circle), which makes the SMACOF iteration a pure function of the distance matrix with no random state. Alternatively, Sugiyama-style layered layout (dagre~\cite{dagre}, ELK~\cite{elk}) is deterministic by construction for a given node and edge ordering.

\subsubsection{Floating-Point Rounding}

\textbf{Problem:} Different CPUs may produce different floating-point results for the same operations due to extended precision, FMA instructions, or rounding mode differences.

\textbf{Solution:} All intermediate values are rounded using round-half-up ($\lfloor v + 0.5 \rfloor$) to two decimal places. Scene coordinates are expressed as fixed-precision decimals, not arbitrary floats.

\subsubsection{Collection Ordering}

\textbf{Problem:} Hash maps and sets have undefined iteration order. Processing elements in arbitrary order produces arbitrary output order.

\textbf{Solution:} Every ordered collection is sorted by an explicit, unambiguous key sequence:
\begin{itemize}
  \item Tracks: by \texttt{index} ascending
  \item Activities: by \texttt{(start\_ordinal, id)} ascending
  \item Milestones: by \texttt{(date\_ordinal, id)} ascending
  \item All other entities: by \texttt{id} ascending (alphabetically)
\end{itemize}
ID uniqueness is guaranteed by schema validation; IDs break all ties.

\subsubsection{Font Metrics}

\textbf{Problem:} System fonts vary across platforms. Text measurement with system fonts produces different label widths on macOS vs.\ Linux vs.\ Windows.

\textbf{Solution:} Themes must specify embedded font files (WOFF2 or equivalent) for all layout-critical fonts. The compiler ships font metrics tables at build time. System fonts are never consulted for text measurement.

\subsubsection{Environment Variables and Locale}

\textbf{Problem:} Environment variables (\texttt{TZ}, \texttt{LANG}, \texttt{LC\_*}) affect date formatting and text processing.

\textbf{Solution:} The compiler uses explicit locale from \texttt{metadata.locale}. Environment variables are ignored. Date formatting uses ISO 8601 internally; display formatting is locale-aware but deterministic given the locale.

\subsubsection{File System State}

\textbf{Problem:} Reading external files (images, fonts) that may change between runs.

\textbf{Solution:} External assets are resolved at IR load time and embedded (base64 data URIs) or checksummed. The Scene IR contains no file paths---only embedded data.

\subsection{The Determinism Verification Protocol}

Determinism is verified at multiple levels:

\subsubsection{Scene-Level Verification}

The \texttt{scene\_hash} (SHA-256 of canonical Scene JSON) is computed after layout. Tests verify:

\begin{lstlisting}[language=javascript,caption={Scene determinism test pattern}]
const scene1 = layoutEngine.compile(ir, theme);
const scene2 = layoutEngine.compile(ir, theme);
expect(scene1.meta.scene_hash).toBe(scene2.meta.scene_hash);
\end{lstlisting}

\subsubsection{SVG-Level Verification}

SVG output is compared byte-for-byte:

\begin{lstlisting}[language=javascript,caption={SVG determinism test pattern}]
const svg1 = svgBackend.render(scene);
const svg2 = svgBackend.render(scene);
expect(svg1).toBe(svg2);  // Exact string equality
\end{lstlisting}

\subsubsection{Raster-Level Verification}

PNG output is compared pixel-by-pixel with zero tolerance (for resvg) or minimal tolerance (for Skia):

\begin{lstlisting}[language=javascript,caption={Raster determinism test pattern}]
const png1 = pngBackend.render(scene);
const png2 = pngBackend.render(scene);
expect(pixelDiff(png1, png2)).toBeLessThan(tolerance);
\end{lstlisting}

\subsubsection{Cross-Platform Verification}

CI runs on macOS, Linux, and Windows. Golden files committed on one platform must match on all platforms.

\subsection{Determinism Across Backend Versions}

Determinism is guaranteed \emph{within} a backend version. When backend versions change (e.g., resvg 0.35 → 0.36), output may change. This is expected and managed by:

\begin{enumerate}
  \item Recording backend version in output metadata
  \item Regenerating golden files when backend versions are bumped
  \item Semantic versioning: backend version bumps that change output are minor or major releases
\end{enumerate}

\subsection{The Exception: Cross-Backend Differences}

Cross-backend pixel identity is \emph{explicitly not promised}. The SVG backend and Skia backend, given the same Scene, produce visually different outputs:

\begin{itemize}
  \item SVG may approximate or omit effects that Skia renders faithfully
  \item Anti-aliasing algorithms differ
  \item Font rendering differs (even with the same font metrics)
\end{itemize}

This is correct behavior. The determinism contract applies \emph{within} a backend, not \emph{across} backends. Cross-backend tests use per-backend golden images.

\section{Animation: Additive and Backend-Conditional}
\label{sec:animation}

Animation is modeled as an \emph{additive}, \emph{declarative}, \emph{backend-conditional} layer on top of the static Scene IR. Backends that support animation (SVG, HTML) honor animation hints; backends that produce static output (PNG, PDF, PPTX) render the resting frame.

\paragraph{Implementation Status.}
Fully implemented with optional \texttt{animation} field in Scene IR. SVG backend emits SMIL \texttt{<animate>} and \texttt{<animateMotion>} elements; raster backends (PNG, Skia) render the byte-identical resting frame. Flow grammar edges support flowing dashes (\texttt{stroke-dashoffset} animation). All animation attributes deterministic and included in SVG hash verification.

\subsection{Design Philosophy}

\paragraph{Additive, not fundamental.}
Animation is not required for a diagram to be valid or useful. The vast majority of use cases (print, static slides, documentation) require no animation. Animation is an optional enhancement for web contexts.

\paragraph{Declarative, not frame-based.}
Animation hints describe \emph{what} should animate, not \emph{how} to render each frame. The implementation uses declarative mechanisms (SVG SMIL, CSS animations) that the browser or renderer interprets. There is no frame-by-frame generation, no video encoding, no GIF pipelines.

\paragraph{Backend-conditional.}
Animation hints are carried in the Scene IR. Backends that support animation render them; backends that don't, ignore them. This mirrors the existing pattern for effects (glow is Skia-only; SVG approximates or omits).

\paragraph{Determinism preserved.}
An animated SVG is still byte-identical markup. The animation is expressed as declarative attributes (\texttt{<animate>}, \texttt{<animateMotion>}, CSS \texttt{@keyframes}), not as different frames. Golden tests compare SVG text, including animation elements.

\subsection{Animation Patterns from the Corpus}

Analysis of the inspiration corpus reveals several animation patterns used in technical explainer infographics:

\begin{description}
  \item[Flowing dashes / marching ants] Arrows or connectors with animated \texttt{stroke-dashoffset}, creating a sense of flow or data movement. Common in flow diagrams and architecture visualizations.
  
  \item[Draw-on] Path strokes that animate from zero length to full length, as if being drawn by hand. Used for emphasis when diagrams appear.
  
  \item[Dot along path] A small circle that travels along a connector path, representing data or signal flow.
  
  \item[Pulse / glow cycle] Elements that pulse (opacity or scale) to draw attention. Used sparingly for emphasis.
  
  \item[Sequential reveal] Elements that appear in sequence (fade-in with delay) rather than all at once. Used for step-by-step explanations.
\end{description}

\subsection{Animation Hint Types}

Animation hints are attached to Scene IR primitives via \texttt{animation\_ref} fields:

\subsubsection{FlowingDashes}

Animates \texttt{stroke-dashoffset} to create flowing/marching effect on lines and paths.

\begin{Verbatim}[frame=single,fontsize=\small]
FlowingDashes {
  target_id: string,              // Line or Path id
  dash_length: number,            // px
  gap_length: number,             // px
  speed_px_per_sec: number,       // flow speed
  direction: 'forward' | 'reverse'
}
\end{Verbatim}

\paragraph{SVG implementation.}
\begin{lstlisting}[language=html,caption={SVG flowing dashes animation}]
<line id="flow-1" stroke-dasharray="10 5">
  <animate attributeName="stroke-dashoffset"
           from="0" to="15" dur="1s"
           repeatCount="indefinite"/>
</line>
\end{lstlisting}

\subsubsection{DrawOn}

Animates a path from fully hidden (\texttt{stroke-dashoffset} = path length) to fully visible.

\begin{Verbatim}[frame=single,fontsize=\small]
DrawOn {
  target_id: string,              // Path id
  duration_ms: number,
  easing: 'linear' | 'ease-in-out' | 'ease-out',
  delay_ms?: number
}
\end{Verbatim}

\subsubsection{DotAlongPath}

Animates a dot (circle) traveling along a path using \texttt{<animateMotion>}.

\begin{Verbatim}[frame=single,fontsize=\small]
DotAlongPath {
  target_id: string,              // Path id to follow
  dot_radius: number,
  dot_fill: string,               // color
  duration_ms: number,
  repeat: number | 'infinite'
}
\end{Verbatim}

\subsubsection{Pulse}

Animates opacity, scale, or fill color in a repeating cycle.

\begin{Verbatim}[frame=single,fontsize=\small]
Pulse {
  target_id: string,
  property: 'opacity' | 'scale' | 'fill',
  from: number | string,
  to: number | string,
  duration_ms: number,
  repeat: number | 'infinite'
}
\end{Verbatim}

\subsubsection{FadeIn}

Animates opacity from 0 to 1, optionally with delay for sequential reveal.

\begin{Verbatim}[frame=single,fontsize=\small]
FadeIn {
  target_id: string,
  duration_ms: number,
  delay_ms?: number,
  easing: 'linear' | 'ease-in-out'
}
\end{Verbatim}

\subsection{Theme-Level Animation Defaults}

Themes may specify default animation behavior:

\begin{lstlisting}[language=yaml,caption={Theme animation configuration}]
animation:
  enabled: true                    # master switch
  default_duration_ms: 500
  default_easing: ease-in-out
  connectors:
    style: flowing-dashes          # none | flowing-dashes | draw-on
    speed_px_per_sec: 30
  appear:
    style: fade-in                 # none | fade-in | draw-on
    stagger_ms: 100                # delay between elements
\end{lstlisting}

\subsection{Backend Behavior}

\begin{center}
\small
\begin{tabular}{lll}
\toprule
\textbf{Backend} & \textbf{Animation Support} & \textbf{Behavior} \\
\midrule
SVG & Full & Emits SMIL \texttt{<animate>} / \texttt{<animateMotion>} \\
HTML & Full & Emits CSS \texttt{@keyframes} and \texttt{animation} properties \\
PNG (resvg) & None & Renders resting frame (t=0 or $t=\infty$) \\
PNG (Skia) & None & Renders resting frame \\
PDF & None & Renders resting frame \\
PPTX & Partial & Ignored (PPTX animation is a separate system) \\
\bottomrule
\end{tabular}
\end{center}

\paragraph{Resting frame definition.}
When animation is ignored, the resting frame is:
\begin{itemize}
  \item For \texttt{FlowingDashes}: solid stroke (no dash animation)
  \item For \texttt{DrawOn}: fully visible stroke
  \item For \texttt{DotAlongPath}: dot at start of path
  \item For \texttt{Pulse}: element at \texttt{from} value
  \item For \texttt{FadeIn}: fully opaque ($t=\infty$ state)
\end{itemize}

\subsection{Scope Discipline: What's Out}

The following animation capabilities are explicitly \textbf{out of scope} for v1:

\begin{description}
  \item[Frame-by-frame generation] No rendering of individual animation frames. No GIF/APNG/WebP animated image output.
  
  \item[Video encoding] No MP4/WebM output. Video requires frame rendering, audio timing, and codec complexity.
  
  \item[Lottie/Bodymovin export] Lottie is a rich animation format that would require a separate export path. Deferred to future.
  
  \item[Interactive animation] No animation triggered by user interaction (hover, click). The diagrams are declarative, not interactive.
  
  \item[Complex sequencing] No timeline-based animation sequencing (keyframes at specific times). Only simple, repeating animations.
\end{description}

\paragraph{Rationale.}
These exclusions preserve the elegance of the small-declarative-spec → crisp-vector pipeline. Frame rendering, video encoding, and Lottie export would require heavy dependencies, break the text-based-determinism property, and add significant complexity. They may be valuable future extensions, but they are not v1.


% ============================================================================
% PART IV — THE DIAGRAM FAMILIES
% ============================================================================

\part{The Diagram Families}
\label{part:families}

\section{The Grammar Concept: Two-IR-Layer Model}
\label{sec:grammar-concept}

This section defines the ``grammar'' abstraction that structures the diagram compiler. A grammar is a self-contained visual vocabulary with its own Domain IR and layout engine(s), all compiling down to the shared Scene IR.

\subsection{What Is a Grammar?}

A \textbf{grammar} in this architecture is:

\begin{quote}
A grammar = (Domain IR, Layout Engine(s), Theme Extensions) that compiles a declarative description of \emph{what to communicate} into Scene IR primitives specifying \emph{what to draw}.
\end{quote}

Each grammar is:
\begin{itemize}
  \item \textbf{Self-contained}: It defines its own IR schema, validation rules, and layout logic.
  \item \textbf{Kernel-dependent}: It consumes kernel services (Scene IR, backends, themes, icons) but does not modify them.
  \item \textbf{Peer to other grammars}: Timeline, flow, and graph grammars are siblings, not nested.
\end{itemize}

% -------------------------------------------------------------------------
\subsection{Grounding in the Grammar of Graphics}
\label{sec:grammar-gog}
% -------------------------------------------------------------------------

\textbf{What the literature says.}
This architecture stands on the shoulders of three decades of
visualization-grammar research.

Wilkinson's \emph{Grammar of Graphics}~\cite{grammarofgraphics2005}
established that any statistical graphic can be decomposed into a small
algebra of orthogonal components: \textbf{data}, \textbf{aesthetic
mappings}, \textbf{geometric marks}, \textbf{statistical transforms},
\textbf{scales}, \textbf{coordinates}, and \textbf{facets}.
The key insight is generativity: a finite, composable set of primitives
describes an infinite variety of charts.
This is precisely the principle underlying the diagram compiler's grammar
abstraction --- a finite set of mark types (box, arrow, label, group,
connector) combines through composition to cover all target diagram classes.

Wickham's layered grammar~\cite{layeredgrammar2010} operationalized Wilkinson
into ggplot2, adding two engineering contributions of direct relevance:
(1)~\textbf{default inference} --- the system fills in unspecified dimensions
(axes, colors, font sizes) from data type and context, so the author
specifies only deviations from sensible defaults; and (2)~\textbf{spec/render
separation} --- a ggplot object is a spec; calling \texttt{ggsave()} triggers
the rendering pipeline.
Both principles are adopted verbatim: Domain IR fields have sensible defaults
inferred by the layout engine, and the render pass is strictly separated from
the compile pass.

Satyanarayan et al.'s Vega-Lite~\cite{vegalite2017} pushed the grammar further:
a concise, JSON-serializable spec compiles to Vega's lower-level dataflow IR,
which backends render to SVG or Canvas.
Vega-Lite's composition algebra (\texttt{layer}, \texttt{hconcat},
\texttt{vconcat}, \texttt{facet}) is the direct ancestor of the composition
layer in Section~\ref{sec:composition}.

\medskip\noindent
\textbf{The diagram extension.}
Where the GoG lineage addresses \emph{statistical graphics}, this architecture
addresses \emph{diagram types} --- swimlane timelines, flow pipelines, node-link
graphs, step-card sequences, and comparison matrices.
These diagrams have no \emph{data/aesthetic-mapping} duality; they have a
\emph{semantic entity / visual layout} duality instead.
A Timeline activity has \texttt{start}, \texttt{end}, and \texttt{track};
the compiler maps these to rectangle position, width, and vertical row ---
an analogous but distinct encoding to GoG's \texttt{field $\to$ channel}
mapping.

% -------------------------------------------------------------------------
\subsection{Agent-Authoring Reliability: Constrained DSL}
\label{sec:grammar-llm}
% -------------------------------------------------------------------------

\textbf{What the literature says.}
A growing body of constrained-generation research establishes that
\emph{grammar-constrained decoding is the essential reliability lever for
LLM-DSL pipelines}, and that small, well-constrained grammars lower the
failure surface dramatically.

Willard \& Louf~\cite{willard2023} show that reformulating LLM generation as
transitions in a finite-state machine derived from a formal grammar (or
regex) eliminates an entire class of syntax failures with negligible
per-token overhead.
Wang et al.~\cite{wang2023grammar} demonstrate that a \emph{minimal grammar
fragment} for a specific diagram instance is more reliable than the full
schema: ``grammar constraints derived from domain-specific knowledge enable
LLMs to generalize from few examples on complex DSL generation tasks.''
Dong et al.'s XGrammar~\cite{dong2024xgrammar} achieves up to $100\times$
speedup over naive per-step CFG execution by separating context-independent
from context-dependent tokens at grammar-compilation time, making
schema-constrained diagram generation practical in production LLM
serving stacks.

Tian et al.~\cite{tian2023chartgpt} studied LLM failure modes for chart spec
generation: LLMs produce syntactically valid but semantically wrong specs,
miss required fields, and use ambiguous field names.
Their mitigations map directly to the compiler's architecture: (1)~grammar
constraints eliminate syntactic failures; (2)~the validation layer (§\ref{sec:two-ir-model}) catches semantic failures; (3)~decomposed
generation (entity identification, then IR emission) handles complex diagrams.

Narechania et al.~\cite{narechania2021nl4dv} showed that targeting a
well-specified JSON grammar (Vega-Lite's schema) rather than a general-purpose
API reduces NL-to-visualization complexity substantially.
For small or on-device models, Ray~\cite{ray2026constraint} characterises the
``constraint tax'' --- hard schema-decoding improves syntactic validity but
can lower semantic accuracy for sub-3B models --- and recommends a
``reason free, constrain late'' strategy.

\medskip\noindent
\textbf{What we recommend.}
The Domain IR JSON Schema (one per grammar) is published as a formal
constraint and submitted to OpenAI Structured Outputs~\cite{openai-structured-outputs},
XGrammar~\cite{dong2024xgrammar}, or llama.cpp
GBNF~\cite{llama2024gbnf} for grammar-constrained generation.
Keeping each Domain IR \emph{small} (few required fields, flat nesting,
minimal enum values) is not a simplification concession --- it is a
reliability design decision grounded in the literature above.
The god-IR (Section~\ref{sec:two-ir-model}) is rejected in part because a
large, polymorphic schema is an agent-authoring failure vector.

\subsection{The Two-IR-Layer Model}
\label{sec:two-ir-model}

The architecture uses \textbf{two layers of intermediate representation}:

\ourdiagram[0.95\linewidth]{two-ir-model}{%
  The two-IR-layer model: each grammar's Domain IR feeds its own Layout Engine,
  all of which compile down to the single shared Scene IR,
  consumed by SVG, Skia, and PPTX rendering backends.}

\paragraph{Domain IRs are diverse, small, and grammar-specific.}
The Timeline IR has tracks, activities, milestones. The Flow IR has nodes, edges, lanes. The Graph IR has vertices, edges, clusters. These entities are \emph{semantically meaningful} within their grammar---an Activity knows it has a start date; a Node knows it has incoming edges.

Domain IRs are kept \textbf{small on purpose}:
\begin{itemize}
  \item A small grammar is reliably emittable by an LLM
  \item A small schema fits in a context window
  \item A small grammar has clear boundaries and fewer edge cases
\end{itemize}

\paragraph{Scene IR is the universal assembly language.}
All Domain IRs compile \emph{down} to the single shared Scene IR. The Scene IR knows nothing about timelines, flows, or graphs---it knows only primitives (Rect, Circle, Line, Path, Text, Group) and effects.

The Scene IR is the \textbf{one integration point}:
\begin{itemize}
  \item Backends consume only Scene IR
  \item Themes style Scene IR primitives
  \item Golden tests can verify Scene IR determinism
  \item Animation hints attach to Scene IR primitives
\end{itemize}

\subsection{Rejecting the God-IR Approach}

An alternative architecture would use a single ``god IR'' that tries to express all diagram types in one schema:

\begin{lstlisting}[language=yaml,caption={Hypothetical god-IR (rejected)}]
# DON'T DO THIS
diagram:
  type: timeline | flow | graph | comparison | ...
  entities:
    - kind: activity | node | cell | stat | ...
      properties: { ... massive union of all possible fields ... }
\end{lstlisting}

This approach is explicitly rejected because:

\begin{description}
  \item[Semantic muddiness] A single schema that means different things in different contexts is harder to validate, harder to understand, and harder to document.
  
  \item[Schema bloat] Every grammar's fields appear in the union, even when irrelevant. An activity has \texttt{start}/\texttt{end}; a graph node does not. The god-IR carries dead weight.
  
  \item[LLM hostility] A large, polymorphic schema with conditional validity rules is difficult for agents to emit correctly. Small, single-purpose schemas are easier.
  
  \item[Brittle evolution] Adding a new grammar requires modifying the central schema, risking regressions in all other grammars.
  
  \item[No clear ownership] Who owns the god-IR? Every grammar must negotiate changes. Separate Domain IRs have clear ownership.
\end{description}

The two-IR-layer model inverts the coupling: Domain IRs are independent; the Scene IR is the shared core. Adding a new grammar does not touch existing grammars' IRs.

\subsection{Composition IR: Above Domain IRs}

For multi-panel posters (Section~\ref{sec:composition}), a Composition IR sits \emph{above} the Domain IRs, not merged into them:

\ourdiagram[0.9\linewidth]{composition-ir}{%
  Composition IR sits above Domain IRs: the Composition Layout Engine dispatches
  each panel to its grammar-specific layout engine, then assembles the resulting
  sub-scenes into a single Unified Scene IR.}

The Composition IR \emph{references} panels, each with its own grammar type and Domain IR. It does not flatten them into a single schema.

\subsection{Grammar Interface Contract}

Every grammar must implement:

\begin{lstlisting}[language=javascript,caption={Grammar interface contract}]
interface Grammar<DomainIR> {
  // Schema
  readonly name: string;               // e.g., 'timeline', 'flow'
  readonly schema: JsonSchema;         // JSON Schema for DomainIR
  
  // Validation
  validate(ir: unknown): ValidationResult;
  
  // Layout
  layout(ir: DomainIR, theme: Theme): Scene;
  
  // Optional: grammar-specific theme extensions
  readonly themeSchema?: JsonSchema;
}
\end{lstlisting}

\paragraph{Schema.}
The JSON Schema for the Domain IR, used for validation and agent generation.

\paragraph{Validate.}
Checks the IR against schema and semantic rules. Returns diagnostics.

\paragraph{Layout.}
The core transformation: Domain IR + Theme → Scene IR. This is where all layout computation happens.

\paragraph{Theme extensions.}
Grammars may extend the theme schema with grammar-specific properties (e.g., \texttt{timeline.trackHeight}, \texttt{flow.nodeSpacing}).

\subsection{Current and Planned Grammars}

\begin{center}
\small
\begin{tabular}{llll}
\toprule
\textbf{Grammar} & \textbf{Domain Entities} & \textbf{Status} & \textbf{Family} \\
\midrule
Timeline & tracks, activities, milestones, sections & Implemented & Timeline/project \\
Flow & nodes, edges, lanes, legends & Implemented & Node-link/graph \\
Sequence & participants, messages, activations, fragments & Implemented & UML/software \\
Tree & root, children (recursive), labels & Implemented & Tree/hierarchy \\
\midrule
Class & classes, relationships, packages & Planned (Tier-1) & UML/software \\
State & states, transitions, composites & Planned (Tier-1) & UML/software \\
ER & entities, relationships, cardinalities & Planned (Tier-1) & UML/software \\
Chart & data, encoding, marks, scales & Planned (Tier-2) & Charts \\
\bottomrule
\end{tabular}
\end{center}

The full 22-type coverage matrix and family taxonomy are detailed in
Section~\ref{sec:family-taxonomy}.

\subsection{The Compounding Payoff}

The kernel/grammar separation creates a compounding payoff:

\begin{itemize}
  \item \textbf{Animation} added to the kernel is inherited by all grammars
  \item \textbf{New backends} (e.g., Keynote) serve all grammars immediately
  \item \textbf{Theme improvements} (new fonts, colors, effects) apply everywhere
  \item \textbf{Determinism verification} at the Scene level covers all grammars
  \item \textbf{Icon registry} additions are available to all grammars
\end{itemize}

Each new grammar starts with a full feature set: determinism, theming, SVG/PNG/PDF/PPTX output, animation support, validation, and agent integration---all for free.
       % Two-IR-layer model; grammar interface
\section{The Five Diagram Families}
\label{sec:family-taxonomy}

The 22 Mermaid diagram types (plus our extensions) are organized into five families.
Each family shares a layout kernel, a visual vocabulary, and a set of Domain IR patterns.
This taxonomy drives architecture decisions: new types within a family reuse the existing
kernel and require primarily a new parser and minor IR extensions.

\subsection{Family Overview}

\begin{table}[h]
\centering\small
\caption{Five diagram families with constituent types and kernel}
\begin{tabularx}{\textwidth}{lXll}
\toprule
\textbf{Family} & \textbf{Diagram Types} & \textbf{Layout Kernel} & \textbf{Status} \\
\midrule
Node-link / Graph & flowchart, C4 (4 variants), architecture, block, requirement, gitGraph, sankey &
  Sugiyama (layered) & Flow built; rest planned \\
\addlinespace
UML / Software & sequence, class, state, erDiagram &
  Grammar-specific & Sequence built; rest Tier-1 \\
\addlinespace
Charts (data$\to$geometry) & pie, xychart, quadrant, radar &
  Grammar-of-graphics & New kernel (Tier-2) \\
\addlinespace
Timeline / Project & gantt, timeline, journey, kanban &
  Track-based & Timeline built; rest planned \\
\addlinespace
Tree / Hierarchy & mindmap, treemap &
  Buchheim--J\"unger--Leipert & Tree built; treemap planned \\
\bottomrule
\end{tabularx}
\label{tab:family-overview}
\end{table}

\subsection{Family 1: Node-Link / Graph}
\label{sec:family-nodelink}

The largest family. All members share the fundamental structure: \textbf{nodes} connected by
\textbf{edges}, laid out by a directed-graph algorithm.

\paragraph{Shared kernel.} The Sugiyama four-phase layered layout
(Section~\ref{sec:flow-grammar}): cycle removal, rank assignment, crossing minimization,
coordinate assignment. Variants (LR, TB, RL, BT) are direction parameters.

\paragraph{Type-specific extensions:}
\begin{description}
  \item[flowchart] Core type. Subgraphs, diverse node shapes, edge labels. \textbf{Built.}
  \item[C4 Context/Container/Component/Dynamic] Bounded contexts as styled node clusters;
        stereotyped relationships. Reuses flow kernel with C4-specific node rendering.
  \item[architecture] Icon-heavy nodes in a layered or grouped topology.
  \item[block] Block diagrams with nested containers---layout variation of flowchart.
  \item[requirement] Requirement nodes + satisfy/derive/refine edges.
  \item[gitGraph] Commit DAG with branch/merge topology. Specialized rank = commit order.
  \item[sankey] Weighted directed edges rendered as variable-width flows.
\end{description}

\subsection{Family 2: UML / Software (Tier-1 Priority)}
\label{sec:family-uml}

The differentiating family. These are the diagrams software teams use daily---and where
PlantUML's dated rendering creates the biggest opportunity.

\paragraph{Sequence diagrams.} Order-deterministic column layout. Participants, messages,
activations, fragments (alt/opt/loop/par/critical/break). \textbf{Built}
(Section~\ref{sec:sequence-grammar}).

\paragraph{Class diagrams.} Classes with attributes/methods, relationships
(inheritance, composition, aggregation, association, dependency), packages/namespaces.
Domain IR: classes[], relationships[]. Layout: hierarchical with inheritance edges
directing rank.

\paragraph{State diagrams.} States, transitions, composite states (nested), start/end
pseudo-states, fork/join bars, history nodes. Domain IR: states[], transitions[].
Layout: layered graph with composite-state nesting.

\paragraph{Entity-Relationship diagrams.} Entities with attributes, relationships with
cardinality (1, 0..1, 1..*, 0..*). Domain IR: entities[], relationships[].
Layout: force-directed or layered with relationship-label placement.

\subsection{Family 3: Charts (Data $\to$ Geometry)}
\label{sec:family-charts-overview}

A genuinely new grammar-of-graphics layer. Unlike the other families, charts map
\textbf{quantitative data} to \textbf{geometric marks} through scales and coordinate
systems. This is the Vega-Lite~\cite{vegalite2017} / Grammar of
Graphics~\cite{grammarofgraphics2005} lineage adapted to our architecture.

See Section~\ref{sec:chart-family} for the full specification.

\paragraph{Types:}
\begin{description}
  \item[pie] Proportional sectors from labeled values.
  \item[xychart] Bar and line charts on Cartesian axes.
  \item[quadrant] Four-quadrant scatter with labeled axes and positioned items.
  \item[radar] Multi-axis radial chart (spider/radar plot).
\end{description}

\subsection{Family 4: Timeline / Project}
\label{sec:family-timeline}

Track-based temporal layouts. All members share the concept of a time axis with items
positioned along it.

\paragraph{Shared kernel.} Track-based layout with configurable time axis, row allocation,
and activity/milestone placement (Section~\ref{sec:timeline-grammar}).

\paragraph{Types:}
\begin{description}
  \item[gantt] Tasks with durations, sections, dependencies (visual arrows), today marker.
        \textbf{Built} (maps to Timeline IR).
  \item[timeline] Sectioned event list on a time axis. \textbf{Built} (maps to Timeline IR).
  \item[journey] User-journey map with tasks scored by satisfaction level.
        Extension: satisfaction score $\to$ color/position mapping.
  \item[kanban] Column-based board (backlog/doing/done). Extension: columns as
        vertical tracks, cards as items.
\end{description}

\subsection{Family 5: Tree / Hierarchy}
\label{sec:family-tree}

Recursive parent-child structures laid out with tidy-tree algorithms.

\paragraph{Shared kernel.} Buchheim--J\"unger--Leipert O(n) tidy-tree layout
(Section~\ref{sec:tree-grammar}).

\paragraph{Types:}
\begin{description}
  \item[mindmap] Indentation-defined hierarchy with styled nodes. \textbf{Built.}
  \item[treemap] Area-proportional nested rectangles from weighted hierarchy.
        Extension: squarified treemap algorithm for the layout kernel.
\end{description}

\subsection{Coverage Roadmap Tiers}

\begin{description}
  \item[Tier 0 --- Wire existing grammars to Mermaid syntax.]
        flowchart, sequence, gantt, timeline, mindmap.
        Work: write parsers for Mermaid syntax $\to$ existing Domain IRs.
        \textbf{5 types. Kernel: already built.}

  \item[Tier 1 --- UML/software line.]
        class, state, erDiagram, C4 (4 variants).
        Work: new Domain IRs + layout extensions + Mermaid parsers.
        \textbf{7 types. High priority; key differentiator.}

  \item[Tier 2 --- Charts.]
        pie, xychart, quadrant, radar.
        Work: new chart grammar-of-graphics kernel + parsers.
        \textbf{4 types. New architecture required.}

  \item[Tier 3 --- Remaining types.]
        sankey, requirement, gitGraph, block, architecture, treemap, kanban, journey, packet.
        Work: mostly parser + IR variants on existing kernels.
        \textbf{9 types. Lower priority; incremental.}
\end{description}
       % 5 families, 22-type coverage matrix
\section{Timeline Grammar: Domain IR}
\label{sec:timeline-grammar}

The Timeline Grammar is the first grammar implemented in the diagram compiler. This section specifies its Domain IR---the declarative representation of \emph{what} a timeline communicates.

\textbf{Note:} This section re-homes and re-contextualizes the original Timeline IR specification as ``Grammar \#1'' in the broader diagram compiler architecture. The IR definition is unchanged; only the framing is updated.

\subsection{Design Principles}

The Timeline Domain IR follows the general grammar design principles (Section~\ref{sec:grammar-concept}), with timeline-specific elaboration:

\begin{enumerate}
  \item \textbf{Make illegal states unrepresentable.} The IR's type system should prevent
        malformed timelines at the structural level, not merely validate them after the fact.
  \item \textbf{Rendering-agnostic.} The IR describes meaning, not pixels. Visual
        decisions---colors, fonts, positioning---belong to the theme engine.
  \item \textbf{Git-friendly.} Stable field ordering, line-oriented YAML syntax, deterministic
        serialization, minimal diff churn.
  \item \textbf{Agent-generatable.} Clear required-vs-optional distinctions, explicit defaults,
        forgiving structure that remains validatable.
  \item \textbf{Orthogonal core.} A small set of primitives; convenience features are optional
        layers, not core complexity.
\end{enumerate}

\subsection{Type Vocabulary}
\label{sec:ir-types}

The IR uses the following type vocabulary consistently throughout this specification:

\begin{table}[h]
\centering
\begin{tabular}{lp{10cm}}
\toprule
\textbf{Type} & \textbf{Description} \\
\midrule
\texttt{string} & UTF-8 text, non-empty unless explicitly noted \\
\texttt{string?} & Optional string (may be omitted or null) \\
\texttt{int} & Integer \\
\texttt{number} & Decimal number \\
\texttt{number[0..1]} & Decimal in closed interval $[0, 1]$ \\
\texttt{bool} & Boolean (\texttt{true}/\texttt{false}) \\
\texttt{date} & A temporal point (see \S\ref{sec:date-model}) \\
\texttt{daterange} & A temporal interval with start and optional end \\
\texttt{id} & Document-unique identifier (kebab-case slug, e.g., \texttt{alpha-release}) \\
\texttt{ref<T>} & Reference to an entity of type \texttt{T} by its \texttt{id} \\
\texttt{list<T>} & Ordered list of elements of type \texttt{T} \\
\texttt{optional<T>} & Value of type \texttt{T} that may be omitted \\
\texttt{enum\{...\}} & One of the listed literal values \\
\texttt{map<K,V>} & Key-value mapping \\
\bottomrule
\end{tabular}
\caption{IR type vocabulary}
\label{tab:type-vocab}
\end{table}

\subsection{Document Structure}
\label{sec:ir-structure}

A Timeline IR document is a single YAML file with the following root structure:

\begin{lstlisting}[language=yaml,caption={Root document structure}]
version: "1.0"
metadata: { ... }
tracks: [ ... ]
groups: [ ... ]       # optional
activities: [ ... ]
milestones: [ ... ]   # optional
annotations: [ ... ]  # optional
sections: [ ... ]     # optional
legend: { ... }       # optional
\end{lstlisting}

\subsubsection{Root Fields}

\begin{table}[h]
\centering
\small
\begin{tabular}{llccp{5cm}}
\toprule
\textbf{Field} & \textbf{Type} & \textbf{Req} & \textbf{Default} & \textbf{Description} \\
\midrule
\texttt{version} & \texttt{string} & Yes & --- & Specification version (semver, e.g., \texttt{"1.0"}) \\
\texttt{metadata} & \texttt{Metadata} & Yes & --- & Document-level metadata \\
\texttt{tracks} & \texttt{list<Track>} & Yes & --- & Swimlanes/rows; at least one required \\
\texttt{groups} & \texttt{list<Group>} & No & \texttt{[]} & Logical groupings of activities \\
\texttt{activities} & \texttt{list<Activity>} & Yes & --- & Time-spanning items; may be empty \\
\texttt{milestones} & \texttt{list<Milestone>} & No & \texttt{[]} & Point-in-time markers \\
\texttt{annotations} & \texttt{list<Annotation>} & No & \texttt{[]} & Callouts, notes, highlights \\
\texttt{sections} & \texttt{list<Section>} & No & \texttt{[]} & Visual/logical document sections \\
\texttt{legend} & \texttt{Legend} & No & auto & Legend mapping; auto-generated if omitted \\
\bottomrule
\end{tabular}
\caption{Root document fields}
\label{tab:root-fields}
\end{table}

\textbf{Invariant:} The document must contain at least one track. Activities and milestones
may both be empty, but a timeline with neither is semantically vacuous (valid but degenerate).


%% ============================================================================
%% METADATA
%% ============================================================================
\subsection{Metadata}
\label{sec:ir-metadata}

The \texttt{metadata} object contains document-level information and temporal framing.

\begin{table}[h]
\centering
\small
\begin{tabular}{llccp{4.5cm}}
\toprule
\textbf{Field} & \textbf{Type} & \textbf{Req} & \textbf{Default} & \textbf{Description} \\
\midrule
\texttt{title} & \texttt{string} & Yes & --- & Primary document title \\
\texttt{subtitle} & \texttt{string?} & No & \texttt{null} & Secondary title or tagline \\
\texttt{author} & \texttt{string?} & No & \texttt{null} & Author or team name \\
\texttt{created} & \texttt{date} & No & now & Document creation date \\
\texttt{updated} & \texttt{date} & No & now & Last modification date \\
\texttt{time\_range} & \texttt{TimeRange} & Yes & --- & Temporal bounds of the timeline \\
\texttt{axis\_unit} & \texttt{AxisUnit} & No & inferred & Granularity of the time axis \\
\texttt{theme} & \texttt{string?} & No & \texttt{null} & Theme identifier (resolved by renderer) \\
\texttt{locale} & \texttt{string?} & No & \texttt{"en"} & BCP-47 locale for date formatting \\
\texttt{today} & \texttt{date?} & No & eval time & Reference date for \texttt{now}, relative dates, and today-marker; see \S\ref{sec:date-anchor} \\
\texttt{fiscal\_year\_start} & \texttt{int [1..12]} & No & \texttt{1} & Calendar month on which the fiscal year begins (1~=~January); see \S\ref{sec:date-anchor} \\
\texttt{description} & \texttt{string?} & No & \texttt{null} & Extended description or notes \\
\bottomrule
\end{tabular}
\caption{Metadata fields}
\label{tab:metadata-fields}
\end{table}

\subsubsection{TimeRange}

The \texttt{time\_range} defines the temporal viewport of the timeline---the span that will
be rendered. Activities or milestones outside this range are valid but may be clipped.

\begin{lstlisting}[language=yaml]
time_range:
  start: 2026-Q1
  end: 2026-Q4
\end{lstlisting}

\begin{table}[h]
\centering
\begin{tabular}{llccp{5cm}}
\toprule
\textbf{Field} & \textbf{Type} & \textbf{Req} & \textbf{Default} & \textbf{Description} \\
\midrule
\texttt{start} & \texttt{date} & Yes & --- & Inclusive start of the viewport \\
\texttt{end} & \texttt{date} & No & open & Inclusive end; omit for open-ended \\
\bottomrule
\end{tabular}
\caption{TimeRange fields}
\label{tab:timerange-fields}
\end{table}

\subsubsection{AxisUnit}

\begin{lstlisting}
enum AxisUnit { day, week, month, quarter, half, year }
\end{lstlisting}

If omitted, the renderer infers axis unit from the time range span. The axis unit affects
visual tick marks and labels, not the precision of individual dates.

\subsubsection{Date Anchor and Fiscal Calendar}
\label{sec:date-anchor}

\paragraph{Date anchor (\texttt{today}).}
The symbolic date \texttt{now} and all relative date offsets (\texttt{+3m}, \texttt{-2w},
etc.) resolve against the \emph{date anchor}, evaluated in this order:
\begin{enumerate}
  \item \texttt{metadata.today} — if present, this value is used exclusively.
  \item \texttt{metadata.created} — fallback when \texttt{today} is absent.
  \item \textbf{Hard error} — if neither is present and any \texttt{now} or relative date
        appears in the document, the document is not deterministically evaluable and
        validation must reject it.
\end{enumerate}

Renderers \textbf{must not} consult the system clock to resolve \texttt{now} or relative
dates. The today-marker annotation (\texttt{type: today-marker}) also uses this anchor.

\textbf{Determinism caveat:} When \texttt{today} is omitted and \texttt{created} serves
as the anchor, output is deterministic (the \texttt{created} date is fixed in the
document). However, the document's visual "current position" does not advance with real
time. Authors and agents \textbf{should} set \texttt{today} explicitly whenever \texttt{now}
or relative dates appear, so that the document's meaning is unambiguous and byte-stable
across environments and rendering runs.

\paragraph{Fiscal calendar (\texttt{fiscal\_year\_start}).}
\texttt{fiscal\_year\_start} specifies the calendar month (1--12) on which the fiscal year
begins. The default is~1 (January), which makes fiscal and calendar years identical.

Fiscal quarter dates (\texttt{FYnn-Qk}) map to calendar dates as follows. \texttt{FY}$nn$
denotes the fiscal year beginning on month \texttt{fiscal\_year\_start} of calendar year
$20nn$. Fiscal quarter $k$ ($k \in \{1,2,3,4\}$) spans three consecutive calendar months
starting at calendar month $m_k$, where:
\[
  m_k = ((\texttt{fiscal\_year\_start} - 1) + (k - 1) \times 3) \bmod 12 + 1
\]
and the calendar year is $20nn$ if $m_k \ge \texttt{fiscal\_year\_start}$, otherwise
$20nn + 1$.

\textbf{Examples with \texttt{fiscal\_year\_start: 1} (default):}
FY26-Q1 = Jan--Mar 2026; FY26-Q2 = Apr--Jun 2026 (identical to calendar quarters).

\textbf{Examples with \texttt{fiscal\_year\_start: 4} (April fiscal year):}
FY26-Q1 = Apr--Jun 2026; FY26-Q2 = Jul--Sep 2026; FY26-Q3 = Oct--Dec 2026;
FY26-Q4 = Jan--Mar 2027.

If any \texttt{FY...} date appears in the document and \texttt{fiscal\_year\_start}
$\neq 1$, authors \textbf{should} set \texttt{fiscal\_year\_start} explicitly to ensure
all renderers produce consistent x-coordinates for fiscal dates.


%% ============================================================================
%% DATE MODEL
%% ============================================================================
\subsection{Date Model}
\label{sec:date-model}

Timelines present temporal information at varying granularities. A product roadmap may show
quarters; a conference schedule shows hours. The IR must represent this heterogeneity without
forcing false precision or losing meaningful vagueness.

\subsubsection{Date Representations}

A \texttt{date} in the IR is one of the following forms:

\begin{table}[h]
\centering
\small
\begin{tabular}{llp{6cm}}
\toprule
\textbf{Form} & \textbf{Example} & \textbf{Semantics} \\
\midrule
ISO date & \texttt{2026-06-09} & Specific calendar day \\
ISO datetime & \texttt{2026-06-09T14:00} & Specific moment (minute precision) \\
Year-month & \texttt{2026-06} & Entire month (June 2026) \\
Year only & \texttt{2026} & Entire year \\
Quarter & \texttt{2026-Q2} & Calendar quarter (Apr--Jun) \\
Half & \texttt{2026-H1} & Calendar half (Jan--Jun) \\
Fiscal quarter & \texttt{FY26-Q2} & Fiscal quarter (calendar mapping per \texttt{fiscal\_year\_start}; see \S\ref{sec:date-anchor}) \\
Relative & \texttt{+3m} & 3 months from the date anchor (\texttt{metadata.today} $\to$ \texttt{metadata.created}; see \S\ref{sec:date-anchor}) \\
Symbolic & \texttt{now} & Resolves to the date anchor (\texttt{metadata.today} $\to$ \texttt{metadata.created}; see \S\ref{sec:date-anchor}) \\
\bottomrule
\end{tabular}
\caption{Date representation forms}
\label{tab:date-forms}
\end{table}

\textbf{Design rationale:} We deliberately support coarse granularity (quarters, halves, years)
as first-class concepts rather than forcing conversion to day-level dates. This preserves
authorial intent: ``Q3 2026'' means the entire quarter, not ``2026-07-01''.

\subsubsection{Uncertain and Open Dates}

For dates that are genuinely unknown, tentative, or open-ended, the IR provides explicit
encodings rather than using null-as-magic:

\begin{table}[h]
\centering
\begin{tabular}{lp{8cm}}
\toprule
\textbf{Value} & \textbf{Semantics} \\
\midrule
\texttt{tbd} & Date is not yet determined; renders as ``TBD'' \\
\texttt{ongoing} & No planned end; activity continues indefinitely \\
\texttt{unknown} & Date exists but is not known to the author \\
\texttt{\textasciitilde 2026-Q3} & Approximate/tentative (prefix tilde) \\
\bottomrule
\end{tabular}
\caption{Uncertain and open date markers}
\label{tab:uncertain-dates}
\end{table}

\textbf{Invariant:} \texttt{tbd}, \texttt{ongoing}, and \texttt{unknown} are valid only in
positions where the schema allows uncertainty (e.g., \texttt{end} of a daterange, not
\texttt{start} of the document's \texttt{time\_range}).

\subsubsection{DateRange}

A \texttt{daterange} represents a temporal interval:

\begin{lstlisting}[language=yaml]
# Explicit start and end
start: 2026-04-01
end: 2026-06-30

# Open-ended (no planned end)
start: 2026-04-01
end: ongoing

# Shorthand for single-granularity span
span: 2026-Q2    # equivalent to start: 2026-04-01, end: 2026-06-30
\end{lstlisting}

\begin{table}[h]
\centering
\begin{tabular}{llccp{5cm}}
\toprule
\textbf{Field} & \textbf{Type} & \textbf{Req} & \textbf{Default} & \textbf{Description} \\
\midrule
\texttt{start} & \texttt{date} & Yes* & --- & Interval start (inclusive) \\
\texttt{end} & \texttt{date|ongoing|tbd} & No & \texttt{ongoing} & Interval end (inclusive); omitting is equivalent to \texttt{end: ongoing} \\
\texttt{span} & \texttt{date} & Yes* & --- & Single-unit shorthand; mutually exclusive with \texttt{start}/\texttt{end} \\
\bottomrule
\end{tabular}
\caption{DateRange fields (*exactly one of \texttt{span} or \texttt{start} (+ optional \texttt{end}) is required; co-presence of \texttt{span} with \texttt{start} or \texttt{end} is a well-formedness error)}
\label{tab:daterange-fields}
\end{table}

\textbf{Invariant:} When both \texttt{start} and \texttt{end} are concrete dates,
\texttt{end >= start}. The \texttt{span} shorthand expands to the natural bounds of its
granularity (e.g., \texttt{span: 2026-Q2} expands to Q2's first and last days).

\textbf{Invariant (exclusivity):} \texttt{span} is \emph{mutually exclusive} with
\texttt{start} and \texttt{end}. A daterange must use exactly one of:
\begin{itemize}
  \item \texttt{span} alone (shorthand form), or
  \item \texttt{start} alone or \texttt{start} + \texttt{end} (explicit form).
\end{itemize}
Co-presence of \texttt{span} with \texttt{start} or \texttt{end} on the same entity is
a well-formedness error (\texttt{SPAN\_START\_CONFLICT}).

\textbf{Open/ongoing semantics:} An Activity or daterange with \texttt{start} specified
but \texttt{end} omitted is an \emph{open/ongoing interval}. This is semantically
identical to writing \texttt{end: ongoing} and is an explicit, valid state (not an
authoring error). Renderers must extend the bar to the canvas right edge and apply
an open-end indicator (e.g., right-pointing chevron).


%% ============================================================================
%% TRACKS
%% ============================================================================
\subsection{Tracks (Swimlanes)}
\label{sec:ir-tracks}

Tracks are horizontal lanes that organize activities visually. Every activity belongs to
exactly one track. Tracks appear in document order (first track at top).

\begin{lstlisting}[language=yaml,caption={Track examples}]
tracks:
  - id: platform
    label: Platform Team
    description: Core infrastructure work

  - id: mobile
    label: Mobile Apps
    color: blue         # optional hint
\end{lstlisting}

\begin{table}[h]
\centering
\small
\begin{tabular}{llccp{4.5cm}}
\toprule
\textbf{Field} & \textbf{Type} & \textbf{Req} & \textbf{Default} & \textbf{Description} \\
\midrule
\texttt{id} & \texttt{id} & Yes & --- & Unique track identifier \\
\texttt{label} & \texttt{string} & Yes & --- & Display label \\
\texttt{description} & \texttt{string?} & No & \texttt{null} & Extended description \\
\texttt{color} & \texttt{string?} & No & theme & Color hint (theme may override) \\
\texttt{group} & \texttt{ref<Group>?} & No & \texttt{null} & Parent group (for nested tracks) \\
\texttt{index} & \texttt{int?} & No & position & Explicit sort index; unique non-negative integer; tracks render top-to-bottom in ascending \texttt{index} order \\
\texttt{collapsed} & \texttt{bool} & No & \texttt{false} & Initial collapsed state (if renderer supports) \\
\bottomrule
\end{tabular}
\caption{Track fields}
\label{tab:track-fields}
\end{table}

\textbf{Invariant:} Track \texttt{id} values must be unique within the document.

%% ============================================================================
%% GROUPS
%% ============================================================================
\subsection{Groups}
\label{sec:ir-groups}

Groups provide logical clustering of activities, tracks, or other groups. They enable
hierarchical organization (e.g., ``Engineering'' containing ``Platform'' and ``Mobile'' tracks)
without conflating logical grouping with visual swimlanes.

\begin{lstlisting}[language=yaml,caption={Group example}]
groups:
  - id: engineering
    label: Engineering
    children:
      - platform    # ref to track
      - mobile      # ref to track

  - id: releases
    label: Release Milestones
    members:
      - alpha-release   # ref to milestone
      - beta-release
\end{lstlisting}

\begin{table}[h]
\centering
\small
\begin{tabular}{llccp{4.5cm}}
\toprule
\textbf{Field} & \textbf{Type} & \textbf{Req} & \textbf{Default} & \textbf{Description} \\
\midrule
\texttt{id} & \texttt{id} & Yes & --- & Unique group identifier \\
\texttt{label} & \texttt{string} & Yes & --- & Display label \\
\texttt{description} & \texttt{string?} & No & \texttt{null} & Extended description \\
\texttt{parent} & \texttt{ref<Group>?} & No & \texttt{null} & Parent group (for nesting) \\
\texttt{children} & \texttt{list<ref<Track|Group>>} & No & \texttt{[]} & Ordered child tracks/groups \\
\texttt{members} & \texttt{list<ref<Activity|Milestone>>} & No & \texttt{[]} & Member entities \\
\texttt{color} & \texttt{string?} & No & theme & Color hint \\
\bottomrule
\end{tabular}
\caption{Group fields}
\label{tab:group-fields}
\end{table}

\textbf{Invariant:} Group hierarchies must be acyclic. A group cannot be its own ancestor.


%% ============================================================================
%% ACTIVITIES
%% ============================================================================
\subsection{Activities}
\label{sec:ir-activities}

Activities are time-spanning items---phases, projects, tasks, or efforts. They are the
primary content of most timelines.

\begin{lstlisting}[language=yaml,caption={Activity examples}]
activities:
  - id: api-redesign
    label: API Redesign
    track: platform
    start: 2026-Q1
    end: 2026-Q2
    status: in-progress
    progress: 0.4

  - id: mobile-launch
    label: Mobile App Launch
    track: mobile
    span: 2026-Q3
    status: planned
    category: launch
\end{lstlisting}

\begin{table}[h]
\centering
\small
\begin{tabular}{llccp{4cm}}
\toprule
\textbf{Field} & \textbf{Type} & \textbf{Req} & \textbf{Default} & \textbf{Description} \\
\midrule
\texttt{id} & \texttt{id} & Yes & --- & Unique activity identifier \\
\texttt{label} & \texttt{string} & Yes & --- & Display label (shown on bar) \\
\texttt{track} & \texttt{ref<Track>} & Yes & --- & Containing track \\
\texttt{start} & \texttt{date} & Yes* & --- & Start date \\
\texttt{end} & \texttt{date|ongoing|tbd} & No & \texttt{ongoing} & End date; omitting is equivalent to \texttt{end: ongoing} (open interval) \\
\texttt{span} & \texttt{date} & Yes* & --- & Single-unit shorthand; mutually exclusive with \texttt{start}/\texttt{end} \\
\texttt{status} & \texttt{Status} & No & \texttt{planned} & Visual communication status \\
\texttt{progress} & \texttt{number[0..1]?} & No & \texttt{null} & Completion fraction \\
\texttt{category} & \texttt{string?} & No & \texttt{null} & Semantic category (for styling) \\
\texttt{description} & \texttt{string?} & No & \texttt{null} & Extended description \\
\texttt{group} & \texttt{ref<Group>?} & No & \texttt{null} & Logical group membership \\
\texttt{url} & \texttt{string?} & No & \texttt{null} & External link \\
\texttt{tags} & \texttt{list<string>} & No & \texttt{[]} & Freeform tags \\
\texttt{metadata} & \texttt{map<string,any>} & No & \texttt{\{\}} & Extension metadata \\
\bottomrule
\end{tabular}
\caption{Activity fields (*exactly one of \texttt{span} or \texttt{start} (+ optional \texttt{end}) is required; co-presence of \texttt{span} with \texttt{start} or \texttt{end} is a well-formedness error)}
\label{tab:activity-fields}
\end{table}

\subsubsection{Status Enum}
\label{sec:status-enum}

The \texttt{status} field communicates \emph{visual} state for presentation purposes.
This is explicitly \textbf{not} a project-management workflow state; it describes what
the audience should understand about the item's current condition.

\begin{lstlisting}
enum Status {
  planned,      # Scheduled but not yet started
  in-progress,  # Currently active
  done,         # Completed
  at-risk,      # May miss target (visual warning)
  blocked,      # Cannot proceed (visual alert)
  cancelled,    # No longer planned
  tentative     # Not yet confirmed
}
\end{lstlisting}

\textbf{Design rationale:} We include \texttt{at-risk} and \texttt{blocked} because executive
timelines frequently need to communicate risk visually. We omit workflow states like
``in review'' or ``approved'' which belong to project-management tools, not visual
communication artifacts.

\subsubsection{Progress}

The \texttt{progress} field is a number in $[0, 1]$ representing completion fraction.
It has no inherent visual meaning---the theme decides whether to render it as a
progress bar, percentage label, or ignore it entirely.

\textbf{Invariant:} If \texttt{progress} is provided, it must be in $[0, 1]$.
\texttt{progress: 1.0} does not imply \texttt{status: done}; these are orthogonal.


%% ============================================================================
%% MILESTONES
%% ============================================================================
\subsection{Milestones}
\label{sec:ir-milestones}

Milestones are point-in-time markers---releases, deadlines, decision points, or events.
Unlike activities, they have no duration.

\begin{lstlisting}[language=yaml,caption={Milestone examples}]
milestones:
  - id: alpha-release
    label: Alpha Release
    date: 2026-03-15
    track: platform
    status: done
    icon: flag
    color: cyan

  - id: board-review
    label: Board Review
    date: 2026-Q2        # first day of Q2
    category: governance
\end{lstlisting}

\begin{table}[h]
\centering
\small
\begin{tabular}{llccp{4.5cm}}
\toprule
\textbf{Field} & \textbf{Type} & \textbf{Req} & \textbf{Default} & \textbf{Description} \\
\midrule
\texttt{id} & \texttt{id} & Yes & --- & Unique milestone identifier \\
\texttt{label} & \texttt{string} & Yes & --- & Display label \\
\texttt{date} & \texttt{date} & Yes & --- & Point in time \\
\texttt{track} & \texttt{ref<Track>?} & No & global & Associated track (or global) \\
\texttt{status} & \texttt{Status} & No & \texttt{planned} & Visual status \\
\texttt{category} & \texttt{string?} & No & \texttt{null} & Semantic category \\
\texttt{icon} & \texttt{string?} & No & theme & Icon hint (diamond, flag, star, etc.) \\
\texttt{color} & \texttt{string?} & No & theme & Color hint (theme may override) \\
\texttt{description} & \texttt{string?} & No & \texttt{null} & Extended description \\
\texttt{group} & \texttt{ref<Group>?} & No & \texttt{null} & Logical group membership \\
\texttt{url} & \texttt{string?} & No & \texttt{null} & External link \\
\texttt{tags} & \texttt{list<string>} & No & \texttt{[]} & Freeform tags \\
\texttt{metadata} & \texttt{map<string,any>} & No & \texttt{\{\}} & Application data; renderers ignore unknown keys \\
\bottomrule
\end{tabular}
\caption{Milestone fields}
\label{tab:milestone-fields}
\end{table}

\textbf{Note:} When \texttt{track} is omitted, the milestone is ``global'' and may be
rendered spanning all tracks (e.g., a vertical line with label at top).

%% ============================================================================
%% ANNOTATIONS
%% ============================================================================
\subsection{Annotations}
\label{sec:ir-annotations}

Annotations add contextual information---callouts, notes, highlights, or today-markers---
without being first-class timeline entities.

\begin{lstlisting}[language=yaml,caption={Annotation examples}]
annotations:
  - type: today-marker
    date: now

  - type: callout
    target: api-redesign
    text: "Critical path item"
    position: above

  - type: bracket
    targets: [alpha-release, beta-release]
    label: "Testing Phase"

  - type: period
    start: 2026-06-01
    end: 2026-06-15
    label: "Code Freeze"
    style: highlight
\end{lstlisting}

\begin{table}[h]
\centering
\small
\begin{tabular}{llccp{4cm}}
\toprule
\textbf{Field} & \textbf{Type} & \textbf{Req} & \textbf{Default} & \textbf{Description} \\
\midrule
\texttt{type} & \texttt{AnnotationType} & Yes & --- & Annotation kind \\
\texttt{id} & \texttt{id?} & No & auto & Optional identifier \\
\texttt{target} & \texttt{ref<Activity|Milestone>?} & Cond & --- & Single target entity \\
\texttt{targets} & \texttt{list<ref<...>>?} & Cond & --- & Multiple targets \\
\texttt{date} & \texttt{date?} & Cond & --- & Point annotation \\
\texttt{start}/\texttt{end} & \texttt{date?} & Cond & --- & Range annotation \\
\texttt{text}/\texttt{label} & \texttt{string?} & No & \texttt{null} & Display text \\
\texttt{position} & \texttt{enum\{above,below,left,right\}?} & No & auto & Hint only \\
\texttt{style} & \texttt{string?} & No & theme & Style hint \\
\bottomrule
\end{tabular}
\caption{Annotation fields (conditional fields depend on type)}
\label{tab:annotation-fields}
\end{table}

\subsubsection{Annotation Types}

\begin{lstlisting}
enum AnnotationType {
  today-marker,  # Vertical line at current date
  callout,       # Note attached to an entity
  bracket,       # Visual bracket spanning entities
  period,        # Highlighted time period
  note,          # Freestanding note
  connector      # Line connecting two entities
}
\end{lstlisting}


%% ============================================================================
%% SECTIONS
%% ============================================================================
\subsection{Sections}
\label{sec:ir-sections}

Sections divide the timeline into logical or visual chapters. They enable multi-page
timelines or distinct phases that should be visually separated.

\begin{lstlisting}[language=yaml,caption={Section example}]
sections:
  - id: current-state
    label: "Current State (Q1-Q2)"
    time_range:
      start: 2026-Q1
      end: 2026-Q2

  - id: future-state
    label: "Future Roadmap (Q3-Q4)"
    time_range:
      start: 2026-Q3
      end: 2026-Q4
\end{lstlisting}

\begin{table}[h]
\centering
\small
\begin{tabular}{llccp{5cm}}
\toprule
\textbf{Field} & \textbf{Type} & \textbf{Req} & \textbf{Default} & \textbf{Description} \\
\midrule
\texttt{id} & \texttt{id} & Yes & --- & Unique section identifier \\
\texttt{label} & \texttt{string} & Yes & --- & Section title \\
\texttt{time\_range} & \texttt{TimeRange?} & No & inherit & Override document time range \\
\texttt{tracks} & \texttt{list<ref<Track>>?} & No & all & Subset of tracks to show \\
\texttt{description} & \texttt{string?} & No & \texttt{null} & Section description \\
\texttt{page\_break} & \texttt{bool} & No & \texttt{false} & Hint for paginated output \\
\bottomrule
\end{tabular}
\caption{Section fields}
\label{tab:section-fields}
\end{table}

%% ============================================================================
%% LEGEND
%% ============================================================================
\subsection{Legend}
\label{sec:ir-legend}

The legend documents the meaning of statuses, categories, and visual conventions used
in the timeline. If omitted, the renderer may auto-generate a legend from used values.

\begin{lstlisting}[language=yaml,caption={Legend example}]
legend:
  show: true
  position: bottom-right
  entries:
    - key: in-progress
      label: In Progress
      description: Currently being worked on
    - key: at-risk
      label: At Risk
      description: May miss planned date
    - key: launch
      label: Launch Event
      category: true
\end{lstlisting}

\begin{table}[h]
\centering
\small
\begin{tabular}{llccp{5cm}}
\toprule
\textbf{Field} & \textbf{Type} & \textbf{Req} & \textbf{Default} & \textbf{Description} \\
\midrule
\texttt{show} & \texttt{bool} & No & \texttt{true} & Whether to render legend \\
\texttt{position} & \texttt{string?} & No & theme & Position hint \\
\texttt{title} & \texttt{string?} & No & \texttt{"Legend"} & Legend title \\
\texttt{entries} & \texttt{list<LegendEntry>} & No & auto & Legend entries \\
\bottomrule
\end{tabular}
\caption{Legend fields}
\label{tab:legend-fields}
\end{table}

\begin{table}[h]
\centering
\small
\begin{tabular}{llccp{5cm}}
\toprule
\textbf{Field} & \textbf{Type} & \textbf{Req} & \textbf{Default} & \textbf{Description} \\
\midrule
\texttt{key} & \texttt{string} & Yes & --- & Status or category key \\
\texttt{label} & \texttt{string} & Yes & --- & Human-readable label \\
\texttt{description} & \texttt{string?} & No & \texttt{null} & Extended description \\
\texttt{category} & \texttt{bool} & No & \texttt{false} & True if this is a category (not status) \\
\bottomrule
\end{tabular}
\caption{LegendEntry fields}
\label{tab:legend-entry-fields}
\end{table}


%% ============================================================================
%% ID AND REFERENCE SYSTEM
%% ============================================================================
\subsection{ID and Reference System}
\label{sec:ir-refs}

\subsubsection{Identifier Format}

All \texttt{id} fields must be:
\begin{itemize}
  \item Non-empty strings
  \item Composed of lowercase letters, digits, and hyphens
  \item Starting with a letter
  \item Unique within their entity type AND globally within the document
\end{itemize}

\textbf{Regex:} \verb|^[a-z][a-z0-9-]*$|

\textbf{Rationale for global uniqueness:} References may appear in contexts where the
entity type is ambiguous (e.g., annotation targets). Global uniqueness eliminates
the need for type prefixes in references.

\begin{lstlisting}[language=yaml,caption={ID examples}]
# Valid IDs
id: api-v2
id: q3-release
id: mobile-team-ios

# Invalid IDs
id: API-v2         # uppercase
id: 2026-release   # starts with digit
id: api_v2         # underscore
\end{lstlisting}

\subsubsection{References}

A reference is a string containing the \texttt{id} of another entity. The IR uses
bare strings for references (not structured objects) for YAML terseness:

\begin{lstlisting}[language=yaml]
activities:
  - id: feature-a
    track: platform      # ref<Track>
    group: engineering   # ref<Group>

annotations:
  - type: callout
    target: feature-a    # ref<Activity>
\end{lstlisting}

\textbf{Invariant:} Every reference must resolve to exactly one entity in the document.
A validator must reject documents with dangling references.

\subsubsection{Reference Namespacing}

Because IDs are globally unique, references are unambiguous without type prefixes.
The schema context determines what types are valid:

\begin{itemize}
  \item \texttt{activity.track}: must reference a Track
  \item \texttt{annotation.target}: may reference Activity or Milestone
  \item \texttt{group.children}: may reference Track or Group
\end{itemize}

A validator checks that the referenced entity has the correct type for its context.

%% ============================================================================
%% WELL-FORMEDNESS AND INVARIANTS
%% ============================================================================
\subsection{Well-Formedness Rules}
\label{sec:ir-invariants}

A Timeline IR document is \textbf{well-formed} if and only if all the following
invariants hold:

\subsubsection{Structural Invariants}

\begin{enumerate}
  \item \textbf{Version present:} \texttt{version} field exists and matches a known
        specification version.
  \item \textbf{Required fields:} All fields marked ``Req: Yes'' are present and non-null.
  \item \textbf{At least one track:} The \texttt{tracks} list is non-empty.
  \item \textbf{Unique IDs:} All \texttt{id} values are globally unique within the document.
  \item \textbf{Valid ID format:} All \texttt{id} values match \verb|^[a-z][a-z0-9-]*$|.
\end{enumerate}

\subsubsection{Referential Invariants}

\begin{enumerate}[resume]
  \item \textbf{References resolve:} Every reference points to an existing entity.
  \item \textbf{Reference types match:} References point to entities of the expected type
        (e.g., \texttt{activity.track} references a Track, not a Group).
  \item \textbf{No circular groups:} The group parent-child graph is acyclic.
\end{enumerate}

\subsubsection{Temporal Invariants}

\begin{enumerate}[resume]
  \item \textbf{Valid time range:} If \texttt{metadata.time\_range} has both start and end,
        then \texttt{end >= start}.
  \item \textbf{Activity duration non-negative:} For activities with concrete start and end,
        \texttt{end >= start}.
  \item \textbf{Date format valid:} All date values conform to the date model
        (\S\ref{sec:date-model}).
  \item \textbf{Activity date-source exclusive:} Every activity must satisfy exactly one of:
        (a)~\texttt{span} is present and neither \texttt{start} nor \texttt{end} is present; or
        (b)~\texttt{start} is present and \texttt{span} is absent (\texttt{end} optional).
        Co-presence of \texttt{span} with \texttt{start} or \texttt{end} is a hard
        well-formedness error: \texttt{SPAN\_START\_CONFLICT}.
  \item \textbf{Date anchor present when required:} If any date field in the document
        contains \texttt{now} or a relative offset (e.g., \texttt{+3m}, \texttt{-2w}),
        then at least one of \texttt{metadata.today} or \texttt{metadata.created} must be
        present. Absence of both is a hard error: \texttt{DATE\_ANCHOR\_MISSING}.
  \item \textbf{Track index values unique:} When \texttt{track.index} is present, all
        supplied values must be unique non-negative integers within the document.
\end{enumerate}

\subsubsection{Value Invariants}

\begin{enumerate}[resume]
  \item \textbf{Progress in range:} If \texttt{progress} is present, $0 \le \text{progress} \le 1$.
  \item \textbf{Status valid:} All \texttt{status} values are members of the Status enum.
  \item \textbf{Axis unit valid:} If present, \texttt{axis\_unit} is a member of AxisUnit enum.
\end{enumerate}

\subsubsection{Soft Warnings (Valid but Suspicious)}

The following are not errors but may indicate authorial mistakes:

\begin{itemize}
  \item Activity or milestone outside \texttt{metadata.time\_range} (will be clipped)
  \item \texttt{progress: 1.0} with \texttt{status: in-progress} (likely stale)
  \item Empty \texttt{activities} and \texttt{milestones} lists (vacuous timeline)
  \item Unused tracks (no activities reference them)
\end{itemize}


%% ============================================================================
%% EXAMPLES
%% ============================================================================
\subsection{Complete Examples}
\label{sec:ir-examples}

The following examples demonstrate realistic Timeline IR documents. Each is a complete,
valid document showcasing different timeline use cases.

\subsubsection{Example 1: Product Roadmap}

A quarterly product roadmap with multiple teams, showing how the IR represents a
typical software planning artifact.

\begin{lstlisting}[language=yaml,caption={Product roadmap example},label={lst:roadmap-example}]
version: "1.0"

metadata:
  title: "Product Roadmap 2026"
  subtitle: "Engineering & Design"
  author: "Product Team"
  created: 2026-06-09
  time_range:
    start: 2026-Q1
    end: 2026-Q4
  axis_unit: quarter
  theme: corporate-blue

tracks:
  - id: platform
    label: Platform
    description: Core infrastructure and APIs

  - id: mobile
    label: Mobile Apps
    description: iOS and Android applications

  - id: design
    label: Design System
    description: Component library and guidelines

groups:
  - id: foundation
    label: Foundation Work
    members: [api-v2, component-lib]

activities:
  - id: api-v2
    label: API v2 Migration
    track: platform
    start: 2026-Q1
    end: 2026-Q2
    status: in-progress
    progress: 0.6
    category: infrastructure

  - id: mobile-redesign
    label: Mobile App Redesign
    track: mobile
    start: 2026-Q2
    end: 2026-Q3
    status: planned
    description: Complete visual refresh

  - id: component-lib
    label: Component Library
    track: design
    start: 2026-Q1
    end: 2026-Q4
    status: in-progress
    progress: 0.3

  - id: analytics
    label: Analytics Integration
    track: platform
    span: 2026-Q3
    status: planned

milestones:
  - id: api-ga
    label: API v2 GA
    date: 2026-06-30
    track: platform
    status: planned
    icon: flag

  - id: app-launch
    label: App Store Launch
    date: 2026-Q4
    track: mobile
    status: planned
    category: launch

annotations:
  - type: today-marker
    date: now

legend:
  show: true
  entries:
    - key: infrastructure
      label: Infrastructure
      category: true
    - key: launch
      label: Launch Event
      category: true
\end{lstlisting}

\subsubsection{Example 2: Executive Program Timeline}

A high-level transformation program timeline suitable for board presentations,
demonstrating status communication and risk visibility.

\begin{lstlisting}[language=yaml,caption={Executive program timeline},label={lst:exec-timeline}]
version: "1.0"

metadata:
  title: "Digital Transformation Program"
  subtitle: "FY26 Executive View"
  author: "PMO"
  created: 2026-06-09
  time_range:
    start: 2026-01-01
    end: 2026-12-31
  axis_unit: month
  theme: executive

tracks:
  - id: strategy
    label: Strategy & Governance

  - id: technology
    label: Technology Modernization

  - id: people
    label: People & Change

activities:
  - id: strategy-define
    label: Strategy Definition
    track: strategy
    start: 2026-01-01
    end: 2026-03-31
    status: done
    progress: 1.0

  - id: vendor-selection
    label: Vendor Selection
    track: technology
    start: 2026-02-01
    end: 2026-04-30
    status: done

  - id: platform-build
    label: Platform Build
    track: technology
    start: 2026-04-01
    end: 2026-09-30
    status: in-progress
    progress: 0.35

  - id: data-migration
    label: Data Migration
    track: technology
    start: 2026-07-01
    end: 2026-11-30
    status: at-risk
    description: Dependency on legacy system documentation

  - id: training
    label: Training Program
    track: people
    start: 2026-06-01
    end: 2026-12-31
    status: planned

  - id: change-mgmt
    label: Change Management
    track: people
    start: 2026-03-01
    end: ongoing
    status: in-progress

milestones:
  - id: board-approval
    label: Board Approval
    date: 2026-03-15
    status: done
    icon: star

  - id: go-live
    label: Go Live
    date: 2026-10-01
    track: technology
    status: planned
    icon: flag

  - id: program-close
    label: Program Close
    date: 2026-12-15
    status: planned

annotations:
  - type: today-marker
    date: now

  - type: callout
    target: data-migration
    text: "Risk: Legacy documentation gaps"
    position: above

  - type: bracket
    targets: [platform-build, data-migration]
    label: "Critical Path"

sections:
  - id: h1
    label: "H1 2026: Foundation"
    time_range:
      start: 2026-01-01
      end: 2026-06-30

  - id: h2
    label: "H2 2026: Execution"
    time_range:
      start: 2026-07-01
      end: 2026-12-31
\end{lstlisting}


\subsubsection{Example 3: Architecture Evolution Timeline}

A technical architecture timeline showing system evolution over multiple years,
demonstrating coarse granularity and technical categorization.

\begin{lstlisting}[language=yaml,caption={Architecture evolution timeline},label={lst:arch-timeline}]
version: "1.0"

metadata:
  title: "Platform Architecture Evolution"
  subtitle: "2026-2028 Technical Roadmap"
  author: "Architecture Team"
  created: 2026-06-09
  time_range:
    start: 2026
    end: 2028
  axis_unit: half
  theme: technical

tracks:
  - id: frontend
    label: Frontend

  - id: backend
    label: Backend Services

  - id: data
    label: Data Platform

  - id: infra
    label: Infrastructure

activities:
  - id: react-migration
    label: React 19 Migration
    track: frontend
    start: 2026-H1
    end: 2026-H2
    status: in-progress
    category: modernization

  - id: micro-frontend
    label: Micro-Frontend Architecture
    track: frontend
    start: 2027-H1
    end: 2027-H2
    status: planned
    category: architecture

  - id: api-gateway
    label: API Gateway
    track: backend
    span: 2026-H1
    status: done
    category: infrastructure

  - id: service-mesh
    label: Service Mesh Adoption
    track: backend
    start: 2026-H2
    end: 2027-H1
    status: planned
    category: infrastructure

  - id: event-driven
    label: Event-Driven Architecture
    track: backend
    start: 2027
    end: 2028
    status: tentative
    category: architecture

  - id: lakehouse
    label: Data Lakehouse
    track: data
    start: 2026-H1
    end: 2027-H1
    status: in-progress
    progress: 0.4
    category: data

  - id: ml-platform
    label: ML Platform
    track: data
    start: 2027
    end: ongoing
    status: tentative
    category: ai

  - id: k8s-migration
    label: Kubernetes Migration
    track: infra
    span: 2026
    status: in-progress
    progress: 0.7
    category: infrastructure

  - id: multi-cloud
    label: Multi-Cloud Strategy
    track: infra
    start: 2027
    end: 2028
    status: planned
    category: infrastructure

milestones:
  - id: legacy-sunset
    label: Legacy System Sunset
    date: 2027-H1
    track: backend
    status: planned
    icon: flag

  - id: cloud-native
    label: Cloud Native Milestone
    date: 2026
    track: infra
    status: planned

annotations:
  - type: period
    start: 2026-H2
    end: 2027-H1
    label: "Transition Period"
    style: highlight

legend:
  entries:
    - key: modernization
      label: Modernization
      category: true
    - key: architecture
      label: Architecture Change
      category: true
    - key: infrastructure
      label: Infrastructure
      category: true
    - key: data
      label: Data Platform
      category: true
    - key: ai
      label: AI/ML
      category: true
\end{lstlisting}


%% ============================================================================
%% EXTENSIBILITY
%% ============================================================================
\subsection{Extensibility}
\label{sec:ir-extensibility}

The IR is designed to be extended without breaking existing consumers.

\subsubsection{Metadata Extension}

Every entity type includes a \texttt{metadata} field (type \texttt{map<string,any>})
for application-specific data:

\begin{lstlisting}[language=yaml]
activities:
  - id: feature-x
    label: Feature X
    track: platform
    span: 2026-Q2
    metadata:
      jira_key: PROJ-1234
      owner: alice@example.com
      priority: high
\end{lstlisting}

Renderers should ignore unrecognized metadata keys. Source integrations may use
metadata to preserve round-trip fidelity with external systems.

\subsubsection{Custom Categories}

The \texttt{category} field accepts arbitrary strings. The theme engine maps
categories to visual styles. New categories can be introduced without IR changes:

\begin{lstlisting}[language=yaml]
activities:
  - id: security-audit
    category: security   # custom category
    # ...

legend:
  entries:
    - key: security
      label: Security Work
      category: true
\end{lstlisting}

\subsubsection{Version Compatibility}

The \texttt{version} field enables forward compatibility:

\begin{itemize}
  \item Consumers should accept documents with higher minor versions (1.1, 1.2, etc.)
        and ignore unknown fields.
  \item Major version changes (2.0) may introduce breaking changes; consumers should
        reject documents with unsupported major versions.
\end{itemize}

%% ============================================================================
%% SERIALIZATION
%% ============================================================================
\subsection{Serialization and Syntax}
\label{sec:ir-serialization}

\subsubsection{Canonical Format}

YAML 1.2 is the canonical surface syntax for Timeline IR documents. Design decisions
prioritizing YAML:

\begin{itemize}
  \item \textbf{Minimal punctuation:} No braces for simple mappings, no quotes for
        most strings.
  \item \textbf{Line-oriented:} Each field on its own line for clean diffs.
  \item \textbf{Stable ordering:} Fields should appear in specification order for
        consistency; validators may enforce this.
  \item \textbf{Comments preserved:} YAML allows comments; tools should preserve them.
\end{itemize}

\subsubsection{Alternative Formats}

Tooling may support additional formats with lossless conversion:

\begin{itemize}
  \item \textbf{JSON:} Direct mapping; verbose but universally parseable.
  \item \textbf{TOML:} Alternative for users who prefer it.
\end{itemize}

The IR is defined abstractly; the YAML syntax is one serialization, not the definition.

\subsubsection{Git Friendliness}

For minimal diff churn:

\begin{enumerate}
  \item Lists use block style (one item per line), not flow style.
  \item Fields appear in consistent order.
  \item Auto-generated fields (like \texttt{updated} timestamps) are optional and
        may be omitted in version-controlled documents.
  \item IDs are stable slugs, not generated UUIDs.
\end{enumerate}

\begin{lstlisting}[language=yaml,caption={Git-friendly style}]
# Good: block style, stable order
activities:
  - id: feature-a
    label: Feature A
    track: platform
    span: 2026-Q2

# Avoid: flow style, hard to diff
activities: [{id: feature-a, label: Feature A, track: platform, span: 2026-Q2}]
\end{lstlisting}

%% ============================================================================
%% VALIDATION
%% ============================================================================
\subsection{Validation}
\label{sec:ir-validation}

A conforming validator must check all invariants in \S\ref{sec:ir-invariants}.
Validation is a function:

\[
\text{validate} : \text{Document} \to \text{Valid} \mid \text{Errors}
\]

\subsubsection{Error Reporting}

Validation errors should include:
\begin{itemize}
  \item Path to the offending field (e.g., \texttt{activities[2].track})
  \item The invariant violated
  \item The actual value
  \item Suggested fix (when determinable)
\end{itemize}

\subsubsection{Strict vs. Lenient Mode}

\begin{itemize}
  \item \textbf{Strict mode:} All invariants enforced; required for rendering.
  \item \textbf{Lenient mode:} Warnings instead of errors for soft issues; useful
        for work-in-progress documents.
\end{itemize}

%% ============================================================================
%% RELATIONSHIP TO OTHER COMPONENTS
%% ============================================================================
\subsection{Relationship to Other Components}
\label{sec:ir-relationships}

The IR is the stable contract between pipeline stages:

\begin{description}
  \item[Ingestion (upstream)] Source adapters (ADO, GitHub, Markdown, etc.) produce
        IR documents. The IR defines the target shape; adapters handle mapping.

  \item[Theme Engine (downstream)] Themes consume IR and produce styled visual
        specifications. Themes interpret \texttt{status}, \texttt{category}, and
        hint fields (\texttt{color}, \texttt{icon}) but cannot require additional
        IR fields.

  \item[Renderer (downstream)] The renderer consumes themed IR (or IR + theme) and
        produces output (SVG, PNG, PDF, PPTX). The IR remains rendering-agnostic;
        pixel-level decisions are not IR concerns.

  \item[Agent Integration] AI agents generate IR directly. The IR's explicit typing,
        clear required/optional fields, and forgiving structure enable reliable
        generation. Agents need not understand rendering.
\end{description}

\textbf{Contract stability:} Downstream components may not require IR fields beyond
this specification. Extension fields (\texttt{metadata}) are pass-through.


%% ============================================================================
%% BUILD VS ADOPT: SURVEY OF EXISTING REPRESENTATIONS
%% (Research: David, 2026-06-10)
%% ============================================================================
\subsection{Build vs.\ Adopt: A Survey of Existing Representations}
\label{sec:ir-build-vs-adopt}

Before committing to a fully greenfield IR, this section surveys existing formats,
grammars, and standards as potential adoption or extension candidates. Every
candidate is assessed against the five criteria derived from the design principles
(\S\ref{sec:ir-structure}):
\begin{enumerate}
  \item \textbf{Visual-communication focus:} roadmap/milestone/swimlane semantics,
        \emph{not} task-scheduling (no dependencies, resource levelling, or
        critical path);
  \item \textbf{Git-friendly:} line-oriented text, stable field order, minimal diff
        churn;
  \item \textbf{LLM-generatable:} clear required/optional schema, simple references,
        forgiving structure;
  \item \textbf{Render/theme separation:} clean boundary between \emph{what} is
        communicated and \emph{how} it looks; and
  \item \textbf{Open licence with community or specification body.}
\end{enumerate}

%% ------------------------------------------------------------------
\subsubsection{Timeline Text Languages}
%% ------------------------------------------------------------------

\paragraph{Markwhen.}
Markwhen~\cite{markwhen} (MIT licence) is the closest existing tool to our
surface-syntax goals. Events are written as \texttt{date: label} or
\texttt{start/end: label} lines; \texttt{group:} blocks provide swimlane-like
structure; \texttt{\#tag} tokens allow lightweight categorisation. The format
is text-native, line-oriented, and git-diffable.

Critical gaps exist, however. There is no first-class status concept
(no equivalent of \texttt{status: at-risk}); granularity is limited to ISO
dates and informal text strings (no \texttt{2026-Q3} or fiscal-quarter
shorthand); render/theme separation is absent---the view container interprets
colour from tags, not a separate theme layer; static SVG, PDF, and PPTX output
are unsupported; and Markwhen does not publish a versioned, machine-verifiable
schema for its internal JSON parse tree, which makes reliable LLM generation
fragile. The editor (markwhen.com) is proprietary; only the parser and view
container are open source.

\paragraph{Mermaid timeline and Gantt.}
The Mermaid \texttt{timeline} type~\cite{mermaiddocs2024} groups events into
sections on a free-text time axis. It has no swimlane concept, no status enum,
no PPTX output, and no render/theme separation (styles are hard-coded in the
renderer). The Mermaid \texttt{gantt} type~\cite{mermaidgantt2024} adds
dependency keywords (\texttt{crit}, \texttt{done}, \texttt{active})---precisely
the scheduling semantics this specification rejects. Neither type exposes a
stable IR that downstream tooling can consume.

\paragraph{PlantUML Gantt.}
PlantUML's \texttt{@startgantt}~\cite{plantumlgantt} is experimental and
scheduling-oriented (dependency declarations, work-remaining percentages). Its
verbose, UML-rooted syntax is not LLM-friendly for presentation timelines.

\paragraph{D2 and Structurizr.}
D2~\cite{d2lang} is a general-purpose diagram DSL with no native timeline type.
Structurizr~\cite{structurizr} is domain-specific to software architecture (C4
model). Neither is a relevant adoption candidate.

%% ------------------------------------------------------------------
\subsubsection{Visualization Grammars}
%% ------------------------------------------------------------------

\paragraph{Vega-Lite.}
Vega-Lite~\cite{vegalite2017} is the most significant analogy in this
category. It demonstrates that a high-level declarative JSON grammar can
cleanly separate \emph{what to visualize} (data, encodings, marks) from
\emph{how to render it} (scales, axes, guides, legends), with actual rendering
delegated to the Vega runtime. This is precisely the WHAT/HOW philosophy that
this IR embodies (see \S\ref{sec:ir-relationships}).

Vega-Lite is, however, a statistical-graphics grammar, not a roadmap grammar.
Encoding a swimlane timeline in Vega-Lite requires manual mark composition:
a \texttt{rect} mark with \texttt{y: track}, \texttt{x: start},
\texttt{x2: end}, plus separate layers for milestones and annotations. There is
no first-class concept of \texttt{tracks}, \texttt{milestones}, \texttt{status},
\texttt{span}, or \texttt{annotations}. The resulting spec is verbose
(100{+} JSON lines for a five-activity roadmap), fragile for LLM generation,
and not human-editable in the YAML sense intended here. Vega-Lite is the
strongest \emph{architectural analogy}---its declarative WHAT/HOW separation is
the direct design precedent for this IR---but cannot be adopted wholesale.

\paragraph{The Grammar of Graphics and ggplot2.}
Wilkinson's grammar~\cite{grammarofgraphics2005} and Wickham's layered
refinement~\cite{layeredgrammar2010} establish the theoretical foundation for
declarative visual grammars: decompose a graphic into data, aesthetics,
geometry, statistics, and scales. The key insight---separating semantic content
from visual treatment---directly motivates this IR's separation of
tracks/activities/status from theme/colours/layout. These are
\emph{conceptual donors}, not adoptable formats.

\paragraph{Observable Plot and Apache ECharts.}
Observable Plot (ISC licence~\cite{observable-plot}) is a concise JavaScript
charting API. It has no native swimlane roadmap type and requires imperative
programming, not declarative text authoring. Apache ECharts has a
\texttt{timeline} component, but it is an animated-playback control that cycles
through multiple chart states over time---not a swimlane roadmap grammar.
Neither is suitable for adoption.

%% ------------------------------------------------------------------
\subsubsection{Timeline Data Models}
%% ------------------------------------------------------------------

\paragraph{Knight Lab TimelineJS.}
TimelineJS~\cite{timelinejs} accepts a flat JSON array of events with
\texttt{start\_date}, \texttt{end\_date}, \texttt{text.headline}, and
\texttt{text.text}. It is storytelling-oriented (narrative slides, embedded
media) with no swimlanes, no status enum, no PPTX output, and no render/theme
separation.

\paragraph{vis-timeline.}
vis-timeline~\cite{vistimeline} is a browser-only interactive widget requiring
JavaScript item objects (\texttt{\{id, content, start, end, group\}}). There is
no declarative text format, no static SVG/PDF/PPTX export, and no theme-layer
concept. Not adoptable.

\paragraph{react-chrono and Timeline Storyteller.}
react-chrono~\cite{react-chrono} (MIT) accepts a JavaScript items array with
\texttt{title}, \texttt{cardTitle}, \texttt{cardDetailedText}; date fields are
plain strings with no semantic granularity and no swimlanes. Microsoft
Research's Timeline Storyteller~\cite{timeline-storyteller} (open-sourced 2019)
accepts a simple \texttt{[\{start, end, label\}]} JSON array---no swimlanes,
no status, no maintained community. Neither is adoptable.

%% ------------------------------------------------------------------
\subsubsection{Temporal and Calendar Standards}
%% ------------------------------------------------------------------

\paragraph{iCalendar (RFC 5545).}
RFC~5545~\cite{ical-rfc5545} defines \textsc{vevent} with properties
\texttt{DTSTART}, \texttt{DTEND}, \texttt{SUMMARY}, \texttt{DESCRIPTION},
\texttt{CATEGORIES}, \texttt{STATUS} (\texttt{TENTATIVE}, \texttt{CONFIRMED},
\texttt{CANCELLED}), and \texttt{URL}. This vocabulary is the outstanding
\emph{donor}: Mark's \texttt{start}/\texttt{end}/\texttt{label}/\texttt{description}/
\texttt{category}/\texttt{status}/\texttt{url} all map directly to iCalendar
terms, establishing prior-art legitimacy for these names
(see Table~\ref{tab:alignment}).

The \textsc{ics} wire format is wholly unsuitable: property-folded lines
(\texttt{DTSTART:20260609T140000Z}), no YAML/JSON structure, and no swimlane,
visual-status, or render-pipeline concept. The \texttt{STATUS} enum covers
scheduling state (\texttt{CONFIRMED}) but not visual-communication urgency
(\texttt{at-risk}, \texttt{blocked}). \textbf{Not adoptable; strongest single
vocabulary donor.}

\paragraph{Schema.org/Event.}
The schema.org \texttt{Event} type~\cite{schemaorg-event} defines
\texttt{name}, \texttt{startDate}, \texttt{endDate}, \texttt{description},
\texttt{url}, and \texttt{eventStatus}. These are the canonical web vocabulary
for events and corroborate Mark's field naming. Schema.org is a
semantic-markup standard (JSON-LD/Microdata), not a visual representation
format. \textbf{Not adoptable; useful vocabulary donor.}

\paragraph{W3C OWL-Time.}
OWL-Time~\cite{owltime} (W3C Recommendation, 2022) formalises temporal
entities as \texttt{Instant} and \texttt{Interval} with object properties
\texttt{hasBeginning}, \texttt{hasEnd}, and \texttt{hasTemporalDuration}.
Crucially, omitting \texttt{hasEnd} explicitly represents an open/ongoing
interval---the precise semantics of the IR's \texttt{end: ongoing}. OWL-Time
implements Allen's thirteen interval relations~\cite{allen1983} and is the
authoritative formal reference for the IR's open-ended activity semantics.

\paragraph{Allen's Interval Algebra.}
Allen's 1983 paper~\cite{allen1983} establishes 13 binary relations between
time intervals: \emph{before}, \emph{meets}, \emph{overlaps}, \emph{starts},
\emph{during}, \emph{finishes}, \emph{equals}, and their converses. The
algebra provides formal vocabulary for the uncertain/open date model: an
activity with \texttt{end: ongoing} participates in the \emph{started-by}
relation with the document's \texttt{time\_range} without requiring a concrete
end date. Not a format; the formal grounding for temporal semantics.

\paragraph{ISO 8601.}
ISO~8601~\cite{iso8601} is already the foundation of the date model
(\S\ref{sec:date-model}). The IR's \texttt{2026-06-09}, \texttt{2026-06}, and
\texttt{2026} forms are ISO~8601; the \texttt{2026-Q2} and \texttt{2026-H1}
extensions follow the same \texttt{YYYY-unit} pattern as a natural extension.
No further alignment is needed.

%% ------------------------------------------------------------------
\subsubsection{Scheduling and PM Formats}
%% ------------------------------------------------------------------

TaskJuggler~\cite{taskjuggler} (GPL~v2), MS Project XML~\cite{msprojectxml},
and Primavera P6 are scheduling grammars built around \emph{computed} dates:
dependency resolution, resource levelling, and critical-path analysis. All are
explicitly out of scope per the Timeline Grammar decision. MS Project XML
additionally fails the git-friendly criterion: the format is non-deterministic
between saves (GUIDs and timestamps change on every write). No vocabulary from
these formats merits borrowing beyond what iCalendar and schema.org already
provide.

%% ------------------------------------------------------------------
\subsubsection{Analogous Document IRs}
%% ------------------------------------------------------------------

The Pandoc AST~\cite{pandoc} is the strongest structural analogy: a typed,
versioned, language-agnostic JSON document IR with a \texttt{pandoc-api-version}
key, a \texttt{meta} block, and a sequence of typed block elements. This
pattern---\texttt{version} + \texttt{metadata} + typed entity lists---maps
directly onto the root structure of \S\ref{sec:ir-structure}. The Dragon
Book~\cite{dragonbook2006} provides theoretical backing for a typed,
pipeline-stable IR as the contract between compilation stages. Both are
\emph{structural analogies}, not adoptable formats.

%% ------------------------------------------------------------------
\subsubsection{Comparison Table}
%% ------------------------------------------------------------------

Table~\ref{tab:build-vs-adopt} scores the seven strongest candidates against
the five criteria. \textbf{Y}~=~good fit; \textbf{P}~=~partial;
\textbf{N}~=~poor fit or not applicable.

\begin{table}[h]
\centering
\small
\begin{tabular}{p{3.2cm}ccccp{1.6cm}c}
\toprule
\textbf{Candidate}
  & \textbf{Vis.}
  & \textbf{Git}
  & \textbf{LLM}
  & \textbf{Sep.}
  & \textbf{Licence}
  & \textbf{Tooling} \\
\midrule
Markwhen~\cite{markwhen}
  & Y & Y & P & N & MIT & P \\
Mermaid timeline~\cite{mermaiddocs2024}
  & P & Y & P & N & MIT & Y \\
iCalendar VEVENT~\cite{ical-rfc5545}
  & N & N & N & N & IETF & Y \\
TimelineJS JSON~\cite{timelinejs}
  & Y & P & Y & N & GPL & P \\
vis-timeline~\cite{vistimeline}
  & Y & N & N & N & MIT & Y \\
Vega-Lite~\cite{vegalite2017}
  & P & P & P & Y & BSD-3 & Y \\
TaskJuggler~\cite{taskjuggler}
  & N & Y & N & N & GPLv2 & P \\
\bottomrule
\end{tabular}
\caption{Build-vs-Adopt scoring. Vis.~=~visual-communication focus (not
scheduling); Git~=~git-friendly text format; LLM~=~LLM-generatable schema;
Sep.~=~render/theme separation; Tooling~=~active community. Y~=~good;
P~=~partial; N~=~poor.}
\label{tab:build-vs-adopt}
\end{table}

No candidate scores \textbf{Y} on all five dimensions. Markwhen---the
closest---fails on render/theme separation, lacks a stable machine-verifiable
schema, and has no PPTX output path.

%% ------------------------------------------------------------------
\subsubsection{Recommendation: \textsc{Build} with Borrowed Vocabulary}
%% ------------------------------------------------------------------

\textbf{Verdict: \textsc{Build} a bespoke IR. No existing format is
adoptable wholesale and no viable extension profile exists.}

Two orthogonal gaps eliminate every candidate:

\begin{enumerate}
  \item \textbf{Semantic gap.} All formats with real communities (Mermaid,
        iCalendar, Vega-Lite) are scheduling-oriented, statistical-graphics-%
        oriented, or calendar-oriented. None is a visual-communication roadmap
        grammar. Defining a ``profile'' of iCalendar or a ``mark recipe'' in
        Vega-Lite would require adding swimlanes, visual status, coarse-grain
        dates, PPTX output, and a theme engine as first-class concepts---at
        which point we have built a new layer regardless.

  \item \textbf{Pipeline gap.} The requirement of a static, multi-backend
        render pipeline (SVG / PDF / PPTX) with a separate theme engine has no
        counterpart in any existing open-source tool. Mermaid, Markwhen, and
        vis-timeline all bake rendering into the format. Vega-Lite delegates to
        a JavaScript runtime, making it unsuitable for PPTX or print-PDF output
        without adaptation that would negate any adoption benefit.
\end{enumerate}

A \textsc{build} decision is therefore justified. It must not be a
wheel-reinvention: the following \textbf{vocabulary and semantic donors} are
explicitly leveraged so that the IR speaks standard language wherever standards
exist.

\begin{description}
  \item[ISO 8601~\cite{iso8601}]
        Date string syntax (\texttt{2026-06-09}, \texttt{2026-Q2}).
        Already used in the date model (\S\ref{sec:date-model});
        no change needed.

  \item[iCalendar RFC~5545~\cite{ical-rfc5545}]
        Field naming: \texttt{DTSTART}~$\to$~\texttt{start},
        \texttt{DTEND}~$\to$~\texttt{end},
        \texttt{SUMMARY}~$\to$~\texttt{label},
        \texttt{DESCRIPTION}~$\to$~\texttt{description},
        \texttt{CATEGORIES}~$\to$~\texttt{tags}/\texttt{category},
        \texttt{URL}~$\to$~\texttt{url}.
        Status vocabulary: \texttt{TENTATIVE} appears verbatim in the Status
        enum; \texttt{CANCELLED} maps to \texttt{cancelled}.

  \item[Schema.org/Event~\cite{schemaorg-event}]
        Confirms \texttt{name}/\texttt{startDate}/\texttt{endDate}/
        \texttt{description}/\texttt{url} as the canonical web vocabulary for
        events. Mark's \texttt{label}/\texttt{start}/\texttt{end} are direct
        equivalents (shorter names preferred for human authoring).

  \item[W3C OWL-Time~\cite{owltime} and Allen's Algebra~\cite{allen1983}]
        Open-interval semantics: omitting \texttt{hasEnd} denotes an ongoing
        interval, validating \texttt{end: ongoing}. The Instant/Interval
        duality formalises the \texttt{Milestone} (instant) vs.\
        \texttt{Activity} (interval) distinction.

  \item[Vega-Lite~\cite{vegalite2017}]
        Architectural precedent: strict WHAT/HOW separation between a
        declarative specification (this IR) and the rendering runtime
        (Theme Engine + Renderer). The pattern of a versioned, typed YAML/JSON
        document as the sole contract between source and output is adopted from
        Vega-Lite's design.

  \item[Markwhen~\cite{markwhen} and Mermaid~\cite{mermaiddocs2024}]
        Surface-syntax precedents: readable event lines, section/group
        structure, tag-based categorisation. These inform future higher-level
        authoring syntaxes that compile down to this IR.

  \item[Pandoc AST~\cite{pandoc}]
        Structural analogy: versioned, typed, language-agnostic document IR with
        \texttt{version} + \texttt{meta} + typed entity lists. Validates the
        root structure (\S\ref{sec:ir-structure}) and the principle of a
        stable, machine-verifiable schema.
\end{description}

%% ------------------------------------------------------------------
\subsubsection{Alignment of Mark's IR with Prior Art}
%% ------------------------------------------------------------------

Table~\ref{tab:alignment} maps Mark's field names to their closest standard
equivalents. Fields marked \emph{Original} have no adequate standard analog
and are justified by the visual-communication use case.

\begin{table}[h]
\centering
\small
\begin{tabular}{p{2.4cm}p{3.0cm}p{2.0cm}p{4.0cm}}
\toprule
\textbf{IR field} & \textbf{iCal / schema.org} & \textbf{OWL-Time} & \textbf{Assessment} \\
\midrule
\texttt{start}, \texttt{end}
  & \texttt{DTSTART/DTEND}; \texttt{startDate/endDate}
  & \texttt{hasBeginning/} \texttt{hasEnd}
  & \textbf{Aligned.} Standard names; shorter preferred for authoring. \\
\texttt{label}
  & \texttt{SUMMARY}; \texttt{name}
  & ---
  & \textbf{Aligned.} Shorter form preferred for YAML. \\
\texttt{description}
  & \texttt{DESCRIPTION}
  & ---
  & \textbf{Aligned.} \\
\texttt{url}
  & \texttt{URL}; \texttt{url}
  & ---
  & \textbf{Aligned.} \\
\texttt{tags}
  & \texttt{CATEGORIES} (multi)
  & ---
  & \textbf{Aligned.} Correct plural form. \\
\texttt{category}
  & \texttt{CATEGORIES} (single)
  & ---
  & \textbf{Aligned.} \\
\texttt{status: tentative}
  & \texttt{STATUS:TENTATIVE}
  & ---
  & \textbf{Aligned.} Verbatim iCalendar value. \\
\texttt{status: cancelled}
  & \texttt{STATUS:CANCELLED}
  & ---
  & \textbf{Aligned.} Lower-case per IR convention. \\
\texttt{status: at-risk}
  & ---
  & ---
  & \textbf{Original.} No standard analog; justified by executive-presentation idiom. \\
\texttt{status: blocked}
  & ---
  & ---
  & \textbf{Original.} No standard analog; justified by visual-communication need. \\
\texttt{end: ongoing}
  & ---
  & omit \texttt{hasEnd}
  & \textbf{Semantically aligned.} Explicit encoding is clearer than omission. \\
\texttt{span}
  & ---
  & ---
  & \textbf{Original.} Convenient shorthand; no standard equivalent. \\
\texttt{track}
  & \texttt{CALENDAR} (weak)
  & ---
  & \textbf{Original.} Markwhen uses \texttt{group}; \texttt{track} is clearer. \\
\texttt{progress}
  & ---
  & ---
  & \textbf{Original.} No standard analog; no change needed. \\
\texttt{Activity} (type)
  & \textsc{vevent} (duration $>$ 0)
  & \texttt{Interval}
  & \textbf{Aligned.} Correct modelling of time-spanning entities. \\
\texttt{Milestone} (type)
  & \textsc{vevent} (zero dur.)
  & \texttt{Instant}
  & \textbf{Aligned.} Correct modelling of point-in-time markers. \\
\bottomrule
\end{tabular}
\caption{Alignment of Mark's IR field names and types with prior-art standards.}
\label{tab:alignment}
\end{table}

\paragraph{Flags for Mark and Leslie (no schema changes required).}
\begin{itemize}
  \item \textbf{No schema changes are recommended.} Mark's field choices are
        well-aligned with iCalendar RFC~5545, schema.org/Event, and OWL-Time.
        The original contributions (\texttt{at-risk}, \texttt{blocked},
        \texttt{span}, \texttt{progress}, \texttt{track}) are justified by the
        visual-communication use case and have no adequate standard equivalent.

  \item \textbf{Document vocabulary donors in the spec preamble.} This section
        serves that purpose. Future contributors should not propose renames (e.g.,
        \texttt{name} for \texttt{label}, \texttt{startDate} for \texttt{start})
        without consulting this survey, since shorter names are deliberate and
        standards-corroborated.

  \item \textbf{Optional surface alias \texttt{type: event}:} schema.org uses
        \texttt{Event} and OWL-Time uses \texttt{Instant} for point-in-time
        items. Exposing \texttt{type: event} as a surface-syntax alias for
        milestones could ease LLM generation without changing the IR schema.
        Flagged for future surface-language design.
\end{itemize}

      % Timeline/project family (built)
\section{Timeline Layout Engine}
\label{sec:timeline-rendering}

This section specifies the layout engine for the Timeline Grammar---the deterministic transformation from Timeline Domain IR to Scene IR. The layout engine is grammar-specific; the rendering backends (Section~\ref{sec:backends}) are shared across all grammars.

\subsection{Layout Pipeline Overview}

The Timeline layout engine defines a \emph{deterministic mapping} from a Timeline IR document and a
theme to Scene IR geometry: positions, dimensions, label placements, and visual
treatments. ``Deterministic'' is not aspirational; it is a hard contract. Two conforming
layout engines, given identical IR and theme, must produce geometrically identical Scene IR~\cite{vegalite2017}.

The engine is organized as an ordered six-phase pipeline. Each phase consumes only outputs
of preceding phases and produces geometry for later phases. The pipeline is \emph{pure}: no
phase reads from the IR after its designated turn, and no phase performs I/O, reads a system
clock, or invokes a random number generator.

\subsection{Determinism Contract}
\label{sec:rendering-determinism}

Two conforming layout engines $R_1$ and $R_2$, given the same IR document $D$ and theme $\theta$,
must produce identical coordinate values, label positions, and visual-treatment assignments.
This contract requires:

\begin{enumerate}
  \item \textbf{Pure function.}
        Output is a pure function of $(D,\,\theta)$. No external state---timestamps, random
        values, environment variables, system locale, or screen resolution---influences any
        computed value.

  \item \textbf{Stable sort keys.}
        Every ordered collection is sorted by an explicit, unambiguous key sequence.
        Collections without semantic ordering sort by \texttt{id} alphabetically:
        \begin{itemize}
          \item Tracks: by \texttt{index} ascending (integer, unique by well-formedness invariant).
          \item Activities within a track: by \texttt{(start\_ordinal,\;id)} ascending.
          \item Milestones: by \texttt{(date\_ordinal,\;id)} ascending.
          \item Annotations, sections, groups: by \texttt{id} ascending.
        \end{itemize}
        The ID values break all ties; global uniqueness is guaranteed by the IR invariant.

  \item \textbf{Fixed rounding.}
        All intermediate floating-point values are rounded using \emph{round-half-up}
        ($\lfloor v + 0.5 \rfloor$). Scene geometry values are expressed to two decimal
        places using the same convention. No other rounding rule is used anywhere in the
        pipeline.

  \item \textbf{Deterministic symbolic date resolution.}
        The symbolic date \texttt{now} and relative dates (\texttt{+3m}, \texttt{-2w})
        resolve to a fixed anchor. The anchor is: \texttt{metadata.today} if present;
        else \texttt{metadata.created}; else a validation error (see Gap~1 in
        \S\ref{sec:rendering-gaps}). The system clock is never consulted.

  \item \textbf{Embedded font metrics.}
        Label-width and line-height computations use metrics from font files bundled with
        the theme. System fonts are not consulted for any layout-critical measurement.
        This guarantees that label placement is identical on macOS, Linux, and Windows.

  \item \textbf{Specification-version governance.}
        The renderer reads \texttt{version} from the IR document and verifies it against a
        supported specification version before beginning layout. A version mismatch is a
        hard error, not a warning.

  \item \textbf{Layered determinism.}
        The determinism contract applies at three levels. (i)~\emph{Scene geometry} is
        always byte-deterministic: the layout pipeline output is bitwise-reproducible
        across runs, platforms, and renderer implementations. (ii)~\emph{Per-backend
        output} is deterministic given a pinned backend version and fixed effect seeds;
        the backend version is an explicit parameter of the contract.
        (iii)~\emph{Cross-backend pixel identity} is explicitly \emph{not} promised: the
        SVG backend and the Raster backend are permitted and expected to produce visually
        different outputs from the same Scene. This is correct behaviour at different
        fidelity tiers, not a defect. See Section~\ref{sec:scene-determinism} for the
        full layered contract.
\end{enumerate}

\paragraph{Consequence.}
If $R_1$ and $R_2$ both satisfy this contract and their outputs diverge, the divergent
renderer has a defect. The specification is the arbiter.

\subsection{Coordinate System}
\label{sec:rendering-coords}

The renderer works in a \emph{logical-pixel} coordinate system. The top-left corner of
the canvas is the origin $(0,\,0)$; $x$ increases rightward, $y$ downward. All layout
computations use integer arithmetic in logical pixels; the Scene/Render IR records
decimal-valued coordinates as its output (see Section~\ref{sec:scene-output}).

\begin{figure}[h]
\centering
\begin{Verbatim}[frame=single,fontsize=\small]
  x=0           x=Hhdr+mL              x=canvas.width
  +--------------+----------------------------------+
  |              |  Axis header (Haxis px)          |  y = mT
  |  Track       +----------------------------------+
  |  header      |  Track[0]  (row_height px)       |
  |  column      |  [activity bars, milestones]     |
  |  (Hhdr px)   +----------------------------------+
  |              |  Track[1]  (row_height px)       |
  |              |    ...                            |
  |              +----------------------------------+
  |              |  Track[N-1]                      |
  +--------------+----------------------------------+
                                         y = canvas.height
\end{Verbatim}
\caption{Canvas coordinate system and layout zones}
\label{fig:canvas-layout}
\end{figure}

\begin{table}[h]
\centering
\small
\begin{tabular}{lp{9.5cm}}
\toprule
\textbf{Symbol} & \textbf{Definition} \\
\midrule
$W$ & \texttt{theme.canvas.width} --- total canvas width in logical pixels \\
$m_L,\,m_R,\,m_T,\,m_B$ & Left, right, top, bottom margins \\
$H_{\mathrm{hdr}}$ & \texttt{theme.track.header\_width} --- left column for track labels \\
$H_{\mathrm{axis}}$ & \texttt{theme.axis.height} --- horizontal time-axis strip \\
$W_{\mathrm{draw}}$ & $W - m_L - m_R - H_{\mathrm{hdr}}$ --- drawable timeline width \\
$H_{\mathrm{draw}}$ & $\displaystyle\sum_{i} h_i + \sum_{i>0} g$ --- total height of all track rows plus gaps \\
$H$ & Computed canvas height: $m_T + H_{\mathrm{axis}} + H_{\mathrm{draw}} + m_B$ \\
\bottomrule
\end{tabular}
\caption{Coordinate system symbols}
\label{tab:coord-symbols}
\end{table}

The canvas \emph{width} is fixed by the theme. The canvas \emph{height} is always computed
from track count and sub-lane expansion; it must never be truncated.

\subsection{Timeline Axis}
\label{sec:rendering-axis}

\subsubsection{Axis Unit and Inference}
\label{sec:axis-inference}

\texttt{metadata.axis\_unit} controls tick granularity. If absent, the renderer infers a
suitable unit from the \texttt{time\_range} span, applying the first matching threshold:

\begin{table}[h]
\centering
\begin{tabular}{ll}
\toprule
\textbf{time\_range span (days, inclusive)} & \textbf{Inferred axis\_unit} \\
\midrule
$\leq 14$  & \texttt{day}     \\
$\leq 91$  & \texttt{week}    \\
$\leq 548$ & \texttt{month}   \\
$\leq 1461$& \texttt{quarter} \\
$\leq 2922$& \texttt{half}    \\
$> 2922$   & \texttt{year}    \\
\bottomrule
\end{tabular}
\caption{Axis unit inference thresholds (first match wins)}
\label{tab:axis-inference}
\end{table}

These thresholds yield 4--26 visible ticks at a 1200-pixel canvas width---the range that
avoids crowding while remaining informative for the executive-roadmap use cases described
in Section~\ref{sec:principles}.

\subsubsection{Date-to-Coordinate Function}
\label{sec:date-to-x}

All horizontal positions derive from a single deterministic function $x(T)$. Let
$T_{\mathrm{ord}}$ denote the \emph{day ordinal} of date $T$: an integer count of days
since 2000-01-01 (epoch day~0). Coarser-precision dates are resolved to their period-start
day before ordinal conversion (see \S\ref{sec:date-coercion}).

\[
  x(T) \;=\; H_{\mathrm{hdr}} + m_L \;+\;
  \left\lfloor
    \dfrac{
      (T_{\mathrm{ord}} - T_{s,\mathrm{ord}}) \;\cdot\; W_{\mathrm{draw}}
    }{
      T_{e,\mathrm{ord}} - T_{s,\mathrm{ord}}
    }
    \;+\; 0.5
  \right\rfloor
\]

where $T_{s,\mathrm{ord}}$ and $T_{e,\mathrm{ord}}$ are the day ordinals of the axis domain
start and end after coercion. The numerator product
$(T_{\mathrm{ord}} - T_{s,\mathrm{ord}}) \cdot W_{\mathrm{draw}}$ is computed before division
to prevent precision loss. For a 1200-pixel canvas spanning ten years (3\,650 days), the
intermediate product is at most $1200 \times 3650 = 4\,380\,000$, well within 32-bit
integer range.

\paragraph{Boundary invariants.}
$x(T_s) = H_{\mathrm{hdr}} + m_L$ exactly;
$x(T_e) = H_{\mathrm{hdr}} + m_L + W_{\mathrm{draw}}$ exactly. Any date at or before $T_s$
maps to the left canvas boundary; any date at or after $T_e$ maps to the right canvas
boundary, before clipping.

\subsubsection{Date Coercion for Bar Edges}
\label{sec:date-coercion}

Dates coarser than \texttt{day} must be coerced to a specific calendar day before ordinal
conversion. The rule depends on whether the date is a \emph{left edge} (start of a bar) or
a \emph{right edge} (end of a bar):

\begin{table}[h]
\centering
\small
\begin{tabular}{lll}
\toprule
\textbf{Date form} & \textbf{As left edge (start)} & \textbf{As right edge (end)} \\
\midrule
\texttt{2026-06-09}  (day)    & 2026-06-09 & 2026-06-09 \\
\texttt{2026-06}     (month)  & 2026-06-01 & 2026-06-30 \\
\texttt{2026}        (year)   & 2026-01-01 & 2026-12-31 \\
\texttt{2026-Q2}     (quarter)& 2026-04-01 & 2026-06-30 \\
\texttt{2026-H1}     (half)   & 2026-01-01 & 2026-06-30 \\
\texttt{FY26-Q2}     (fiscal) & First day of fiscal Q2 & Last day of fiscal Q2 \\
\bottomrule
\end{tabular}
\caption{Date coercion to day boundaries by edge role}
\label{tab:date-coercion}
\end{table}

\textbf{Fiscal dates} require a fiscal-year start month that is currently absent from the IR
contract (see Gap~2 in \S\ref{sec:rendering-gaps}).

The same left-edge/right-edge coercion applies to \texttt{time\_range.start} (left) and
\texttt{time\_range.end} (right) when establishing the axis domain.

\subsubsection{Tick and Gridline Placement}
\label{sec:tick-placement}

\begin{enumerate}
  \item Enumerate all \texttt{axis\_unit} period-start dates within $[T_s,\,T_e]$ inclusive.
  \item For each period boundary $T_k$: compute $x_k = x(T_k)$ per \S\ref{sec:date-to-x}.
  \item Place a tick mark of height \texttt{theme.axis.tick\_height}~px at $x_k$, at the
        bottom of the axis header band.
  \item Place a vertical gridline at $x_k$ spanning from the top of Track[0] to the bottom
        of Track[$N-1$], styled per \texttt{theme.axis.gridline\_*} properties.
  \item Place a tick label centered at $x_k$, \texttt{theme.axis.tick\_label\_offset}~px
        below the tick. Format using \texttt{metadata.locale} and the label rules in
        Table~\ref{tab:tick-label-density}. If a label would overlap its neighbor (right
        edge $>$ neighbor's left edge), suppress the label for this tick only.
\end{enumerate}

\begin{table}[h]
\centering
\small
\begin{tabular}{lll}
\toprule
\textbf{axis\_unit} & \textbf{Primary label content} & \textbf{Secondary label (every $N$ ticks)} \\
\midrule
\texttt{day}     & Day-of-month number         & Month name (every 7) \\
\texttt{week}    & Week-start day              & Month name (every 4) \\
\texttt{month}   & Month abbreviation          & Year number (every 12) \\
\texttt{quarter} & Q1 / Q2 / Q3 / Q4           & Year number (every 4) \\
\texttt{half}    & H1 / H2                     & Year number (every 2) \\
\texttt{year}    & Year number                 & None \\
\bottomrule
\end{tabular}
\caption{Tick label content rules per axis unit}
\label{tab:tick-label-density}
\end{table}

\paragraph{Week ticks.}
For \texttt{axis\_unit = week}, week boundaries are ISO~8601 Monday starts regardless of
locale. The locale affects label string formatting only (language, order), never tick positions.

\paragraph{Section bands.}
If \texttt{sections[]} is non-empty, the theme may render alternating column shading between
section boundaries. Band edges coincide with section start/end dates coerced per
Table~\ref{tab:date-coercion}. Section bands are painted at the lowest z-order, below
gridlines and all content elements.

\subsection{The Six-Phase Layout Algorithm}
\label{sec:layout-algorithm}

The renderer executes geometry computation in exactly six phases in the order given below.
Each phase's output is immutable input to subsequent phases. No phase may invoke logic from
a later phase.

\begin{lstlisting}[caption={Layout pipeline overview}]
Phase 1 | AXIS COMPUTATION
        |   Input:  IR metadata (time_range, axis_unit, today-anchor)
        |   Output: axis domain [Ts, Te], axis_unit, tick_positions[]

Phase 2 | TRACK PLACEMENT
        |   Input:  tracks[] sorted by index
        |   Output: track.y_top[], track.height[] (provisional)

Phase 3 | ACTIVITY GEOMETRY
        |   Input:  activities[], Phase-1 axis, Phase-2 track positions
        |   Output: activity.{x_left, x_right, y, lane, width}
        |           track.height[] (final, expanded for sub-lanes)

Phase 4 | MILESTONE GEOMETRY
        |   Input:  milestones[], Phase-1 axis, Phase-3 final track heights
        |   Output: milestone.{x_center, y_center, stack_index}

Phase 5 | SECTIONS AND ANNOTATIONS
        |   Input:  sections[], annotations[], all Phase 1-4 geometry
        |   Output: section.{x_left, x_right}, annotation.{x, y}

Phase 6 | LABEL COLLISION RESOLUTION
        |   Input:  all geometry from Phases 3-5
        |   Output: final label positions for milestones
\end{lstlisting}

\subsubsection{Phase 1: Axis Computation}
\label{sec:phase1}

\begin{lstlisting}[caption={Phase 1: Axis computation}]
1. Resolve today-anchor:
     if metadata.today is set : anchor = metadata.today
     elif metadata.created is set : anchor = metadata.created
     else : ERROR("Cannot resolve symbolic/relative dates: no today/created")

2. Resolve time_range.start:
     a. Expand 'now' -> anchor; expand '+3m' -> anchor + 3 months; etc.
     b. Coerce to period start (left-edge rule, Table 5.3).  -> T_s

3. Resolve time_range.end (same expansion; coerce to period end -> T_e).
     ASSERT T_e >= T_s; else ERROR.

4. Infer axis_unit if absent:
     span_days = T_e_ord - T_s_ord
     Apply Table 5.2 thresholds in order; assign axis_unit.

5. Compute W_draw = canvas.width - mL - mR - theme.track.header_width.

6. Enumerate tick positions:
     T_k = all period-start dates of axis_unit in [T_s, T_e] inclusive.
     For each T_k: x_k = x(T_k) per Section 5.3.2.
     Compute tick label strings using metadata.locale.
\end{lstlisting}

\subsubsection{Phase 2: Track Placement}
\label{sec:phase2}

\begin{lstlisting}[caption={Phase 2: Track placement}]
1. Sort tracks by index ascending.
   ASSERT all index values are unique non-negative integers.

2. y_cursor = mT + Haxis

3. For i in 0 .. N_tracks-1:
     track[i].y_top  = y_cursor
     track[i].height = theme.track.row_height     -- provisional
     y_cursor       += theme.track.row_height + theme.track.row_gap

4. Record group bracket extents after Phase 3 finalises track heights.
\end{lstlisting}

Track heights set here are provisional. Phase~3 may expand them when overlapping activities
require sub-lanes. Downstream phases must use the Phase-3-finalised heights.

\subsubsection{Phase 3: Activity Geometry}
\label{sec:phase3}

\paragraph{Sub-lane assignment.}

\begin{lstlisting}[caption={Phase 3: Activity geometry and greedy sub-lane assignment}]
For each track T (in index order):
  1. A = activities where A.track == T.id.
  2. Resolve each A.start_ord = ordinal(coerce_left(A.start)).
     Resolve each A.end_x per special-end-date rules below.
  3. Sort A by (start_ord, id) ascending (stable sort).
  4. Greedy sub-lane assignment:
       lane_end_x = [-infinity]
       For each A_i in sort order:
         k = smallest index where lane_end_x[k] <= x(A_i.start_ord)
         if no such k exists:
           if len(lane_end_x) < theme.track.max_sub_lanes:
             append -infinity to lane_end_x; k = new last index
           else:
             k = theme.track.max_sub_lanes - 1   -- cap; emit WARNING
         A_i.lane    = k
         lane_end_x[k] = A_i.end_x
  5. n_lanes = max(A.lane) + 1
  6. if n_lanes > 1:
       track[T].height = theme.track.row_height
                       + (n_lanes - 1) * theme.track.sub_lane_height
  7. For each A_i:
       A_i.y = track[T].y_top
             + theme.activity.bar_top_margin
             + A_i.lane * theme.track.sub_lane_height
       A_i.x_left  = x(A_i.start_ord)
       A_i.x_right = A_i.end_x
       logical_w   = A_i.x_right - A_i.x_left
       if logical_w < theme.activity.min_width:
         mid         = (A_i.x_left + A_i.x_right) / 2   -- round-half-up
         A_i.x_left  = mid - theme.activity.min_width / 2
         A_i.x_right = A_i.x_left + theme.activity.min_width
       A_i.width = A_i.x_right - A_i.x_left
\end{lstlisting}

\paragraph{Special end-date resolution.}
\begin{itemize}
  \item \texttt{end} absent (omitted): treated identically to \texttt{ongoing}.
  \item \texttt{ongoing}: \texttt{end\_x} $=$ $H_{\mathrm{hdr}} + m_L + W_{\mathrm{draw}}$
        (right canvas boundary); append right-chevron ($\rightarrow$) decoration.
  \item \texttt{tbd}: \texttt{end\_x} $=$ \texttt{x\_left} $+$
        \texttt{theme.activity.tbd\_extension\_px}
        (default: $2 \times \overline{\Delta x_{\text{tick}}}$, the mean tick spacing in pixels);
        render rightward extension with dashed border at 50\% opacity;
        place ``TBD'' label inside extension zone.
  \item \texttt{unknown}: same geometry as \texttt{tbd}; use \texttt{unknown} visual
        treatment from theme status map.
  \item \texttt{\textasciitilde{}date}: coerce date normally to ordinal; \texttt{end\_x}
        $= x(\text{date\_ord})$; flag right edge as approximate for gradient-fade rendering.
\end{itemize}

\paragraph{Label placement (deterministic three-way rule).}
The rule is applied exactly once per activity; no iteration:
\begin{enumerate}
  \item \textbf{Inside}: if \texttt{A.width} $\geq$ \texttt{theme.activity.label\_inside\_min\_width},
        render label horizontally centered in bar, vertically centered on bar height, in
        \texttt{theme.activity.label\_color\_inside}.
  \item \textbf{Right}: render label starting at $\texttt{A.x\_right} + 4$~px in
        \texttt{theme.activity.label\_color\_outside}.
        If label right edge $>$ $H_{\mathrm{hdr}} + m_L + W_{\mathrm{draw}} - 4$, proceed to
        step~3.
  \item \textbf{Left}: render label ending at $\texttt{A.x\_left} - 4$~px.
        If label left edge $<$ $H_{\mathrm{hdr}} + m_L$, truncate label to
        \texttt{theme.activity.label\_truncate\_chars} characters $+$
        ``\ldots'' (U+2026) and re-attempt step~2.
\end{enumerate}

\paragraph{Progress indicator.}
If \texttt{activity.progress} is set: render a filled strip of height
\texttt{theme.activity.progress\_bar\_height}~px at the bottom of the bar, width $=
\lfloor \texttt{A.width} \cdot \texttt{A.progress} + 0.5 \rfloor$~px, inset inside the
bar border.

\subsubsection{Phase 4: Milestone Geometry}
\label{sec:phase4}

\paragraph{Placement and stacking.}

\begin{lstlisting}[caption={Phase 4: Milestone placement and stacking}]
For each milestone M (sorted by date_ord, id):
  1. x_center = x(M.date_ordinal).
  2. Determine base_y:
       if M.track is set:
         base_y = track[M.track].y_top + track[M.track].height / 2
       else (cross-track):
         base_y = mT + Haxis + H_draw / 2

  3. Stacking within (effective_track_id, x_center) group:
       Sort group members by id ascending.
       Assign stack_index k = 0, 1, 2, ...
       M.y_center = base_y + k * theme.milestone.stack_offset_y

  4. Diamond vertices (integer arithmetic):
       d = theme.milestone.diamond_size
       top    = (x_center,     M.y_center - d)
       right  = (x_center + d, M.y_center    )
       bottom = (x_center,     M.y_center + d)
       left   = (x_center - d, M.y_center    )
\end{lstlisting}

The diamond is a four-vertex polygon with vertices computed directly. No
\texttt{transform="rotate(...)"} attribute is used; rotation-transform-dependent rendering
differences across SVG viewers are avoided.

\paragraph{Cross-track milestones.}
When \texttt{track} is null, the milestone is drawn at the vertical centre of the full
rendered area. Its diamond size is \texttt{theme.milestone.diamond\_size} (not scaled);
visual emphasis comes from the spanning vertical extent of the label, which may be increased
by a theme multiplier \texttt{theme.milestone.cross\_track\_label\_scale} (default: 1.2).

\subsubsection{Phase 5: Sections and Annotations}
\label{sec:phase5}

\paragraph{Section bands.}
For each section $S$ (sorted by start date, then id):
\begin{itemize}
  \item $x_{\text{left}}(S) = x(\text{coerce\_left}(S.\texttt{start}))$.
  \item $x_{\text{right}}(S) = x(\text{coerce\_right}(S.\texttt{end}))$.
  \item Band spans vertically from $m_T + H_{\mathrm{axis}}$ to
        $m_T + H_{\mathrm{axis}} + H_{\mathrm{draw}}$.
  \item Z-order: sections paint first (below gridlines, activities, milestones, labels).
\end{itemize}
Overlapping sections are valid; a later section in sort order paints over an earlier one.

\paragraph{Annotation placement.}
\begin{enumerate}
  \item Sort annotations by \texttt{id} ascending.
  \item For each annotation: compute anchor as centre of target entity's bounding box.
  \item Place annotation box in the first quadrant that keeps it entirely within the canvas.
        Preference order: top-right, top-left, bottom-right, bottom-left.
        If no quadrant fits, use top-right and clip to canvas boundary.
  \item Draw connector (line/elbow) from annotation box edge to anchor,
        styled per \texttt{theme.annotation.connector\_style}.
\end{enumerate}

\subsubsection{Phase 6: Label Collision Resolution}
\label{sec:phase6}

Activity labels are placed definitively in Phase~3 (no iteration). Milestone labels require
a collision-resolution pass because multiple milestones at similar $x$ positions can produce
overlapping labels.

\begin{lstlisting}[caption={Phase 6: Milestone label collision resolution}]
1. Collect all milestone label bounding boxes (x_label, y_label, width, height, id).
2. Sort by (x_label, y_label, id) ascending (stable).
3. Scan pass (bounded by N iterations, N = label count):
     For each consecutive pair (L_i, L_{i+1}):
       if L_i.right_edge > L_{i+1}.left_edge
          AND |L_i.y_label - L_{i+1}.y_label| < label_height:
         L_{i+1}.y_label += theme.milestone.label_stack_offset
4. Repeat step 3 until no overlapping pairs remain OR N iterations reached.
   If overlaps persist after N iterations: truncate labels to
   theme.milestone.label_max_width px, accepting visual overlap.
5. Clamp all final y_label values to [mT, mT + Haxis + H_draw + mB].
\end{lstlisting}

Because the input sort and every shift decision are deterministic, two renderers executing
this algorithm on identical inputs produce identical label placements.

\subsection{Scene / Render IR --- Output of the Pipeline}
\label{sec:scene-output}

The six-phase layout pipeline does not target any output format directly. Its output is the
\textbf{Scene / Render IR}: a byte-deterministic, backend-agnostic record of every drawing
primitive, resolved coordinate, visual treatment, and effect request. The Scene is described
fully in Section~\ref{sec:scene-ir}; its role here is to name what the pipeline produces
and to record what each phase contributes.

Every value computed in Phases 1--6 becomes a field of a Scene drawing primitive or a
top-level Scene property. The Scene is a pure data record: it carries no rendering
instructions specific to SVG, rasterisation, or PowerPoint. Rendering backends (SVG,
Raster, PPTX, HTML) consume the Scene and emit their respective formats; they are described
in Section~\ref{sec:outputs}.

\paragraph{What the Scene preserves from the pipeline.}
\begin{itemize}
  \item \textbf{Phase 1} outputs: axis domain $[T_s,\,T_e]$, axis unit, tick positions
        $x_k$, and tick label strings.
  \item \textbf{Phases 2--3} outputs: all activity bar coordinates
        $(x_{\text{left}},\,x_{\text{right}},\,y,\,\text{width})$, lane assignments,
        and final track heights.
  \item \textbf{Phase 4} outputs: all milestone vertex coordinates and stack indices.
  \item \textbf{Phase 5} outputs: section band bounds
        $(x_{\text{left}},\,x_{\text{right}})$, annotation positions, and connector
        endpoints.
  \item \textbf{Phase 6} outputs: final label positions after collision resolution.
  \item \textbf{Theme-resolved visual treatments}: fill colours, stroke styles, opacity,
        corner radii, pattern references, and effect requests with their fallback policies
        (see Section~\ref{sec:schema-fidelity}).
\end{itemize}

The Scene does not carry the IR document or theme; it is the fully resolved,
self-contained rendering specification. A backend needs only the Scene and its own
capability profile to produce output.

\subsection{Edge Cases}
\label{sec:edge-cases}

Every case below is normative. A renderer that produces different behaviour for any case
has a defect. The summary table gives the complete catalogue; extended descriptions follow
for the most structurally complex cases.

\begin{longtable}{@{}p{3.2cm}p{3.4cm}p{5.8cm}@{}}
\caption{Edge-case definitions (normative)}
\label{tab:edge-cases} \\
\toprule
\textbf{Case} & \textbf{Condition} & \textbf{Rendering Rule} \\
\midrule
\endfirsthead
\multicolumn{3}{c}{\small\textit{Table~\ref{tab:edge-cases} continued}} \\[2pt]
\toprule
\textbf{Case} & \textbf{Condition} & \textbf{Rendering Rule} \\
\midrule
\endhead
\midrule
\multicolumn{3}{r}{\small\textit{continued \ldots}} \\
\endfoot
\bottomrule
\endlastfoot
Zero-duration activity &
  \texttt{start == end} or \texttt{span} that resolves to a single point &
  Render bar of \texttt{min\_width}~px centred at $x(\text{start})$. Never render 0-width. Label placed right (outside). \\[4pt]
Overlapping activities (same track) &
  Two activities in same track whose date ranges overlap &
  Greedy sub-lane assignment (Phase 3). Track height expands. Cap at \texttt{max\_sub\_lanes}. \\[4pt]
Open interval (\texttt{ongoing}) &
  \texttt{end: ongoing} or \texttt{end} absent &
  Bar extends to right canvas boundary. Right-chevron decoration. No overflow beyond canvas edge. \\[4pt]
TBD end date &
  \texttt{end: tbd} &
  Bar extends \texttt{tbd\_extension\_px} beyond \texttt{x\_left}. Dashed, 50\% opacity. ``TBD'' label inside extension. \\[4pt]
Both dates TBD &
  \texttt{start: tbd} or equivalent &
  Full-track-width hatched block at 40\% opacity. Activity label $+$ ``(TBD)''. \\[4pt]
Unknown start &
  \texttt{start: unknown} &
  Left edge placed at $x(T_s)$. Left-fade gradient over \texttt{approx\_fade\_px}. \\[4pt]
Unknown end &
  \texttt{end: unknown} &
  Same geometry as \texttt{tbd}; rendered with \texttt{unknown} visual treatment from theme. \\[4pt]
Approximate start &
  \texttt{start: \textasciitilde{}date} &
  Geometry at nominal date; left edge rendered with gradient fade of \texttt{approx\_fade\_px}. \\[4pt]
Approximate end &
  \texttt{end: \textasciitilde{}date} &
  Geometry at nominal date; right edge rendered with gradient fade. \\[4pt]
Activity fully before range &
  $\text{end} < T_s$ &
  Not rendered. Renderer warning. Activity still appears in legend. \\[4pt]
Activity fully after range &
  $\text{start} > T_e$ &
  Not rendered. Renderer warning. Activity still appears in legend. \\[4pt]
Activity partially before range &
  $\text{start} < T_s$, $\text{end} \geq T_s$ &
  Bar clipped to $[x(T_s),\;x(\text{end})]$. Left-clip indicator on clipped edge. \\[4pt]
Activity partially after range &
  $\text{start} \leq T_e$, $\text{end} > T_e$ &
  Bar clipped to $[x(\text{start}),\;x(T_e)]$. Right-clip indicator on clipped edge. \\[4pt]
Very short span &
  $x(\text{end}) - x(\text{start}) < \text{min\_width}$ after rounding &
  Render at \texttt{min\_width}; centre on logical midpoint (round-half-up). Label right (outside). \\[4pt]
Very long span &
  Bar occupies $> 80\%$ of $W_{\mathrm{draw}}$ &
  Render normally. Prefer inside label placement if label fits within bar. \\[4pt]
Simultaneous milestones (same track) &
  Two milestones with equal \texttt{date} on the same \texttt{track} &
  Stack vertically by \texttt{stack\_offset\_y}, sorted by \texttt{id} ascending (Phase 4). \\[4pt]
Simultaneous milestones (cross-track) &
  Two milestones with \texttt{track: null} and equal \texttt{date} &
  Stack at full-timeline vertical centre, sorted by \texttt{id} ascending. \\[4pt]
Milestone outside range &
  Milestone date $< T_s$ or $> T_e$ &
  Not rendered. Renderer warning. Appears in legend. \\[4pt]
Empty track &
  Track has no activities and no milestones &
  Row rendered at \texttt{row\_height}. Label visible. Body empty. Never collapsed. \\[4pt]
Sub-lane cap exceeded &
  Overlapping activities require $>$ \texttt{max\_sub\_lanes} lanes &
  Excess activities placed in the last lane (visible overlap). Renderer warning. \\
\end{longtable}

\paragraph{Overlapping activities: sub-lane assignment in detail.}

The greedy algorithm assigns each activity (in stable \texttt{(start\_ord,\;id)} order) to the
lowest-indexed lane where it does not overlap any previously placed activity. Overlap is
defined as $x_{\text{left}}(A_j) < x_{\text{right}}(A_i)$ for activities $A_i$, $A_j$ in
the same lane. Because sort order is deterministic and each assignment is a pure function of
the current \texttt{lane\_end\_x} state, all renderers produce identical lane assignments
for identical inputs.

\paragraph{Simultaneous milestones in detail.}

When milestones share an $(\texttt{track},\,x_{\text{center}})$ pair, they are stacked downward
from the nominal track centre. The first milestone in \texttt{id}-ascending order is at the
base position (stack index 0). Subsequent milestones are offset by
$k \times \texttt{theme.milestone.stack\_offset\_y}$. Stacking is permitted to overflow into
the inter-track gap; milestones must \emph{not} be silently suppressed if they cannot fit
within track height.

\paragraph{Clip indicators.}

An activity bar clipped at the axis boundary receives a \emph{clip indicator}: a small
``/\!/'' mark (two angled lines, $4 \times 8$~px) drawn in the activity's fill colour at
full opacity, positioned at the clipped edge centre. The indicator signals to the reader
that the activity continues beyond the visible axis range.

\paragraph{Minimum-width semantics.}

\texttt{theme.activity.min\_width} (default: 4~px) prevents activities from being rendered
invisibly narrow at coarse zoom levels. A zero-duration activity and a very-short-span
activity both produce a bar of exactly \texttt{min\_width} px. The two cases are
\emph{visually indistinguishable}; distinguishing them would require a separate icon or
tooltip, which is a theme-level option outside the core rendering contract.

\subsection{Known IR Gaps}
\label{sec:rendering-gaps}

The following gaps in Mark's IR contract (\S\ref{sec:ir}) affect rendering determinism.
They are flagged here for Mark and Leslie to resolve; this spec does \emph{not} unilaterally
extend the IR. Renderers encountering these gaps should apply the documented fallback and
emit a validation warning.

\begin{description}
  \item[Gap 1 --- \texttt{metadata.today} field missing.]
    Leslie's decision record mandates an explicit \texttt{today} field in the IR for
    deterministic resolution of the \texttt{now} symbolic date and the today-marker
    feature. Mark's current \texttt{metadata} object does not include this field.
    \textbf{Proposed resolution:} add \texttt{today} (type \texttt{date}, optional) to
    \texttt{metadata}. Renderer fallback chain:
    \texttt{metadata.today} $\to$ \texttt{metadata.created} $\to$ validation error.

  \item[Gap 2 --- \texttt{metadata.fiscal\_year\_start} missing.]
    The \texttt{FY26-Q2} fiscal-date format requires knowing the fiscal-year start month.
    Two renderers assuming different organisations' fiscal calendars (US federal: October;
    UK: April; common corporate: January) will produce different $x$-coordinates for
    identical fiscal dates.
    \textbf{Proposed resolution:} add \texttt{fiscal\_year\_start} (type \texttt{int},
    range 1--12, default 1 = January) to \texttt{metadata}.

  \item[Gap 3 --- Relative date anchor undefined.]
    Relative dates (\texttt{+3m}, \texttt{-2w}) are described in Mark's contract as
    resolving from ``context'', which is insufficient. A renderer must know the anchor date
    to compute an ordinal.
    \textbf{Proposed resolution:} define that relative dates use the same anchor as Gap~1
    (\texttt{today} $\to$ \texttt{created} $\to$ error).

  \item[Gap 4 --- Omitted \texttt{end} semantics ambiguous.]
    Mark's Activity contract marks \texttt{end} as optional but does not state whether an
    absent \texttt{end} is equivalent to \texttt{ongoing} or a validation error. This
    rendering model treats omitted \texttt{end} as \texttt{ongoing} (see the open-interval
    row in Table~\ref{tab:edge-cases}). The IR spec should make this explicit.
\end{description}
             % Timeline rendering specifics
\section{Flow Grammar: Domain IR}
\label{sec:flow-grammar}

The Flow Grammar is the second grammar implemented in the diagram compiler and the \textbf{template grammar} for all future grammar additions. Where the Timeline Grammar (Section~\ref{sec:timeline-grammar}) proved the kernel's Scene IR contract with a time-axis-centric vocabulary, the Flow Grammar proves that the kernel is genuinely grammar-agnostic: it supports node-link topologies, pipeline sequences, and process flows---fundamentally different visual semantics---without kernel modifications.

% ============================================================================
\subsection{Scope \& Intent}
\label{sec:flow-scope}
% ============================================================================

\subsubsection{What a Flow Diagram IS}

A \textbf{flow diagram} in this grammar is a \emph{directed node-link diagram} representing:

\begin{itemize}
  \item \textbf{Pipelines/sequences}: ordered chains of processing steps (e.g., CI/CD, RAG retrieval, ETL).
  \item \textbf{Process flows}: multi-path directed graphs with branching and joining (e.g., decision trees, agent loops, state machines).
  \item \textbf{Architecture topologies}: layered or grouped node-link structures (e.g., microservice communication, data flow).
\end{itemize}

These correspond to the corpus taxonomy's \emph{node-link/topology} and \emph{pipeline/sequence} kinds (Section~\ref{sec:corpus-taxonomy}). The animated-arrow variant (flowing dashes along connectors) is a first-class pattern in this grammar.

\subsubsection{What a Flow Diagram is NOT}

\begin{description}
  \item[Not a comparison/matrix.] Tabular grid layouts with cells (Section~\ref{sec:corpus-comparison-matrix}) belong to the Comparison grammar. Flows have directed edges; comparisons do not.
  \item[Not a timeline.] The timeline grammar's entities have temporal semantics (\texttt{start}, \texttt{end}, \texttt{track}). Flow nodes have no time axis.
  \item[Not a general graph.] The Graph grammar (Section~\ref{sec:graph-grammar}) covers undirected networks, trees with Reingold-Tilford layout, and hub-and-spoke topologies. The Flow grammar is restricted to \emph{directed} structures with a declared layout intent---it does not attempt general force-directed placement.
\end{description}

\subsubsection{Modeling Boundary}

The Flow grammar covers diagrams where:
\begin{enumerate}
  \item Entities are discrete \textbf{nodes} (processing steps, decisions, data stores).
  \item Relationships are \textbf{directed edges} (data flow, control flow, sequence).
  \item Layout admits a natural \textbf{directional reading order} (left-to-right, top-to-bottom).
\end{enumerate}

Diagrams without directed edges, or without a dominant reading direction, are out of scope for this grammar.

% ============================================================================
\subsection{Flow Domain IR}
\label{sec:flow-domain-ir}
% ============================================================================

The Flow Domain IR is the small, grammar-specific representation that compiles \emph{down} to the shared Scene IR (Section~\ref{sec:scene-ir}). It is \textbf{not} a ``god-IR''---it contains only the constructs necessary to describe directed node-link diagrams with optional grouping.

\subsubsection{Document Structure}

A Flow IR document is a single YAML/JSON object with the following root structure:

\begin{lstlisting}[language=yaml,caption={Flow IR root document structure}]
version: "1.0"
metadata:
  title: string             # required
  subtitle: string          # optional; default null
  theme: string             # optional; default "default-flow"
flow:
  direction: left-to-right | top-to-bottom  # required
  nodes: [...]              # required; list<FlowNode>; may be empty
  edges: [...]              # optional; list<FlowEdge>; default []
  groups: [...]             # optional; list<FlowGroup>; default []
\end{lstlisting}

\subsubsection{Root Fields}

\begin{table}[h]
\centering
\small
\begin{tabular}{llccp{4.8cm}}
\toprule
\textbf{Field} & \textbf{Type} & \textbf{Req} & \textbf{Default} & \textbf{Description} \\
\midrule
\texttt{version} & \texttt{string} & Yes & --- & Spec version (semver, e.g., \texttt{"1.0"}) \\
\texttt{metadata} & \texttt{FlowMetadata} & Yes & --- & Document-level metadata \\
\texttt{flow} & \texttt{FlowDefinition} & Yes & --- & The flow graph definition \\
\bottomrule
\end{tabular}
\caption{Flow IR root fields}
\label{tab:flow-root-fields}
\end{table}

\subsubsection{FlowNode}
\label{sec:flow-node}

A node is a discrete entity in the flow: a processing step, decision point, data store, or any labeled box.

\begin{table}[h]
\centering
\small
\begin{tabular}{llccp{4.5cm}}
\toprule
\textbf{Field} & \textbf{Type} & \textbf{Req} & \textbf{Default} & \textbf{Description} \\
\midrule
\texttt{id} & \texttt{id} & Yes & --- & Document-unique identifier (kebab-case) \\
\texttt{label} & \texttt{string} & Yes & --- & Display text for the node \\
\texttt{shape} & \texttt{enum\{rect, rounded-rect, circle, diamond, stadium\}} & No & \texttt{rounded-rect} & Visual shape \\
\texttt{icon} & \texttt{string?} & No & \texttt{null} & Icon registry key \\
\texttt{status} & \texttt{enum\{default, active, success, warning, error, muted\}} & No & \texttt{default} & Semantic status (mapped to color by theme) \\
\texttt{description} & \texttt{string?} & No & \texttt{null} & Tooltip/secondary text \\
\texttt{group} & \texttt{ref<FlowGroup>?} & No & \texttt{null} & Group membership (alternative to declaring in group.nodes) \\
\bottomrule
\end{tabular}
\caption{FlowNode fields}
\label{tab:flow-node-fields}
\end{table}

\paragraph{Identity.}
Node \texttt{id} is the primary key within a document. It must be unique across all nodes. References in edges and groups use this \texttt{id}. Two nodes with the same \texttt{id} are a validation error.

\paragraph{Ordering.}
Nodes are ordered by their position in the \texttt{nodes} list. This order is the \textbf{canonical order} used by the layout engine for tie-breaking (determinism) and by serialization for stable output.

\subsubsection{FlowEdge}
\label{sec:flow-edge}

An edge represents a directed relationship between two nodes.

\begin{table}[h]
\centering
\small
\begin{tabular}{llccp{4.2cm}}
\toprule
\textbf{Field} & \textbf{Type} & \textbf{Req} & \textbf{Default} & \textbf{Description} \\
\midrule
\texttt{source} & \texttt{ref<FlowNode>} & Yes & --- & Source node id \\
\texttt{target} & \texttt{ref<FlowNode>} & Yes & --- & Target node id \\
\texttt{source\_port} & \texttt{enum\{auto, top, right, bottom, left\}} & No & \texttt{auto} & Anchor point on source node \\
\texttt{target\_port} & \texttt{enum\{auto, top, right, bottom, left\}} & No & \texttt{auto} & Anchor point on target node \\
\texttt{label} & \texttt{string?} & No & \texttt{null} & Edge annotation text \\
\texttt{style} & \texttt{enum\{solid, dashed, dotted\}} & No & \texttt{solid} & Stroke style \\
\texttt{animated} & \texttt{bool} & No & \texttt{false} & Enable flowing-dash animation \\
\texttt{directedness} & \texttt{enum\{directed, bidirectional, undirected\}} & No & \texttt{directed} & Arrow rendering \\
\bottomrule
\end{tabular}
\caption{FlowEdge fields}
\label{tab:flow-edge-fields}
\end{table}

\paragraph{Identity.}
Edges are identified by position in the \texttt{edges} list (index). Unlike nodes, edges do not carry an explicit \texttt{id}---their identity is their list position. This list position is the canonical order for deterministic processing.

\paragraph{Directedness.}
\begin{description}
  \item[\texttt{directed}] Arrow from source to target (default).
  \item[\texttt{bidirectional}] Arrows on both ends.
  \item[\texttt{undirected}] No arrowheads (connector line only).
\end{description}

\paragraph{Port resolution.}
When \texttt{source\_port} or \texttt{target\_port} is \texttt{auto}, the layout engine assigns the port based on the relative positions of source and target nodes. For left-to-right flows, \texttt{auto} resolves to \texttt{right} on source and \texttt{left} on target for forward edges.

\subsubsection{FlowGroup}
\label{sec:flow-group}

Groups are optional visual containers (lanes, clusters, swim-lanes) that partition nodes.

\begin{table}[h]
\centering
\small
\begin{tabular}{llccp{4.5cm}}
\toprule
\textbf{Field} & \textbf{Type} & \textbf{Req} & \textbf{Default} & \textbf{Description} \\
\midrule
\texttt{id} & \texttt{id} & Yes & --- & Document-unique identifier \\
\texttt{label} & \texttt{string} & Yes & --- & Group heading text \\
\texttt{nodes} & \texttt{list<ref<FlowNode>>} & No & \texttt{[]} & Node ids in this group (augments per-node \texttt{group} field) \\
\texttt{style} & \texttt{enum\{lane, cluster, outline\}} & No & \texttt{lane} & Visual rendering style \\
\bottomrule
\end{tabular}
\caption{FlowGroup fields}
\label{tab:flow-group-fields}
\end{table}

\paragraph{Group membership resolution.}
A node belongs to a group if \emph{either} it appears in the group's \texttt{nodes} list \emph{or} it has a \texttt{group} field referencing the group's \texttt{id}. Both mechanisms are equivalent; using both for the same node--group pair is valid (idempotent). A node may belong to at most one group; dual membership is a validation error.

\paragraph{Ungrouped nodes.}
Nodes not assigned to any group are rendered in the main flow area, outside any group container. They participate in layout normally.

\subsubsection{Layout Intent}
\label{sec:flow-layout-intent}

The \texttt{flow.direction} field declares the dominant layout direction:

\begin{description}
  \item[\texttt{left-to-right}] Nodes are ranked left-to-right; edges flow rightward. This is the default for pipeline/sequence diagrams.
  \item[\texttt{top-to-bottom}] Nodes are ranked top-to-bottom; edges flow downward. This suits vertical process flows.
\end{description}

The declared direction is a \textbf{hard constraint} on the layout engine: it determines layer-assignment orientation in the Sugiyama algorithm and reading order in the linear-sequence layout.

% ============================================================================
\subsection{Layout Contract}
\label{sec:flow-layout-contract}
% ============================================================================

The Flow grammar's layout engine converts the Domain IR into Scene IR coordinates. The choice of algorithm depends on graph topology, but all algorithms must satisfy the \textbf{determinism contract}: identical IR in $\rightarrow$ byte-identical Scene out.

\subsubsection{Layout Family Selection}

\begin{center}
\small
\begin{tabular}{lp{4.5cm}p{5.5cm}}
\toprule
\textbf{Topology} & \textbf{Algorithm} & \textbf{When Applied} \\
\midrule
Linear chain & Linear sequence layout & All nodes have in-degree $\le 1$ and out-degree $\le 1$ (a simple path or set of disjoint paths). \\
DAG & Sugiyama layered~\cite{sugiyama1981} & Directed acyclic graph with branching/joining. \\
Cyclic & Sugiyama with cycle removal & Cycles detected; feedback arcs reversed before layering. \\
\bottomrule
\end{tabular}
\end{center}

\paragraph{Topology detection is automatic.}
The layout engine inspects the graph structure and selects the appropriate algorithm. The author does not choose the algorithm---only the \texttt{direction}. This is a key difference from the Graph grammar (Section~\ref{sec:graph-grammar}), which requires explicit layout selection.

\subsubsection{Sugiyama Layered Layout (DAGs and Cyclic Flows)}

The primary layout algorithm for non-trivial flows follows the four-phase Sugiyama framework~\cite{sugiyama1981}:

\begin{enumerate}
  \item \textbf{Cycle removal.} Detect cycles via DFS. Reverse a minimal feedback arc set to produce a DAG. Reversed edges are rendered with a visual back-edge indicator (curved path routed behind the main flow).
  
  \item \textbf{Layer assignment.} Assign each node to a layer (rank) using network simplex~\cite{gansner1993dot}, minimizing total edge span. The \texttt{direction} field determines whether layers are columns (left-to-right) or rows (top-to-bottom).
  
  \item \textbf{Crossing minimization.} Barycenter heuristic iterated over adjacent layer pairs (24 sweeps, per standard practice~\cite{sugiyama1981}). Ties broken by canonical node list order (determinism).
  
  \item \textbf{Coordinate assignment.} Brandes \& K\"opf~\cite{brandesKopf2001} linear-time algorithm: four candidate alignments, median selected. Guarantees at most two bends per long edge.
\end{enumerate}

\paragraph{Edge routing.}
Edges are routed orthogonally where the Brandes--K\"opf coordinate assignment produces aligned segments. For non-orthogonal cases (long spans, back-edges), edges use smooth cubic B\'ezier curves. The choice is deterministic: given fixed node positions, edge routing is a pure function.

Orthogonal routing follows the principles of Tamassia's topology--shape--metrics framework~\cite{tamassia1987}: edges consist of horizontal and vertical segments only, with bends minimized by the coordinate-assignment phase. This produces clean, architecture-diagram-style connectors.

\subsubsection{Linear Sequence Layout (Pipelines)}

When the flow graph is a simple chain (every node has at most one predecessor and one successor), the layout engine uses a trivial linear placement:

\begin{itemize}
  \item Nodes placed at equal intervals along the primary axis (determined by \texttt{direction}).
  \item Edges are straight horizontal (LTR) or vertical (TTB) connectors.
  \item Groups (if present) span the extent of their member nodes.
\end{itemize}

This produces clean, minimal pipeline diagrams without the overhead of Sugiyama layer assignment.

\subsubsection{Determinism Mandate}

\textbf{All flow layouts are fully deterministic.} Specifically:

\begin{enumerate}
  \item No randomized initialization. No force-directed simulation. No RNG.
  \item Tie-breaking uses the canonical node order (list position in the IR).
  \item Crossing minimization uses a fixed sweep count (24), not convergence-based termination.
  \item The same IR document, processed by the same engine version, produces byte-identical Scene IR on every execution.
\end{enumerate}

If a force-like aesthetic is ever required for a flow diagram (not anticipated for v1), the implementation \textbf{must} use stress majorization~\cite{gansnerStressMaj2004} with a deterministic initial layout (nodes on a circle in canonical order) and a fixed iteration count. Raw Fruchterman-Reingold~\cite{fruchterman1991} with random initialization is prohibited.

\paragraph{Stress majorization rationale.}
Stress majorization (SMACOF) is a pure function of the graph-distance matrix given a fixed initial configuration. It converges monotonically (no oscillation), producing identical results across runs~\cite{gansnerStressMaj2004}. This is the only acceptable ``force-like'' family for a deterministic compiler.

% ============================================================================
\subsection{Lowering to the Scene IR}
\label{sec:flow-lowering}
% ============================================================================

The flow layout engine emits Scene IR primitives (Section~\ref{sec:scene-ir}). This section specifies the mapping from flow constructs to Scene primitives. \textbf{No new Scene IR primitives are required}---the existing kernel is sufficient.

\subsubsection{Node Lowering}

\begin{center}
\small
\begin{tabular}{lp{9cm}}
\toprule
\textbf{Flow Shape} & \textbf{Scene IR Primitive(s)} \\
\midrule
\texttt{rect} & \texttt{Rect\{corner\_radius: 0\}} \\
\texttt{rounded-rect} & \texttt{Rect\{corner\_radius: theme.flow.nodeRadius\}} \\
\texttt{circle} & \texttt{Circle} \\
\texttt{diamond} & \texttt{Path} (four-point rhombus via \texttt{M/L/L/L/Z} commands) \\
\texttt{stadium} & \texttt{Rect\{corner\_radius: height/2\}} (pill shape) \\
\bottomrule
\end{tabular}
\end{center}

Each node emits a \texttt{Group} containing:
\begin{enumerate}
  \item The shape primitive (background/border).
  \item A \texttt{Text} primitive for the label (centered within the shape).
  \item Optionally, an \texttt{Image} primitive for the icon (positioned above or left of the label, per theme).
\end{enumerate}

\paragraph{Status $\rightarrow$ color.}
The node's \texttt{status} field is resolved by the theme to \texttt{fill} and \texttt{stroke} colors on the shape primitive. The Domain IR does \emph{not} carry color values---color is always theme-resolved.

\subsubsection{Edge Lowering}

Each edge emits:

\begin{enumerate}
  \item A \texttt{Path} primitive for the connector line (routed by the layout engine). Stroke style (\texttt{solid}, \texttt{dashed}, \texttt{dotted}) maps directly to the Scene Path's \texttt{stroke\_style} field.
  
  \item An arrowhead: a small filled \texttt{Path} (triangle) at the target end for \texttt{directed} edges; at both ends for \texttt{bidirectional}; omitted for \texttt{undirected}.
  
  \item Optionally, a \texttt{Text} primitive for the edge label, positioned at the midpoint of the path (offset perpendicular to the edge direction to avoid overlap).
\end{enumerate}

\paragraph{Animated edges.}
When \texttt{animated: true}, the edge's \texttt{Path} primitive receives a \texttt{FlowingDashes} animation hint (Section~\ref{sec:scene-animation}):

\begin{Verbatim}[frame=single,fontsize=\small]
FlowingDashes {
  target_id: <edge_path_id>,
  dash_length: theme.flow.animDashLength,     // default 8px
  gap_length:  theme.flow.animGapLength,      // default 4px
  speed_px_per_sec: theme.flow.animSpeed,     // default 40
  direction: 'forward'
}
\end{Verbatim}

SVG/HTML backends render this as \texttt{stroke-dashoffset} animation. Raster backends (PNG/Skia) render the resting frame: a static dashed stroke.

\subsubsection{Group Lowering}

Groups emit:
\begin{enumerate}
  \item A background \texttt{Rect} spanning the bounding box of all member nodes (with padding from theme).
  \item A \texttt{Text} primitive for the group label (positioned at the top of the rect).
  \item All member nodes rendered as children of a \texttt{Group} primitive (for logical containment, not visual nesting---nodes are positioned absolutely).
\end{enumerate}

\paragraph{No new primitives needed.}
The flow grammar maps entirely to existing Scene IR primitives: \texttt{Rect}, \texttt{Circle}, \texttt{Path}, \texttt{Text}, \texttt{Image}, and \texttt{Group}. No kernel changes are required.

\textbf{Open question for Barbara/Mark:} If future flow diagrams require \emph{nested} groups (groups within groups), should the Scene IR \texttt{Group} primitive support recursive clip regions, or should nesting be flattened at lowering time? Current spec assumes single-level groups only. Nested groups are deferred to v2.

% ============================================================================
\subsection{Edge Cases \& Total Semantics}
\label{sec:flow-edge-cases}
% ============================================================================

Every possible input must have a defined meaning or produce a clear validation error. This section pins down the edge cases.

\subsubsection{Cycles}

\textbf{Cycles are valid.} A flow may contain cycles (e.g., retry loops, feedback arcs in agent pipelines). The layout engine handles cycles by:

\begin{enumerate}
  \item Detecting strongly-connected components via Tarjan's algorithm.
  \item Selecting a minimal feedback arc set (edges to reverse for layering). Selection is deterministic: among equally-scored candidates, the edge appearing \emph{later} in the canonical edge list is chosen for reversal.
  \item Reversed edges are rendered as \emph{back-edges}: routed with a longer path (curved or stepped) that visually indicates ``going back'' in the flow direction.
\end{enumerate}

The resulting diagram remains readable: forward edges flow in the declared direction; back-edges are visually distinct.

\subsubsection{Self-Loops}

\textbf{Self-loops are valid.} An edge where \texttt{source == target} is rendered as a loopback connector: a curved path that exits the node from one port and re-enters at an adjacent port (e.g., exits \texttt{right}, loops around, enters \texttt{top}). The specific port pair is determined by the flow direction and node position.

\subsubsection{Multi-Edges}

\textbf{Multiple edges between the same pair of nodes are valid.} Each edge is rendered as a separate path, offset perpendicular to the primary path direction to avoid visual overlap. Offset distance is \texttt{theme.flow.multiEdgeGap} (default: 6px). Edges are offset in canonical order (list position).

\subsubsection{Disconnected / Orphan Nodes}

\textbf{Nodes with no edges are valid.} Orphan nodes (zero in-degree and zero out-degree) are placed in a designated ``orphan region''---below the main flow (for LTR) or to the right (for TTB)---sorted by canonical order. They participate in group membership normally.

\textbf{Disconnected components} (multiple disjoint subgraphs) are each laid out independently using the appropriate algorithm, then arranged sequentially along the secondary axis (vertical for LTR, horizontal for TTB).

\subsubsection{Missing Target/Source ID}

\textbf{Validation error.} If an edge references a node \texttt{id} that does not exist in the \texttt{nodes} list, the document fails validation with a diagnostic: \texttt{"Edge at index N references unknown node id 'X' in field 'source|target'"}. The layout engine does not process invalid documents.

\subsubsection{Empty Flow}

\textbf{Valid but degenerate.} A flow with zero nodes and zero edges produces a Scene with only the canvas and metadata (no drawing primitives). This is analogous to the timeline grammar's empty-activities case. A lint warning is emitted: \texttt{"Flow contains no nodes; output will be empty."}.

\subsubsection{Duplicate Node IDs}

\textbf{Validation error.} Two nodes with the same \texttt{id} fail validation: \texttt{"Duplicate node id 'X' at indices M and N."}.

\subsubsection{Edge Referencing Same Source and Target (Self-Loop)}

Covered above (\S self-loops). This is valid, not an error.

% ============================================================================
\subsection{Worked Example: RAG Pipeline}
\label{sec:flow-worked-example}
% ============================================================================

This example is drawn from the corpus (rag-vectorless-explainer, Section~\ref{sec:corpus-taxonomy}): a simplified Retrieval-Augmented Generation pipeline.

\subsubsection{Flow Domain IR}

\begin{lstlisting}[language=yaml,caption={RAG pipeline --- Flow Domain IR}]
version: "1.0"
metadata:
  title: "RAG Retrieval Pipeline"
  theme: dark-flow
flow:
  direction: left-to-right
  nodes:
    - id: query
      label: "User Query"
      shape: stadium
      icon: message-circle
      status: active
    - id: embed
      label: "Embed Query"
      shape: rounded-rect
      icon: cpu
    - id: retrieve
      label: "Vector Search"
      shape: rounded-rect
      icon: database
    - id: rerank
      label: "Re-rank"
      shape: rounded-rect
      icon: filter
    - id: generate
      label: "Generate Answer"
      shape: rounded-rect
      icon: sparkles
      status: success
  edges:
    - source: query
      target: embed
      style: solid
    - source: embed
      target: retrieve
      style: solid
      animated: true
    - source: retrieve
      target: rerank
      style: dashed
      label: "top-k"
    - source: rerank
      target: generate
      style: solid
      animated: true
    - source: generate
      target: query
      style: dotted
      label: "follow-up"
      directedness: directed
  groups:
    - id: retrieval
      label: "Retrieval Stage"
      nodes: [embed, retrieve, rerank]
      style: lane
\end{lstlisting}

\subsubsection{Lowering to Scene IR (Prose Description)}

The layout engine processes this document as follows:

\begin{enumerate}
  \item \textbf{Topology detection.} The edge \texttt{generate $\rightarrow$ query} creates a cycle. The engine identifies this as a feedback arc and reverses it for layering purposes.

  \item \textbf{Layer assignment.} Network simplex assigns five layers (left-to-right): L0=\texttt{query}, L1=\texttt{embed}, L2=\texttt{retrieve}, L3=\texttt{rerank}, L4=\texttt{generate}. The reversed back-edge \texttt{generate $\rightarrow$ query} spans 4 layers.

  \item \textbf{Crossing minimization.} With a single node per layer, no crossings exist. The barycenter pass completes trivially.

  \item \textbf{Coordinate assignment.} Nodes are placed at equal horizontal intervals. The ``Retrieval Stage'' group's bounding box spans L1--L3.

  \item \textbf{Scene IR emission.} The resulting Scene contains:
  \begin{itemize}
    \item 5 \texttt{Group} primitives (one per node), each containing a \texttt{Rect} (stadium/rounded-rect) + \texttt{Text} (label) + \texttt{Image} (icon).
    \item 5 \texttt{Path} primitives (edges). The \texttt{embed$\rightarrow$retrieve} and \texttt{rerank$\rightarrow$generate} paths carry \texttt{FlowingDashes} animation hints. The \texttt{generate$\rightarrow$query} back-edge is routed as a curved path below the main flow.
    \item 5 small \texttt{Path} primitives (arrowheads on directed edges).
    \item 2 \texttt{Text} primitives for edge labels (``top-k'' and ``follow-up'').
    \item 1 \texttt{Rect} (group background for ``Retrieval Stage'') + 1 \texttt{Text} (group label).
    \item Canvas: dimensions computed from content bounding box + padding.
    \item Meta: \texttt{scene\_hash} computed from the serialized element list.
  \end{itemize}
\end{enumerate}

\paragraph{Determinism verification.}
Running this IR through the layout engine twice produces an identical \texttt{scene\_hash}. The cycle-handling, layer assignment, and coordinate assignment are all pure functions of the input and canonical ordering.

% ============================================================================
\subsection{Determinism \& Opt-In Discipline}
\label{sec:flow-determinism}
% ============================================================================

The flow grammar respects the project's sacred determinism contract (Section~\ref{sec:determinism}):

\begin{description}
  \item[Byte-identical output.] The same Flow IR + same engine version + same theme = identical Scene IR (and therefore identical SVG). No exceptions.
  
  \item[No-change defaults.] All optional fields (\texttt{shape}, \texttt{status}, \texttt{style}, \texttt{animated}, \texttt{directedness}, ports, groups) have explicit defaults. Omitting optional fields produces the same output as specifying the default value explicitly.
  
  \item[New tokens are opt-in.] Future additions to the Flow IR (e.g., new shape values, new edge styles) must have no-change defaults: existing documents that do not use new features must produce byte-identical output after an engine upgrade.
  
  \item[Canonical ordering.] List position in \texttt{nodes} and \texttt{edges} is the tie-breaking order for all layout decisions. Reordering nodes in the IR may change layout (this is intentional---order is semantic).
  
  \item[Animation does not break determinism.] The \texttt{animated} flag adds a declarative \texttt{FlowingDashes} hint to the Scene. The SVG markup is byte-identical regardless of whether the animation is ``playing'' --- animation is CSS/SMIL, not frame-by-frame.
\end{description}

% ============================================================================
\subsection{Open Questions (Deferred to Mark/Barbara)}
\label{sec:flow-open-questions}
% ============================================================================

The following design points are flagged for specialist review before the schema is finalized:

\begin{description}
  \item[Mark (IR schema detail):]
  \begin{itemize}
    \item Exact JSON Schema for the Flow Domain IR (enum values, string patterns, min/max constraints).
    \item Whether \texttt{FlowEdge} should carry an explicit \texttt{id} field for referencing in composition or annotation contexts.
    \item Port model: is the 5-value enum (\texttt{auto, top, right, bottom, left}) sufficient, or should named custom ports be supported (useful for nodes with multiple outgoing paths)?
    \item Validation rule set: exhaustive list of semantic validation rules beyond schema conformance.
  \end{itemize}
  
  \item[Barbara (rendering semantics):]
  \begin{itemize}
    \item Self-loop routing: exact curve parameters (radius, port pair selection heuristic).
    \item Back-edge rendering: curve style (cubic B\'ezier, stepped, or arc) and visual offset distance.
    \item Multi-edge offset geometry: perpendicular offset vs.\ curvature-based separation.
    \item Group visual style details: border radius, padding, header positioning rules.
    \item Edge label collision avoidance with nodes/other labels: post-layout adjustment rules.
  \end{itemize}
\end{description}

\paragraph{Kernel change flag.}
No kernel (Scene IR) changes are required for the flow grammar as specified. If nested groups or custom port models later prove necessary, those would require kernel-level changes needing Barbara and Mark sign-off before adoption.
          % Node-link family: Flow (built)
\section{Sequence Grammar: Domain IR}
\label{sec:sequence-grammar}

The Sequence Grammar is the third grammar implemented in the diagram compiler and the \textbf{first de-risked new grammar}: its layout is deterministic-by-construction---no graph auto-layout algorithm is required. Where the Flow Grammar (Section~\ref{sec:flow-grammar}) requires Sugiyama layered layout~\cite{sugiyama1981} with cycle removal, crossing minimization, and coordinate assignment (Section~\ref{sec:layout-engines}), the Sequence Grammar's placement is fully determined by two ordered lists: participants declared left-to-right, and messages declared top-to-bottom. This makes it the lowest-risk grammar to add after Timeline and Flow.

\paragraph{Implementation Status.}
Fully implemented in \texttt{packages/core/src/grammars/sequence/} with participant/message/activation/fragment IR compilation. Supports all fragment kinds (alt with multi-guard, loop, opt, par), self-messages with curved paths, and stereotype rendering (actor, boundary, control, entity, database). SequenceTheme with light and ByteByteGo variants.

% ============================================================================
\subsection{Scope \& Intent}
\label{sec:sequence-scope}
% ============================================================================

\subsubsection{What a Sequence Diagram IS}

A \textbf{sequence diagram} in this grammar is a \emph{time-ordered interaction diagram} representing:

\begin{itemize}
  \item \textbf{Participants}: Named entities (actors, objects, services, boundaries) declared in a fixed horizontal order.
  \item \textbf{Messages}: Ordered communications between participants---synchronous calls, asynchronous signals, and return/reply messages---flowing downward along vertical lifelines.
  \item \textbf{Lifelines}: Vertical dashed lines descending from each participant, representing the entity's existence over time.
  \item \textbf{Activations} (optional): Thin rectangles on lifelines indicating processing spans.
  \item \textbf{Fragments} (optional): Labeled rectangles spanning a range of messages, representing control-flow semantics (loop, alt, opt, par).
\end{itemize}

The visual lineage is UML 2.x Sequence Diagrams~\cite{uml25} and ITU-T Message Sequence Charts~\cite{itu-msc}. The grammar captures the essential structural semantics---participant ordering plus message ordering---without the full UML behavioral metamodel.

\subsubsection{What a Sequence Diagram is NOT}

\begin{description}
  \item[Not a free node-link flow.] Flows (Section~\ref{sec:flow-grammar}) have nodes at arbitrary graph positions connected by routed edges. Sequence diagrams have a rigid two-axis structure: participants span x, messages span y. There is no node-link topology to lay out.
  \item[Not a timeline.] The timeline grammar's entities have temporal semantics (\texttt{start}, \texttt{end}, \texttt{track}) on a calibrated time axis. Sequence diagrams use \emph{ordinal} vertical position (message 1, message 2, \ldots)---not absolute time.
  \item[Not a general graph.] No force-directed, Sugiyama, or tree-layout algorithm is involved. Placement is tabular.
\end{description}

\subsubsection{Modeling Boundary}

The Sequence Grammar covers diagrams where:
\begin{enumerate}
  \item Entities are \textbf{participants} arranged in a fixed horizontal order.
  \item Interactions are \textbf{messages} between participants, with a total order defining vertical position.
  \item Time is \textbf{ordinal} (row index), not metric.
\end{enumerate}

Diagrams requiring metric time (Gantt-like durations) belong to the Timeline grammar. Diagrams requiring arbitrary node placement belong to Flow or Graph grammars.

\subsubsection{Why De-Risked}

\textbf{Deterministic layout is trivial.} Given $n$ participants and $m$ messages:
\begin{itemize}
  \item Participant $i$ is placed at $x_i$ determined solely by declared order and measured label widths (no optimization).
  \item Message $j$ is placed at $y_j = \text{headerHeight} + j \times \text{rowHeight}$ (no crossing minimization, no coordinate assignment).
  \item No iterative algorithm, no RNG, no convergence criterion.
\end{itemize}

Contrast with the Flow grammar's reliance on four-phase Sugiyama layout (Section~\ref{sec:flow-layout-contract}) where even tie-breaking heuristics must be carefully pinned for determinism. Here, placement is a pure arithmetic function of declaration order---the ``hard problem'' of \S\ref{sec:layout-engines} simply does not apply.

% ============================================================================
\subsection{Sequence Domain IR}
\label{sec:sequence-domain-ir}
% ============================================================================

The Sequence Domain IR is the small, grammar-specific representation that compiles \emph{down} to the shared Scene IR (Section~\ref{sec:scene-ir}). It is \textbf{not} a ``god-IR''---it contains only the constructs necessary to describe participant-message interaction diagrams.

\subsubsection{Document Structure}

A Sequence IR document is a single YAML/JSON object with the following root structure:

\begin{lstlisting}[language=yaml,caption={Sequence IR root document structure}]
version: "1.0"
metadata:
  title: string             # required
  subtitle: string          # optional; default null
  theme: string             # optional; default "default-sequence"
sequence:
  participants: [...]       # required; list<Participant>; min 1
  messages: [...]           # required; list<Message>; may be empty
  activations: [...]        # optional; list<Activation>; default []
  fragments: [...]          # optional; list<Fragment>; default []
\end{lstlisting}

\subsubsection{Root Fields}

\begin{table}[h]
\centering
\small
\begin{tabular}{llccp{4.8cm}}
\toprule
\textbf{Field} & \textbf{Type} & \textbf{Req} & \textbf{Default} & \textbf{Description} \\
\midrule
\texttt{version} & \texttt{string} & Yes & --- & Spec version (semver, e.g., \texttt{"1.0"}) \\
\texttt{metadata} & \texttt{SeqMetadata} & Yes & --- & Document-level metadata \\
\texttt{sequence} & \texttt{SeqDefinition} & Yes & --- & The sequence definition \\
\bottomrule
\end{tabular}
\caption{Sequence IR root fields}
\label{tab:seq-root-fields}
\end{table}

\subsubsection{Participant}
\label{sec:seq-participant}

A participant is a named entity with a vertical lifeline.

\begin{table}[h]
\centering
\small
\begin{tabular}{llccp{4.5cm}}
\toprule
\textbf{Field} & \textbf{Type} & \textbf{Req} & \textbf{Default} & \textbf{Description} \\
\midrule
\texttt{id} & \texttt{id} & Yes & --- & Document-unique identifier (kebab-case) \\
\texttt{label} & \texttt{string} & Yes & --- & Display text in the header box \\
\texttt{kind} & \texttt{enum\{actor, object, boundary, control, entity, database\}} & No & \texttt{object} & Participant stereotype (affects header icon/shape) \\
\texttt{description} & \texttt{string?} & No & \texttt{null} & Tooltip/secondary text \\
\bottomrule
\end{tabular}
\caption{Participant fields}
\label{tab:seq-participant-fields}
\end{table}

\paragraph{Identity.}
Participant \texttt{id} is the primary key. It must be unique across all participants. Messages reference participants by this \texttt{id}. Two participants with the same \texttt{id} are a validation error.

\paragraph{Declared order.}
Participants are placed left-to-right in the order they appear in the \texttt{participants} list. This order is the \textbf{canonical horizontal order}---the sole determinant of x-position. Reordering participants in the IR changes the diagram (this is intentional; order is semantic).

\paragraph{Kind semantics.}
The \texttt{kind} field controls the visual stereotype of the participant header:
\begin{description}
  \item[\texttt{actor}] Stick-figure icon above the label (UML actor).
  \item[\texttt{object}] Rectangular box with label (default, most common).
  \item[\texttt{boundary}] Box with a vertical bar on the left edge.
  \item[\texttt{control}] Box with a circular arrow icon.
  \item[\texttt{entity}] Box with a horizontal underline.
  \item[\texttt{database}] Cylinder (DB drum) icon above the label.
\end{description}

The \texttt{kind} does NOT affect layout---only the header's visual representation. All participants occupy the same column width regardless of kind.

\subsubsection{Message}
\label{sec:seq-message}

A message is a directed communication between two participants (or a self-message within one participant).

\begin{table}[h]
\centering
\small
\begin{tabular}{llccp{4.0cm}}
\toprule
\textbf{Field} & \textbf{Type} & \textbf{Req} & \textbf{Default} & \textbf{Description} \\
\midrule
\texttt{from} & \texttt{ref<Participant>} & Yes & --- & Source participant id \\
\texttt{to} & \texttt{ref<Participant>} & Yes & --- & Target participant id \\
\texttt{label} & \texttt{string} & Yes & --- & Message annotation text \\
\texttt{order} & \texttt{integer} & Yes & --- & Explicit sequence index ($\ge 0$); determines y-position \\
\texttt{kind} & \texttt{enum\{sync, async, reply\}} & No & \texttt{sync} & Message kind (controls arrowhead style) \\
\bottomrule
\end{tabular}
\caption{Message fields}
\label{tab:seq-message-fields}
\end{table}

\paragraph{Identity.}
Messages are identified by their \texttt{order} field---the explicit sequence index. The \texttt{order} is the sole determinant of vertical position. Messages with lower \texttt{order} values appear higher in the diagram.

\paragraph{Total order contract.}
The \texttt{order} field imposes a \textbf{total order} on messages. This is the deterministic backbone of the layout:
\begin{itemize}
  \item Messages are sorted by \texttt{order} (ascending) before layout.
  \item Message at rank $k$ (0-indexed after sorting) is placed at $y_k = \text{headerHeight} + k \times \text{rowHeight}$.
  \item If two messages share the same \texttt{order} value, the tie is broken by list position in the \texttt{messages} array (stable sort). See edge cases (\S\ref{sec:seq-edge-cases}).
\end{itemize}

\paragraph{Self-message.}
When \texttt{from == to}, the message is a self-message: a loopback arrow on a single lifeline. It occupies one message row like any other message, but is rendered as a small rightward loop returning to the same lifeline (see \S\ref{sec:seq-lowering}).

\paragraph{Message kind semantics.}
\begin{description}
  \item[\texttt{sync}] Synchronous call. Rendered as a solid line with a \textbf{filled} triangular arrowhead at the target.
  \item[\texttt{async}] Asynchronous signal. Rendered as a solid line with an \textbf{open} (line-only) arrowhead at the target.
  \item[\texttt{reply}] Return/response message. Rendered as a \textbf{dashed} line with an open arrowhead at the target.
\end{description}

\subsubsection{Activation}
\label{sec:seq-activation}

An activation represents a processing span on a participant's lifeline---the period during which the participant is actively handling a request.

\begin{table}[h]
\centering
\small
\begin{tabular}{llccp{4.5cm}}
\toprule
\textbf{Field} & \textbf{Type} & \textbf{Req} & \textbf{Default} & \textbf{Description} \\
\midrule
\texttt{participant} & \texttt{ref<Participant>} & Yes & --- & Lifeline on which the activation appears \\
\texttt{from\_order} & \texttt{integer} & Yes & --- & Message order at which activation begins \\
\texttt{to\_order} & \texttt{integer} & Yes & --- & Message order at which activation ends \\
\bottomrule
\end{tabular}
\caption{Activation fields}
\label{tab:seq-activation-fields}
\end{table}

\paragraph{Semantics.}
An activation bar is drawn as a thin filled rectangle on the participant's lifeline, spanning from the y-position of the message with \texttt{order == from\_order} to the y-position of the message with \texttt{order == to\_order}. The bar is centered on the lifeline.

\paragraph{Validation.}
\texttt{from\_order} must be $\le$ \texttt{to\_order}. Both values must reference \texttt{order} values that exist in the \texttt{messages} list (not necessarily on messages involving this participant---activations can span any message range). A zero-height activation (\texttt{from\_order == to\_order}) is valid: it renders as a minimal-height bar at that message row.

\subsubsection{Fragment}
\label{sec:seq-fragment}

A fragment represents a combined fragment (UML 2.x interaction fragment)---a labeled region spanning a range of messages that carries control-flow semantics.

\begin{table}[h]
\centering
\small
\begin{tabular}{llccp{4.0cm}}
\toprule
\textbf{Field} & \textbf{Type} & \textbf{Req} & \textbf{Default} & \textbf{Description} \\
\midrule
\texttt{kind} & \texttt{enum\{loop, alt, opt, par, critical, break\}} & Yes & --- & Fragment operator \\
\texttt{label} & \texttt{string} & Yes & --- & Guard condition or description text \\
\texttt{from\_order} & \texttt{integer} & Yes & --- & First message order in the span \\
\texttt{to\_order} & \texttt{integer} & Yes & --- & Last message order in the span \\
\texttt{participants} & \texttt{list<ref<Participant>>?} & No & all & Subset of participants the fragment spans horizontally \\
\bottomrule
\end{tabular}
\caption{Fragment fields}
\label{tab:seq-fragment-fields}
\end{table}

\paragraph{Semantics.}
A fragment is rendered as a labeled rounded rectangle spanning:
\begin{itemize}
  \item Vertically: from the y-position of \texttt{from\_order} to the y-position of \texttt{to\_order} (with padding).
  \item Horizontally: from the leftmost to the rightmost participant in \texttt{participants} (or all participants if omitted), with padding.
\end{itemize}

The \texttt{kind} keyword is rendered in a small tab in the upper-left corner of the rectangle (e.g., ``\texttt{loop}'', ``\texttt{alt}'').

\paragraph{Nesting.}
Fragments may nest: a fragment with \texttt{from\_order}=2, \texttt{to\_order}=5 may contain another with \texttt{from\_order}=3, \texttt{to\_order}=4. Nested fragments are rendered inside their parent with visual inset. Overlapping but non-nested fragments (partial overlap) are a validation error.

% ============================================================================
\subsection{Layout Contract (Deterministic-by-Construction)}
\label{sec:sequence-layout-contract}
% ============================================================================

The Sequence Grammar's layout is \textbf{fully determined by two ordered lists}: participants (horizontal) and messages (vertical). No optimization, heuristic, or iterative algorithm is involved.

\subsubsection{Participant Placement (x-axis)}

Given $n$ participants $p_1, p_2, \ldots, p_n$ in declared order:

\begin{enumerate}
  \item Compute the header width for each participant: $w_i = \text{measureText}(p_i.\text{label}) + 2 \times \text{headerPadX}$.
  \item Apply a minimum column width: $\text{colW}_i = \max(w_i, \text{minColWidth})$.
  \item Place participant centers at cumulative positions:
    \begin{align*}
      x_1 &= \text{marginLeft} + \text{colW}_1 / 2 \\
      x_i &= x_{i-1} + \text{colW}_{i-1}/2 + \text{colGap} + \text{colW}_i/2 \quad (i > 1)
    \end{align*}
\end{enumerate}

\paragraph{Determinism.}
Placement is a pure function of the participant list and theme geometry tokens (\texttt{headerPadX}, \texttt{minColWidth}, \texttt{colGap}, \texttt{marginLeft}). Same input $\rightarrow$ identical pixel positions. No force-directed layout, no crossing minimization, no Sugiyama---contrast with the Flow Grammar's four-phase Sugiyama pipeline (Section~\ref{sec:flow-layout-contract}).

\subsubsection{Lifelines (y-axis skeleton)}

Each participant $p_i$ emits a vertical dashed line from the bottom of its header box to the bottom of the diagram:
$$
  \text{lifeline}_i: (x_i, y_{\text{headerBottom}}) \longrightarrow (x_i, y_{\text{diagramBottom}})
$$

The lifeline is purely decorative; it does not participate in layout computation.

\subsubsection{Message Placement (y-axis)}

Given $m$ messages sorted by \texttt{order} (ties broken by list position), message at rank $k$ (0-indexed) is placed at:
$$
  y_k = y_{\text{headerBottom}} + \text{firstMsgGap} + k \times \text{rowHeight}
$$

Each message is rendered as a horizontal (or angled) arrow from $x_{\text{from}}$ to $x_{\text{to}}$ at height $y_k$, with its label text positioned above the arrow.

\paragraph{Self-messages.}
A self-message ($\text{from} = \text{to}$) is rendered as a small rightward loop: an arrow that exits the lifeline to the right, descends by a fixed \texttt{selfMsgHeight}, and returns to the same lifeline. It still occupies exactly one message row at $y_k$.

\subsubsection{Activation Bars}

An activation bar for participant $p_i$ spanning orders $[a, b]$ is placed as a thin filled rectangle:
$$
  \text{rect}: (x_i - \text{barHalfW},\; y_{\text{rank}(a)}) \longrightarrow (x_i + \text{barHalfW},\; y_{\text{rank}(b)})
$$

where $\text{rank}(a)$ is the y-position of the message with that order value.

\subsubsection{Fragment Rectangles}

A fragment spanning orders $[a, b]$ over participants $[p_l, p_r]$ is placed as a rounded rectangle:
\begin{align*}
  x_{\text{left}}  &= x_l - \text{colW}_l/2 - \text{fragPadX} \\
  x_{\text{right}} &= x_r + \text{colW}_r/2 + \text{fragPadX} \\
  y_{\text{top}}   &= y_{\text{rank}(a)} - \text{fragPadY} \\
  y_{\text{bot}}   &= y_{\text{rank}(b)} + \text{fragPadY}
\end{align*}

The operator keyword tab is positioned at $(x_{\text{left}} + \text{tabInset}, y_{\text{top}} + \text{tabInset})$.

\subsubsection{Contrast with Flow Grammar Layout}

\begin{center}
\small
\begin{tabular}{lp{5cm}p{5cm}}
\toprule
\textbf{Aspect} & \textbf{Flow Grammar} & \textbf{Sequence Grammar} \\
\midrule
Layout family & Sugiyama 4-phase (layer assignment, crossing min., coordinate assignment) & Arithmetic placement (declared order $\rightarrow$ position) \\
Determinism mechanism & Fixed sweep count, canonical tie-breaking & Trivially deterministic (no optimization) \\
Risk & Cycle removal heuristics, convergence control & None---placement is a closed-form function \\
New algorithms needed & Yes (Sugiyama, network simplex) & No \\
\bottomrule
\end{tabular}
\end{center}

% ============================================================================
\subsection{Lowering to the Scene IR}
\label{sec:seq-lowering}
% ============================================================================

The sequence layout engine emits Scene IR primitives (Section~\ref{sec:scene-ir}). \textbf{No new Scene IR primitives are required}---the existing kernel is sufficient.

\subsubsection{Participant Headers}

Each participant header emits:
\begin{enumerate}
  \item A \texttt{Rect} primitive for the header box (or a \texttt{Path} for non-rectangular kinds like \texttt{database} cylinder).
  \item A \texttt{Text} primitive for the participant label, centered within the box.
  \item For \texttt{kind: actor}: an additional small \texttt{Path} (stick-figure) above the label.
\end{enumerate}

\subsubsection{Lifelines}

Each lifeline emits a single \texttt{Line} (or \texttt{Path}) primitive:
\begin{Verbatim}[frame=single,fontsize=\small]
Line {
  x1: x_i, y1: headerBottom,
  x2: x_i, y2: diagramBottom,
  stroke_style: dashed,
  stroke_width: theme.sequence.lifelineWidth   // default 1px
}
\end{Verbatim}

\subsubsection{Messages}

Each message emits:
\begin{enumerate}
  \item A \texttt{Line} or \texttt{Path} primitive for the connector:
    \begin{description}
      \item[\texttt{sync}] Solid stroke.
      \item[\texttt{async}] Solid stroke.
      \item[\texttt{reply}] Dashed stroke.
    \end{description}
  \item An arrowhead \texttt{Path}:
    \begin{description}
      \item[\texttt{sync}] Filled triangular arrowhead (solid, closed polygon).
      \item[\texttt{async}] Open arrowhead (two lines forming a ``V'').
      \item[\texttt{reply}] Open arrowhead (two lines forming a ``V''), on a dashed connector.
    \end{description}
  \item A \texttt{Text} primitive for the message label, positioned above the connector midpoint.
\end{enumerate}

\paragraph{Self-message rendering.}
A self-message emits a \texttt{Path} with three segments: rightward horizontal, downward vertical (height = \texttt{selfMsgHeight}), leftward horizontal returning to the lifeline. The arrowhead is placed at the return point.

\textbf{Open question for Barbara:} The exact curve geometry for self-messages (rounded corners vs.\ sharp right angles vs.\ a smooth arc) is a rendering-detail decision. The spec defines the logical shape (exit right, descend, return left); Barbara owns the curve parameterization.

\subsubsection{Activation Bars}

Each activation emits a single \texttt{Rect} primitive:
\begin{Verbatim}[frame=single,fontsize=\small]
Rect {
  x: x_i - barHalfW,  y: y_rank(from_order),
  width: 2 * barHalfW,
  height: y_rank(to_order) - y_rank(from_order),
  fill: theme.sequence.activationFill,
  stroke: theme.sequence.activationStroke
}
\end{Verbatim}

\subsubsection{Fragments}

Each fragment emits:
\begin{enumerate}
  \item A \texttt{Rect} primitive (rounded corners, no fill or light fill per theme) for the frame.
  \item A small \texttt{Rect} + \texttt{Text} for the operator keyword tab in the upper-left corner.
  \item A \texttt{Text} primitive for the guard label, positioned after the keyword tab.
\end{enumerate}

\paragraph{No new kernel primitives needed.}
The entire sequence grammar maps to existing Scene IR primitives: \texttt{Rect}, \texttt{Line}/\texttt{Path}, \texttt{Text}, and \texttt{Group}. No kernel changes are required. If future extensions (e.g., destruction markers, creation messages with offset lifeline starts) prove necessary, those would require kernel review---\textbf{flagged for Barbara/Mark sign-off} rather than assumed.

\paragraph{Animation.}
Animation semantics are out of scope for this grammar specification. Per the project's additive animation model (Section~\ref{sec:scene-animation}), future animation hints (e.g., sequential message fade-in) can be layered onto the Scene IR without modifying the domain IR.

% ============================================================================
\subsection{Edge Cases \& Total Semantics}
\label{sec:seq-edge-cases}
% ============================================================================

Every possible input must have a defined meaning or produce a clear validation error.

\subsubsection{Self-Message}

\textbf{Valid.} A message where \texttt{from == to} is rendered as a loopback on the participant's lifeline (see \S\ref{sec:seq-lowering}). It occupies one row like any other message.

\subsubsection{Message Referencing an Undeclared Participant}

\textbf{Validation error.} If a message's \texttt{from} or \texttt{to} references a participant \texttt{id} not in the \texttt{participants} list, the document fails validation: \texttt{"Message at order N references unknown participant id 'X' in field 'from|to'."}

\subsubsection{Duplicate Order Indices}

\textbf{Valid (tie-broken).} If two or more messages share the same \texttt{order} value, they are placed at consecutive y-positions in the order they appear in the \texttt{messages} array (stable sort). A lint warning is emitted: \texttt{"Messages at indices M and N share order value K; list position used as tie-breaker."}

This ensures total determinism even with duplicate orders---but authors are encouraged to use unique order values for clarity.

\subsubsection{Participant with No Messages}

\textbf{Valid.} A participant that is never referenced by any message still appears in the diagram with its header box and lifeline. It occupies its declared column position. A lint warning is emitted: \texttt{"Participant 'X' has no messages; lifeline will be empty."}

\subsubsection{Empty Diagram}

\textbf{Partially valid.} A sequence with at least one participant but zero messages produces a diagram showing only participant headers and lifelines (no message arrows). A sequence with zero participants is a validation error: \texttt{"Sequence must declare at least one participant."}

\subsubsection{Nested Fragments}

\textbf{Valid.} Fragment $A$ may fully contain fragment $B$ ($A.\text{from\_order} \le B.\text{from\_order}$ and $B.\text{to\_order} \le A.\text{to\_order}$). Nested fragments are rendered with visual inset (each nesting level offsets inward by \texttt{fragNestInset} pixels). Rendering order: outermost fragments first (painter's algorithm), innermost on top.

\textbf{Partial overlap is invalid.} If fragment $A$ and fragment $B$ overlap but neither fully contains the other ($A.\text{from\_order} < B.\text{from\_order} < A.\text{to\_order} < B.\text{to\_order}$), validation fails: \texttt{"Fragments at indices M and N partially overlap; fragments must nest or be disjoint."}

\textbf{Open question for Barbara:} Maximum nesting depth for readable rendering. The spec imposes no hard limit; Barbara may recommend a soft limit (e.g., 3--4 levels) with a lint warning.

\subsubsection{Async Message with No Corresponding Reply}

\textbf{Valid.} An \texttt{async} message with no subsequent \texttt{reply} from the target back to the source is perfectly valid---it represents a fire-and-forget signal. No implicit reply is generated.

\subsubsection{Activation Referencing Non-Existent Order}

\textbf{Validation error.} If \texttt{from\_order} or \texttt{to\_order} does not match any message's \texttt{order} value, validation fails: \texttt{"Activation for participant 'X' references non-existent order value N."}

% ============================================================================
\subsection{Worked Example: Token-Based REST API Authentication}
\label{sec:seq-worked-example}
% ============================================================================

This example models a simple token-based REST API authentication flow: a client requests a token, the server returns it, then the client uses the token in a subsequent authenticated request.

\subsubsection{Sequence Domain IR}

\begin{lstlisting}[language=yaml,caption={REST API auth --- Sequence Domain IR}]
version: "1.0"
metadata:
  title: "Token-Based REST API Authentication"
  theme: default-sequence
sequence:
  participants:
    - id: client
      label: "Client"
      kind: actor
    - id: server
      label: "Auth Server"
      kind: object
  messages:
    - from: client
      to: server
      label: "POST /auth/token (credentials)"
      order: 1
      kind: sync
    - from: server
      to: client
      label: "200 OK { token: 'jwt...' }"
      order: 2
      kind: reply
    - from: client
      to: server
      label: "GET /api/data (Authorization: Bearer jwt...)"
      order: 3
      kind: sync
    - from: server
      to: client
      label: "200 OK { data: [...] }"
      order: 4
      kind: reply
  activations:
    - participant: server
      from_order: 1
      to_order: 2
    - participant: server
      from_order: 3
      to_order: 4
\end{lstlisting}

\subsubsection{Lowering to Scene IR (Detailed)}

The layout engine processes this document as follows:

\begin{enumerate}
  \item \textbf{Participant placement.} Two participants in declared order: Client at $x_1$, Auth Server at $x_2$. With \texttt{colGap}=80px and measured label widths, suppose $x_1 = 120$px, $x_2 = 320$px.

  \item \textbf{Message placement.} Four messages sorted by order (1, 2, 3, 4). With \texttt{rowHeight}=50px and \texttt{firstMsgGap}=30px:
  \begin{itemize}
    \item Message 1 at $y = 30 + 0 \times 50 = 30$px below header.
    \item Message 2 at $y = 30 + 1 \times 50 = 80$px below header.
    \item Message 3 at $y = 30 + 2 \times 50 = 130$px below header.
    \item Message 4 at $y = 30 + 3 \times 50 = 180$px below header.
  \end{itemize}

  \item \textbf{Scene IR emission.} The resulting Scene contains:
  \begin{itemize}
    \item \textbf{2 header Groups}: each a \texttt{Rect} + \texttt{Text}. Client additionally has a stick-figure \texttt{Path} (kind=actor).
    \item \textbf{2 lifeline Lines}: dashed vertical lines from header bottom to diagram bottom, at $x_1$ and $x_2$.
    \item \textbf{4 message Lines}: horizontal arrows from $x_1$ to $x_2$ (messages 1, 3) and from $x_2$ to $x_1$ (messages 2, 4). Messages 2 and 4 have \texttt{stroke\_style: dashed} (reply kind).
    \item \textbf{4 arrowhead Paths}: filled triangles for messages 1, 3 (sync); open V-shapes for messages 2, 4 (reply).
    \item \textbf{4 label Texts}: positioned above each message arrow at the midpoint.
    \item \textbf{2 activation Rects}: thin filled rectangles on the server lifeline ($x_2$), spanning $[y_1, y_2]$ and $[y_3, y_4]$.
    \item \textbf{Canvas}: width computed from rightmost participant + margin; height from last message row + footer margin.
    \item \textbf{Meta}: \texttt{scene\_hash} computed from serialized element list.
  \end{itemize}
\end{enumerate}

\paragraph{Primitive count.}
Total Scene primitives: 2 Groups (headers) + 2 Lines (lifelines) + 4 Lines (messages) + 4 Paths (arrowheads) + 4 Texts (labels) + 2 Rects (activations) + canvas + meta = 18 drawing primitives. All are existing kernel types---no extensions required.

\paragraph{Determinism verification.}
Running this IR through the layout engine twice produces an identical \texttt{scene\_hash}. Participant placement is a function of declaration order and label widths; message placement is a function of order values. No randomness, no iteration, no heuristic.

% ============================================================================
\subsection{Determinism \& Opt-In Discipline}
\label{sec:seq-determinism}
% ============================================================================

The Sequence Grammar respects the project's sacred determinism contract (Section~\ref{sec:determinism}):

\begin{description}
  \item[Byte-identical output.] The same Sequence IR + same engine version + same theme = identical Scene IR (and therefore identical SVG). No exceptions.

  \item[No-change defaults.] All optional fields (\texttt{kind}, \texttt{description}, \texttt{activations}, \texttt{fragments}) have explicit defaults. Omitting optional fields produces the same output as specifying the default value explicitly.

  \item[New tokens are opt-in.] Future additions to the Sequence IR (e.g., new participant kinds, destruction markers, creation messages) must have no-change defaults: existing documents that do not use new features must produce byte-identical output after an engine upgrade.

  \item[Canonical ordering.] Declaration order in \texttt{participants} determines x-position; \texttt{order} field in messages determines y-position; list position is the final tie-breaker. These are the \emph{only} inputs to layout.

  \item[No iterative algorithms.] Unlike the Flow Grammar's Sugiyama layout, the Sequence Grammar uses no iterative or heuristic algorithms. Placement is a closed-form arithmetic function. This is the strongest possible determinism guarantee.
\end{description}

% ============================================================================
\subsection{Open Questions (Deferred to Mark/Barbara)}
\label{sec:seq-open-questions}
% ============================================================================

The following design points are flagged for specialist review before the schema is finalized:

\begin{description}
  \item[Mark (IR schema detail):]
  \begin{itemize}
    \item Exact JSON Schema for the Sequence Domain IR (enum values, string patterns, min/max constraints on \texttt{order}).
    \item Whether \texttt{Message} should carry an optional explicit \texttt{id} field for cross-referencing in composition or annotation contexts.
    \item Validation rule completeness: the exhaustive set of semantic checks beyond schema conformance (e.g., fragment overlap detection algorithm, activation range validation).
    \item Whether \texttt{order} should be required or whether implicit ordering (list position as order) is acceptable as an alternative mode.
  \end{itemize}

  \item[Barbara (rendering semantics):]
  \begin{itemize}
    \item Self-message curve geometry: rounded corners, smooth arc, or sharp right-angle bends. Radius parameters.
    \item Fragment nesting depth: recommended maximum for readability; visual inset scaling per level.
    \item Participant header rendering for non-\texttt{object} kinds: exact icon geometries for actor (stick-figure), boundary, control, entity, database stereotypes.
    \item Arrowhead sizing: scale relative to stroke width or fixed pixel size.
    \item Activation bar width: fixed or proportional to lifeline stroke width.
  \end{itemize}
\end{description}

\paragraph{Kernel change flag.}
No kernel (Scene IR) changes are required for the Sequence Grammar as specified. If future extensions (destruction markers with ``X'' terminators, participant creation mid-diagram with offset lifeline starts, or ``found/lost'' messages from/to diagram boundaries) prove necessary, those would require kernel-level review needing Barbara and Mark sign-off before adoption.
      % UML/software family: Sequence (built)
\section{Tree Grammar: Domain IR}
\label{sec:tree-grammar}

The Tree Grammar is the fourth grammar in the diagram compiler and the \textbf{second de-risked grammar}: its layout is the Buchheim--J{\"u}nger--Leipert tidy-tree algorithm~\cite{buchheim2002}---deterministic, $O(n)$, and fully solved. Where the Flow Grammar (Section~\ref{sec:flow-grammar}) requires four-phase Sugiyama layered layout with cycle removal and crossing minimization (Section~\ref{sec:layout-engines}), and the Sequence Grammar (Section~\ref{sec:sequence-grammar}) sidesteps layout entirely via arithmetic placement, the Tree Grammar occupies a middle ground: it \emph{does} require a real layout algorithm, but one that is provably deterministic and linear-time for its constrained input class (rooted trees).

% ============================================================================
\subsection{Scope \& Intent}
\label{sec:tree-scope}
% ============================================================================

\subsubsection{What a Tree Diagram IS}

A \textbf{tree diagram} in this grammar is a \emph{node-link rendering of a rooted tree}: a connected, acyclic graph with exactly one root node and where every non-root node has exactly one parent. Trees represent containment, decomposition, and hierarchical relationships:

\begin{itemize}
  \item \textbf{Document structure}: document $\rightarrow$ chapters $\rightarrow$ sections $\rightarrow$ subsections.
  \item \textbf{Organizational hierarchies}: CEO $\rightarrow$ VPs $\rightarrow$ directors $\rightarrow$ teams.
  \item \textbf{Decision trees}: root question $\rightarrow$ branches at each decision point.
  \item \textbf{File systems}: root directory $\rightarrow$ subdirectories $\rightarrow$ files.
  \item \textbf{Taxonomies and classification}: kingdom $\rightarrow$ phylum $\rightarrow$ class $\rightarrow$ order.
\end{itemize}

The visual form is a \emph{tidy tree}: nodes placed at discrete depth levels with sibling groups spread horizontally, connected by parent-to-child edges. This is the canonical ``org chart'' or ``tree view'' layout.

\subsubsection{What a Tree Diagram is NOT}

\begin{description}
  \item[Not a general directed graph.] Flows (Section~\ref{sec:flow-grammar}) support cycles, multi-parent nodes, and arbitrary DAG topologies. Trees require exactly one parent per non-root node---no shared children, no cycles.
  \item[Not a sequence diagram.] Sequence diagrams (Section~\ref{sec:sequence-grammar}) have a rigid participant$\times$message grid structure. Trees have variable branching and depth.
  \item[Not a timeline.] Timelines have a calibrated time axis. Tree depth levels are ordinal (level 0, 1, 2\ldots), not temporal.
  \item[Not a mind map.] Mind maps are typically radial, free-form, and allow cross-links. This grammar produces strictly hierarchical tidy-tree layouts with no cross-edges.
\end{description}

\subsubsection{Modeling Boundary}

The Tree Grammar covers diagrams where:
\begin{enumerate}
  \item Entities are \textbf{nodes} in a containment/decomposition hierarchy.
  \item Relationships are strictly \textbf{parent$\rightarrow$child} (one parent per node, one root).
  \item Layout admits a natural \textbf{level-based reading} (root at top, children below).
\end{enumerate}

Diagrams with shared children (a node having two parents) belong to the Flow grammar (DAG). Diagrams with no clear hierarchy belong to the Graph grammar's force-directed or hub-spoke layouts.

\subsubsection{Why De-Risked}

\textbf{Tidy-tree layout is deterministic and $O(n)$.} The Buchheim--J{\"u}nger--Leipert algorithm~\cite{buchheim2002} (correcting Walker~\cite{walker1990}, itself extending Reingold--Tilford~\cite{reingold1981}) computes compact, non-overlapping node positions in a single linear-time pass:

\begin{itemize}
  \item No NP-hard subproblems (contrast: crossing minimization in Sugiyama is NP-complete~\cite{garey1983}).
  \item No iterative convergence (contrast: force-directed methods).
  \item No randomized initialization.
  \item Fully deterministic given a fixed sibling order.
  \item Proven $O(n)$ time and space.
\end{itemize}

This makes Tree the lowest-risk graph-layout grammar after Sequence (which requires no layout algorithm at all). The ``hard problem'' of Section~\ref{sec:layout-engines} is bypassed by the structural constraint: trees are a strict, easier subclass of general graphs.

% ============================================================================
\subsection{Tree Domain IR}
\label{sec:tree-domain-ir}
% ============================================================================

The Tree Domain IR is the small, grammar-specific representation that compiles \emph{down} to the shared Scene IR (Section~\ref{sec:scene-ir}). It is \textbf{not} a ``god-IR''---it contains only the constructs necessary to describe rooted-tree hierarchies.

\subsubsection{Document Structure}

A Tree IR document is a single YAML/JSON object with the following root structure:

\begin{lstlisting}[language=yaml,caption={Tree IR root document structure}]
version: "1.0"
metadata:
  title: string             # required
  subtitle: string          # optional; default null
  theme: string             # optional; default "default-tree"
tree:
  root: TreeNode            # required; the single root node
\end{lstlisting}

\subsubsection{Root Fields}

\begin{table}[h]
\centering
\small
\begin{tabular}{llccp{4.8cm}}
\toprule
\textbf{Field} & \textbf{Type} & \textbf{Req} & \textbf{Default} & \textbf{Description} \\
\midrule
\texttt{version} & \texttt{string} & Yes & --- & Spec version (semver, e.g., \texttt{"1.0"}) \\
\texttt{metadata} & \texttt{TreeMetadata} & Yes & --- & Document-level metadata \\
\texttt{tree} & \texttt{TreeDefinition} & Yes & --- & The tree definition (contains the root) \\
\bottomrule
\end{tabular}
\caption{Tree IR root fields}
\label{tab:tree-root-fields}
\end{table}

\subsubsection{TreeNode}
\label{sec:tree-node}

A \texttt{TreeNode} is the sole structural construct. Trees are recursive: each node may contain child nodes.

\begin{table}[h]
\centering
\small
\begin{tabular}{llccp{4.0cm}}
\toprule
\textbf{Field} & \textbf{Type} & \textbf{Req} & \textbf{Default} & \textbf{Description} \\
\midrule
\texttt{id} & \texttt{id} & Yes & --- & Document-unique identifier (kebab-case) \\
\texttt{label} & \texttt{string} & Yes & --- & Display text for the node \\
\texttt{children} & \texttt{list<TreeNode>} & No & \texttt{[]} & Ordered child nodes \\
\texttt{kind} & \texttt{string?} & No & \texttt{null} & Semantic kind hint (e.g., \texttt{"chapter"}, \texttt{"person"}, \texttt{"decision"}) \\
\texttt{icon} & \texttt{string?} & No & \texttt{null} & Icon registry key \\
\texttt{collapsed} & \texttt{bool} & No & \texttt{false} & Hint: render subtree as collapsed (elided) \\
\texttt{description} & \texttt{string?} & No & \texttt{null} & Tooltip or secondary text \\
\bottomrule
\end{tabular}
\caption{TreeNode fields}
\label{tab:tree-node-fields}
\end{table}

\paragraph{Canonical representation: children-list.}
The \textbf{canonical} tree representation uses an embedded \texttt{children} list on each node. This is chosen over a flat-list-with-parent-ref representation for three reasons:
\begin{enumerate}
  \item \textbf{Natural authoring}: trees are written top-down; nesting children inside parents mirrors the mental model.
  \item \textbf{Structural guarantee}: a well-formed nested document is automatically a valid tree (no cycles, no orphans, exactly one root). Validation is trivial---only duplicate \texttt{id} checks are needed.
  \item \textbf{Sibling order is explicit}: the list order of \texttt{children} defines left-to-right placement. No separate ordering field is required.
\end{enumerate}

\paragraph{Alternative: flat list with parent refs.}
An alternative flat representation (a list of nodes where each non-root carries a \texttt{parent} ref) is a valid \emph{input convenience}. If supported, the validator \textbf{normalizes} it to the canonical children-list form before layout:
\begin{enumerate}
  \item Verify exactly one node has no \texttt{parent} (the root).
  \item Verify all \texttt{parent} refs resolve to existing node ids.
  \item Verify no cycles (topological sort succeeds).
  \item Group nodes by parent; sibling order = list order among nodes sharing the same parent.
\end{enumerate}

\textbf{Open question for Mark}: Whether to support the flat parent-ref form as an alternative input mode, or require the canonical children-list form exclusively. The spec defines the canonical form; the schema may optionally accept both.

\paragraph{Identity.}
Node \texttt{id} is the primary key across the entire document. It must be globally unique (across all levels of nesting). Two nodes with the same \texttt{id}---regardless of depth---are a validation error.

\paragraph{Sibling order.}
Children are placed left-to-right in the order they appear in the \texttt{children} list. This order is the \textbf{canonical horizontal order} and the sole determinant of sibling arrangement. Reordering children in the IR changes the diagram (this is intentional; order is semantic).

\paragraph{Kind, icon, collapsed: semantic hints only.}
Per the grammar$=$semantics / theme$=$style principle (Section~\ref{sec:grammar-concept}):
\begin{itemize}
  \item \texttt{kind} is a semantic tag (``chapter'', ``person'', ``folder'') that a \texttt{TreeTheme} \emph{may} use to select visual treatment (e.g., different fill colors or border styles per kind). The grammar assigns no visual meaning to kind values---that is a theme concern.
  \item \texttt{icon} references the shared icon registry. Placement and sizing are theme-driven.
  \item \texttt{collapsed} is a rendering hint: when \texttt{true}, the subtree rooted at this node is visually elided (e.g., replaced by an expander indicator). The node itself is still rendered; its children are hidden. This is a \emph{hint}---static backends render it as a visual marker; interactive backends may allow expansion.
\end{itemize}

None of these fields affect layout geometry of \emph{other} nodes. A collapsed subtree occupies zero layout space for its hidden children (the collapsed node itself retains its measured size).

% ============================================================================
\subsection{Layout Contract (Deterministic Tidy-Tree)}
\label{sec:tree-layout-contract}
% ============================================================================

The Tree Grammar's layout engine converts the Domain IR into Scene IR coordinates using the layered tidy-tree algorithm family. The layout is \textbf{fully deterministic}: identical IR in $\rightarrow$ byte-identical Scene out.

\subsubsection{Algorithm: Buchheim--J{\"u}nger--Leipert (2002)}

The layout follows the Reingold--Tilford~\cite{reingold1981} / Walker~\cite{walker1990} / Buchheim--J{\"u}nger--Leipert~\cite{buchheim2002} lineage:

\begin{description}
  \item[Depth assignment.] Each node is assigned to a \textbf{depth level} $d$: the root at $d=0$, its children at $d=1$, and so on. Depth determines the y-coordinate: $y = d \times \text{levelHeight}$.
  
  \item[Horizontal placement.] Siblings are spread horizontally such that:
  \begin{itemize}
    \item No two node bounding boxes overlap (minimum horizontal separation = \texttt{siblingGap}).
    \item Subtrees of adjacent siblings do not overlap (the algorithm walks contours of left and right subtrees, advancing a ``thread'' pointer to detect conflicts).
    \item A parent is centered above its children (the midpoint of the leftmost and rightmost child x-positions).
  \end{itemize}
  
  \item[Complexity.] $O(n)$ time and $O(n)$ space, where $n$ is the number of nodes~\cite{buchheim2002}. The thread-pointer technique avoids Walker's~\cite{walker1990} quadratic worst case.
  
  \item[Determinism.] The algorithm is a pure function of the tree structure and sibling order. No randomness, no iteration count, no convergence criterion. Same input $\rightarrow$ same output, always.
\end{description}

\subsubsection{Orientation}

The primary (default) orientation is \textbf{top-down}: root at the top of the canvas, children below. The depth axis maps to y (increasing downward); the sibling axis maps to x.

\textbf{Left-to-right} orientation (root at left, children to the right) is supported as a theme/layout option---it is the same algorithm with axes transposed. Additional orientations (bottom-up, right-to-left) follow the same transposition pattern.

\textbf{Note}: Orientation is a layout/theme option, not a grammar-level semantic field. The IR describes structure; orientation is presentation. A \texttt{TreeTheme} declares which orientation to use.

\subsubsection{Node Sizing}

Each node's bounding box is computed as:
$$
  w_i = \text{measureText}(\text{node}_i.\text{label}) + 2 \times \text{nodePadX} + \text{iconWidth}_i
$$
$$
  h_i = \text{lineHeight} + 2 \times \text{nodePadY}
$$

All sizing parameters (\texttt{nodePadX}, \texttt{nodePadY}, \texttt{iconWidth}, \texttt{lineHeight}) are \textbf{theme tokens}. The grammar does not prescribe pixel values---only the formula. This ensures that node size is a pure function of label content + theme geometry.

\subsubsection{Spacing Parameters (Theme Tokens)}

\begin{table}[h]
\centering
\small
\begin{tabular}{lp{7cm}}
\toprule
\textbf{Token} & \textbf{Meaning} \\
\midrule
\texttt{levelHeight} & Vertical distance between depth levels (center-to-center) \\
\texttt{siblingGap} & Minimum horizontal gap between adjacent sibling bounding boxes \\
\texttt{subtreeGap} & Minimum horizontal gap between adjacent subtrees (may differ from \texttt{siblingGap}) \\
\texttt{nodePadX}, \texttt{nodePadY} & Internal padding within a node box \\
\texttt{marginTop}, \texttt{marginLeft} & Canvas margins \\
\bottomrule
\end{tabular}
\caption{Tree layout theme tokens (geometry)}
\label{tab:tree-layout-tokens}
\end{table}

\subsubsection{Contrast with Flow Grammar's Sugiyama}

\begin{center}
\small
\begin{tabular}{lp{5cm}p{5cm}}
\toprule
\textbf{Aspect} & \textbf{Flow Grammar (Sugiyama)} & \textbf{Tree Grammar (Tidy-Tree)} \\
\midrule
Input class & General directed graphs (DAGs, cyclic with FAS removal) & Rooted trees only (one parent per node, no cycles) \\
Phases & 4 (cycle removal, layer assignment, crossing min., coordinate assignment) & 1 (recursive bottom-up + top-down adjustment) \\
Crossing minimization & NP-complete; barycenter heuristic, 24 sweeps & Not applicable---trees have no edge crossings \\
Complexity & $O(|V|^2)$ for crossing minimization & $O(n)$ \\
Determinism mechanism & Fixed sweep count + canonical tie-breaking & Inherently deterministic (pure recursion over fixed structure) \\
\bottomrule
\end{tabular}
\end{center}

Trees are a strict subclass of DAGs. The tidy-tree algorithm exploits this constraint to eliminate all hard subproblems that Sugiyama must address.

% ============================================================================
\subsection{Lowering to the Scene IR}
\label{sec:tree-lowering}
% ============================================================================

The tree layout engine emits Scene IR primitives (Section~\ref{sec:scene-ir}). \textbf{No new Scene IR primitives are required}---the existing kernel is sufficient.

\subsubsection{Node Lowering}

Each \texttt{TreeNode} emits a \texttt{Group} containing:

\begin{enumerate}
  \item A \texttt{Rect} primitive for the node box. Shape is theme-driven: default is rounded rectangle (\texttt{corner\_radius: theme.tree.nodeRadius}). The theme may select different shapes per \texttt{kind} (e.g., circles for leaf nodes, diamonds for decision nodes).
  \item A \texttt{Text} primitive for the node label, centered within the box.
  \item Optionally, an \texttt{Image} primitive for the icon (positioned left of label or above, per theme).
\end{enumerate}

\paragraph{Collapsed nodes.}
When \texttt{collapsed: true}, the node is rendered normally but:
\begin{itemize}
  \item Its children are \emph{not} emitted to the Scene IR (zero primitives for the subtree).
  \item An expander indicator is added: a small \texttt{Path} or \texttt{Text} glyph (e.g., ``$+$'' or ``\ldots'') positioned below or beside the node, per theme.
\end{itemize}

\subsubsection{Edge Lowering}

Each parent$\rightarrow$child relationship emits a connector:

\begin{description}
  \item[Orthogonal elbow (default).] A \texttt{Path} with three segments: vertical from parent bottom-center down to the midpoint between levels, horizontal to the child's x-position, then vertical down to the child's top-center. This produces the classic ``org chart'' connector style.
  
  \item[Straight line.] A single \texttt{Line} from parent bottom-center to child top-center. Simpler but less readable for wide trees.
  
  \item[Curved.] A cubic B\'ezier \texttt{Path} from parent to child. Produces smooth, flowing connectors.
\end{description}

The choice among these styles is a \textbf{theme option} (\texttt{theme.tree.edgeStyle: elbow | straight | curved}). The grammar does not prescribe edge rendering---only that a visual connector exists for each parent-child pair.

\paragraph{Edge directionality.}
Edges are rendered \textbf{without arrowheads} by default (the top-down reading direction is implicit). Arrowheads are a theme option for cases where directionality needs emphasis.

\subsubsection{Scene IR Mapping Summary}

\begin{center}
\small
\begin{tabular}{lp{8cm}}
\toprule
\textbf{Tree Construct} & \textbf{Scene IR Primitive(s)} \\
\midrule
Node box & \texttt{Rect} (rounded, theme-styled) \\
Node label & \texttt{Text} (centered in box) \\
Node icon & \texttt{Image} (optional, positioned per theme) \\
Parent$\rightarrow$child edge & \texttt{Path} (elbow/straight/curved, per theme) \\
Collapsed indicator & \texttt{Path} or \texttt{Text} (expander glyph) \\
Canvas & Computed from tree bounding box + margins \\
\bottomrule
\end{tabular}
\end{center}

\paragraph{No new kernel primitives needed.}
The entire Tree Grammar maps to existing Scene IR primitives: \texttt{Rect}, \texttt{Text}, \texttt{Path}, \texttt{Line}, \texttt{Image}, and \texttt{Group}. No kernel changes are required. If future extensions (e.g., cross-links between tree nodes for showing relationships outside the hierarchy) prove necessary, those would require kernel-level review---\textbf{flagged for Barbara/Mark sign-off} rather than assumed.

\paragraph{Styling deferred to TreeTheme.}
Consistent with the \texttt{SequenceTheme} and \texttt{FlowTheme} precedent, all visual styling (colors, font sizes, border widths, edge styles, node shapes per kind, spacing) is controlled by a \texttt{TreeTheme} type. The grammar specifies structure; the theme specifies presentation.

% ============================================================================
\subsection{Edge Cases \& Total Semantics}
\label{sec:tree-edge-cases}
% ============================================================================

Every possible input must have a defined meaning or produce a clear validation error.

\subsubsection{Multiple Roots (Forest)}

\textbf{Rejected.} The Tree Grammar requires exactly one root. A document with multiple roots (e.g., a flat list with two nodes having no parent) is a validation error: \texttt{"Tree must have exactly one root node; found N root candidates."}

\textbf{Rationale}: A forest (multiple disjoint trees) is a fundamentally different layout problem---it requires deciding how to arrange separate trees relative to each other. This is a composition concern. Authors who need a forest should use the Composition layer (Section~\ref{sec:composition}) with multiple Tree panels.

\subsubsection{Cycles}

\textbf{Rejected.} A cycle in the tree structure (node $A$ is a descendant of node $B$ which is a descendant of $A$) is a validation error: \texttt{"Cycle detected involving node 'X'; tree structures must be acyclic."}

In the canonical children-list representation, cycles are structurally impossible (a node cannot appear as its own ancestor in a nested JSON/YAML document). Cycles can only arise in the flat parent-ref alternative input form, where validation must perform a topological sort to detect them.

\subsubsection{Orphan Nodes (Unresolved Parent)}

\textbf{Rejected.} In the flat parent-ref input form, if a node references a parent \texttt{id} that does not exist, validation fails: \texttt{"Node 'X' references non-existent parent 'Y'."}

In the canonical children-list form, orphans are structurally impossible---every node is either the root or a child of another node.

\subsubsection{Single-Node Tree}

\textbf{Valid.} A tree with only a root node (no children) is valid. It produces a single node box centered on the canvas. This is the minimal valid tree.

\subsubsection{Very Wide Trees (Many Siblings)}

\textbf{Valid.} A node with many children (e.g., 50+ siblings) produces a wide diagram. The layout engine does not impose a hard limit on sibling count. The tidy-tree algorithm handles this correctly---it simply places all siblings in a row. A lint warning may be emitted for extreme cases: \texttt{"Node 'X' has N children; consider restructuring for readability."} (Threshold is theme-configurable, e.g., $>20$.)

\subsubsection{Very Deep Trees}

\textbf{Valid.} A tree with many levels (e.g., 20+ depth) produces a tall diagram. No hard limit is imposed. The $O(n)$ algorithm handles deep trees efficiently. As with wide trees, a lint warning may be emitted for extreme depth.

\subsubsection{Duplicate IDs}

\textbf{Validation error.} Two nodes anywhere in the tree with the same \texttt{id} fail validation: \texttt{"Duplicate node id 'X' at paths 'P1' and 'P2'."} (Paths indicate the nesting location for debugging.)

\subsubsection{Empty Children List}

\textbf{Valid.} A node with \texttt{children: []} is a leaf node. This is the common case for terminal nodes and is not degenerate.

% ============================================================================
\subsection{Worked Example: Document Structure Hierarchy}
\label{sec:tree-worked-example}
% ============================================================================

This example models a document hierarchy: a root ``Document'' node with three chapters, each containing a few sections. This is drawn from the corpus's document-structure diagram kind.

\subsubsection{Tree Domain IR}

\begin{lstlisting}[language=yaml,caption={Document hierarchy --- Tree Domain IR}]
version: "1.0"
metadata:
  title: "Document Structure"
  theme: default-tree
tree:
  root:
    id: doc
    label: "Document"
    kind: root
    children:
      - id: ch1
        label: "Chapter 1: Introduction"
        kind: chapter
        children:
          - id: s1-1
            label: "1.1 Background"
            kind: section
          - id: s1-2
            label: "1.2 Motivation"
            kind: section
      - id: ch2
        label: "Chapter 2: Methods"
        kind: chapter
        children:
          - id: s2-1
            label: "2.1 Data Collection"
            kind: section
          - id: s2-2
            label: "2.2 Analysis"
            kind: section
          - id: s2-3
            label: "2.3 Validation"
            kind: section
      - id: ch3
        label: "Chapter 3: Results"
        kind: chapter
        children:
          - id: s3-1
            label: "3.1 Findings"
            kind: section
\end{lstlisting}

\subsubsection{Layout Computation}

The layout engine processes this document as follows:

\begin{enumerate}
  \item \textbf{Depth assignment.} Three levels: \texttt{doc} at $d=0$; \texttt{ch1}, \texttt{ch2}, \texttt{ch3} at $d=1$; all sections at $d=2$.
  
  \item \textbf{Node sizing.} Each node's width is measured from its label text + padding. Suppose (with default theme): \texttt{doc}$=180$px, chapters$\approx220$px, sections$\approx200$px wide; all$=36$px tall.
  
  \item \textbf{Tidy-tree placement (Buchheim et al.).}
  \begin{itemize}
    \item Leaf nodes (sections) are placed with \texttt{siblingGap}$=20$px between them.
    \item Each chapter is centered above its children.
    \item Subtree gaps ensure \texttt{ch1}'s rightmost section and \texttt{ch2}'s leftmost section maintain at least \texttt{subtreeGap}$=40$px separation.
    \item The root \texttt{doc} is centered above the three chapter nodes.
  \end{itemize}
  
  \item \textbf{Vertical positions.} With \texttt{levelHeight}$=100$px: $y_{\text{doc}}=50$, $y_{\text{chapters}}=150$, $y_{\text{sections}}=250$.
\end{enumerate}

\subsubsection{Lowering to Scene IR}

The resulting Scene contains:

\begin{itemize}
  \item \textbf{10 node Groups}: each containing a \texttt{Rect} (rounded, theme-styled) + \texttt{Text} (label). The root node may have a distinct fill color (theme resolves \texttt{kind: root} to a highlight fill).
  
  \item \textbf{9 edge Paths}: one for each parent$\rightarrow$child relationship. With the default \texttt{elbow} edge style, each path has three segments (vertical-horizontal-vertical).
  
  \item \textbf{Canvas}: width computed from the leftmost section's left edge to the rightmost section's right edge, plus margins. Height computed from top margin to depth-2 bottom edge, plus footer margin.
  
  \item \textbf{Meta}: \texttt{scene\_hash} computed from the serialized element list.
\end{itemize}

\paragraph{Primitive count.}
Total Scene primitives: 10 Groups (containing 10 Rects + 10 Texts) + 9 Paths (edges) + canvas + meta = 29 drawing primitives. All are existing kernel types---no extensions required.

\paragraph{Determinism verification.}
Running this IR through the Buchheim layout engine twice produces an identical \texttt{scene\_hash}. Node positions are a pure function of the tree structure, sibling order, and theme geometry tokens. No randomness, no iteration, no heuristic.

% ============================================================================
\subsection{Determinism \& Theme/Semantics Split}
\label{sec:tree-determinism}
% ============================================================================

The Tree Grammar respects the project's sacred determinism contract (Section~\ref{sec:determinism}) and the grammar$=$semantics / theme$=$style principle:

\begin{description}
  \item[Byte-identical output.] The same Tree IR + same engine version + same theme = identical Scene IR (and therefore identical SVG). No exceptions.
  
  \item[No-change defaults.] All optional fields (\texttt{kind}, \texttt{icon}, \texttt{collapsed}, \texttt{description}, \texttt{children}) have explicit defaults. Omitting optional fields produces the same output as specifying the default value explicitly.
  
  \item[New tokens are opt-in.] Future additions to the Tree IR (e.g., new semantic hint fields, annotation support) must have no-change defaults: existing documents that do not use new features must produce byte-identical output after an engine upgrade.
  
  \item[Canonical ordering.] Sibling order in \texttt{children} lists is the sole determinant of horizontal arrangement. Reordering children changes the diagram (intentional---order is semantic).
  
  \item[Deterministic algorithm.] Buchheim--J{\"u}nger--Leipert is a pure function over a tree structure. No iteration count, no sweep count, no RNG. This provides the strongest possible determinism guarantee short of trivial arithmetic (Sequence Grammar).
  
  \item[Grammar is semantic-only.] The Tree Domain IR carries \emph{structure} (parent-child hierarchy, sibling order) and \emph{semantic hints} (\texttt{kind}, \texttt{icon}, \texttt{collapsed}). It does \textbf{not} carry:
  \begin{itemize}
    \item Colors, fills, or strokes (theme tokens).
    \item Font sizes or families (theme tokens).
    \item Node shapes (theme resolves \texttt{kind} to a shape).
    \item Edge styles (theme option: elbow/straight/curved).
    \item Spacing values (theme geometry tokens).
    \item Orientation (theme layout option).
  \end{itemize}
  All visual presentation is controlled by a \texttt{TreeTheme} type, following the \texttt{SequenceTheme} and \texttt{FlowTheme} precedent.
\end{description}

% ============================================================================
\subsection{Open Questions (Deferred to Mark/Barbara)}
\label{sec:tree-open-questions}
% ============================================================================

The following design points are flagged for specialist review before the schema is finalized:

\begin{description}
  \item[Mark (IR schema detail):]
  \begin{itemize}
    \item \textbf{Canonical form}: Should the schema accept \emph{only} the children-list form, or also support a flat parent-ref alternative input? If both, specify the normalization algorithm and which takes precedence on conflict.
    \item \textbf{Forest handling}: Confirm rejection of multiple roots. If future requirements demand forest support, specify it as a separate grammar or as a composition-layer concern.
    \item \textbf{Validation invariants}: Exhaustive list of semantic checks beyond schema conformance (duplicate id detection across nested levels, maximum depth/breadth lint thresholds).
    \item \textbf{Kind enum}: Should \texttt{kind} be a free string or a closed enum? Free string is more extensible (themes can map arbitrary kinds); closed enum enables schema-level validation.
    \item \textbf{Node \texttt{id} format}: Confirm kebab-case requirement. Nested ids may benefit from path-based namespacing (e.g., \texttt{ch1/s1-1}) vs.\ global flat namespace.
  \end{itemize}
  
  \item[Barbara (rendering semantics):]
  \begin{itemize}
    \item \textbf{Edge routing style}: Default elbow geometry---exact midpoint calculation, corner radius for rounded elbows, minimum segment length before defaulting to straight.
    \item \textbf{Collapsed-node handling}: Visual indicator design (``$+$'' glyph, ellipsis, count badge showing hidden children). Interactive vs.\ static rendering semantics.
    \item \textbf{TreeTheme token surface}: Complete list of theme tokens (colors, typography, geometry, edge style) needed for the initial default theme and at least one showcase theme.
    \item \textbf{Kind$\rightarrow$shape mapping}: Default visual mappings for common kind values (e.g., \texttt{person}$\rightarrow$circle, \texttt{folder}$\rightarrow$rounded-rect with icon). How many built-in kind mappings should the default theme provide?
    \item \textbf{Label overflow}: Behavior when label text exceeds available node width---truncate with ellipsis, wrap to multiple lines, or auto-expand node width (current spec assumes auto-expand via \texttt{measureText}).
  \end{itemize}
\end{description}

\paragraph{Kernel change flag.}
No kernel (Scene IR) changes are required for the Tree Grammar as specified. If future extensions (cross-links between tree nodes, animated expand/collapse transitions, or multi-root forest layout) prove necessary, those would require kernel-level review needing Barbara and Mark sign-off before adoption.
          % Tree/hierarchy family: Tree (built)
\section{The Chart Family: A Grammar-of-Graphics Layer}
\label{sec:chart-family}

The chart family is the one genuinely new architectural component required by the
Mermaid-superset strategy. Unlike node-link, timeline, and tree families (which map
\emph{structural} entities to \emph{positioned shapes}), charts map \textbf{quantitative
data} to \textbf{geometric marks} through scales and coordinate systems.

This section specifies the chart grammar drawing on the Grammar of
Graphics~\cite{grammarofgraphics2005} and Vega-Lite~\cite{vegalite2017} traditions,
adapted to our two-IR-layer architecture.

\subsection{Design Principles for Charts}

\begin{enumerate}
  \item \textbf{Data $\to$ geometry, not data $\to$ pixels.} The Domain IR specifies data
        and encoding; the layout engine computes geometry; the theme styles it.
  \item \textbf{Mermaid-compatible syntax.} Each chart type's Mermaid syntax parses
        faithfully to our Chart IR.
  \item \textbf{Deterministic layout.} All scale computations, tick placement, and mark
        positioning are closed-form (no iteration, no randomness).
  \item \textbf{Theme-driven aesthetics.} Colors, fonts, gridlines, and labels are
        theme-controlled. The IR carries data and encoding, never style.
\end{enumerate}

\subsection{Chart Domain IR}

\begin{lstlisting}[language=javascript,caption={Chart Domain IR shape (TypeScript)}]
interface ChartDocument {
  version: string;
  chartType: 'pie' | 'xy' | 'quadrant' | 'radar';
  title?: string;
  data: ChartData;
  encoding: ChartEncoding;
  config?: ChartConfig;
}

interface ChartData {
  values: Array<Record<string, number | string>>;
}

interface ChartEncoding {
  // Type-specific encoding channels
  x?: FieldEncoding;      // xy charts
  y?: FieldEncoding;      // xy charts
  color?: FieldEncoding;  // categorical color
  size?: FieldEncoding;   // mark size (radar axis values)
  theta?: FieldEncoding;  // pie angle
  label?: FieldEncoding;  // text labels
}

interface FieldEncoding {
  field: string;          // key into data values
  type: 'quantitative' | 'nominal' | 'ordinal';
  scale?: ScaleConfig;
}
\end{lstlisting}

\subsection{Chart Types}

\subsubsection{Pie Chart}

Maps categorical values to angular sectors. Mermaid syntax:

\begin{Verbatim}[fontsize=\small]
pie title Pet Adoption
    "Dogs" : 386
    "Cats" : 85
    "Rats" : 15
\end{Verbatim}

\noindent Domain IR encoding: \texttt{theta} channel maps value to sector angle;
\texttt{color} channel maps label to fill color (theme palette).

\subsubsection{XY Chart (Bar/Line)}

Cartesian charts with categorical or quantitative axes. Mermaid syntax:

\begin{Verbatim}[fontsize=\small]
xychart-beta
    title "Monthly Revenue"
    x-axis [Jan, Feb, Mar, Apr, May]
    y-axis "Revenue (USD)" 0 --> 5000
    bar [1200, 2400, 3100, 2800, 4200]
    line [1000, 2200, 2900, 2600, 4000]
\end{Verbatim}

\noindent Domain IR: \texttt{x} encodes categories or values; \texttt{y} encodes quantitative
measure; mark type (bar/line) stored in encoding config.

\subsubsection{Quadrant Chart}

Four-quadrant scatter plot with labeled axes and positioned items. Mermaid syntax:

\begin{Verbatim}[fontsize=\small]
quadrantChart
    title Reach and engagement
    x-axis Low Reach --> High Reach
    y-axis Low Engagement --> High Engagement
    quadrant-1 We should expand
    quadrant-2 Need to promote
    Campaign A: [0.3, 0.6]
    Campaign B: [0.45, 0.23]
\end{Verbatim}

\noindent Domain IR: \texttt{x} and \texttt{y} encode position on [0,1] scales;
quadrant labels are metadata.

\subsubsection{Radar Chart}

Multi-axis radial chart. Mermaid syntax (proposed):

\begin{Verbatim}[fontsize=\small]
radar
    title Skills Assessment
    axes: [Speed, Reliability, Comfort, Safety, Efficiency]
    "Model A": [4, 3, 2, 5, 4]
    "Model B": [2, 5, 4, 3, 3]
\end{Verbatim}

\noindent Domain IR: axes define radial dimensions; each series maps values to axis positions.

\subsection{Layout Engine}

The chart layout engine is shared across all four chart types:

\begin{enumerate}
  \item \textbf{Scale computation.} From data domain (min/max or categories) to pixel range,
        using linear, band, or radial scales. All scales are closed-form.
  \item \textbf{Axis/grid generation.} Tick positions, labels, and gridlines computed
        deterministically from scale + theme config.
  \item \textbf{Mark placement.} Data points mapped through scales to Scene IR primitives:
        Rect (bars), Path (lines, pie sectors, radar polygons), Circle (scatter points).
  \item \textbf{Label placement.} Deterministic label positioning with collision avoidance
        (priority-based, not iterative).
\end{enumerate}

\subsection{Relation to Vega-Lite}

The chart grammar draws on Vega-Lite's decomposition (data, transform, mark, encoding,
scale, axis, legend) but is deliberately simpler:

\begin{itemize}
  \item No data transforms (filter, aggregate, window)---data arrives pre-aggregated.
  \item No interaction or selection.
  \item No composition operators (layer/concat/facet)---composition is handled by the
        poster layer (Part~\ref{part:composition}).
  \item Subset of mark types: rect, line, arc, point, text.
\end{itemize}

This simplification keeps the chart Domain IR small enough for reliable agent generation
while covering the four Mermaid chart types completely.

\subsection{Implementation Status}

\textbf{Not yet built.} The chart family is the Tier-2 deliverable. It requires:
\begin{itemize}
  \item New Chart Domain IR and JSON Schema
  \item Scale computation library (linear, band, radial)
  \item Chart layout engine compiling to Scene IR
  \item Mermaid parsers for pie, xychart, quadrant, radar syntax
  \item Chart-specific theme extensions (axis styling, gridlines, palettes)
\end{itemize}

The existing Scene IR and backends (SVG, PNG) require no changes---charts compile to the
same Rect/Path/Text/Circle primitives used by all other families.
          % Charts family: grammar-of-graphics layer (NEW)

% ============================================================================
% PART V — AESTHETICS & THEMING
% ============================================================================

\part{Aesthetics \& Theming}
\label{part:aesthetics}

\section{Theme Architecture}
\label{sec:themes}

A theme is \emph{one coherent design system}, defined once at a general, component-agnostic
level. It is \emph{not} an attribute owned separately by each diagram type; it is an
authority that every component conforms to. Each component---flow, sequence, timeline, tree,
class diagram, ER diagram, and so on---provides a \emph{specific implementation} of that
general system: it knows how to realise the theme's tokens in its own geometric terms. The
theme knows nothing about flowcharts or Gantt bars; the component knows nothing about brand
palettes or spacing scales. The theme sets the contract; the component fulfils it.

\begin{quote}
\textbf{Confirmed design principle (binding):} A theme must be a coherent system applicable
to many (ideally all) components. You cannot conceive of a theme as an attribute of each
component as a separate system. A theme is defined once in general terms; each component
provides its own specific implementation of that general definition.
\end{quote}

Themes also drive \textbf{geometry, layout, and routing}---not just colour. A theme defines
widths, fonts, weights, spacing rhythms, density levels, and shape language. The layout
engine consults the theme to make rendering, layout, and routing decisions. The old rule that
``a theme may not alter geometry'' is explicitly retired: it contradicts both this model and
the existing code, where the timeline's \texttt{spineSpacing} and \texttt{nodeWrap} tokens
already drive layout behaviour.

\paragraph{Implementation status.}
As of 2026-06-14 the compiler ships 16 grammar-specific \texttt{theme.ts} files (one per
component: \texttt{FlowTheme}, \texttt{SequenceTheme}, \texttt{ClassTheme},
\texttt{TreeTheme}, etc.) plus the timeline's \texttt{ResolvedTheme} (which covers 14 named
timeline themes). These are separate, incompatible vocabularies---the fragmentation this
section resolves. The migration path is to generalise \texttt{ResolvedTheme} upward into a
shared general contract, then adapt each per-component file to consume that contract.

% ──────────────────────────────────────────────────────────────────────────────
\subsection{Current Reality: Per-Component Fragmentation}
\label{sec:themes-fragmentation}

Today every grammar family carries an independent styling vocabulary:

\begin{itemize}
  \item \textbf{Timeline:} \texttt{packages/core/src/themes/types.ts} ---
        \texttt{ResolvedTheme} with \texttt{canvas}, \texttt{typography},
        \texttt{axis}, \texttt{track}, \texttt{milestone}, \texttt{serpentine},
        \texttt{roadmap}, and layout-affecting flags such as
        \texttt{spineSpacing} and \texttt{nodeWrap}. Fourteen named themes:
        \texttt{consulting}, \texttt{executive}, \texttt{product}, \texttt{release},
        \texttt{minimal}, \texttt{showcase}, \texttt{bytebytego}, \texttt{ai-timeline},
        and others.
  \item \textbf{Flow:} \texttt{FlowTheme}---orientation, edge routing style, node
        geometry, font sizes.
  \item \textbf{Sequence:} \texttt{SequenceTheme}---participant render mode (box vs.\
        card), row heights, activation bar widths, fragment styling.
  \item \textbf{Sixteen further grammars} (class, ER, state, tree, C4, chart,
        gitGraph, journey, sankey, kanban, requirement, block, packet, architecture)
        each with their own isolated token structs.
\end{itemize}

No named theme (``consulting'', ``executive'') spans more than one grammar today. There is no
shared vocabulary for palette roles, type scale, or spacing rhythm across components.
This is the problem the new model resolves.

% ──────────────────────────────────────────────────────────────────────────────
\subsection{The Themes $\times$ Components Matrix}
\label{sec:themes-matrix}

The new model is a two-dimensional system:

\begin{description}
  \item[Theme (the row)] A general design language. It defines palette by role, type
        scale, spacing rhythm, density, and shape language---nothing component-specific.
        Adding a new theme adds a new row.
  \item[Component (the column)] A diagram grammar. It knows how to realise the general
        theme in its own geometric terms. Adding a new component adds a new column.
  \item[Intersection ($T \times C$)] How theme $T$ renders component $C$. This is
        computed by $C$'s implementation consuming $T$'s general tokens; it is not
        stored as a separate artefact.
\end{description}

\begin{table}[h]
\centering
\small
\begin{tabularx}{\textwidth}{l|XXXXXX}
\toprule
\textbf{Theme $\downarrow$ / Component $\rightarrow$} &
  \textbf{Timeline} & \textbf{Flow} & \textbf{Sequence} &
  \textbf{Tree} & \textbf{Class} & \textbf{\ldots} \\
\midrule
Consulting  & \checkmark & \checkmark & \checkmark & \checkmark & \checkmark & \ldots \\
Executive   & \checkmark & \checkmark & \checkmark & \checkmark & \checkmark & \ldots \\
ByteByteGo  & \checkmark & \checkmark & \checkmark & \checkmark & \checkmark & \ldots \\
Minimal     & \checkmark & \checkmark & \checkmark & \checkmark & \checkmark & \ldots \\
\textit{New theme}  & \checkmark & \checkmark & \checkmark & \checkmark & \checkmark & \ldots \\
\bottomrule
\end{tabularx}
\caption{Themes $\times$ Components matrix. A theme is the row authority; a component is the
column implementer. A \checkmark\ means the component's implementation consumes that theme's
tokens. Add a theme once---every component renders it. Add a component once (implement the
contract)---it inherits every existing theme for free.}
\label{tab:themes-matrix}
\end{table}

This is the inverse of the fragmented model. In the fragmented model each component owns its
theme, so adding a new theme requires editing every component; adding a new component
requires writing a new, isolated theme. In the matrix model each direction is O(1) in
authoring work: a new theme is one new general-contract file; a new component is one new
implementation that automatically supports all existing themes.

% ──────────────────────────────────────────────────────────────────────────────
\subsection{The General Theme Contract}
\label{sec:themes-contract}

The general theme contract is the primary design artefact of the theming system. It is an
abstract, component-agnostic token vocabulary expressive enough to drive any component's
geometry, routing, and look---with zero component-specific fields. Every component's
implementation consumes this vocabulary; nothing outside it may be required for a correct
rendering~\cite{dtcg2024}.

\begin{quote}
\textbf{General-contract rule (binding):} The general theme vocabulary MUST be rich enough
that a component's layout engine can ask the theme for spacing, density, routing style, and
palette roles and receive a deterministic answer. No component-specific token may leak into
the general contract. Components may \emph{interpret} tokens in their own context; they may
not \emph{demand} tokens not in the contract.
\end{quote}

The vocabulary is organized into six domains; the specific field names are subject to
refinement but the domain coverage is fixed:

\subsubsection{Palette}
\label{sec:contract-palette}

The general contract carries \textbf{two distinct palette systems}. Structural
components consume the \emph{role palette}; data and chart components consume the
\emph{data palette}. Both systems are part of the general contract.

\paragraph{Role palette (semantic/structural).}
Structural roles are expressed by name, not raw hex. Raw hex values are an
implementation detail of a specific theme instance.

\begin{lstlisting}[language=yaml,caption={General theme contract: role palette (decided)}]
palette:
  # Surface roles
  surface:         "#FFFFFF"   # primary background
  surface-raised:  "#F8F8F8"  # cards, panels
  surface-overlay: "#0A0A0A"  # dark overlay / inverse background
  surface-panel:   "#F4F4F4"  # lane / panel / track-header background (distinct from surface)

  # Content roles
  ink:             "#111111"   # primary text
  ink-muted:       "#666666"   # secondary / de-emphasized text
  ink-inverse:     "#FFFFFF"   # text on dark backgrounds
  ink-panel:       "#333333"   # text/ink rendered on panel or lane backgrounds

  # Brand / accent
  accent:          "#1F497D"   # primary brand colour
  accent-muted:    "#6A92BF"   # lighter accent (inactive, planned)
  accent-strong:   "#0D2B4E"   # darker accent (active, in-progress)

  # Border
  border:          "#DDDDDD"   # default stroke / divider
  border-strong:   "#999999"   # emphasized boundary

  # Muted fill (background regions, disabled states)
  muted:           "#F0F0F0"
  muted-strong:    "#CCCCCC"

  # Semantic status roles (components map these to IR status values)
  status-success:  { fill: "#4CAF50", stroke: "#388E3C" }
  status-warning:  { fill: "#D97706", stroke: "#A35A04" }
  status-error:    { fill: "#DC2626", stroke: "#A81919" }
  status-info:     { fill: "#1F497D", stroke: "#0D2B4E" }

  # IR workflow states (super-set of status roles above)
  status-planned:    { fill: "#6A92BF", stroke: "#1F497D" }
  status-active:     { fill: "#1F497D", stroke: "#0D2B4E" }
  status-done:       { fill: "#888888", stroke: "#555555", opacity: 0.85 }
  status-cancelled:  { fill: "#CCCCCC", stroke: "#999999", opacity: 0.60 }
  status-uncertain:  { fill: "#CCCCCC", stroke: "#999999", opacity: 0.70 }

  # Optional fill pattern applied to workflow-state renders
  status-pattern: "solid"  # solid | hatch | dashed-border
\end{lstlisting}

\paragraph{Data palette (quantities and sequences).}
Chart and data components (pie, XY chart, kanban, gitGraph, radar, sankey) consume the
data palette. It provides three sub-systems covering the full range of quantitative
encodings~\cite{dtcg2024}.

\begin{lstlisting}[language=yaml,caption={General theme contract: data palette (decided)}]
data-palette:
  # Categorical sequence — ordered colours for pie slices, chart series,
  # kanban columns, branch colours. Use by index (0-based); themes may
  # supply any length >= 6.
  categorical:
    - "#1F497D"   # 0 — primary
    - "#6A92BF"   # 1
    - "#4CAF50"   # 2
    - "#D97706"   # 3
    - "#9C27B0"   # 4
    - "#00BCD4"   # 5

  # Sequential ramp — for ordered quantitative encodings (sankey
  # throughput, heat intensity, radar fill). Low → high.
  sequential:
    low:  "#EEF4FB"
    mid:  "#6A92BF"
    high: "#0D2B4E"

  # Diverging ramp — for values around a meaningful midpoint
  # (positive/negative, above/below target). Negative → zero → positive.
  diverging:
    negative: "#DC2626"
    zero:     "#F5F5F5"
    positive: "#4CAF50"
\end{lstlisting}

\subsubsection{Typography}
\label{sec:contract-typography}

Typography specifies family and fallback, a \emph{type scale} (not a list of fixed sizes),
and a weight set. Components consume the scale by step name, not by pixel value.

\begin{lstlisting}[language=yaml,caption={General theme contract: typography block (proposed shape)}]
typography:
  family:    "Helvetica Neue"
  fallback:  "Arial, Helvetica, sans-serif"
  # Embedded font files for deterministic metric computation.
  files:
    regular: "fonts/HelveticaNeue-Regular.woff2"
    bold:    "fonts/HelveticaNeue-Bold.woff2"

  # Type scale: step names → point sizes.
  # Components reference steps by name, not by pt value.
  scale:
    xs:   8    # pt — axis ticks, dense labels
    sm:   9    # pt — captions, secondary labels
    base: 11   # pt — body text, activity labels
    md:   13   # pt — subtitles, section headers
    lg:   18   # pt — document title
    xl:   24   # pt — hero / keynote headline

  # Weight palette
  weights:
    regular: 400
    medium:  500
    semibold: 600
    bold:    700
\end{lstlisting}

\subsubsection{Spacing / Rhythm Scale}
\label{sec:contract-spacing}

A spacing scale defines the rhythm unit and derived named steps. Components consult the
scale as \emph{advice}, not as a binding pixel-identical contract: each component maps
step names onto its own geometry as it sees fit, preserving layout judgement while sharing
a common rhythm. The scale is not normalised across components.

\begin{lstlisting}[language=yaml,caption={General theme contract: spacing block (advisory)}]
spacing:
  unit: 8           # px — base rhythm unit
  # Named steps (advisory; components map these to their own geometry)
  steps:
    xxs:  2   # 0.25u
    xs:   4   # 0.5u
    sm:   8   # 1u
    md:  16   # 2u
    lg:  24   # 3u
    xl:  32   # 4u
    xxl: 48   # 6u
\end{lstlisting}

\subsubsection{Density}
\label{sec:contract-density}

Density is a discrete named level that components use to select between compact and
comfortable configurations. Three levels are defined---no more, no less. The layout engine
consults density to choose row heights, inter-element gaps, and label truncation thresholds
(e.g.\ whether secondary labels are shown at all).

\begin{lstlisting}[language=yaml,caption={General theme contract: density (decided)}]
density: "normal"   # compact | normal | comfortable
\end{lstlisting}

A \texttt{compact} theme instructs every component to select its smallest safe row height,
tightest inter-element gap, and most aggressive label truncation. A \texttt{comfortable}
theme selects the most generous variants. Components map the three density levels to their
own geometry constants; density is not a continuous scalar.

\subsubsection{Shape Language}
\label{sec:contract-shape}

Shape language defines how geometric objects feel. Components use these tokens when
deciding corner radius, node padding, stroke scale, and connector style.

\begin{lstlisting}[language=yaml,caption={General theme contract: shape block (proposed shape)}]
shape:
  corner-radius:  0       # px — 0 = square; >0 = rounded
  node-padding:   8       # px — internal padding inside node shapes
  stroke-scale:   1.0     # multiplier on all stroke widths (0.5 = hairline; 2.0 = bold)
  connector-style: "elbow" # straight | elbow | curved | orthogonal
  marker-shape: "circle"  # circle | diamond | square — milestone / marker glyph
\end{lstlisting}

\subsubsection{Effects and Motion}
\label{sec:contract-effects}

Effects declare the fidelity tier and optional art-effect tokens. These gate which rendering
backend capability profile is required (see Section~\ref{sec:fidelity-tiers}).

\begin{lstlisting}[language=yaml,caption={General theme contract: effects block (proposed shape)}]
effects:
  fidelity: 1        # 0=Minimal | 1=Crisp | 2=Polished | 3=Showcase
  drop-shadow:  { enabled: false }
  glow:         { enabled: false }
  motion:       { enabled: false }   # SMIL animation opt-in (Section~14)
\end{lstlisting}

\subsubsection{The Timeline ResolvedTheme as the Generalisation Seed}
\label{sec:contract-seed}

The timeline's \texttt{ResolvedTheme}
(\texttt{packages/core/src/themes/types.ts}) is the closest existing artefact to the
general contract. It already covers canvas, typography (with weights), axis geometry, track
geometry, and layout-affecting flags
(\texttt{spineSpacing}, \texttt{nodeWrap}, \texttt{maxSubLanes}). It demonstrates that a
theme can---and must---drive layout decisions, not merely colour.

The migration plan is to \emph{generalise upward}: extract the conceptual domains (palette
roles, type scale, spacing rhythm, density, shape language) from \texttt{ResolvedTheme}
into the shared general contract, then recast the timeline-specific fields (axis height,
track header width, serpentine row spacing) as the timeline component's interpretation of
those general tokens~\cite{vegalite2017}.

% ──────────────────────────────────────────────────────────────────────────────
\subsection{Theme Drives Geometry, Layout, and Routing}
\label{sec:themes-drives-layout}

Themes are not a post-processing colour filter. The layout engine queries the theme
\emph{during} layout computation. Representative examples from the existing codebase:

\begin{itemize}
  \item \textbf{\texttt{spineSpacing: 'time'} vs.\ \texttt{'even'}} --- controls whether
        vertical-spine milestone nodes are placed proportionally to elapsed time or at
        uniform y-intervals. This changes node $y$-coordinates, not just colour.
  \item \textbf{\texttt{nodeWrap: 'over-under'}} --- routes the horizontal timeline spine
        around each on-axis node as a semicircular arc (alternating above / below),
        producing entirely different path geometry compared to the default straight line.
  \item \textbf{\texttt{density}} --- compact density reduces \texttt{row\_height} and
        \texttt{sub\_lane\_height}, shrinking the canvas and shifting all track and
        milestone $y$-coordinates.
  \item \textbf{\texttt{shape.corner-radius}} --- a flow diagram component may use this
        token to choose between rectangular and rounded node shapes, which affects label
        area computation and thus label wrap point.
  \item \textbf{\texttt{connector-style: 'orthogonal'}} --- instructs the flow / sequence
        layout engine to use orthogonal edge routing rather than diagonal or curved,
        changing every edge path.
\end{itemize}

\begin{quote}
\textbf{Layout-engine rule (binding):} A layout engine MUST consult the active theme for
any geometry decision that varies across themes. It MUST NOT hard-code values that the
general contract expresses. Theme-driven geometry is still \emph{closed-form}: the theme
supplies constants; the engine computes coordinates deterministically from those constants.
No iterative solver or random sampling is introduced.
\end{quote}

% ──────────────────────────────────────────────────────────────────────────────
\subsection{Reach: Theme, Layout, and New Components}
\label{sec:themes-reach}

Drawing the right boundary between theme, layout argument, and new grammar is the
architectural decision that keeps the system from proliferating components unnecessarily.
The rule is stated at the IR boundary:

\begin{description}
  \item[\emph{New component (grammar)}] Only when the diagram requires \emph{new semantic
        data} that the existing IR cannot carry. A new grammar is a last resort.
  \item[\emph{Layout argument}] When the same data can be positioned differently without
        needing new IR fields. A layout mode is a named argument passed at render time,
        not a new grammar.
  \item[\emph{Theme tokens}] Look-and-feel within a layout: colour, typography, spacing,
        shape, effects. A theme produces a visually distinct output from the same IR +
        same layout.
\end{description}

\paragraph{Worked example: the timeline as six diagrams.}

The timeline IR is one grammar. A single set of entries---each with a date, label, status,
and optional icon---can be rendered six ways that look and behave radically differently,
purely via \texttt{\{layout: $L$, theme: $T$\}}:

\begin{table}[h]
\centering
\small
\begin{tabularx}{\textwidth}{llX}
\toprule
\textbf{Layout argument} & \textbf{Theme} & \textbf{Visual result} \\
\midrule
\texttt{vertical-spine}  & consulting   & Spine-and-card with navy accent, tidy whitespace \\
\texttt{horizontal-arc}  & executive    & Arc path with serif labels, subtle quarter shading \\
\texttt{roadmap}         & product      & Horizontal Gantt-style bars, dense information \\
\texttt{serpentine}      & showcase     & Gradient boustrophedon path, premium glow effects \\
\texttt{gantt}           & release      & Traffic-light Gantt, monospace labels, weekly grid \\
\texttt{timeline-columns}& minimal      & Column-per-period grid, greyscale, LaTeX-safe \\
\bottomrule
\end{tabularx}
\caption{One timeline IR; six visual outputs via \{layout, theme\} arguments. No new grammar
is required for any of these. Same data, radically different look and behaviour.}
\label{tab:reach-example}
\end{table}

The rule of thumb: if the data the IR already carries is sufficient to drive the new visual,
use a layout argument and/or theme tokens. Only introduce a new grammar when the IR would
need to grow new fields to support it (e.g., swimlane headers, dependency arrows, resource
calendars---these are genuinely new semantic data, not layout variations).

% ──────────────────────────────────────────────────────────────────────────────
\subsection{Mermaid Configuration Surface}
\label{sec:themes-mermaid-surface}

All user-facing surface is Mermaid or Mermaid-like. Theme selection, token overrides,
layout selection, and density are expressed via Mermaid-native configuration mechanisms,
so a user can remain in standard Mermaid syntax while accessing the extended vocabulary.

\paragraph{Frontmatter (\texttt{---}~\ldots~\texttt{---}).}

\begin{lstlisting}[language=yaml,caption={Theme and layout selection via Mermaid YAML frontmatter}]
---
theme: executive          # named theme from the registry
layout: vertical-spine    # layout family for this diagram
density: compact          # compact | normal | comfortable
tokens:                   # inline token overrides (scoped to this diagram)
  palette.accent: "#8B0000"
  typography.scale.base: 12
---
timeline
  2026-Q1 : Project kick-off
  2026-Q2 : Beta release
  2026-Q3 : GA
\end{lstlisting}

\paragraph{Init directive (\texttt{\%\%\{init:\ \ldots\}\%\%}).}

For compatibility with tools that strip frontmatter, the same surface is available via
Mermaid's existing init directive hook:

\begin{lstlisting}[language=yaml,caption={Theme selection via Mermaid init directive}]
%%{init: {"theme": "executive", "layout": "vertical-spine", "density": "compact"}}%%
timeline
  2026-Q1 : Project kick-off
\end{lstlisting}

\paragraph{Surface contract.}
Mermaid is the floor: any standard Mermaid document is valid input, and the compiler
respects Mermaid's built-in theme names (\texttt{default}, \texttt{dark}, \texttt{neutral},
\texttt{forest}, \texttt{base}) as aliases to our nearest equivalent. Our named themes and
token overrides are the ceiling---available but never required. A document with no theming
directives renders under the default theme~\cite{mermaiddocs2024}.

% ──────────────────────────────────────────────────────────────────────────────
\subsection{Determinism Under Theme-Driven Layout}
\label{sec:themes-determinism}

Theme-driven geometry does not compromise determinism. The constraint is:

\begin{quote}
\textbf{Determinism rule (binding):} Themes supply \emph{constants}; engines compute
coordinates \emph{closed-form} from those constants. Same IR + same theme + same layout
argument $\to$ byte-identical output. Themes may change \emph{which} constants are
supplied; they may not introduce randomness, iteration-count sensitivity, or
non-deterministic solver state~\cite{vegalite2017}.
\end{quote}

This is already the case for the existing layout-driving tokens: \texttt{spineSpacing} and
\texttt{nodeWrap} are constants supplied to the layout engine; the engine's coordinate
computation is a closed-form function of those constants. Extending this discipline to the
general contract requires only that every new token be a scalar, enumerated value, or
named step---not a runtime-computed value.

Effect tokens (\texttt{drop-shadow}, \texttt{glow}) degrade gracefully per backend; the
\emph{geometry} under any effect is always identical (see Section~\ref{sec:fidelity-tiers}).

% ──────────────────────────────────────────────────────────────────────────────
\subsection{Fidelity Tiers and Backend Effect Profiles}
\label{sec:fidelity-tiers}

Fidelity tiers define what class of visual richness a theme requests and therefore which
rendering backend is sufficient to honour it fully. A backend that cannot satisfy the
requested tier applies the fallback policies declared in each effect's
\texttt{fallback-policy} field and degrades gracefully; \emph{geometry is never affected}.

\begin{table}[h]
\centering
\small
\begin{tabularx}{\textwidth}{lXXXX}
\toprule
\textbf{Tier} & \textbf{Name} & \textbf{Effect classes} & \textbf{Min.\ backend} & \textbf{Examples} \\
\midrule
0 & Minimal &
  No effects; solid fills, patterns, and opacity only &
  SVG backend &
  Minimal \\[4pt]
1 & Crisp &
  Linear/radial gradients, diagonal hatch, dashed borders, clip paths &
  SVG backend &
  Consulting, Release \\[4pt]
2 & Polished &
  Drop shadows, soft glow; native PPTX effects; SVG safe-filter set (determinism caveat) &
  SVG (caveat) or Raster &
  Executive, Product \\[4pt]
3 & Showcase &
  Bloom, cloud/atmosphere layers, noise textures, gradient meshes, volumetric compositing &
  Raster backend required &
  Showcase, ByteByteGo \\
\bottomrule
\end{tabularx}
\caption{Fidelity tier definitions, effect classes, minimum required backend, and example themes}
\label{tab:fidelity-tiers}
\end{table}

\paragraph{Degradation policy.}
A Tier-3 cloud-layer effect targeting the SVG backend:
\begin{enumerate}
  \item The SVG backend checks capability against \texttt{cloud-layer}: ``not available.''
  \item It reads the effect's \texttt{fallback-policy}: \texttt{embed-raster}.
  \item It invokes the Raster backend for the cloud-layer group only, receives a PNG, and
        embeds it as a \texttt{<image>} element. The rest of the diagram remains clean
        vector.
  \item If the Raster backend is also unavailable, the policy chain falls to
        \texttt{omit}: the cloud layer is silently dropped; no geometry changes.
\end{enumerate}

\paragraph{Geometry invariant under degradation.}
Fallback decisions never alter coordinates, dimensions, or label positions. The Scene
geometry is the invariant shared across all backends and all fidelity levels.

% ──────────────────────────────────────────────────────────────────────────────
\subsection{Theme Inheritance and Extension}
\label{sec:theme-inheritance}

Themes support single-level inheritance via the \texttt{extends} key. A child theme
specifies only the tokens that differ from its parent; all others inherit. This enables
branded variants (\texttt{acme-consulting} extending \texttt{consulting}) with minimal
duplication.

\begin{lstlisting}[language=yaml,caption={Custom theme extending Consulting (general-contract form)}]
theme-id:      "acme-corp"
display-name:  "ACME Corporate"
extends:       "consulting"

palette:
  accent:       "#8B0000"   # ACME brand red replaces navy
  accent-strong: "#5A0000"

typography:
  family:   "Futura"
  fallback: "Century Gothic, sans-serif"
  files:
    regular: "fonts/Futura-Book.woff2"
    bold:    "fonts/Futura-Bold.woff2"
\end{lstlisting}

Inheritance rules:
\begin{enumerate}
  \item Properties not specified in the child are taken from the parent exactly.
  \item \texttt{palette} and nested maps are merged entry-by-entry: a child overrides
        only the entries it declares; unspecified entries inherit from parent.
  \item A font family change requires supplying the corresponding embedded font files.
        A theme that names a font without providing files is malformed.
  \item Inheritance depth is exactly one level. Chains
        (\texttt{A} extends \texttt{B} extends \texttt{C}) are not supported.
  \item The general contract must be fully satisfiable from the resolved (merged) theme.
        A child that leaves a required contract field unresolved is malformed.
\end{enumerate}

% ──────────────────────────────────────────────────────────────────────────────
\subsection{Built-in Named Themes}
\label{sec:builtin-themes}

The compiler ships six built-in named themes that address the primary use cases identified
across David's research and the existing timeline implementation. Each is a complete
realisation of the general contract; each component's implementation produces its own
visual interpretation of the same theme tokens~\cite{thinkcell,mermaid2023}.

\begin{description}
  \item[\texttt{consulting} (Tier 1 — Crisp)] McKinsey/BCG-style. Black, white, single navy
        accent. Heavy typography, generous whitespace, square corners, no decorative chrome.
        \texttt{density: normal}. Optimised for A4/Letter print and A3 poster output.

  \item[\texttt{executive} (Tier 2 — Polished)] Corporate boardroom. Navy/slate on white,
        subtle quarter shading, 4~px corner radius on bars, serif heading font. Status
        communicated by colour $+$ icon. \texttt{density: normal}. 16:9 slide target.

  \item[\texttt{product} (Tier 2 — Polished)] Developer-friendly; information-dense. Per-
        track colour coding. Category colours over status colours; status via small icon.
        \texttt{density: compact}. 1400~px canvas.

  \item[\texttt{release} (Tier 1 — Crisp)] Engineering release calendar. Traffic-light
        palette (green/amber/red). Monospace font. Today marker most prominent element.
        \texttt{density: compact}. Optimised for CI/CD dashboard display.

  \item[\texttt{minimal} (Tier 0 — Minimal)] Maximum whitespace, minimum decoration.
        Single grey fill regardless of status; pattern differentiates status. Print-
        optimised; renders correctly in greyscale. \texttt{density: comfortable}.

  \item[\texttt{showcase} (Tier 3 — Showcase)] Premium keynote aesthetic. Gradient fills,
        drop shadows on nodes, optional atmospheric cloud layer. Full fidelity requires
        Raster backend. \texttt{density: normal}. 1920~px canvas.
\end{description}

\paragraph{Cross-component example.}
Consider the \texttt{consulting} theme applied to three component types simultaneously in a
composition poster:
\begin{itemize}
  \item \textbf{Timeline component:} \texttt{spineSpacing: time}, \texttt{accent} drives
        activity bars, \texttt{corner-radius: 0} produces square bars, \texttt{scale.base}
        = 11~pt for activity labels.
  \item \textbf{Flow component:} \texttt{accent} drives node fills, \texttt{ink} drives
        edge labels, \texttt{corner-radius: 0} produces rectangular nodes, \texttt{scale.sm}
        = 9~pt for edge labels.
  \item \textbf{Sequence component:} \texttt{accent} drives participant header fills,
        \texttt{ink-muted} drives lifeline colour, \texttt{scale.base} = 11~pt for
        message labels.
\end{itemize}
All three diagrams share the same navy accent, the same type scale, the same square shape
language, and the same generous whitespace rhythm---because all three consume the same
\texttt{consulting} general-contract tokens. The poster is visually coherent without any
per-component style co-ordination.

% ──────────────────────────────────────────────────────────────────────────────
\subsection{The Three-Tier Token Architecture}
\label{sec:token-tiers}

Token vocabulary is organized into three tiers following the standard design-token
pattern~\cite{dtcg2024}. The tier boundaries are firm: references flow
\emph{upward only}; the general contract is never contaminated by component-specific
concerns.

\begin{description}
  \item[Tier~1 — Primitives] Raw, un-themed values: colour ramps (HSL ladders, raw hex
        families), typeface family names, the spacing base unit. These are the raw
        material from which Tier-2 tokens are assembled. Components \emph{never} reference
        Tier-1 tokens directly; themes build Tier-2 values from them.

  \item[Tier~2 — Semantic tokens (the general contract)] Role palette, data palette,
        type scale and weights, spacing steps, density, shape language (corner radius,
        node padding, stroke scale, connector style), and effects/fidelity tier.
        \textbf{This is the interface every component implements against.} The full Tier-2
        shape is specified in Section~\ref{sec:themes-contract} and illustrated in the
        listings above.

  \item[Tier~3 — Component tokens] Namespaced \texttt{components.\textit{name}.\textit{token}}
        (e.g.\ \texttt{components.serpentine.turnRadius},
        \texttt{components.sankey.ribbonOpacity},
        \texttt{components.gitgraph.mergeCurveTension}). Each Tier-3 token is
        \emph{derived} from a Tier-2 token by the component's theme-binding as its default.
        A concrete theme \emph{may} override any Tier-3 token as an optional fine-tuning
        step; it is never required to do so.
\end{description}

\begin{quote}
\textbf{Binding invariants (tier contract):}
\begin{itemize}
  \item The general contract (Tier~2) \textbf{NEVER} references a Tier-3 component token.
        The contract is component-agnostic by construction.
  \item Components reference \textbf{upward only}: Tier-3 tokens derive from Tier-2
        defaults; a component may not require a Tier-2 token that does not exist in the
        contract.
  \item A theme \textbf{may} reach downward into a \texttt{components.\textit{name}}
        namespace as an optional override. This is the escape hatch for per-component
        fine-tuning; it is never mandatory.
\end{itemize}
\end{quote}

\subsubsection{Tier-3 Example: serpentine turn radius}

\begin{lstlisting}[language=yaml,caption={Tier-3 token deriving from Tier-2 with optional theme override}]
# ── Tier 2 (general contract, theme-authored) ──────────────────────────
shape:
  corner-radius: 4       # px — general shape language

# ── Tier 3 (component binding, inside serpentine layout engine) ─────────
# Default derivation (computed at bind time, not stored in the theme file):
#   components.serpentine.turnRadius = shape.corner-radius * 6
#   → 4 * 6 = 24 px (sufficient arc to avoid a sharp hairpin)

# ── Optional Tier-3 override in a concrete theme (e.g. "showcase") ──────
components:
  serpentine:
    turnRadius: 40        # override: generous arc for showcase aesthetic
  sankey:
    ribbonOpacity: 0.55   # override: slightly more opaque ribbons
  gitgraph:
    mergeCurveTension: 0.4  # override: softer merge-curve
\end{lstlisting}

The \texttt{components} block is optional in every theme. A theme that omits it entirely
is valid and fully correct: all Tier-3 tokens resolve to their Tier-2 derivation defaults.
Only themes with a strong visual rationale for per-component fine-tuning need to supply it.

% ──────────────────────────────────────────────────────────────────────────────
\subsection{Contract Vocabulary Summary}
\label{sec:contract-summary}

The Tier-2 contract in full, showing all six domains together:

\begin{lstlisting}[language=yaml,caption={Full Tier-2 general contract skeleton (decided)}]
# ── Palette ───────────────────────────────────────────────────────────────────
palette:            # role palette — structural components
  surface: …        # see §12.3 for full role list
  surface-panel: …  # lane / panel / track-header background
  ink: …
  ink-panel: …      # text rendered on panel backgrounds
  accent: …
  border: …
  muted: …
  status-success: { fill: …, stroke: … }
  status-warning: { fill: …, stroke: … }
  status-error:   { fill: …, stroke: … }
  status-info:    { fill: …, stroke: … }
  # workflow states: status-planned, status-active, status-done,
  #                  status-cancelled, status-uncertain
  status-pattern: "solid"  # solid | hatch | dashed-border

data-palette:       # data palette — chart/quantity components
  categorical: [ …, …, …, …, …, … ]   # min 6 entries, ordered
  sequential:  { low: …, mid: …, high: … }
  diverging:   { negative: …, zero: …, positive: … }

# ── Typography ────────────────────────────────────────────────────────────────
typography:
  family: …
  fallback: …
  files: { regular: …, bold: … }
  scale: { xs: 8, sm: 9, base: 11, md: 13, lg: 18, xl: 24 }   # pt
  weights: { regular: 400, medium: 500, semibold: 600, bold: 700 }

# ── Spacing (advisory) ────────────────────────────────────────────────────────
spacing:
  unit: 8   # px — base rhythm unit
  steps: { xxs: 2, xs: 4, sm: 8, md: 16, lg: 24, xl: 32, xxl: 48 }

# ── Density ───────────────────────────────────────────────────────────────────
density: "normal"   # compact | normal | comfortable

# ── Shape language ────────────────────────────────────────────────────────────
shape:
  corner-radius: 0
  node-padding:  8
  stroke-scale:  1.0
  connector-style: "elbow"  # straight | elbow | curved | orthogonal
  marker-shape: "circle"    # circle | diamond | square

# ── Effects / fidelity ────────────────────────────────────────────────────────
effects:
  fidelity: 1   # 0=Minimal | 1=Crisp | 2=Polished | 3=Showcase
  drop-shadow: { enabled: false }
  glow:        { enabled: false }
  motion:      { enabled: false }
\end{lstlisting}

% ──────────────────────────────────────────────────────────────────────────────
\subsection{Contract Adoption Plan}
\label{sec:contract-adoption}

\subsubsection{Proof Set: Validate the Contract First}

Before migrating any production component, the contract is validated by a spike: implement
the \texttt{executive} theme against three components chosen to exercise the entire
contract surface jointly.

\begin{description}
  \item[\texttt{flowchart}] Node-link boxes, edge routing, orientation. Exercises: role
        palette, shape language (corner radius, connector style), type scale, density,
        spacing advisory steps.

  \item[\texttt{sequence}] Text/compartment geometry, activation bars, fragments, lifelines.
        Exercises: role palette, type scale, density (row heights, secondary label
        visibility), shape language (node padding, corner radius).

  \item[\texttt{xychart}] Axes, gridlines, series bars/lines, tick labels. Exercises: the
        \emph{data palette} (categorical sequence + sequential ramp), type scale, spacing
        rhythm. This is the critical third leg: the other two components do not touch the
        data palette, so xychart is required to prove the full contract surface.
\end{description}

The spike succeeds when all three diagrams render coherently under the single
\texttt{executive} theme token set without any component-specific hacks. A passing spike
commits the Tier-2 vocabulary as final~\cite{dtcg2024}.

\subsubsection{Migration Order}

After the proof spike, the 16 existing per-component \texttt{theme.ts} files are migrated
in the following order:

\begin{enumerate}
  \item \textbf{Generalise the timeline \texttt{ResolvedTheme} upward} into the Tier-2
        contract (\texttt{packages/core/src/themes/types.ts}). The timeline is the deepest
        existing layout-driven theme and serves as the reference implementation; all
        subsequent migrations align to the shape it establishes.

  \item \textbf{Node-link family:} \texttt{class}, \texttt{state}, \texttt{ER}, \texttt{C4},
        \texttt{requirement}, \texttt{block}, \texttt{architecture}. These share the same
        fundamental node-edge visual vocabulary and are highest value after the proof set.

  \item \textbf{Remaining charts:} \texttt{pie}, \texttt{quadrant}, \texttt{radar}.
        These consume the data palette and validate categorical and diverging ramp usage
        beyond xychart.

  \item \textbf{Timeline-family adoption of the contract.} The timeline component
        re-consumes the now-general Tier-2 contract, keeping its Tier-3 tokens
        (\texttt{components.serpentine.*}, \texttt{components.roadmap.*}) as derived
        overrides.

  \item \textbf{Specialised components:} \texttt{sankey}, \texttt{gitGraph},
        \texttt{journey}, \texttt{kanban}, \texttt{mindmap}, \texttt{packet}.
        These have the most Tier-3 component-specific tokens and are migrated last
        once the Tier-3 override mechanism is proven by earlier steps.
\end{enumerate}

\begin{quote}
\textbf{Migration invariant:} Each step is independently shippable. No step requires a
downstream step to complete. At every point the compiler is in a valid state: un-migrated
components continue to consume their legacy \texttt{theme.ts}; migrated components consume
the general contract. Both coexist until Step~5 completes.
\end{quote}

% ──────────────────────────────────────────────────────────────────────────────
\subsection{Implementation Status (2026)}
\label{sec:contract-impl-status}

The contract described in this section is fully implemented. The Tier-2
\texttt{ThemeContract} interface and the \texttt{executive} reference theme reside in
\texttt{packages/core/src/theme-contract/}. Every one of the 21 diagram components now
adopts the contract via a per-component binding: for most grammars the binding is
\texttt{grammars/\textit{<component>}/contract-binding.ts}; the timeline uses
\texttt{themes/contract-binding.ts}. Adoption is \textbf{opt-in}: a theme name registered
in \texttt{CONTRACT\_THEMES} resolves via the binding; all legacy per-component named themes
are unchanged and their rendered output is byte-identical to pre-contract builds. The
\texttt{executive} theme has been verified to render all 21 diagram types as one coherent
design system (navy accent, Georgia serif, white surface).

\begin{itemize}
  \item \textbf{Precedence rule:} Legacy component theme names always win. A contract
        theme name must not duplicate a name already claimed by a legacy per-component
        theme; the \texttt{CONTRACT\_THEMES} resolver never overrides a legacy route.
  \item \textbf{Contract location:}
        \texttt{packages/core/src/theme-contract/}
        (\texttt{ThemeContract} Tier-2 interface $+$ \texttt{executive} concrete theme).
  \item \textbf{Binding pattern:}
        \texttt{grammars/\textit{<component>}/contract-binding.ts} (20 grammars)
        and \texttt{themes/contract-binding.ts} (timeline) --- 21 bindings total.
\end{itemize}
                % Theme architecture + schema
\section{The Aesthetic Bar: Beating Mermaid Visually}
\label{sec:aesthetics}

Aesthetics is not a finishing coat applied after architecture---it is a \textbf{first-class
pillar} of the product. The explicit goal: every diagram produced by this compiler must look
demonstrably better than its Mermaid equivalent. ``Pro,'' ``neat,'' ``coherent'' are
requirements, not aspirations.

\subsection{Why Aesthetics Is Architecture}

Mermaid is ubiquitous because of its ecosystem integration, not its visual quality.
Users accept its output because it's free and frictionless---not because it looks good.
The aesthetic gap is the primary differentiation lever:

\begin{description}
  \item[Mermaid] Documentation-grade. Fixed styles. Poor typography. No design system.
        Functional but never beautiful.
  \item[This compiler] Presentation-grade. Themeable. Proper typography. Coherent design
        language. Output you'd put in a board deck or technical blog.
\end{description}

This is not subjective polish---it is the \textbf{primary reason users choose this tool
over Mermaid}. If the output doesn't look materially better, there is no product.

\subsection{The Design System}

Every diagram shares a coherent visual language:

\begin{description}
  \item[Typography] System-level font stack with proper weight hierarchy
        (title/heading/body/label/caption). Consistent sizing ratios across all diagram types.
        No all-caps labels. No tiny illegible text.

  \item[Spacing and rhythm] Consistent padding, margins, and gaps derived from a base unit.
        Generous whitespace. Visual breathing room. The ``no clutter'' principle
        (Section~\ref{principle:minimal-clutter}) is enforced structurally.

  \item[Color palette] Curated palettes with proper contrast ratios (WCAG AA minimum).
        Semantic color mapping: status colors, category colors, emphasis colors.
        Dark and light modes as first-class variants, not afterthoughts.

  \item[Shape language] Rounded corners, consistent border weights, soft shadows where
        appropriate. No harsh black-on-white wireframe aesthetic.

  \item[Edge/connector style] Smooth curves, proper arrowheads, consistent stroke widths.
        Routing that avoids overlaps where the layout algorithm permits.
\end{description}

\subsection{Default Theme Craft}

The default theme must be beautiful \emph{out of the box}. Users who never configure a theme
should still produce output that surpasses Mermaid. Design commitments:

\begin{itemize}
  \item Modern, clean aesthetic inspired by technical publications (not wireframes)
  \item Light theme: warm grays, blue accent, Inter/system-ui stack
  \item Dark theme: deep navy/charcoal, blue-teal accent, proper contrast
  \item Professional appearance suitable for slides, blog posts, documentation
  \item No garish colors, no 1990s UML aesthetic, no pixelated edges
\end{itemize}

\subsection{Brand Themes and Coherence}

The theme system enables organization-specific brand themes:

\begin{itemize}
  \item Themes can define complete visual identities (colors, fonts, shapes, spacing)
  \item Single inheritance: branded themes extend base themes, overriding selectively
  \item Cross-diagram coherence: a theme applied to a poster produces visually unified
        output across all constituent diagram types
  \item Built-in exemplars: \texttt{professional-light}, \texttt{professional-dark},
        \texttt{bytebytego} (inspired by the popular technical illustration style)
\end{itemize}

\subsection{Grammar = Semantics; Theme = Style}

The architectural discipline that enables aesthetic excellence:

\begin{quote}
\emph{The Domain IR carries what to communicate. The theme carries how it looks.
Never mix them. A theme change cannot alter diagram meaning. A content change
cannot alter diagram style.}
\end{quote}

This separation means:
\begin{itemize}
  \item The same diagram looks professional in \texttt{professional-light} and playful
        in a hypothetical \texttt{sketch} theme---without any IR changes.
  \item Agents author content without making aesthetic decisions.
  \item Designers craft themes without understanding grammar semantics.
  \item New grammars inherit the full theme system and look coherent from day one.
\end{itemize}

\subsection{Aesthetic Verification}

Beauty is hard to test, but we enforce baselines:

\begin{itemize}
  \item \textbf{Golden-image regression}: visual output is diffed against known-good
        baselines. Any change to rendering requires explicit approval.
  \item \textbf{Contrast ratio checks}: automated WCAG AA verification for all
        text-on-background combinations in built-in themes.
  \item \textbf{Gallery examples}: every grammar ships with gallery outputs demonstrating
        the aesthetic bar. These serve as both documentation and regression targets.
  \item \textbf{Side-by-side comparisons}: marketing/documentation includes Mermaid vs.\
        ours comparisons for every supported diagram type.
\end{itemize}

\subsection{Implementation Status}

\textbf{Built:} The theme system is fully implemented. Five timeline themes (including dark
and ByteByteGo variants), two flow themes, two sequence themes, two tree themes, and two
composition themes are shipped. All enforce the grammar=semantics/theme=style discipline.

\textbf{In progress:} Expanding the aesthetic standard to the new diagram types as they
ship. The UML/software line (Tier-1) will establish the ``modern UML'' visual standard
that differentiates from PlantUML's dated look.
            % The aesthetic bar; beating Mermaid visually

% ============================================================================
% PART VI — COMPOSITION
% ============================================================================

\part{Composition}
\label{part:composition}

\section{Composition: Multi-Panel Posters}
\label{sec:composition}

The Composition layer arranges multiple grammar outputs into a single poster or infographic. This is how the multi-section corpus images (10x-coding-playbook, archon-harness-builder, rag-vectorless-explainer) are assembled.

\subsection{Design Philosophy}

\paragraph{Composition is a thin layer.}
The Composition layer does not define new visual primitives. It arranges existing grammar outputs (each a Scene IR) into a larger canvas.

\paragraph{Each panel is an independent grammar.}
A panel might be a timeline, a flow diagram, a comparison table, or a stat callout. Each panel has its own Domain IR and is compiled independently.

\paragraph{The composition specifies arrangement, not content.}
Content is defined in each panel's Domain IR. The Composition IR specifies panel positions, sizes, and inter-panel elements (dividers, titles).

\subsection{Composition IR}

\begin{lstlisting}[language=yaml,caption={Composition IR}]
version: "1.0"
metadata:
  title: "The 10x AI Coding Playbook"
  subtitle: "Everything is infrastructure for one thing: parallel development."
  theme: dark-poster
  canvas:
    width: 1200
    height: auto               # computed from panels
composition:
  layout: stack                # stack | grid | custom
  gap: 40                      # px between panels
  panels:
    - id: reframe
      grammar: timeline
      title: "THE REFRAME"
      ir:                      # inline Domain IR
        tracks:
          - id: solo
            label: "2x = solo + AI"
            # ...
    - id: infrastructure
      grammar: flow
      title: "THE INFRASTRUCTURE"
      ir:
        direction: left-to-right
        nodes:
          # ...
    - id: catch
      grammar: flow
      title: "THE CATCH"
      ir:
        # ...
    - id: payoff
      grammar: timeline
      title: "THE PAYOFF"
      ir:
        # ...
  chrome:                      # editorial elements
    header:
      logo: { src: "icon.png", position: top-left }
      title_style: large-centered
    footer:
      badges: ["GitHub", "Linear", "Jira"]
    dividers: between-panels   # none | between-panels | after-each
\end{lstlisting}

\subsection{Layout Modes}

\subsubsection{Stack Layout}

Panels arranged vertically, each taking full width. Most common for long-form infographics.

\begin{Verbatim}[frame=single,fontsize=\small]
+----------------------------------+
|  Panel 1 (full width)            |
+----------------------------------+
|  Panel 2 (full width)            |
+----------------------------------+
|  Panel 3 (full width)            |
+----------------------------------+
\end{Verbatim}

\subsubsection{Grid Layout}

Panels arranged in a grid (rows × columns). Useful for dashboards or comparison layouts.

\begin{lstlisting}[language=yaml,caption={Grid layout specification}]
composition:
  layout: grid
  grid:
    columns: 2
    rows: auto                 # computed from panel count
  panels:
    - { id: p1, span: 1 }      # 1 cell
    - { id: p2, span: 1 }
    - { id: p3, span: 2 }      # 2 cells (full row)
\end{lstlisting}

\subsubsection{Custom Layout}

Explicit (x, y, width, height) for each panel. Maximum control, most verbose.

\begin{lstlisting}[language=yaml,caption={Custom layout specification}]
composition:
  layout: custom
  panels:
    - id: p1
      x: 0
      y: 0
      width: 600
      height: 300
    - id: p2
      x: 600
      y: 0
      width: 600
      height: 300
\end{lstlisting}

\subsection{Compilation Process}

\begin{enumerate}
  \item \textbf{Parse Composition IR}: Validate structure, resolve grammar references.
  
  \item \textbf{Compile each panel}: For each panel, invoke the appropriate grammar's layout engine to produce a Scene IR.
  
  \item \textbf{Compute panel dimensions}: If not explicit, compute from panel Scene dimensions.
  
  \item \textbf{Arrange panels}: Apply layout mode to compute (x, y) for each panel.
  
  \item \textbf{Merge Scenes}: Combine all panel Scenes into a single composition Scene, translating coordinates.
  
  \item \textbf{Add chrome}: Render header, footer, dividers, badges.
  
  \item \textbf{Output}: The final Scene IR (a single scene containing all panels).
\end{enumerate}

\subsection{Panel References}

Panels can reference external IR files instead of inline definitions:

\begin{lstlisting}[language=yaml,caption={Panel with external IR reference}]
panels:
  - id: timeline-panel
    grammar: timeline
    title: "Project Roadmap"
    ir_file: "./roadmap.timeline.yaml"  # external file
\end{lstlisting}

This enables:
\begin{itemize}
  \item Reusing the same diagram in multiple compositions
  \item Separate version control for each panel's IR
  \item Incremental updates (change one panel, not the whole poster)
\end{itemize}

\subsection{Chrome Elements}

\subsubsection{Header}

\begin{itemize}
  \item \textbf{Logo}: positioned at top-left or top-right
  \item \textbf{Title}: large centered text
  \item \textbf{Subtitle}: smaller text below title
\end{itemize}

\subsubsection{Section Titles}

Each panel can have a title rendered above it:

\begin{lstlisting}[language=yaml,caption={Panel title configuration}]
panels:
  - id: infrastructure
    title: "THE INFRASTRUCTURE"
    title_style: section-header   # theme-defined style
\end{lstlisting}

\subsubsection{Dividers}

Horizontal lines between panels:

\begin{lstlisting}[language=yaml,caption={Divider options}]
chrome:
  dividers: between-panels       # none | between-panels | after-each
  divider_style:
    color: "#333"
    thickness: 1
    margin: 20
\end{lstlisting}

\subsubsection{Footer}

\begin{itemize}
  \item \textbf{Badges}: pill-shaped labels (``GitHub'', ``MIT License'')
  \item \textbf{Author}: avatar + handle
  \item \textbf{Branding}: company logo or URL
\end{itemize}

\subsection{Determinism in Composition}

Composition inherits determinism from its constituent grammars:

\begin{itemize}
  \item Each panel's Scene is deterministic (by grammar contract)
  \item Panel arrangement is deterministic (explicit layout)
  \item Chrome rendering is deterministic (explicit configuration)
\end{itemize}

The composition's \texttt{scene\_hash} is computed from the merged Scene, enabling golden testing of complete posters.

\subsection{Limitations}

\begin{description}
  \item[No cross-panel connectors] Edges cannot span from a node in panel A to a node in panel B. Panels are visually independent.
  
  \item[No responsive layout] Panel sizes are computed once; there is no responsive reflow for different canvas widths.
  
  \item[No interactive linking] Clicking a panel does not navigate or filter. Output is static.
\end{description}

These limitations are intentional: they keep the Composition layer simple and deterministic. Cross-panel connections and responsive layout would significantly complicate the architecture.

% =============================================================================
% §30b — Cross-Diagram Node Linking
% =============================================================================
% Specifies the mechanism for drawing edges between nodes that live in
% different cells of a poster/composition.  Superset-only.  Spec-first (2026).
% =============================================================================

\section{Cross-Diagram Node Linking}
\label{sec:cross-diagram-links}

% ─────────────────────────────────────────────────────────────────────────────
\subsection{Concept and Motivation}
% ─────────────────────────────────────────────────────────────────────────────

A poster assembled from independent grammar cells
(Section~\ref{sec:composition}) is, by default, a set of \emph{panels side by
side}: each cell compiles and renders autonomously with no structural
relationship to its neighbours.  Cross-diagram node linking removes that
constraint.  A \texttt{link} statement in the poster DSL draws a directed or
undirected edge between a node in one cell's diagram and a node in another
cell's diagram, turning the poster into a \textbf{linked multi-view system
diagram}.

\begin{description}
  \item[Presentation-layer assertion.]
        A cross-diagram link does not modify either diagram's Domain IR or
        Scene IR.  It is an annotation on the composition itself, resolved
        after every cell has compiled independently.

  \item[Superset-only.]
        The \texttt{link} statement is a new top-level poster keyword with no
        Mermaid equivalent (Section~\ref{sec:superset-extensions}).  A
        standard Mermaid renderer will reject a poster document that contains
        \texttt{link} statements; this is the intended signal that the file
        requires our compiler.

  \item[Additive and non-fatal.]
        Unresolvable links (unknown cell, unknown node, non-linkable diagram
        type) emit a warning and are skipped.  The poster and all valid links
        still render.
\end{description}

% ─────────────────────────────────────────────────────────────────────────────
\subsection{Link Syntax}
\label{sec:cdl-syntax}
% ─────────────────────────────────────────────────────────────────────────────

A \texttt{link} statement is a poster-level declaration written alongside
\texttt{cell} definitions.  Its abstract grammar is:

\begin{lstlisting}[language=,caption={Cross-diagram link statement — abstract grammar}]
link <cellAddr>.<nodeId> <edge> <cellAddr>.<nodeId> [: "label"]

<cellAddr>  ::= "[" row "," col "]"   (* bracket form, 0-indexed *)
             |  colLetters rowNumber  (* Excel form, e.g. A1, B2, AA3 *)

<nodeId>    ::= (* the node identifier as authored in that cell's diagram source *)

<edge>      ::= "-->"                 (* directed arrow *)
             |  "---"                 (* plain undirected *)
             |  "-.-"                 (* dashed undirected *)
             |  "-.."                 (* dashed undirected, alternate *)
             |  "-.->"|"-.->"        (* dashed directed arrow *)
\end{lstlisting}

Cell addresses reuse the dual addressing scheme introduced for the
\texttt{poster} DSL and specified in Section~\ref{sec:composition-posters}:
the bracket form \texttt{[row,col]} is 0-indexed (row first); the Excel form
\texttt{A1}, \texttt{B2}, \ldots{} uses bijective base-26 column letters with
1-indexed row numbers.  Both forms may be mixed freely within a single poster.

The optional \texttt{: "label"} suffix attaches a text label rendered at
the midpoint of the overlay edge, styled by the composition theme's
\texttt{overlayLabelFont} token.

\paragraph{Worked example.}

The following poster specifies an Order Fulfilment system with four cells and
three cross-diagram links that connect nodes across cells:

\begin{lstlisting}[language=,caption={Cross-diagram links in a four-cell poster}]
---
theme: executive
layout: grid 2x2
---
poster "Order Fulfilment System"

  cell A1: flowchart LR
    recv[Receive Order] --> valid[Validate]
    valid --> pay[Payment]
    pay --> pick[Warehouse Pick] --> ship[Dispatch]

  cell B1: classDiagram
    class OrderService
    class PaymentGateway
    class WarehouseSystem
    OrderService --> PaymentGateway
    OrderService --> WarehouseSystem

  cell A2: stateDiagram-v2
    [*] --> Pending
    Pending --> Processing
    Processing --> Shipped
    Shipped --> [*]

  cell B2: sequenceDiagram
    Customer->>OrderService: placeOrder()
    OrderService->>PaymentGateway: charge()
    PaymentGateway-->>OrderService: receipt

link A1.pay --> B1.PaymentGateway : "handled by"
link A1.pick -.-> A2.Processing : "triggers"
link B1.OrderService --- A2.Pending : "creates"
\end{lstlisting}

\texttt{link A1.pay --> B1.PaymentGateway} reads: ``draw a directed arrow
from the node with id \texttt{pay} in cell \texttt{A1} to the class named
\texttt{PaymentGateway} in cell \texttt{B1}''.  The cell addresses
\texttt{A1} and \texttt{B1} are the Excel equivalents of \texttt{[0,0]}
and \texttt{[0,1]}.

% ─────────────────────────────────────────────────────────────────────────────
\subsection{The Node-Anchor Registry}
\label{sec:cdl-anchor-registry}
% ─────────────────────────────────────────────────────────────────────────────

\paragraph{The core engineering prerequisite.}

Grammar layout engines today compute precise node positions internally but
\textbf{discard them}: the output of \texttt{buildFlowScene}, \texttt{buildTreeScene},
and their siblings is a bare \texttt{Scene} containing only flattened
rendering primitives.  For example, the flow layout engine builds a
\texttt{PlacedNode} record for every node (carrying \texttt{x}, \texttt{y},
\texttt{w}, \texttt{h}, \texttt{cx}, \texttt{cy}, and port coordinates) but
returns only the rendered \texttt{Scene}; there is no \texttt{nodeId
$\to$ position} registry in the output type.  \textbf{Making node positions
addressable from the composition layer is the single prerequisite that must be
resolved before any cross-diagram link can be drawn.}

\subsubsection{Candidate Implementations}

Two approaches exist for surfacing node anchors to the composition layer.

\paragraph{Candidate (a) — Populate \texttt{GroupPrimitive.id} and walk the
placed scene.}

The \texttt{GroupPrimitive} in the Scene IR carries an optional \texttt{id?:\
string} field (see \texttt{packages/core/src/scene.ts}).  Each grammar's
layout engine could wrap each node's primitives in a \texttt{GroupPrimitive}
tagged with that node's id.  At composition time the link layer would walk the
placed (already translated-and-scaled) scene, compute the axis-aligned
bounding box of each named group, and use it as the anchor.

\medskip
\textit{Pros:} No change to grammar function return types; \texttt{GroupPrimitive.id}
already exists in the IR.

\medskip
\textit{Cons:}
\begin{itemize}
  \item Recovering a bounding box requires walking the full scene tree for each
        node query (O($n$) primitives per lookup for a scene with $n$ primitives).
  \item The walk must be performed after \texttt{translateAndScale}, meaning
        coordinate reconstruction happens at the composition layer, not the
        grammar layer.
  \item No port information (which side of the node to attach an edge to) is
        recoverable from a bounding box without additional metadata.
  \item Requires every linkable grammar to wrap each node in a named
        \texttt{GroupPrimitive}, which is not currently the case and is a
        fragile convention: a grammar that emits node primitives without a
        wrapping group would silently produce empty anchors.
  \item Not reusable for hit-testing or agent node-addressing without
        additional work.
\end{itemize}

\paragraph{Candidate (b) — Sidecar \texttt{anchors} record on the render
result.}

Each grammar's layout function is extended to return a pair
\texttt{\{ scene: Scene, anchors: NodeAnchorRegistry \}} in addition to the
Scene.  The registry maps each node id to its bounding box and optional
connection ports, expressed in the cell's \textbf{local coordinate space}
(i.e., before composition-level translation and scaling).  Grammars that opt
in populate the registry from data they already hold during layout (the flow
grammar's \texttt{PlacedNode} is a direct source); grammars that have not yet
opted in return an empty registry, and links referencing their nodes degrade
gracefully.

\begin{lstlisting}[language=javascript,caption={NodeAnchorRegistry --- type specification (pseudo-TypeScript)}]
/** Port positions on the four cardinal sides of a node's bounding box. */
type CardinalSide = 'N' | 'S' | 'E' | 'W';

/**
 * Anchor for one node: bounding-box in local cell coordinates,
 * plus optional explicit port attachment points.
 * If ports is absent, the link layer derives port midpoints from
 * (x,y,w,h) at rendering time.
 */
interface NodeAnchor {
  x: number;                                   // bounding-box top-left x
  y: number;                                   // bounding-box top-left y
  w: number;                                   // bounding-box width
  h: number;                                   // bounding-box height
  ports?: Partial<Record<CardinalSide, { x: number; y: number }>>;
}

/** Maps diagram-local node id -> anchor in local cell coordinates. */
type NodeAnchorRegistry = Record<string, NodeAnchor>;

/**
 * Extended render result returned by linkable grammars.
 * Non-linkable grammars continue to return Scene only; the composition
 * layer treats an absent registry as {}.
 */
interface RenderResult {
  scene: Scene;
  anchors: NodeAnchorRegistry;   // empty ({}) for non-linkable grammars
}
\end{lstlisting}

\textit{Pros:} Explicit, typed contract authored by the grammar that knows its
node positions best; no scene walking; carries port information; each grammar
opts in independently; composable independently of the scene.

\medskip
\textit{Cons:} Requires changing grammar function return types; migration is
incremental (one grammar at a time).

\subsubsection{Recommended Implementation: Sidecar Anchors (Candidate b)}

\textbf{The recommended approach is the sidecar} \textbf{\texttt{NodeAnchorRegistry}} \textbf{(candidate b).}  The rationale:

\begin{enumerate}
  \item The grammar layout engine already owns the exact node geometry.  For
        the flow grammar, \texttt{PlacedNode.x}, \texttt{.y}, \texttt{.w},
        \texttt{.h} and the port coordinates \texttt{.rx}/\texttt{.lx}/\texttt{.cx}/\texttt{.by}
        are computed during Phase 4 (coordinate assignment) and available
        before scene emission.  Populating the registry is a two-line loop over
        the placed-node map --- no new computation required.
  \item The composition layer already tracks each cell's
        \texttt{(dx, dy, scale)} from the \texttt{translateAndScale} call in
        \texttt{layoutComposition} (see \texttt{packages/core/src/composition/layout.ts}).
        Transforming local anchors to poster coordinates is exact and trivial:
        \[
          x_{\mathrm{poster}} = x_{\mathrm{local}} \cdot s + d_x, \quad
          y_{\mathrm{poster}} = y_{\mathrm{local}} \cdot s + d_y, \quad
          w_{\mathrm{poster}} = w_{\mathrm{local}} \cdot s, \quad
          h_{\mathrm{poster}} = h_{\mathrm{local}} \cdot s
        \]
        where $s$ is the uniform scale factor and $(d_x, d_y)$ is the
        translation offset already computed by the grid layout pass.
  \item Candidate (a) would require reconstructing the same geometry from the
        scene tree, at the cost of an O($n$)-walk on the flattened primitive
        list, after the coordinate transform has already been applied.
  \item The opt-in migration path means zero risk to existing grammars:
        a grammar that has not yet emitted anchors returns \texttt{\{\}};
        links referencing it degrade to a WARN and are skipped
        (Section~\ref{sec:cdl-semantics}).
\end{enumerate}

\paragraph{Transformation to poster coordinates.}

After each cell's \texttt{RenderResult} is produced, the composition layer
performs the anchor transformation:

\begin{lstlisting}[language=javascript,caption={Anchor transformation to poster coordinates (pseudo-code)}]
function transformAnchors(
  registry: NodeAnchorRegistry,
  dx: number,    // cell's x offset in poster (from translateAndScale)
  dy: number,    // cell's y offset in poster
  scale: number, // uniform scale factor applied to the sub-scene
): NodeAnchorRegistry {
  const out: NodeAnchorRegistry = {};
  for (const [id, anchor] of Object.entries(registry)) {
    const pa: NodeAnchor = {
      x: anchor.x * scale + dx,
      y: anchor.y * scale + dy,
      w: anchor.w * scale,
      h: anchor.h * scale,
    };
    if (anchor.ports) {
      pa.ports = {};
      for (const [side, pt] of Object.entries(anchor.ports)) {
        pa.ports[side as CardinalSide] = {
          x: pt.x * scale + dx,
          y: pt.y * scale + dy,
        };
      }
    }
    out[id] = pa;
  }
  return out;
}
\end{lstlisting}

The resulting poster-coordinate registries from all cells are stored alongside
the placed-cell rectangles for use by the overlay routing pass.

\paragraph{Reusability argument.}

The node-anchor registry is valuable well beyond cross-diagram linking and
should be built regardless of whether the linking feature ships in the first
increment:

\begin{itemize}
  \item \textbf{Hit-testing:} An interactive SVG export can route mouse events
        to the correct diagram node by name, enabling click/hover handlers.
  \item \textbf{Tooltips:} Structured hover labels (from the Domain IR's node
        metadata) can be attached to the exact node bounding box.
  \item \textbf{Agent node-addressing:} The agent integration surface
        (Section~\ref{sec:agent-integration}) currently addresses diagrams by
        their IR structure.  A published anchor registry lets an agent
        reference ``cell A1, node \texttt{pay}'' by logical id rather than by
        pixel coordinates --- a stable, round-trip-safe reference.
  \item \textbf{Accessibility:} ARIA \texttt{role="img"} subtrees can carry
        \texttt{aria-label} attributes derived from each node's anchor id and
        IR label, improving screen-reader navigation of complex diagrams.
\end{itemize}

% ─────────────────────────────────────────────────────────────────────────────
\subsection{Overlay Routing}
\label{sec:cdl-overlay-routing}
% ─────────────────────────────────────────────────────────────────────────────

\paragraph{Overlay layer.}

Cross-diagram edges are rendered as a synthetic pass executed \emph{after} all
cell sub-scenes have been embedded.  The overlay primitives (line/path) are
appended to the poster's primitive list last, placing them above all cell
content in painter's-algorithm order.  No new Scene IR primitive type is
required: overlay edges are \texttt{PathPrimitive} (directed, with arrowhead)
or \texttt{LinePrimitive} (plain) constructed from the same kernel as
intra-diagram edges.

\paragraph{Port selection.}

Given source anchor $A$ (with cardinal ports $A_N$, $A_S$, $A_E$, $A_W$) and
target anchor $B$, the link layer selects the port pair $(p_A, p_B)$ that
minimises the Euclidean distance $\|p_A - p_B\|$.  Tie-breaking order: E
before W before S before N (favouring left-to-right reading direction).  When
explicit ports are absent from the registry, the four cardinal midpoints are
derived from the bounding box:
\[
  A_N = (x + w/2,\; y), \quad
  A_S = (x + w/2,\; y+h), \quad
  A_E = (x+w,\; y+h/2), \quad
  A_W = (x,\; y+h/2).
\]

\paragraph{Phase 1 routing: straight and single-elbow.}

The initial implementation uses two routing strategies selected by the
relative cell positions:

\begin{description}
  \item[Same row, adjacent columns] Straight horizontal segment from $p_A$
        to $p_B$, passing through the inter-cell gap.  This is the most common
        case (left-to-right flow) and produces clean, readable results with no
        routing logic.

  \item[Different rows or non-adjacent columns] Single-elbow (L-shaped) route:
        a horizontal segment from $p_A$ to a midpoint column, then a vertical
        segment to $p_B$'s row, then a horizontal segment to $p_B$.  The
        elbow column is placed at the midpoint between the two cell centres.
        All segments run through the inter-cell gap space and do not overlap
        cell bodies for standard grid layouts.
\end{description}

\paragraph{Phase 2 hardening: orthogonal routing.}

For complex poster layouts where the straight/elbow route would cross a cell
body, the composition engine can activate orthogonal routing:

\begin{enumerate}
  \item Build a set of \emph{forbidden rectangles} from each placed cell's
        bounding box, expanded outward by the theme's \texttt{gap} / 2.
  \item Compute the rectilinear shortest path from $p_A$ to $p_B$ that avoids
        all forbidden rectangles, using the inter-cell gap channels as routing
        channels.
  \item If no route exists (e.g., a cell is surrounded on all sides), emit a
        WARNING and fall back to the Phase 1 route drawn on top of the
        obstructing cell (readable but visually overlapping).
\end{enumerate}

Phase 2 routing is a \emph{hardening step}; Phase 1 covers the vast majority
of practical poster layouts and is the only mode required for the first
increment.

\paragraph{Dimension guard.}

The overlay routing pass applies dimension guards consistent with the
project's overall guard mindset
(Section~\ref{sec:composition-edge-cases}):

\begin{itemize}
  \item If the route must cross more than three cell bodies (unreachable via
        gap channels), emit \texttt{WARNING: link \textit{A.nodeId}
        $\to$ \textit{B.nodeId}: route crosses N cell bodies; overlay may be
        unreadable.}  Still render best-effort.
  \item If both endpoints resolve to the same cell, emit
        \texttt{WARNING: intra-cell link \textit{A.nodeId} $\to$
        \textit{A.nodeId2}: use an intra-diagram edge instead.}  The link is
        skipped (a cross-diagram link within one cell is semantically
        incoherent).
  \item If the computed route length (sum of segment lengths) exceeds three
        times the poster canvas diagonal, emit a WARNING.  This guards against
        degenerate posters where two cells are placed at extreme corners of a
        very large canvas.
\end{itemize}

% ─────────────────────────────────────────────────────────────────────────────
\subsection{Linkable-Type Rules}
\label{sec:cdl-linkable-types}
% ─────────────────────────────────────────────────────────────────────────────

A diagram type is \textbf{linkable} if and only if its layout engine can
produce a non-empty \texttt{NodeAnchorRegistry} --- that is, it authors stable
ids for its visual nodes and can map those ids to bounding boxes.

\begin{table}[h]
\centering\small
\caption{Linkable diagram types (v1)}
\begin{tabularx}{\textwidth}{llX}
\toprule
\textbf{Grammar} & \textbf{Addressable unit} & \textbf{Node id source} \\
\midrule
\texttt{flowchart} / \texttt{graph} & Node & Author-assigned id (e.g.\ \texttt{A}, \texttt{pay}) \\
\texttt{classDiagram} & Class & Class name as written \\
\texttt{stateDiagram-v2} & State & State name / explicit id \\
\texttt{erDiagram} & Entity & Entity name as written \\
\texttt{C4Context} / \texttt{C4Container} / \texttt{C4Component} & Person, System, Container, Component & Alias string \\
\texttt{block-beta} & Block & Block id \\
\texttt{architecture-beta} & Service, Group & Id as written \\
\texttt{mindmap} & Node & Label-derived kebab-slug \\
\texttt{gitGraph} & Commit, Branch & Commit id or branch name \\
\bottomrule
\end{tabularx}
\label{tab:cdl-linkable}
\end{table}

\begin{table}[h]
\centering\small
\caption{Non-linkable diagram types (v1) --- excluded or deferred}
\begin{tabularx}{\textwidth}{llX}
\toprule
\textbf{Grammar} & \textbf{Reason excluded} & \textbf{Degradation} \\
\midrule
\texttt{pie} & No stable node ids; data is label+value pairs & WARN + skip link \\
\texttt{sankey} & Edge-based flow; no persistent node ids & WARN + skip link \\
\texttt{quadrantChart} & Continuous 2-D axes; data points have no ids & WARN + skip link \\
\texttt{radar} & Series labels, not structural nodes & WARN + skip link \\
\texttt{packet-beta} & Byte-field labels, not structural nodes & WARN + skip link \\
\texttt{xychart-beta} & Continuous axis; no node abstraction & WARN + skip link \\
\texttt{timeline} & Events are date-keyed, not node-keyed & WARN + skip link \\
\texttt{sequenceDiagram} & Participants are stable identifiers; deferred to v2 & WARN + skip link \\
\texttt{gantt} & Bars are task-keyed; deferred to v2 & WARN + skip link \\
\bottomrule
\end{tabularx}
\label{tab:cdl-non-linkable}
\end{table}

\paragraph{Rule statement.}

\begin{quote}
A \texttt{link} endpoint is resolvable if and only if (1) the referenced cell
address is valid, (2) the cell's diagram type appears in
Table~\ref{tab:cdl-linkable}, and (3) the specified \texttt{nodeId} is present
in that cell's \texttt{NodeAnchorRegistry} after compilation.  Any endpoint
that fails any condition is \textbf{unresolvable}: the compiler emits one
\texttt{WARNING} per unresolvable endpoint, and the entire \texttt{link}
statement is omitted from the overlay.  A link with one valid endpoint and one
unresolvable endpoint is wholly omitted (a dangling edge has no sensible
routing target).
\end{quote}

% ─────────────────────────────────────────────────────────────────────────────
\subsection{Semantics and Degradation Contract}
\label{sec:cdl-semantics}
% ─────────────────────────────────────────────────────────────────────────────

\begin{quote}
\textbf{Cross-diagram link contract.}
Cross-diagram links are \emph{presentation-layer assertions on the
composition}.  They do not modify, extend, or reference any individual
diagram's Domain IR or Scene IR; each cell compiles as if \texttt{link}
statements did not exist.  Cross-diagram links are \emph{superset-only}: a
standard Mermaid renderer will reject a poster document containing
\texttt{link} statements (the keyword is outside Mermaid's grammar), which is
the intended portability signal.  All link resolution --- cell lookup, node
anchor lookup, port selection, and route computation --- occurs at composition
time, after every cell has compiled independently.  An unresolvable link
(unknown cell address, unknown node id within a valid cell, non-linkable
diagram type, or empty anchor registry) emits \textbf{exactly one
\texttt{WARNING} per unresolvable endpoint} and the link is \emph{silently
omitted}.  The poster renders in full; all other valid links render normally.
\textbf{Cross-diagram links are never fatal.}
\end{quote}

\paragraph{Mermaid portability.}

As with the \texttt{poster} keyword itself (Section~\ref{sec:superset-extensions}),
a file containing \texttt{link} statements cannot be rendered by a standard
Mermaid renderer.  This is intentional and clearly signals the superset
requirement.  Individual cell bodies (the diagram source embedded in each
\texttt{cell} block) remain standard Mermaid syntax and can be extracted and
rendered independently by any Mermaid renderer.

\paragraph{IR integrity.}

Because cross-diagram links are resolved entirely at composition time and
rendered as overlay primitives, they require no changes to any grammar's
Domain IR schema, to the Scene IR primitive types, or to the determinism
contract.  The \texttt{sceneHash} of any individual cell's sub-scene is
byte-identical to what it would be without the \texttt{link} statement.

\paragraph{Agent addressing.}

The \texttt{nodeId} token in a \texttt{link} statement is the same identifier
an agent would use when addressing a node via the structured IR path
(Section~\ref{sec:agent-integration}).  This alignment is intentional: the
anchor registry (Section~\ref{sec:cdl-anchor-registry}) is the shared
mechanism that makes both the cross-diagram link syntax and agent
node-addressing stable and consistent.

% ─────────────────────────────────────────────────────────────────────────────
\subsection{Engineering Prerequisites}
\label{sec:cdl-prereqs}
% ─────────────────────────────────────────────────────────────────────────────

The following work items must be completed before any cross-diagram link can
be drawn.  They are listed in dependency order.

\begin{enumerate}
  \item \textbf{Define \texttt{NodeAnchorRegistry} and \texttt{RenderResult}
        types} in \texttt{packages/core/src/scene.ts} (or a new
        \texttt{scene-anchors.ts}).  This is a pure type addition; no
        behaviour changes.

  \item \textbf{Extend \texttt{buildFlowScene}} to return
        \texttt{RenderResult}.  The flow grammar already computes
        \texttt{PlacedNode} records; emitting the registry requires only
        a loop over the placed-node map at the end of Phase 4.  This is
        the \emph{reference implementation} that verifies the approach
        end-to-end.

  \item \textbf{Extend linkable grammar layout functions} (one per grammar,
        per Table~\ref{tab:cdl-linkable}) to emit anchors.  Each extension is
        independent and can be shipped incrementally.  Non-extended grammars
        return \texttt{anchors: \{\}}.

  \item \textbf{Extend the composition layer} (\texttt{layoutComposition} in
        \texttt{packages/core/src/composition/layout.ts}) to (a) collect anchor
        registries from each cell's \texttt{RenderResult}, (b) apply the
        coordinate transform, and (c) build the poster-coordinate anchor map.

  \item \textbf{Implement the link parser} in the poster DSL
        (\texttt{packages/core/src/frontend/mermaid/poster.ts}): recognise
        \texttt{link} statements and store them alongside the \texttt{cells}
        array in \texttt{PosterDocument}.

  \item \textbf{Implement the overlay routing pass} after the cell embed loop
        in \texttt{layoutComposition}: resolve endpoints, select ports, compute
        routes, and append overlay primitives.

  \item \textbf{Implement graceful degradation}: WARN on unresolvable
        endpoints, skip the link, continue.
\end{enumerate}

Items 1 and 5 are pure additions.  Items 2--4 and 6--7 are incremental
extensions to existing functions with no breaking changes to current callers
(the \texttt{Scene} return type is replaced by \texttt{RenderResult} for
linkable grammars; callers that ignore the \texttt{anchors} field are
unaffected).

% ─────────────────────────────────────────────────────────────────────────────
\subsection*{Cross-references}
% ─────────────────────────────────────────────────────────────────────────────

\begin{itemize}
  \item Composition / Posters (poster DSL, cell addressing):
        Section~\ref{sec:composition} and Section~\ref{sec:composition-posters}
  \item Dual cell addressing (bracket and Excel forms):
        Section~\ref{sec:composition-posters}
  \item Scene IR primitives (\texttt{GroupPrimitive.id},
        \texttt{PathPrimitive}): Section~\ref{sec:scene-ir} and
        \texttt{packages/core/src/scene.ts}
  \item \texttt{translateAndScale} coordinate helper:
        Section~\ref{sec:composition-embed} (paragraph \ref{para:translate-scale-helper})
  \item Superset extension mechanisms (\texttt{poster} keyword,
        superset-only features): Section~\ref{sec:superset-extensions}
  \item Agent integration surface (IR-as-agent-API, node-addressing):
        Section~\ref{sec:agent-integration}
  \item Dimension guard mindset: Section~\ref{sec:composition-edge-cases}
\end{itemize}
  % Cross-diagram node linking (superset-only)

% ============================================================================
% PART VII — ARCHITECTURE & PACKAGING
% ============================================================================

\part{Architecture \& Packaging}
\label{part:architecture}

\section{Layered Architecture}
\label{sec:architecture}

This section defines the layered architecture of the diagram compiler, emphasizing the clean separation between kernel, grammars, and composition.

\subsection{The Three Layers}

\ourdiagram[0.95\linewidth]{three-layers}{%
  The three-layer architecture: the Composition layer depends on peer Grammar modules,
  each Grammar depends on the Kernel.
  The Kernel (Scene IR, Backends, Themes, Determinism, Icons, Lint) is grammar-agnostic.}

\ourdiagram[0.95\linewidth]{architecture}{%
  The two-frontend, two-IR compilation pipeline.
  Mermaid DSL (\texttt{.mmd}) and structured IR (YAML/JSON) both feed into a Domain IR,
  which the layout engine transforms into a Scene IR consumed by SVG, PNG, and Skia backends.}

\subsection{Layer 1: The Kernel}
\label{sec:arch-kernel}

The kernel is the shared infrastructure that all grammars depend on. It is grammar-agnostic: it knows nothing about timelines, flows, or graphs.

\subsubsection{Kernel Components}

\begin{description}
  \item[Scene IR (§\ref{sec:scene-ir})] The universal primitive set (Rect, Circle, Line, Path, Text, Group) and effect definitions. The render contract.
  
  \item[Rendering Backends (§\ref{sec:backends})] SVG, PNG (resvg), Skia, PPTX backends. Consume Scene IR, produce output.
  
  \item[Theme System (§\ref{sec:themes})] Theme schema, theme loading, inheritance, property resolution. Grammar-agnostic but grammars may extend.
  
  \item[Determinism Contract (§\ref{sec:determinism})] Rounding rules, sort ordering, no-random, no-clock guarantees.
  
  \item[Deterministic Font Metrics] Embedded font metric tables for layout-critical text measurement.
  
  \item[Icon Registry] Named icons (SVG paths) available to all grammars.
  
  \item[Asset/Image Loader] Resolves and embeds external images as data URIs.
  
  \item[Layout Helpers] Common layout primitives: text measurement, box collision, orthogonal routing.
  
  \item[Lint/Quality-Gate Framework] Validation diagnostics, severity levels, suggestion generation.
  
  \item[Animation Hints (§\ref{sec:animation})] Declarative animation attached to Scene primitives.
\end{description}

\subsubsection{Kernel Contract}

The kernel makes the following guarantees to grammars:

\begin{enumerate}
  \item \textbf{Grammar-agnosticism}: The kernel never imports from, or references, any grammar module.
  
  \item \textbf{Determinism}: All kernel operations are deterministic given the same inputs.
  
  \item \textbf{Backend uniformity}: All backends consume the same Scene IR interface.
  
  \item \textbf{Theme uniformity}: Theme properties are resolved identically for all grammars.
\end{enumerate}

\subsection{Layer 2: Grammars}
\label{sec:arch-grammars}

Grammars are peer modules that sit atop the kernel. Each grammar:

\begin{itemize}
  \item Defines its own Domain IR schema
  \item Implements one or more layout engines (Domain IR → Scene IR)
  \item May extend the theme schema with grammar-specific properties
  \item Exposes validation and compilation entry points
\end{itemize}

\subsubsection{Grammar Independence}

Grammars are \emph{independent} of each other:

\begin{itemize}
  \item The Timeline grammar does not import from the Flow grammar
  \item Adding a new grammar does not modify existing grammars
  \item Each grammar can be versioned and released independently
\end{itemize}

\subsubsection{Grammar-Kernel Dependency}

All grammars depend on the kernel:

\begin{lstlisting}[language=javascript,caption={Grammar dependency structure}]
// packages/timeline/package.json
{
  "dependencies": {
    "@diagram-compiler/kernel": "^1.0.0"
  }
}

// packages/flow/package.json
{
  "dependencies": {
    "@diagram-compiler/kernel": "^1.0.0"
  }
}
\end{lstlisting}

\subsection{Layer 3: Composition}
\label{sec:arch-composition}

The Composition layer sits above grammars. It:

\begin{itemize}
  \item References panels, each specifying a grammar type and Domain IR
  \item Invokes grammar compilation for each panel
  \item Arranges panel Scenes into a unified poster Scene
  \item Adds editorial chrome (headers, footers, dividers)
\end{itemize}

\subsubsection{Composition Dependencies}

Composition depends on grammars (to compile panels) and kernel (for chrome rendering):

\begin{lstlisting}[language=javascript,caption={Composition dependency structure}]
// packages/composition/package.json
{
  "dependencies": {
    "@diagram-compiler/kernel": "^1.0.0",
    "@diagram-compiler/timeline": "^1.0.0",
    "@diagram-compiler/flow": "^1.0.0",
    "@diagram-compiler/comparison": "^1.0.0"
    // ... all available grammars
  }
}
\end{lstlisting}

\subsection{Dependency Rules}

\begin{center}
\small
\begin{tabular}{lcccc}
\toprule
\textbf{Module} & \textbf{→ Kernel} & \textbf{→ Grammars} & \textbf{→ Composition} & \textbf{→ Other Grammars} \\
\midrule
Kernel & — & ✗ & ✗ & ✗ \\
Grammar (any) & ✓ & — & ✗ & ✗ \\
Composition & ✓ & ✓ & — & — \\
\bottomrule
\end{tabular}
\end{center}

\paragraph{Key rule:} The kernel must not import anything grammar-specific. This is enforced by linting.

\subsection{Module Boundaries}

\subsubsection{What Lives in the Kernel}

\begin{itemize}
  \item Scene IR types and serialization
  \item All rendering backends
  \item Theme schema (core properties)
  \item Theme loader and resolver
  \item Icon registry
  \item Font metric tables
  \item Validation framework (DiagnosticSeverity, ValidationResult)
  \item Determinism utilities (rounding, sorting)
  \item Animation hint types
\end{itemize}

\subsubsection{What Lives in Grammars}

\begin{itemize}
  \item Domain IR types (e.g., Activity, Track, Milestone for Timeline)
  \item Domain IR schema (JSON Schema)
  \item Validation rules (semantic validation beyond schema)
  \item Layout engine(s)
  \item Grammar-specific theme extensions
\end{itemize}

\subsubsection{What Lives in Composition}

\begin{itemize}
  \item Composition IR types
  \item Composition IR schema
  \item Panel arrangement algorithms (stack, grid, custom)
  \item Chrome rendering (headers, footers, dividers)
  \item Grammar dispatcher (routes panel to correct grammar)
\end{itemize}

\subsection{API Surface}

Each layer exposes a clean API:

\subsubsection{Kernel API}

\begin{lstlisting}[language=javascript,caption={Kernel public API}]
// Scene IR
export type Scene, DrawingPrimitive, EffectDefinition, AnimationHint;
export function sceneHash(scene: Scene): string;

// Backends
export function renderSvg(scene: Scene): string;
export function renderPng(scene: Scene): Buffer;
export function renderPdf(scene: Scene): Buffer;
export function renderPptx(scene: Scene): Buffer;

// Themes
export function loadTheme(name: string): Theme;
export function resolveTheme(base: Theme, overrides: Partial<Theme>): Theme;

// Validation
export type Diagnostic, ValidationResult;

// Icons
export function getIcon(name: string): SvgPath | undefined;

// Font metrics
export function measureText(text: string, font: FontSpec): TextMetrics;
\end{lstlisting}

\subsubsection{Grammar API}

Each grammar exposes the same interface shape:

\begin{lstlisting}[language=javascript,caption={Grammar public API (Timeline example)}]
// @diagram-compiler/timeline
export type TimelineIR;  // Domain IR type
export const schema: JsonSchema;
export function validate(ir: unknown): ValidationResult;
export function compile(ir: TimelineIR, theme: Theme): Scene;
\end{lstlisting}

\subsubsection{Composition API}

\begin{lstlisting}[language=javascript,caption={Composition public API}]
// @diagram-compiler/composition
export type CompositionIR;
export const schema: JsonSchema;
export function validate(ir: unknown): ValidationResult;
export function compile(ir: CompositionIR, theme: Theme): Scene;
\end{lstlisting}

\section{Packaging \& Repository Strategy}
\label{sec:packaging}

This section defines the packaging strategy for the diagram compiler: how modules are organized in the repository, how they are published, and the incremental path from the current timeline-only codebase to the full grammar family.

\subsection{Strategic Principle}

The diagram compiler lives \textbf{beside timeline in the same monorepo on a shared kernel}---not inside timeline's IR (which would bloat it), not in a separate repo (which would duplicate the engine and drift).

\begin{figure}[h]
\centering
\begin{Verbatim}[frame=single,fontsize=\small]
packages/
  kernel/          @diagram-compiler/kernel
  timeline/        @diagram-compiler/timeline
  flow/            @diagram-compiler/flow        (future)
  graph/           @diagram-compiler/graph       (future)
  comparison/      @diagram-compiler/comparison  (future)
  composition/     @diagram-compiler/composition (future)
  cli/             @diagram-compiler/cli
  mcp/             @diagram-compiler/mcp
\end{Verbatim}
\caption{Target package structure}
\label{fig:package-structure}
\end{figure}

\subsection{Compounding Payoff}

This structure creates compounding value:

\begin{itemize}
  \item \textbf{Animation} added to the kernel is inherited by all grammars
  \item \textbf{New backends} serve all grammars immediately
  \item \textbf{Theme improvements} apply everywhere
  \item \textbf{Determinism tests} cover all grammars
  \item \textbf{Icon additions} are available to all grammars
  \item \textbf{Security patches} to the kernel propagate to all grammars
\end{itemize}

A new grammar starts with the full feature set: determinism, theming, SVG/PNG/PDF/PPTX output, animation support, validation---all for free.

\subsection{The Incremental Path}

The current codebase (\texttt{packages/core}) is timeline-only. Extracting the kernel requires care to avoid a premature big-bang refactor. The recommended path:

\subsubsection{Phase 0: Draw the Seam (Current)}

\textbf{Goal:} Establish the kernel/grammar boundary \emph{inside} \texttt{packages/core}, without restructuring packages.

\textbf{Actions:}
\begin{enumerate}
  \item Create \texttt{src/kernel/} directory within \texttt{packages/core}
  \item Move kernel components (Scene IR, backends, themes, icons, font metrics) into \texttt{src/kernel/}
  \item Move timeline-specific components (Timeline IR, layout engine) into \texttt{src/timeline/}
  \item Add lint rule: \texttt{src/kernel/} may not import from \texttt{src/timeline/}
  \item Verify all tests pass
\end{enumerate}

\textbf{Outcome:} The kernel/timeline boundary is enforced by convention and linting, but no packages are split. The monorepo structure is unchanged.

\subsubsection{Phase 1: Prove Grammar-Agnosticism}

\textbf{Goal:} Build the first new grammar (Flow) as a kernel-only consumer, proving the architecture.

\textbf{Actions:}
\begin{enumerate}
  \item Create \texttt{src/flow/} within \texttt{packages/core}
  \item Implement Flow grammar (IR, validation, layout) using only \texttt{src/kernel/} imports
  \item Add Flow grammar tests
  \item Verify the kernel imports nothing from \texttt{src/timeline/} or \texttt{src/flow/}
\end{enumerate}

\textbf{Outcome:} Two grammars (Timeline, Flow) coexist, both consuming the same kernel. The architecture is validated.

\subsubsection{Phase 2: Extract Packages}

\textbf{Goal:} Split into separate npm packages when publishing/team boundaries justify it.

\textbf{Actions:}
\begin{enumerate}
  \item Create \texttt{packages/kernel/} from \texttt{packages/core/src/kernel/}
  \item Create \texttt{packages/timeline/} from \texttt{packages/core/src/timeline/}
  \item Create \texttt{packages/flow/} from \texttt{packages/core/src/flow/}
  \item Update \texttt{packages/cli/} and \texttt{packages/mcp/} to depend on the new packages
  \item Publish as \texttt{@diagram-compiler/kernel}, \texttt{@diagram-compiler/timeline}, etc.
  \item Deprecate the old \texttt{@timeline-compiler/core} package (or keep as umbrella re-export)
\end{enumerate}

\textbf{Trigger:} Phase 2 is triggered when:
\begin{itemize}
  \item A third grammar is added (justifies the extraction overhead)
  \item OR external contributors need to develop grammars independently
  \item OR the monorepo becomes unwieldy
\end{itemize}

\textbf{Outcome:} Clean package boundaries, independent versioning, clear ownership.

\subsubsection{Phase 3: Grammar Marketplace (Future)}

\textbf{Goal:} Enable community-contributed grammars outside the core monorepo.

\textbf{Actions:}
\begin{enumerate}
  \item Publish kernel with stable API guarantees (semantic versioning)
  \item Document grammar authoring guide
  \item Create template repository for new grammars
  \item Establish discovery mechanism (registry, awesome-list)
\end{enumerate}

\textbf{Outcome:} Third-party grammars that interoperate with the kernel and can be used in compositions.

\subsection{Package Naming}

\begin{center}
\small
\begin{tabular}{ll}
\toprule
\textbf{Package} & \textbf{npm Name} \\
\midrule
Kernel & \texttt{@diagram-compiler/kernel} \\
Timeline Grammar & \texttt{@diagram-compiler/timeline} \\
Flow Grammar & \texttt{@diagram-compiler/flow} \\
Graph Grammar & \texttt{@diagram-compiler/graph} \\
Comparison Grammar & \texttt{@diagram-compiler/comparison} \\
Stat Grammar & \texttt{@diagram-compiler/stat} \\
Composition & \texttt{@diagram-compiler/composition} \\
CLI & \texttt{@diagram-compiler/cli} \\
MCP Server & \texttt{@diagram-compiler/mcp} \\
\bottomrule
\end{tabular}
\end{center}

\textbf{Note:} The existing \texttt{@timeline-compiler/*} packages may be retained as aliases or deprecated with migration guidance.

\subsection{Versioning Strategy}

\begin{description}
  \item[Kernel] Follows semantic versioning strictly. Breaking changes require major version bump.
  
  \item[Grammars] Version independently. May have different version numbers.
  
  \item[Composition] Version tracks the highest grammar version it supports.
  
  \item[CLI/MCP] Version tracks the kernel version.
\end{description}

\paragraph{Compatibility matrix.}
A grammar at version X is compatible with kernel versions [X, X+1). Major kernel versions may break grammars; grammars must be updated.

\subsection{Monorepo Tooling}

The monorepo uses:

\begin{description}
  \item[pnpm workspaces] For package management and linking
  \item[TypeScript project references] For incremental builds and cross-package type checking
  \item[Changesets] For versioning and changelog generation
  \item[Turborepo or Nx] (optional) For build caching and task orchestration
\end{description}

\subsection{CI/CD Considerations}

\begin{itemize}
  \item \textbf{Test matrix:} All grammars tested against kernel changes
  \item \textbf{Publish order:} Kernel publishes before grammars
  \item \textbf{Golden images:} Per-grammar golden test suites
  \item \textbf{Cross-platform:} macOS, Linux, Windows in CI
\end{itemize}

\section{Graph Auto-Layout: The Hard Problem}
\label{sec:layout-engines}

Graph auto-layout is the hardest problem in the diagram compiler. This section is honest about the difficulty, prescribes practical solutions, and identifies risks.

\subsection{Why Graph Layout Is Hard}

General 2D graph layout is research-grade:

\begin{itemize}
  \item \textbf{NP-hard subproblems}: Minimizing edge crossings in general graphs is NP-complete~\cite{garey1983}.
  \item \textbf{Aesthetic trade-offs}: No single algorithm optimizes all criteria (crossing minimization, edge length uniformity, node distribution, symmetry).
  \item \textbf{Non-determinism risk}: Many algorithms use randomized initialization or iterative convergence, which threatens the determinism contract.
  \item \textbf{Label placement}: Text labels add collision constraints that classic algorithms don't handle.
\end{itemize}

\subsection{The Constraint as a Feature}

The architecture embraces the constraint: \textbf{a small, opinionated grammar is both shippable AND reliably LLM-emittable}. Rather than attempting general graph layout, the compiler supports opinionated layouts per diagram sub-kind:

\begin{center}
\small
\begin{tabular}{llp{7cm}}
\toprule
\textbf{Sub-Kind} & \textbf{Algorithm} & \textbf{Constraints / Trade-offs} \\
\midrule
Tree & Reingold-Tilford~\cite{reingold1981} & Single root, no cycles. Clean, compact. \\
DAG (pipeline) & Sugiyama~\cite{sugiyama1981} & Directed acyclic. Layers + crossing minimization. \\
Hub-and-spoke & Radial custom & Single hub. Spokes arranged radially. \\
Left-to-right flow & Linear ranking & Nodes ranked by distance from sources. \\
Swimlane & Constrained grid & Nodes assigned to lanes; x by sequence. \\
\bottomrule
\end{tabular}
\end{center}

\paragraph{Explicit layout type.}
The Graph grammar requires the author to specify the layout type:

\begin{lstlisting}[language=yaml,caption={Explicit layout type in Graph IR}]
graph:
  layout: tidy-tree     # required: tidy-tree | dag | hub-spoke | swimlane
  vertices: [...]
  edges: [...]
\end{lstlisting}

This is a feature, not a limitation: the author (human or agent) declares the graph's structure, and the compiler uses an appropriate algorithm. Ambiguous ``auto-detect'' is out of scope.

\subsection{Algorithm Details}

\subsubsection{Tidy-Tree (Reingold-Tilford)}

For rooted trees: parent-child hierarchies with no cross-links.

\begin{description}
  \item[Algorithm] Reingold-Tilford (1981)~\cite{reingold1981}, extended by Walker (1990) for arbitrary trees.
  \item[Complexity] O(n) time.
  \item[Determinism] Fully deterministic given node order.
  \item[Implementation] Well-documented; available in dagre, ELK~\cite{elk}.
\end{description}

\subsubsection{DAG (Sugiyama Layered Layout)}

For directed acyclic graphs: pipelines, dependency graphs, workflow stages.

\begin{description}
  \item[Algorithm] Sugiyama, Tagawa \& Toda (1981)~\cite{sugiyama1981}: (1) cycle removal, (2) layer assignment, (3) crossing minimization, (4) x-coordinate assignment.
  \item[Complexity] O(|V| + |E|) for layers, O(|V|²) for crossing minimization heuristics.
  \item[Determinism] Deterministic with fixed layer-assignment and barycenter ordering.
  \item[Implementation] dagre~\cite{dagre}, ELK~\cite{elk}, Graphviz dot~\cite{graphviz}.
\end{description}

\subsubsection{Hub-and-Spoke (Radial)}

For star topologies: central hub with peripheral spokes.

\begin{description}
  \item[Algorithm] Custom radial placement: hub at center, spokes at equal angular intervals.
  \item[Complexity] O(n).
  \item[Determinism] Fully deterministic with sorted spoke order.
\end{description}

\subsubsection{Force-Directed (General Graphs)}

For graphs with no clear hierarchy or structure.

\begin{description}
  \item[Algorithm] Fruchterman-Reingold (1991)~\cite{fruchterman1991} or Eades spring model (1984)~\cite{eades1984}.
  \item[Complexity] O(|V|² × iterations).
  \item[Determinism] \textbf{Risk}: iterative simulation can converge to different local minima. Mitigations:
  \begin{itemize}
    \item Initial positions from sorted node IDs (deterministic starting point)
    \item Perturbations seeded by \texttt{scene\_hash}
    \item Fixed iteration count (no convergence-based termination)
  \end{itemize}
  \item[Implementation] d3-force, sigma.js, custom.
\end{description}

\paragraph{Force-directed caveat.}
Even with mitigations, force-directed layout is the least reliable for determinism. Small changes to nodes/edges can cause large layout changes. It is recommended for exploratory use, not production artifacts.

\subsection{Practical Layout Engines}

Two battle-tested open-source layout engines are recommended:

\subsubsection{ELK (Eclipse Layout Kernel)}
\label{sec:elk}

\begin{description}
  \item[Project] Eclipse Layout Kernel~\cite{elk}
  \item[License] EPL-2.0
  \item[Algorithms] Layered (Sugiyama), tree, force, radial, and more
  \item[Integration] Java library; JavaScript port (elkjs) via emscripten
  \item[Determinism] Designed for determinism; used in production diagramming tools
\end{description}

\subsubsection{dagre}
\label{sec:dagre}

\begin{description}
  \item[Project] dagre~\cite{dagre}
  \item[License] MIT
  \item[Algorithms] Sugiyama layered layout only
  \item[Integration] Pure JavaScript; widely used (Mermaid uses dagre)
  \item[Determinism] Deterministic for DAGs
\end{description}

\subsection{Build vs.\ Adopt Decision}

\begin{center}
\small
\begin{tabular}{lp{5cm}p{5cm}}
\toprule
\textbf{Option} & \textbf{Pros} & \textbf{Cons} \\
\midrule
Build custom & Full control; no dependencies; tailored to Scene IR & High effort; bugs; maintenance burden \\
Adopt ELK & Battle-tested; multiple algorithms; deterministic & Large dependency (elkjs ~2MB); EPL license \\
Adopt dagre & Small; MIT; widely adopted & DAG only; no tree/radial \\
Hybrid & dagre for DAGs; custom for simple cases & Two codepaths; integration complexity \\
\bottomrule
\end{tabular}
\end{center}

\paragraph{Recommendation.}
Start with dagre for DAG/pipeline layouts (already proven in Mermaid). Implement custom radial layout for hub-and-spoke. Evaluate ELK for tree layout. Build custom force-directed only if needed.

\subsection{Risk Assessment}

\begin{description}
  \item[Layout quality] Auto-layout may produce suboptimal results. Users may want manual position hints.
  
  \item[Determinism violations] Force-directed layout is risky. Recommend DAG/tree for production; force for drafts.
  
  \item[Label collisions] Standard algorithms don't handle labels. Post-processing required.
  
  \item[Performance] Large graphs (>1000 nodes) may be slow. Define reasonable limits.
\end{description}

\subsection{Layout Hints (Escape Hatch)}

For cases where auto-layout fails, the IR supports position hints:

\begin{lstlisting}[language=yaml,caption={Position hints in Graph IR}]
graph:
  layout: dag
  vertices:
    - id: start
      position: { x: 0, y: 0 }   # hint: override auto-layout for this node
    - id: end
      position: { x: 400, y: 0 }
  edges: [...]
\end{lstlisting}

Position hints are optional. When present, auto-layout respects them; when absent, auto-layout computes positions.

\section{Agent Integration}
\label{sec:agent-integration}

AI agents are first-class users of Timeline Compiler. They are not an afterthought wired to
a human-facing CLI; they are a primary consumer of the IR, the primary beneficiary of a
published JSON Schema, and the primary audience for the MCP server surface. This section
designs the full agent integration path: from raw inputs through ingestion to a valid IR,
through validation and repair, and through round-trip editing without loss of provenance.

The fundamental unit of agent work is the \emph{ingestion task}: given a \textbf{data source}
(work items, prose, structured export) and a \textbf{natural-language prompt}, produce a valid
Timeline IR. The prompt is a first-class input---it determines what is selected, how it is
framed, which entities are promoted to milestones, and which are dropped entirely. An agent
that ignores the prompt and imports everything it can find is not doing its job.


%%============================================================================
\subsection{The Ingestion Contract}
\label{sec:ingestion-contract}
%%============================================================================

The ingestion contract defines what the ingestion path is permitted to do with the inputs it
receives. Table~\ref{tab:ingestion-contract} captures the four categories of ingestion
decision. Violations of this contract---particularly silent data loss or silent inference---
are the most common source of incorrect or misleading timelines.

\begin{table}[h]
\centering
\small
\begin{tabular}{lp{4.5cm}p{6cm}}
\toprule
\textbf{Category} & \textbf{Definition} & \textbf{Examples} \\
\midrule
\textbf{Assumed} & Facts taken from the prompt or schema without source evidence & Title derived from prompt instruction; version \texttt{"1.0"} \\
\addlinespace
\textbf{Inferred} & Derived from source data by a deterministic rule & \texttt{axis\_unit} from date span; \texttt{time\_range} from min/max activity dates; track from \texttt{AreaPath} \\
\addlinespace
\textbf{Defaulted} & Omitted from IR, relying on documented schema default & \texttt{status: planned} omitted when all items are future; \texttt{locale} omitted \\
\addlinespace
\textbf{Rejected} & Source data discarded, reason logged & Bugs dropped when prompt says ``features and milestones only''; items outside prompted time window; duplicate or zero-duration events with no label \\
\bottomrule
\end{tabular}
\caption{Ingestion decision categories}
\label{tab:ingestion-contract}
\end{table}

\subsubsection{The Prompt as First-Class Input}
\label{sec:prompt-as-input}

The prompt directs every non-trivial ingestion decision. Table~\ref{tab:prompt-controls}
lists what the prompt controls and what the ingestion layer must extract from it before
touching the source data.

\begin{table}[h]
\centering
\small
\begin{tabular}{lp{9cm}}
\toprule
\textbf{Prompt Signal} & \textbf{Ingestion Effect} \\
\midrule
Time horizon (\textit{``Q1--Q3 2026''}) & Sets or constrains \texttt{time\_range}; items outside window are rejected \\
Entity filter (\textit{``features and milestones only''}) & Restricts work item types included; bugs, tasks dropped \\
Track grouping (\textit{``group by team''}) & Maps source field (e.g., \texttt{AreaPath}) to tracks \\
Emphasis (\textit{``highlight at-risk items''}) & Promotes \texttt{status} field; may add annotation \\
Framing (\textit{``executive summary''}) & Reduces label verbosity; promotes high-level items \\
Title / subtitle & Populates \texttt{metadata.title} / \texttt{metadata.subtitle} \\
Milestone promotion (\textit{``mark GA as a milestone''}) & Converts matching activity to milestone entity \\
\bottomrule
\end{tabular}
\caption{Prompt signals and their ingestion effects}
\label{tab:prompt-controls}
\end{table}

\subsubsection{What Ingestion May Not Do}

\begin{itemize}
  \item \textbf{Must not compute dates.} If a source item has no date, it must use
        \texttt{tbd}, not invent one. Timeline Grammar does not compute schedules
        (see \S\ref{sec:scope}).
  \item \textbf{Must not collapse ambiguous status silently.} An ADO work item in state
        ``On Hold'' that has no clear IR equivalent must log a rejection warning, not
        silently map to \texttt{planned}.
  \item \textbf{Must not import the whole backlog.} The prompt defines the editorial
        scope. Importing everything that exists and hoping the renderer trims it is a
        contract violation.
  \item \textbf{Must not generate non-stable IDs.} IDs derived from sequential numbers
        (\texttt{act-1}, \texttt{act-2}) are permitted only as a last resort. IDs should
        be derived from the item's title slug for git-friendliness (see
        \S\ref{sec:round-trip}).
\end{itemize}


%%============================================================================
\subsection{Ingestion Workflows}
\label{sec:ingestion-workflows}
%%============================================================================

The following subsections describe four concrete ingestion flows, each with a source
excerpt, a prompt, the resulting IR YAML, and an explicit account of what the prompt
selected, inferred, and rejected.

%% ---------------------------------------------------------------------------
\subsubsection{Azure DevOps Work Items}
\label{sec:ingest-ado}
%% ---------------------------------------------------------------------------

\textbf{Source format.} The ADO Work Items REST API \cite{ado-workitems} returns JSON
objects. The fields relevant to timeline ingestion are listed in
Table~\ref{tab:ado-field-mapping}.

\begin{table}[h]
\centering
\small
\begin{tabular}{lll}
\toprule
\textbf{ADO Field} & \textbf{IR Target} & \textbf{Notes} \\
\midrule
\texttt{System.Title} & \texttt{activity.label} / \texttt{milestone.label} & Verbatim or
  prompt-shortened \\
\texttt{System.State} & \texttt{activity.status} & Via Table~\ref{tab:ado-state-map} \\
\texttt{System.WorkItemType} & \texttt{activity.category} / entity type & Feature $\to$
  activity; Milestone $\to$ milestone \\
\texttt{System.IterationPath} & \texttt{span} or \texttt{start}/\texttt{end} & E.g.
  \texttt{Platform\textbackslash 2026\textbackslash Q2} $\to$ \texttt{span: 2026-Q2} \\
\texttt{System.AreaPath} & \texttt{track} & Leaf segment becomes track label/id \\
\texttt{Microsoft.VSTS.Scheduling.TargetDate} & \texttt{milestone.date} & Used when
  promoting to milestone \\
\texttt{System.Tags} & \texttt{activity.tags} & Passed through \\
\texttt{System.Id} & \texttt{metadata.ado\_id} & Provenance (not rendered) \\
\bottomrule
\end{tabular}
\caption{ADO work item field mapping}
\label{tab:ado-field-mapping}
\end{table}

\begin{table}[h]
\centering
\small
\begin{tabular}{ll}
\toprule
\textbf{ADO State} & \textbf{IR Status} \\
\midrule
New          & \texttt{planned} \\
Active       & \texttt{in-progress} \\
Resolved     & \texttt{done} \\
Closed       & \texttt{done} \\
Removed      & \texttt{cancelled} \\
On Hold      & \texttt{blocked} \\
\textit{(other)} & \texttt{planned} (warn: unmapped state) \\
\bottomrule
\end{tabular}
\caption{ADO state to IR status mapping}
\label{tab:ado-state-map}
\end{table}

\textbf{Source excerpt.}

\begin{lstlisting}[language=yaml,caption={ADO work item JSON (abridged)}]
# ADO API response (workItems array, fields selected)
- id: 1024
  fields:
    System.Title: "API Gateway Migration"
    System.State: "Active"
    System.WorkItemType: "Feature"
    System.IterationPath: "Platform\\2026\\Q2"
    System.AreaPath: "Platform Team"
- id: 1025
  fields:
    System.Title: "Mobile App 3.0 GA"
    System.State: "New"
    System.WorkItemType: "Milestone"
    System.AreaPath: "Mobile Team"
    Microsoft.VSTS.Scheduling.TargetDate: "2026-09-15"
- id: 1031
  fields:
    System.Title: "Fix login timeout bug"
    System.State: "Active"
    System.WorkItemType: "Bug"
    System.AreaPath: "Platform Team"
\end{lstlisting}

\textbf{Prompt.} \textit{``Show the Q1--Q3 2026 Platform and Mobile roadmap. Include
features and milestones only. Group by team. Drop bugs and tasks.''}

\textbf{Resulting IR.}

\begin{lstlisting}[language=yaml,caption={IR produced from ADO work items}]
version: "1.0"
metadata:
  title: Platform & Mobile Roadmap
  subtitle: Q1-Q3 2026
  time_range:
    start: 2026-Q1
    end: 2026-Q3
  axis_unit: quarter   # inferred: 3-quarter window -> quarter granularity

tracks:
  - id: platform
    label: Platform Team
  - id: mobile
    label: Mobile Team

activities:
  - id: api-gateway-migration
    label: API Gateway Migration
    track: platform
    span: 2026-Q2        # inferred from IterationPath "Platform\2026\Q2"
    status: in-progress   # mapped from ADO state "Active"
    category: feature
    metadata:
      source: ado
      ado_id: 1024
      ado_iteration: "Platform\\2026\\Q2"

milestones:
  - id: mobile-app-3-ga
    label: Mobile App 3.0 GA
    date: 2026-09-15    # from TargetDate field
    track: mobile
    status: planned      # mapped from ADO state "New"
    metadata:
      source: ado
      ado_id: 1025
\end{lstlisting}

\textbf{What the prompt did.}
\begin{itemize}
  \item \textbf{Selected:} Items 1024 (Feature) and 1025 (Milestone). Item 1031 (Bug)
        was rejected because the prompt said ``features and milestones only''.
  \item \textbf{Framed:} \texttt{time\_range} set to Q1--Q3 2026 from prompt instruction.
        Item 1025's target date (2026-09-15) falls within Q3, so it is included.
  \item \textbf{Inferred:} Tracks derived from \texttt{AreaPath} leaf segments. \texttt{axis\_unit}
        inferred as \texttt{quarter} from the three-quarter window.
  \item \textbf{Dropped:} Bug work item 1031 (not in filter). Had the time range excluded
        any item, it would be dropped with a logged rejection.
\end{itemize}


%% ---------------------------------------------------------------------------
\subsubsection{Natural Language Prose}
\label{sec:ingest-prose}
%% ---------------------------------------------------------------------------

When the source is free-form text---a planning document, a briefing note, an email
thread---the agent must extract temporal structure from prose. This is the most
ambiguity-prone ingestion path; the contract requires explicit handling of every point
of uncertainty.

\textbf{Source text.}

\begin{lstlisting}[caption={Natural language source}]
We are planning to migrate the data warehouse in the first half of 2026.
The new customer portal is targeted for Q3. Both initiatives are owned
by the Engineering org. An executive review is scheduled for 15 July 2026.
The legacy decommission date is TBD, pending the warehouse cutover.
\end{lstlisting}

\textbf{Prompt.} \textit{``Show this as a two-track roadmap: Infrastructure and Product.
Mark the executive review as a milestone. Include the decommission even though its date
is not fixed.''}

\textbf{Resulting IR.}

\begin{lstlisting}[language=yaml,caption={IR produced from natural language prose}]
version: "1.0"
metadata:
  title: 2026 Engineering Roadmap
  time_range:
    start: 2026-H1
    end: 2026-Q3
  axis_unit: quarter  # inferred: H1+Q3 window, quarterly detail

tracks:
  - id: infrastructure
    label: Infrastructure
  - id: product
    label: Product

activities:
  - id: data-warehouse-migration
    label: Data Warehouse Migration
    track: infrastructure
    span: 2026-H1
    status: planned

  - id: customer-portal
    label: New Customer Portal
    track: product
    span: 2026-Q3
    status: planned

  - id: legacy-decommission
    label: Legacy Decommission
    track: infrastructure
    start: tbd          # explicit: date not yet determined
    status: planned

milestones:
  - id: executive-review
    label: Executive Review
    date: 2026-07-15
    status: planned
\end{lstlisting}

\textbf{What the prompt did.}
\begin{itemize}
  \item \textbf{Structured:} Track assignments (Infrastructure vs.\ Product) came from the
        prompt. The source text only named ``Engineering org''.
  \item \textbf{Promoted:} The executive review was elevated to a milestone (not an
        activity) by explicit prompt instruction.
  \item \textbf{Included with uncertainty:} Legacy decommission included despite no
        fixed date; \texttt{start: tbd} preserves the fact on the timeline without
        inventing a date.
  \item \textbf{Inferred:} \texttt{time\_range} derived from the earliest (\texttt{2026-H1}
        start) and latest (\texttt{2026-07-15}) concrete dates. \texttt{axis\_unit}
        set to \texttt{quarter} as the coarsest granularity that fits the items.
\end{itemize}

\textbf{Ambiguity resolution.} ``First half of 2026'' is unambiguously \texttt{2026-H1}
(ISO half, per the date model). ``Q3'' without a year is inferred as \texttt{2026-Q3}
from context (all other dates are 2026). If the year were ambiguous, the agent must
surface a warning rather than guess.


%% ---------------------------------------------------------------------------
\subsubsection{GitHub Projects}
\label{sec:ingest-github}
%% ---------------------------------------------------------------------------

GitHub Projects (ProjectV2) \cite{github-projects,github-graphql-projectv2} supports
custom field types that carry roadmap-relevant data. The ingestion mapping is
summarised in Table~\ref{tab:github-field-mapping}.

\begin{table}[h]
\centering
\small
\begin{tabular}{lll}
\toprule
\textbf{GitHub ProjectV2 Field} & \textbf{IR Target} & \textbf{Notes} \\
\midrule
\texttt{title} (built-in) & \texttt{activity.label} / \texttt{milestone.label} & Verbatim or shortened \\
\texttt{status} (select field) & \texttt{activity.status} & Label mapped to IR enum \\
\texttt{DATE} custom field & \texttt{start} / \texttt{milestone.date} & Named field determines role \\
\texttt{ITERATION} custom field & \texttt{span} & Iteration title maps to quarter \\
\texttt{labels} & \texttt{activity.category} (first) / \texttt{activity.tags} & First label to category \\
\texttt{milestone} (linked) & \texttt{milestones} entry & Promoted to IR milestone \\
\texttt{repository} & \texttt{track} & When prompt says ``group by repo'' \\
\texttt{assignees} & \texttt{metadata.assignees} & Not rendered; preserved for provenance \\
\texttt{number} (issue \#) & \texttt{metadata.github\_issue} & Provenance link \\
\bottomrule
\end{tabular}
\caption{GitHub ProjectV2 field mapping}
\label{tab:github-field-mapping}
\end{table}

The ITERATION field type is the most useful for quarterly roadmaps: iteration titles
like ``2026 Q2'' map directly to \texttt{span: 2026-Q2}, making the ingestion path
from GitHub Projects to a quarterly roadmap nearly lossless.

\textbf{Gap note.} GitHub Projects does not have a first-class concept analogous to
``work item type'' (Feature vs.\ Bug vs.\ Task). Filtering by type requires label or
custom field conventions, which vary per team. The prompt must specify the filter
criterion explicitly.


%% ---------------------------------------------------------------------------
\subsubsection{Mermaid Timeline Syntax}
\label{sec:ingest-mermaid}
%% ---------------------------------------------------------------------------

Mermaid's \texttt{timeline} diagram \cite{mermaiddocs2024,mermaid2023} is the most
common structured timeline format already embedded in team documentation. Its syntax
maps cleanly to IR sections and activities.

\textbf{Source.}
\begin{lstlisting}[caption={Mermaid timeline source}]
timeline
    title Platform Evolution 2026
    section H1 2026
        API Gateway : Q1 : Migration complete
        Auth Refresh : Q2 : New OAuth2 provider
    section H2 2026
        Mobile 3.0 GA : Q3
        Legacy Cutover : Q4 : Final decommission
\end{lstlisting}

\textbf{Resulting IR.}
\begin{lstlisting}[language=yaml,caption={IR produced from Mermaid timeline syntax}]
version: "1.0"
metadata:
  title: Platform Evolution 2026
  time_range:
    start: 2026-Q1
    end: 2026-Q4
  axis_unit: quarter

tracks:
  - id: platform
    label: Platform      # default single-track when Mermaid has no swimlanes

sections:
  - id: h1-2026
    label: H1 2026
    time_range:
      start: 2026-Q1
      end: 2026-Q2

  - id: h2-2026
    label: H2 2026
    time_range:
      start: 2026-Q3
      end: 2026-Q4

activities:
  - id: api-gateway
    label: API Gateway
    track: platform
    span: 2026-Q1
    description: "Migration complete"
    status: planned

  - id: auth-refresh
    label: Auth Refresh
    track: platform
    span: 2026-Q2
    description: "New OAuth2 provider"
    status: planned

  - id: mobile-3-ga
    label: Mobile 3.0 GA
    track: platform
    span: 2026-Q3
    status: planned

  - id: legacy-cutover
    label: Legacy Cutover
    track: platform
    span: 2026-Q4
    description: "Final decommission"
    status: planned
\end{lstlisting}

\textbf{Mapping notes.}
Mermaid \texttt{timeline} does not have a swimlane concept; all events are rendered in
a single horizontal flow grouped by section. When promoting to IR, a default single track
is created unless the prompt specifies a multi-track structure. Mermaid's colon-separated
description field maps to \texttt{description}. Because Mermaid \texttt{gantt}
\cite{mermaidgantt2024} uses a different syntax (task dependencies, progress), Mermaid
gantt ingestion requires a separate mapping pass; that mapping is out of scope for the
initial IR but is an extension point (see \S\ref{sec:scope}).


%%============================================================================
\subsection{Reliable IR Generation by Agents}
\label{sec:agent-ir-generation}
%%============================================================================

LLM-based agents generate IR reliably when three conditions hold: (1) the target schema
is machine-readable and published, (2) the generation call is constrained to valid
output shapes, and (3) the IR's structural choices minimise the surface area for
generation errors.

%% ---------------------------------------------------------------------------
\subsubsection{JSON Schema Publication}
\label{sec:json-schema}
%% ---------------------------------------------------------------------------

The IR schema must be published as a JSON Schema document \cite{openai-structured-outputs}.
This is a binding requirement established by the positioning analysis
(\S\ref{sec:comparison}): without a machine-readable schema, agent generation is
unreliable. The schema is a first-class deliverable of this project, versioned alongside
the spec.

The top-level JSON Schema structure mirrors the IR document structure directly:

\begin{lstlisting}[language=yaml,caption={Abbreviated JSON Schema for the IR (YAML representation)}]
# $schema: https://json-schema.org/draft/2020-12/schema
# $id: https://timeline-compiler.dev/schema/v1/timeline.json
title: Timeline IR
type: object
required: [version, metadata, tracks, activities]
properties:
  version:
    type: string
    description: "Specification version, e.g. \"1.0\""

  metadata:
    type: object
    required: [title, time_range]
    properties:
      title:       { type: string }
      subtitle:    { type: string }
      time_range:
        type: object
        required: [start]
        properties:
          start: { "$ref": "#/$defs/Date" }
          end:   { "$ref": "#/$defs/Date" }
      axis_unit:   { "$ref": "#/$defs/AxisUnit" }
      theme:       { type: string }
      locale:      { type: string }

  tracks:
    type: array
    minItems: 1
    items: { "$ref": "#/$defs/Track" }

  activities:
    type: array
    items: { "$ref": "#/$defs/Activity" }

  milestones:
    type: array
    items: { "$ref": "#/$defs/Milestone" }

  # groups, annotations, sections, legend all optional arrays

$defs:
  ID:
    type: string
    pattern: "^[a-z][a-z0-9-]*$"

  Date:
    type: string
    description: "ISO date, quarter (2026-Q2), half (2026-H1), relative (+3m),
                  symbolic (now), or uncertain token (tbd, ongoing, unknown,
                  or tilde-prefix ~2026-Q2)"

  AxisUnit:
    type: string
    enum: [day, week, month, quarter, half, year]

  Status:
    type: string
    enum: [planned, in-progress, done, at-risk, blocked, cancelled, tentative]

  Track:
    type: object
    required: [id, label]
    properties:
      id:    { "$ref": "#/$defs/ID" }
      label: { type: string }
      color: { type: string }
      order: { type: integer, minimum: 0 }

  Activity:
    type: object
    required: [id, label, track]
    oneOf:
      - required: [start]    # start + optional end
      - required: [span]     # span shorthand
    properties:
      id:       { "$ref": "#/$defs/ID" }
      label:    { type: string }
      track:    { type: string }   # bare string ref to Track.id
      start:    { "$ref": "#/$defs/Date" }
      end:      { "$ref": "#/$defs/Date" }
      span:     { "$ref": "#/$defs/Date" }
      status:   { "$ref": "#/$defs/Status", default: planned }
      progress: { type: number, minimum: 0, maximum: 1 }
      category: { type: string }
      group:    { type: string }
      description: { type: string }
      metadata: { type: object }

  Milestone:
    type: object
    required: [id, label, date]
    properties:
      id:     { "$ref": "#/$defs/ID" }
      label:  { type: string }
      date:   { "$ref": "#/$defs/Date" }
      track:  { type: string }
      status: { "$ref": "#/$defs/Status", default: planned }
      category: { type: string }
      icon:     { type: string }
      description: { type: string }
      metadata: { type: object }
\end{lstlisting}

The schema is published at a versioned URL. Minor version bumps to the IR may add
optional fields; the schema handles this via \texttt{additionalProperties: true} for
\texttt{metadata} maps and by allowing unknown top-level fields on minor revisions.

\paragraph{New cite key needed.} The JSON Schema specification itself does not have a
cite key in \texttt{references.bib}. A key \texttt{json-schema2020} pointing to the
IETF JSON Schema draft (2020-12) should be added by David.


%% ---------------------------------------------------------------------------
\subsubsection{Structured Output Constraints}
\label{sec:structured-outputs}
%% ---------------------------------------------------------------------------

Platforms that support constrained generation \cite{openai-structured-outputs} should
use the published JSON Schema as the output schema for the IR generation call. This
eliminates entire classes of generation errors:

\begin{itemize}
  \item Required fields can never be missing when the schema enforces \texttt{required}.
  \item ID format is enforced by the \texttt{pattern} constraint on the \texttt{ID}
        definition, preventing spaces, uppercase, or other invalid characters.
  \item Status values are constrained to the enum; the model cannot hallucinate
        ``completed'' or ``active''.
  \item References are just bare strings---the model emits \texttt{track: "platform"},
        not a nested object, which is trivially correct.
\end{itemize}

When structured output is not available, the agent prompt must include the JSON Schema
inline and instruct the model to validate its own output before returning it.

%% ---------------------------------------------------------------------------
\subsubsection{IR Design Choices That Help Agents}
\label{sec:ir-agent-friendly}
%% ---------------------------------------------------------------------------

Several deliberate IR design choices reduce the error rate for LLM-generated documents:

\begin{description}
  \item[Optional fields default to safe values.] \texttt{status} defaults to
        \texttt{planned}; \texttt{progress} defaults to \texttt{null}. An agent that
        omits these fields produces a valid document. Agents only need to emit fields
        they have evidence for.

  \item[\texttt{span} shorthand reduces reasoning load.] Instead of requiring an agent
        to compute both a \texttt{start} and \texttt{end} date for a quarter-length
        activity, it can emit \texttt{span: 2026-Q2} and let the compiler expand it.
        This avoids off-by-one boundary errors (is Q2 end June 30 or July 1?).

  \item[References are bare strings.] \texttt{track: platform} is far easier for an
        LLM to emit correctly than a nested reference object. The schema context (field
        name) determines the target type; no wrapper object is needed.

  \item[IDs are kebab-case slugs.] Agents derive IDs from labels:
        ``API Gateway Migration'' $\to$ \texttt{api-gateway-migration}. This is a
        simple and reliable transformation. UUIDs would be safe but opaque; sequential
        numbers would collide with re-generation.

  \item[No dependency scheduling.] Since the IR is a Timeline Grammar (not a Task
        Scheduling Grammar), there is no dependency resolution. Each activity is
        independently valid. Agents do not need to compute a topological order or
        check resource constraints.
\end{description}


%%============================================================================
\subsection{Validation Pipeline}
\label{sec:validation-pipeline}
%%============================================================================

Validation is layered. Each layer is a gate: failure at an earlier layer produces
errors that block later layers. This design prevents cryptic failures (e.g., a
reference-resolution error caused by a silently malformed ID).

\begin{table}[h]
\centering
\small
\begin{tabular}{clp{4cm}p{4.5cm}}
\toprule
\textbf{Layer} & \textbf{Name} & \textbf{What it checks} & \textbf{Failure mode} \\
\midrule
1 & Syntactic parse       & Valid YAML or JSON; no parse errors & Hard error: stop \\
2 & Schema conformance    & Required fields present; types match; enum values valid; ID regex; \texttt{oneOf} for start/span & Hard error: list all violations \\
3 & Well-formedness       & Referential integrity; ID uniqueness; date ordering; progress bounds; group acyclicity & Hard error: list all violations \\
4 & Render-readiness      & See \S\ref{sec:rendering} (Barbara) & Error or warning depending on severity \\
5 & Semantic advisory     & Degenerate documents; empty \texttt{activities} and \texttt{milestones}; items outside \texttt{time\_range} & Warning only \\
\bottomrule
\end{tabular}
\caption{Validation pipeline layers}
\label{tab:validation-layers}
\end{table}

\textbf{Layer 4 dependency.} Render-readiness validation---checking that the IR can
produce a legible visual output---is defined and owned by Barbara in \S\ref{sec:rendering}.
This includes checks such as overlapping activities on the same track that cannot be
visually separated, labels that exceed track width at the chosen axis unit, and milestone
dates that fall outside the document's \texttt{time\_range}. The validation pipeline
calls Barbara's render-readiness checks as Layer 4; this section does not duplicate them.

\subsubsection{Errors vs.\ Warnings}

\begin{description}
  \item[Error] Prevents rendering. The document is invalid and the renderer must refuse
        to proceed. Errors come from Layers 1--3 and from Layer 4 severity ``fatal''.
        The validator must return \emph{all} errors in a single pass, not just the first.
  \item[Warning] The document is valid but the output may be suboptimal. Warnings come
        from Layer 4 (advisory) and Layer 5. Renderers proceed and may annotate the
        output (e.g., a clipped activity) or log the warning and continue.
\end{description}


%%============================================================================
\subsection{Error Handling and Agent Repair Cycle}
\label{sec:error-repair}
%%============================================================================

\subsubsection{Error Message Format}

Every validation error is \textbf{path-anchored}: it identifies the exact location in
the document, the nature of the violation, and a suggested fix. This design enables
an agent to mechanically apply repairs without re-reading the whole document.

The error message format is:

\begin{lstlisting}[caption={Validation error message format}]
ERROR  <path>: <description>
       Expected: <constraint>
       Found:    <actual value>
       Fix:      <suggested correction>
\end{lstlisting}

\textbf{Concrete error examples:}

\begin{lstlisting}[caption={Example validation errors with suggested fixes}]
# Missing required field
ERROR  activities[0].id: required field 'id' is missing
       Expected: string matching ^[a-z][a-z0-9-]*$
       Found:    (absent)
       Fix:      add id: api-gateway-migration

# Invalid ID format
ERROR  activities[1].id: "API Migration" does not match ^[a-z][a-z0-9-]*$
       Expected: kebab-case slug
       Found:    "API Migration"
       Fix:      id: api-migration

# Unresolved track reference
ERROR  activities[2].track: "eng" references a non-existent track
       Expected: one of [platform, mobile, data]
       Found:    "eng"
       Fix:      track: platform  (or add a track with id: eng)

# Invalid status value
ERROR  activities[3].status: "active" is not a valid Status
       Expected: one of [planned, in-progress, done, at-risk,
                          blocked, cancelled, tentative]
       Found:    "active"
       Fix:      status: in-progress

# Date ordering violation
ERROR  activities[4]: end precedes start
       Expected: end >= start
       Found:    start: 2026-Q3, end: 2026-Q1
       Fix:      swap dates, or use span: 2026-Q3 if this is
                 a single-quarter activity

# Progress out of range
ERROR  activities[5].progress: 1.4 is outside [0, 1]
       Expected: number in [0, 1]
       Found:    1.4
       Fix:      progress: 1.0

# Duplicate ID
ERROR  milestones[1].id: "platform-launch" is already used by activities[2].id
       Expected: globally unique identifier
       Found:    duplicate
       Fix:      id: platform-launch-ms  (or another unique slug)

# Missing date source (neither start nor span present)
ERROR  activities[6]: no temporal anchor defined
       Expected: exactly one of: (start [+ optional end]) or span
       Found:    neither start nor span present
       Fix:      add  span: 2026-Q2  or  start: 2026-04-01
\end{lstlisting}

\subsubsection{Agent Repair Cycle}
\label{sec:repair-cycle}

The validation-error-repair loop is the core of reliable agent IR generation. The
cycle is:

\begin{enumerate}
  \item \textbf{Generate.} Agent generates an IR document, optionally using structured
        output constraints (\S\ref{sec:structured-outputs}).
  \item \textbf{Validate.} The IR is submitted to the \texttt{validate\_timeline} MCP
        tool (\S\ref{sec:mcp-tools}). All layers are run; all errors collected.
  \item \textbf{Report.} Errors are returned to the agent as a structured list with
        path, description, and suggested fix.
  \item \textbf{Repair.} The agent applies targeted fixes. Because errors are
        path-anchored and include suggested fixes, repairs are surgical: the agent
        patches \texttt{activities[2].track} rather than regenerating the entire document.
  \item \textbf{Re-validate.} The patched document is re-submitted. This loop repeats
        until validation passes with no errors (warnings may remain).
  \item \textbf{Render.} Once validation passes, the agent calls \texttt{render\_timeline}
        to produce output.
\end{enumerate}

In practice, structured output constraints eliminate most Layer~2 errors on the first
generation pass. Layers~3 (referential integrity, date ordering) are the most common
source of residual errors for agents generating multi-track documents.


%%============================================================================
\subsection{MCP Server}
\label{sec:mcp-server}
%%============================================================================

The MCP server \cite{mcp-spec} exposes Timeline Compiler as a tool-callable service for
AI agents. It is a first-class distribution target, not an afterthought (see the
distribution architecture in \S\ref{sec:distribution}). The server hosts four tools.

%% ---------------------------------------------------------------------------
\subsubsection{Tool Contracts}
\label{sec:mcp-tools}
%% ---------------------------------------------------------------------------

\paragraph{\texttt{validate\_timeline}}

Runs the full validation pipeline (Layers 1--5) against a provided IR document.
Returns all errors and warnings in structured form.

\begin{lstlisting}[language=yaml,caption={\texttt{validate\_timeline} tool contract}]
tool: validate_timeline
description: >
  Validate a Timeline IR document. Runs syntactic, schema, well-formedness,
  render-readiness, and semantic advisory checks. Returns structured errors
  and warnings with path anchors and suggested fixes.

input:
  ir:
    type: object | string   # IR as parsed object or raw YAML/JSON string
    required: true
  stop_at_layer:
    type: integer           # optional; 1-5; default: 5 (run all layers)
    required: false

output:
  valid:     boolean        # true iff zero errors (warnings do not block)
  errors:
    - path:    string       # e.g. "activities[2].track"
      code:    string       # machine-readable error code, e.g. UNRESOLVED_REF
      message: string       # human-readable description
      fix:     string       # suggested correction
  warnings:
    - path:    string
      code:    string
      message: string
  layers_run: integer       # how many layers were executed before stopping
\end{lstlisting}

\paragraph{\texttt{render\_timeline}}

Renders a valid Timeline IR to a visual output format. Returns the rendered output
as a base64-encoded string or a file path (for local MCP deployments).

\begin{lstlisting}[language=yaml,caption={\texttt{render\_timeline} tool contract}]
tool: render_timeline
description: >
  Render a Timeline IR document to SVG, PNG, PDF, or PPTX. The IR must
  be valid (validate_timeline returns valid: true) before rendering.
  Rendering is deterministic: identical IR + theme + format always
  produces bit-identical output.

input:
  ir:
    type: object | string   # IR as parsed object or raw YAML/JSON string
    required: true
  format:
    type: string
    enum: [svg, png, pdf, pptx]
    default: svg
    required: false
  theme:
    type: string            # theme identifier, e.g. "default", "corporate"
    default: "default"
    required: false
  width:
    type: integer           # output width in pixels (PNG/PDF only)
    required: false

output:
  format:  string           # echo of requested format
  data:    string           # base64-encoded output bytes
  mime_type: string         # e.g. "image/svg+xml"
  warnings:                 # render-layer warnings (does not block output)
    - message: string
\end{lstlisting}

\paragraph{\texttt{describe\_schema}}

Returns the full JSON Schema for the IR and field-level documentation. Agents call
this to bootstrap IR generation or to recover from schema errors.

\begin{lstlisting}[language=yaml,caption={\texttt{describe\_schema} tool contract}]
tool: describe_schema
description: >
  Return the JSON Schema for the Timeline IR, plus human-readable
  documentation for each field. Use this to understand what fields are
  required, their types, allowed values, and defaults.

input:
  version:
    type: string            # IR version to describe; default: latest
    required: false
  entity:
    type: string            # optional: limit to one entity, e.g. "Activity"
    required: false

output:
  version:     string       # IR version described
  schema:      object       # full JSON Schema document
  docs:
    - entity:  string       # entity name
      fields:
        - name:        string
          type:        string
          required:    boolean
          default:     string | null
          description: string
          example:     string
\end{lstlisting}

\paragraph{\texttt{suggest\_time\_range}}

A convenience tool for agents that need to derive a sensible \texttt{time\_range} and
\texttt{axis\_unit} from a collection of candidate activities before composing the full
document. This is especially useful when ingesting data from sources that lack explicit
time windows.

\begin{lstlisting}[language=yaml,caption={\texttt{suggest\_time\_range} tool contract}]
tool: suggest_time_range
description: >
  Given a list of activity date hints (start, end, span values), suggest
  appropriate time_range and axis_unit values for the IR metadata. Does not
  require a complete IR document.

input:
  dates:
    type: array             # list of date strings, e.g. ["2026-Q1", "2026-09-15"]
    items: { type: string }
    required: true
  prefer_unit:
    type: string            # optional hint: "quarter", "month", etc.
    required: false

output:
  time_range:
    start: string
    end:   string
  axis_unit: string         # recommended AxisUnit
  rationale: string         # explanation of the choice
\end{lstlisting}

\subsubsection{MCP Server Deployment}

The MCP server is deployed in two modes:

\begin{description}
  \item[Local server] Runs as a subprocess alongside the agent environment. Started
        via the CLI (\texttt{timeline mcp-server --port 3000}) or as an npm/pip
        package. Suitable for development, CI, and agent runtimes on the same host.
  \item[Hosted endpoint] A cloud-deployed instance of the MCP server. Suitable for
        cloud-native agent runtimes where a local subprocess is impractical.
\end{description}

Both modes expose identical tool contracts. The local server operates on the local
filesystem (output file paths are local); the hosted server returns base64-encoded
output in all cases.


%%============================================================================
\subsection{Round-Trip and Iterative Editing}
\label{sec:round-trip}
%%============================================================================

A Timeline IR document lives in version control alongside the source data it was derived
from. Humans and agents edit it directly. Re-syncing against an updated source must not
silently overwrite editorial decisions made in the IR. This section designs the mechanisms
that make round-trip editing safe.

%% ---------------------------------------------------------------------------
\subsubsection{Source Provenance via Metadata}
\label{sec:provenance}
%% ---------------------------------------------------------------------------

Every entity produced by an ingestion workflow includes a \texttt{metadata} block that
records the source of record:

\begin{lstlisting}[language=yaml,caption={Provenance metadata block (ADO example)}]
activities:
  - id: api-gateway-migration
    label: API Gateway Migration
    track: platform
    span: 2026-Q2
    status: in-progress
    metadata:
      source: ado
      ado_id: 1024
      ado_revision: 42           # ADO ETag for change detection
      ado_area: "Platform Team"
      ado_iteration: "Platform\\2026\\Q2"
      ingested_at: 2026-06-09
\end{lstlisting}

\begin{lstlisting}[language=yaml,caption={Provenance metadata block (GitHub Projects example)}]
activities:
  - id: customer-portal-redesign
    label: Customer Portal Redesign
    track: product
    span: 2026-Q3
    status: planned
    metadata:
      source: github
      github_project_id: PVT_kwDOA12345
      github_issue: 214
      github_url: "https://github.com/acme/portal/issues/214"
      ingested_at: 2026-06-09
\end{lstlisting}

The \texttt{source} key in \texttt{metadata} is a reserved provenance marker. The
ingester writes it; the renderer ignores it. Re-sync logic uses it to correlate IR
entities back to their source records.

\subsubsection{Re-Sync Strategy}
\label{sec:resync}

When source data is updated (e.g., a work item's state changes from Active to Resolved),
the re-sync operation must:

\begin{enumerate}
  \item \textbf{Match by source ID.} Use \texttt{metadata.ado\_id} or
        \texttt{metadata.github\_issue} as the stable foreign key, not the IR \texttt{id}.
        The IR \texttt{id} is a human-edited slug that may have been renamed.
  \item \textbf{Update only source-mapped fields.} Fields that the ingestion workflow
        maps from the source (e.g., \texttt{status}, \texttt{span} derived from
        IterationPath) are overwritten. Fields the human may have edited in the IR
        (\texttt{label}, \texttt{category}, \texttt{description}, \texttt{color},
        \texttt{progress}) are \emph{not} overwritten.
  \item \textbf{Preserve the IR \texttt{id}.} The kebab-case slug must not be
        regenerated on re-sync. Once set, it is stable.
  \item \textbf{Log what changed.} The re-sync operation must produce a structured
        change log (entity id, field, old value, new value) before writing, enabling
        human review.
  \item \textbf{Reject destructive deletions.} If a source item is deleted or removed
        from scope, the re-sync must flag it as a warning rather than silently removing
        the IR entity. A human or agent must confirm removal.
\end{enumerate}

\subsubsection{Stable IDs and Git-Friendly YAML}
\label{sec:stable-ids}

The combination of stable kebab-case IDs and line-oriented YAML serialisation produces
IR documents that diff well under version control:

\begin{lstlisting}[caption={Git diff of an IR update (status change + progress update)}]
  activities:
    - id: api-gateway-migration
      label: API Gateway Migration
      track: platform
      span: 2026-Q2
-     status: in-progress
-     progress: 0.4
+     status: done
+     progress: 1.0
      metadata:
        source: ado
        ado_id: 1024
\end{lstlisting}

Key properties of the YAML serialisation that preserve git-friendliness:

\begin{itemize}
  \item Fields within each entity are emitted in a \textbf{canonical order} (required
        fields first, optional fields in schema definition order). Reordering fields
        does not change semantics but does generate spurious diffs; canonical order
        prevents this.
  \item \textbf{No auto-generated IDs or timestamps} in non-\texttt{metadata} fields.
        The IR ID is derived from content and frozen; \texttt{ingested\_at} in
        \texttt{metadata} only changes when the entity is re-ingested from source.
  \item \textbf{Consistent quoting conventions.} Strings that are safe as bare YAML
        scalars are not quoted; strings with special characters use double quotes.
        Serialisers must apply this consistently to avoid quote-flip diffs.
  \item \textbf{Lists maintain insertion order}, not alphabetical. New entities appended
        to the end of a list produce a single-line diff at the bottom.
\end{itemize}

\subsubsection{Human Edits and Incremental Refinement}

A human author may open the IR YAML in any editor, add annotations, adjust labels,
add a milestone the ingester missed, or change a \texttt{status} value to reflect a
decision made after ingestion. These edits are valid IR operations and must survive
re-sync (per the strategy above).

An agent may similarly refine an IR: for example, an editorial agent might re-read the
document, notice that five activities in one track are unlabelled, and add descriptions.
Because the agent is editing a YAML file with stable IDs and a well-defined schema, this
is a safe, auditable operation. The MCP \texttt{validate\_timeline} tool can be called
after any edit to confirm the document remains valid.

The editing workflow is:

\begin{lstlisting}[caption={Iterative refinement loop (human or agent)}]
1. ingest(source, prompt)   -> roadmap.yaml          # initial ingestion
2. validate(roadmap.yaml)   -> valid, 0 errors        # confirm valid
3. edit(roadmap.yaml)       -> roadmap.yaml (modified) # human/agent edit
4. validate(roadmap.yaml)   -> valid, 0 errors        # re-confirm
5. render(roadmap.yaml, svg) -> roadmap.svg            # deterministic output
6. (commit roadmap.yaml + roadmap.svg to git)
\end{lstlisting}

Step~6 commits both the IR source and the rendered output. Because rendering is
deterministic, the SVG in the repository is always reproducible from the YAML. The YAML
is the source of truth; the SVG is a derived artefact that can be regenerated on demand.


%%============================================================================
\subsection{IR Gaps and Open Questions}
\label{sec:ir-gaps}
%%============================================================================

During the design of this section, three gaps or ambiguities in the IR contract were
identified. These are flagged here for Leslie (reviewer) and Mark to reconcile; they
have \emph{not} been silently resolved in this section's examples.

\begin{enumerate}
  \item \textbf{\texttt{today} field missing from metadata.} Leslie's scope decisions
        specify that the \texttt{now} symbolic date must be evaluated from an explicit
        \texttt{today} field in the IR (not from the system clock), in order to preserve
        determinism. However, the \texttt{metadata} table in Mark's IR contract does not
        include a \texttt{today} field. Until this is resolved, the annotation type
        \texttt{today-marker} with \texttt{date: now} is non-deterministic. Suggested
        resolution: add \texttt{today: date?} to the \texttt{metadata} schema; if absent,
        render-time clock is used and a warning is emitted.

  \item \textbf{Track sort field: \texttt{index} vs.\ \texttt{order}.} Mark's binding
        contract (well-formedness invariant~14) refers to a \texttt{track.index} field
        with the rule ``track index order values unique and non-negative''. The
        \texttt{04-ir.tex} specification defines the same field as \texttt{order}
        (type \texttt{int?}). The canonical field name should be resolved to one term;
        this section uses \texttt{order} to match the spec, but the contract should
        be updated to match.

  \item \textbf{\texttt{span} vs.\ \texttt{start}+\texttt{end} mutual exclusion.}
        The IR allows either \texttt{span} (shorthand) or \texttt{start} plus optional
        \texttt{end}, but does not specify the behaviour when both are present (e.g.,
        \texttt{span: 2026-Q2} and \texttt{start: 2026-05-01} coexist in the same
        activity). Should this be a schema error, or should one take precedence?
        For agent generation, structured output can prevent this from occurring, but the
        validation specification needs an explicit ruling. Suggested resolution: treat
        co-presence of \texttt{span} and \texttt{start}/\texttt{end} as a schema error.
\end{enumerate}
     % Agent surface: MCP/SDK/CLI

% ============================================================================
% PART VIII — ROADMAP & MVP
% ============================================================================

\part{Roadmap \& MVP}
\label{part:roadmap}

\section{Coverage Roadmap}
\label{sec:roadmap}

The product roadmap is organized around the \textbf{tiered coverage strategy}: achieving
full Mermaid compatibility in priority order, then extending beyond Mermaid.

\subsection{Strategy: Mermaid-Complete First, Net-New After}

The sequencing discipline:
\begin{enumerate}
  \item Cover all 22 Mermaid diagram types (Tiers 0--3).
  \item Establish the aesthetic bar for each type (beat Mermaid's rendering).
  \item Only then invest in net-new diagram types with no Mermaid equivalent.
\end{enumerate}

This ensures the migration story (``drop-in replacement with better output'') is complete
before pursuing differentiation that requires user education.

\subsection{Tier 0: Wire Existing Grammars to Mermaid Syntax}

\textbf{Types:} flowchart, sequence, gantt, timeline, mindmap (5 types).

\textbf{Work required:}
\begin{itemize}
  \item Mermaid syntax parsers for each type (PEG/Nearley grammars)
  \item CST $\to$ Domain IR lowering for each parser
  \item Compatibility test suite (parse Mermaid examples $\to$ verify output)
  \item Frontmatter/directive handling in parser layer
\end{itemize}

\textbf{Kernel work:} None. All four layout engines are built and tested.

\textbf{Exit criteria:} Standard Mermaid examples from \texttt{mermaid.js.org} documentation
parse and render with our compiler, producing visually superior output.

\subsection{Tier 1: The UML/Software Line}

\textbf{Types:} class, state, erDiagram, C4Context, C4Container, C4Component, C4Dynamic
(7 types).

\textbf{Work required:}
\begin{itemize}
  \item New Domain IRs: ClassDiagram IR, StateDiagram IR, ER IR
  \item C4 IR (extends Flow IR with stereotype/boundary semantics)
  \item Layout: class/state use layered layout (reuse Sugiyama kernel); ER uses
        force-directed or layered
  \item Mermaid syntax parsers for all UML types
  \item UML-specific theme extensions (stereotypes, compartments, cardinality labels)
  \item ``Modern UML'' visual standard (beating PlantUML's dated rendering)
\end{itemize}

\textbf{Kernel work:} Minor extensions to node rendering (compartmented boxes, stereotype
headers, cardinality decorators on edges).

\textbf{Exit criteria:} Real-world UML diagrams from popular open-source projects render
beautifully. Side-by-side with PlantUML shows clear aesthetic superiority.

\subsection{Tier 2: The Chart Family}

\textbf{Types:} pie, xychart, quadrant, radar (4 types).

\textbf{Work required:}
\begin{itemize}
  \item New Chart Domain IR and JSON Schema (Section~\ref{sec:chart-family})
  \item Scale computation library (linear, band, radial)
  \item Chart layout engine compiling to Scene IR
  \item Mermaid parsers for all four chart syntaxes
  \item Chart-specific theme extensions (axis styling, gridlines, palettes, legends)
\end{itemize}

\textbf{Kernel work:} The chart layout engine is genuinely new architecture---the only
family that requires a new kernel rather than reusing an existing one.

\textbf{Exit criteria:} All four chart types render from Mermaid syntax with proper axes,
legends, and responsive sizing. Output quality comparable to simple Vega-Lite charts.

\subsection{Tier 3: Remaining Types}

\textbf{Types:} sankey, requirement, gitGraph, block, architecture, treemap, kanban,
journey, packet (9 types).

\textbf{Work required per type:}
\begin{itemize}
  \item Mermaid syntax parser
  \item Domain IR definition (most are minor variants of existing IRs)
  \item Layout integration (most reuse node-link, tree, or timeline kernels)
  \item Type-specific rendering (sankey flow widths, git branch colors, etc.)
\end{itemize}

\textbf{Kernel work:} Minimal. Sankey requires weighted-edge rendering; treemap requires
squarified rectangle packing; others reuse existing layout families.

\textbf{Exit criteria:} All 22 Mermaid diagram types render correctly.

\subsection{MVP Definition (Reframed)}

The MVP is defined as:

\begin{description}
  \item[Tier 0 complete] All five existing grammars accessible via Mermaid syntax parsers.
  \item[Aesthetic bar met] Output demonstrably superior to Mermaid for all Tier-0 types.
  \item[Agent IR path working] Structured IR input for all Tier-0 types with published
        JSON Schemas.
  \item[CLI + library] \texttt{timeline render input.mmd -o output.svg} works.
  \item[Composition] Multi-diagram posters from superset DSL.
\end{description}

The MVP answers: \emph{Can users drop in Mermaid files and get better output, with zero
migration cost, plus the agent IR and composition capabilities Mermaid lacks?}

\subsection{Post-MVP: Net-New Diagram Types}

After Mermaid-complete coverage, new diagram types with no Mermaid equivalent:
\begin{itemize}
  \item \textbf{Architecture Decision Records} (ADR visualization)
  \item \textbf{Dependency graphs} (package/module dependency visualization)
  \item \textbf{Infrastructure diagrams} (cloud architecture beyond C4)
  \item \textbf{Data pipeline diagrams} (ETL/streaming topology)
\end{itemize}

These are post-Tier-3 and driven by user demand, not pre-planned.

\subsection{Implementation Status Summary}

\begin{table}[h]
\centering\small
\caption{Current implementation status vs.\ roadmap tiers}
\begin{tabular}{llll}
\toprule
\textbf{Component} & \textbf{Tier} & \textbf{Status} & \textbf{Tests} \\
\midrule
Timeline grammar + themes   & 0 & Built & 551 \\
Flow grammar + themes       & 0 & Built & 33 \\
Sequence grammar + themes   & 0 & Built & 37 \\
Tree grammar + themes       & 0 & Built & 26 \\
Composition layer           & --- & Built & 25 \\
Animation (SMIL)            & --- & Built & incl.\ in flow \\
Schema validation           & --- & Built & 13 \\
Mermaid DSL parsers         & 0 & \textbf{Planned} & --- \\
UML/software line           & 1 & Planned & --- \\
Chart family                & 2 & Planned & --- \\
Remaining types             & 3 & Planned & --- \\
\bottomrule
\end{tabular}
\label{tab:implementation-status}
\end{table}

\textbf{Total tests passing:} 795 (all grammars + composition + schema + CLI).
All existing goldens byte-identical across runs.
               % Tier 0-3 coverage roadmap
\input{sections/51-distribution}
\input{sections/53-oss-strategy}
\input{sections/55-target-outputs}

% ----------------------------------------------------------------------------
% BIBLIOGRAPHY
% ----------------------------------------------------------------------------

\bibliographystyle{plainnat}
\bibliography{references}

\end{document}
